% 高一63_吴淞中学_10月月考.tex
%   —— 高一上数学“2026-2027 年吴淞中学高一上 10 月月考”（讲评版）
% 试卷版：xelatex 高一63_吴淞中学_10月月考.tex
% 答案版：xelatex "\def\isanswer{1}% 高一63_吴淞中学_10月月考.tex
%   —— 高一上数学“2026-2027 年吴淞中学高一上 10 月月考”（讲评版）
% 试卷版：xelatex 高一63_吴淞中学_10月月考.tex
% 答案版：xelatex "\def\isanswer{1}% 高一63_吴淞中学_10月月考.tex
%   —— 高一上数学“2026-2027 年吴淞中学高一上 10 月月考”（讲评版）
% 试卷版：xelatex 高一63_吴淞中学_10月月考.tex
% 答案版：xelatex "\def\isanswer{1}% 高一63_吴淞中学_10月月考.tex
%   —— 高一上数学“2026-2027 年吴淞中学高一上 10 月月考”（讲评版）
% 试卷版：xelatex 高一63_吴淞中学_10月月考.tex
% 答案版：xelatex "\def\isanswer{1}\input{高一63_吴淞中学_10月月考.tex}"
% 紧凑版：xelatex "\def\iscompact{1}\input{高一63_吴淞中学_10月月考.tex}"
%
% 原始材料：微信公众号“魔都数学加油站”文章，作者破晓，连载“27学年高一63”；
% 原文标题“27学年高一63——2026-2027年上海市吴淞中学高一上10月月考”；
% 原文链接：https://mp.weixin.qq.com/s/TFQ4TLnt_l66t7t51es4OQ。
% 页面标注发布时间 2026-10-16 21:44（上海）。
% 正文为 10 张 PNG 页（`原图/高一63_吴淞中学_10月月考/01..10.png`）：
% 第 06、10 页 2560×3620，第 01、02、03、04、05、07、08、09 页 4762×6735
% （**同一篇各页分辨率不同**），均按页序读取。10 页全为试卷正文页（卷首标题在第 1 页页顶，
% 卷末解答题第 21 题解析止于第 10 页），无封面页、目录页、二维码或宣传页。
% 粉色斜排“魔都数学加油站”为卷面水印，与公众号页脚一并未录入。
%
% 卷首标题：照卷面印作“2026--2027 年吴淞中学高一上 10 月月考”（公众号标题中的“上海市”
% 卷面标题行没有，本稿照卷面），年份区间按全稿惯例排 en dash。原卷只有一行标题、无副标题。
%
% 篇幅与题号：原卷自第 3 题起编号（第 1、2 题不在本卷内）。填空题 3—12（10 题）、
% 选择题 13—16（4 题）、解答题 17—21（5 题），共 19 题。本稿沿用原节与原题号，
% 只改排为全稿统一的大节样式（原卷小节标题为“一、填空题”“二、选择题”“三、解答题”）。
%
% 讲评版：原卷每题下直接印【答案】或【解析】。**只印【答案】的 2 题**（第 4、7 题），
% 本稿只给 \ans{}；其余 17 题里 16 题都印【解析】（按全稿口径这类题不另提 \ans{}——第 13、15、16 题
% 的结论写在解析末句（“故选 B”等），亦照录在解析内；第 14 题则既无【答案】也无【解析】段。**另有 1 题答案印在题干括号内**——
% 第 14 题（“(  D  )”，且该题无【答案】/【解析】段），本稿按 高一37 第 13 题成例把括号留空、
% 答案另提为 \ans{$D$}。第 16 题题干括号为空但印有【解析】，不算括号内答案。
%
% 副标题：原卷未印副标题。本稿按“知识模块”口径标作“集合与常用逻辑用语、不等式”：
% 19 题中集合与常用逻辑用语 12 题（第 3、4、6、7、9、10、11、12、13、15、18、21 题），
% 不等式 7 题（含等式（韦达定理、恒等式）与不等式性质、基本不等式、二次不等式、恒成立：
% 第 5、8、14、16、17、19、20 题），两类均超过三分之一，故并标。口径同 高一59。
%
% 与原卷的排版差异：
%   ① 原卷每题下直接印【答案】/【解析】，本稿按全稿惯例将答案与解析置于答案版，试卷版保持干净；
%   ② 原卷未印副标题，本稿按“知识模块”口径补出；
%   ③ 原卷填空横线约四个字宽，本稿按全稿惯例统一排 \blankline[4.4em]；
%   ④ 原卷实数集排斜体 $R$（第 9、10、19 题题干与解析）或粗体 $\mathbf{R}$（第 18 题题干），
%      整数集排斜体 $Z$（第 12 题解析），自然数集排斜体 $N$（第 20 题解析、第 21 题题干
%      “$(n\geqslant2,n\in N)$”）或粗体 $\mathbf{N}$（第 21 题 (3) 题干）、斜体 $N^*$
%      （第 21 题 (3) 解析），本稿按全稿惯例统一改用 $\R$、$\Z$、$\N$、$\N^*$
%      （同 高一14 排版差异④：记号归一，粗体／斜体不作区分，$N$ 与 $N^*$ 的区别照原卷保留）；
%      空集记号原卷作 $\varnothing$，照录；
%   ⑤ 原卷在数学式内多处用半角逗号（如 $\{y\mid y=[2x]+[3x],-1\leqslant x\leqslant2\}$、
%      $m=-1,0,1$、$A=\{x\mid-2<x<15\}$,$B=\cdots$），均照录未强改；
%      句末句点按全稿惯例排西文句点；“和积平衡集”自定义名词照录；
%   ⑥ 第 12 题解析里的“法一：”“法二：”标号原卷排**红色**（全卷共 2 处），
%      同 高一53“情况1：”成例，本稿排 \stepm{法一：}{…}、\stepm{法二：}{…}
%      （加粗着色、不改黑色基调），冒号照原卷排全角；
%   ⑦ 第 17—21 题的子问原卷各自另起一行，本稿按全稿惯例用 enumerate 的 (1)(2)(3)；
%      解析中的子问标号一律用 \stepm{(1)}{…}；第 19 题解析中的①②分情形，按全稿惯例作
%      \stepm{①}{…}、\stepm{②}{…}；
%   ⑧ 第 14 题印在题干括号内的答案移入 \ans{}，见上；
%   ⑨ 每道解答题原卷未留作答空白，本稿按全稿惯例加 \ansspace；
%   ⑩ 选择题的作答括号原卷排半角“(　　)”（第 13、15、16 题括号内空白、第 14 题括号内印 D），
%      本稿按全稿惯例排 \mbox{（\quad）}。
%
% 原卷笔误/存疑（均照录未改，正文对应处另有 % 注）：
%   ① 第 21 题题干“$(n\geqslant2$,$n\in N)$”后连印两个全角逗号“，，”（衍字，
%      8 倍放大确认，非水印干扰），本稿照录；
%   ② 第 21 题 (2) 题干将“写出以 $a_1$、$a_2$ 为根的一个一元二次方程”与“用反证法证明：
%      $a_1$、$a_2$ 至少有一个大于 2”两项任务并作一句、之间仅以逗号相接（“……表示），
%      可尝试用反证法证明：……”），读来像抄并两小问，原卷如此，照录未改。
%
% 插图：无（全卷 10 页无任何图形、无“如图”字样）。
%
\documentclass[a4paper,11pt]{article}
\usepackage{common}

\begin{document}

\examtitlest[3]{2026--2027 年吴淞中学高一上 10 月月考}{集合与常用逻辑用语、不等式}

\bigsec{一、填空题}

\begin{enumerate}
\setcounter{enumi}{2}

\item 满足 $A\subseteq\{1,2,3,4\}$ 的集合 $A$ 共有\blankline[4.4em]个.

\ifanswer
\solbegin
\step{由题意，集合 $A$ 是集合 $\{1,2,3,4\}$ 的子集，共有 $2^4=16$ 个.}
\solend
\fi

\item 陈述句“$a>1$ 或 $b>3$”的否定形式为\blankline[4.4em].

\ifanswer
\ans{$a\leqslant1$ 且 $b\leqslant3$}
\fi

\item 已知方程 $3x^2+4x-4=0$ 的两个根为 $x_1$、$x_2$，则 $x_1^2+x_2^2=$\blankline[4.4em].

\ifanswer
\solbegin
\step{对于方程 $3x^2+4x-4=0$，因为 $\Delta=16-4\times3\times(-4)=64>0$，}
\step{所以方程 $3x^2+4x-4=0$ 的两个不相等的实根为 $x_1$、$x_2$，}
\step{由韦达定理可得 $x_1+x_2=-\dfrac43$，$x_1x_2=-\dfrac43$，}
\step{所以 $x_1^2+x_2^2=(x_1+x_2)^2-2x_1x_2=\dfrac{16}9+\dfrac83=\dfrac{40}9$.}
\solend
\fi

\item 设集合 $A=\{y\mid y=x^2+1\}$,$B=\{y\mid y=-x^2+2x+3\}$，则 $A\cap B=$\blankline[4.4em].

\ifanswer
\solbegin
\step{因为 $y=x^2+1\geqslant1$,$y=-x^2+2x+3=-(x-1)^2+4\leqslant4$，}
\step{所以 $A=\{y\mid y\geqslant1\}$,$B=\{y\mid y\leqslant4\}$，得 $A\cap B=\{y\mid1\leqslant y\leqslant4\}$.}
\solend
\fi

\item 若 $x\in\{1,x^2\}$，则 $x=$\blankline[4.4em].

\ifanswer
\ans{$0$}
\fi

\item 已知等式 $2x^2+3x+5=a(2x+1)(x+1)+c$，则 $a+c=$\blankline[4.4em].

\ifanswer
\solbegin
\step{因为 $2x^2+3x+5=a(2x+1)(x+1)+c$ 恒成立，}
\step{所以 $2x^2+3x+5=2ax^2+3ax+a+c$ 恒成立，则 $\begin{cases}2=2a\\ 3=3a\\ 5=a+c\end{cases}$，解得 $a+c=5$.}
\solend
\fi

\item 已知 $a\in\R$，设 $p:x<2$，$q:x<5a-3$，且 $p$ 是 $q$ 的必要非充分条件，则 $a$ 的取值范围是\blankline[4.4em].

\ifanswer
\solbegin
\step{设 $p:x<2$，$q:x<5a-3$，且 $p$ 是 $q$ 的必要非充分条件，}
\step{则 $\{x\mid x<5a-3\}\subset\{x\mid x<2\}$，所以 $5a-3<2$，即 $a<1$.}
\solend
\fi

\item 已知 $a\in\R$，集合 $A=\{x\mid ax^2+3x-1=0\}$，若集合 $A$ 中只有一个元素，则 $a$ 的取值范围是\blankline[4.4em].

\ifanswer
\solbegin
\step{当 $a=0$ 时，则有 $\{x\mid ax^2+3x-1=0\}=\{x\mid3x-1=0\}=\left\{\dfrac13\right\}$，符合题意；}
\step{当 $a\neq0$ 时，由题意得 $\Delta=9+4a=0$，解得 $a=-\dfrac94$. 符合题意；}
\step{综上所述，实数 $a$ 的取值集合为 $\left\{0,-\dfrac94\right\}$.}
\solend
\fi

\item 设集合 $A=\{x\mid x^2+2(m+1)x+m^2-1=0,m\in\R\}$，且 $A\subseteq\{0,-4\}$，则实数 $m$ 的取值范围是\blankline[4.4em].

\ifanswer
\solbegin
\step{由题意得 $A=\varnothing$ 或 $A=\{0\}$ 或 $A=\{-4\}$ 或 $A=\{0,-4\}$，}
\step{当 $A=\varnothing$ 时，方程 $x^2+2(m+1)x+m^2-1=0$ 无实根，}
\step{则 $\Delta=[2(m+1)]^2-4\times1\times(m^2-1)=8m+8<0$，解得 $m<-1$，}
\step{当 $A=\{0\}$ 时，方程 $x^2+2(m+1)x+m^2-1=0$ 有一个实根 $0$，}
\step{则 $\begin{cases}\Delta=8m+8=0\\ m^2-1=0\end{cases}$，解得 $m=-1$，}
\step{当 $A=\{-4\}$ 时，方程 $x^2+2(m+1)x+m^2-1=0$ 有一个实根 $-4$，}
\step{则 $\begin{cases}\Delta=8m+8=0\\ 16-8(m+1)+m^2-1=0\end{cases}$，无解，}
\step{当 $A=\{0,-4\}$ 时，方程 $x^2+2(m+1)x+m^2-1=0$ 有两个实根 $-4$ 和 $0$，}
\step{则 $\begin{cases}\Delta=8m+8>0\\ 16-8(m+1)+m^2-1=0\\ m^2-1=0\end{cases}$，解得 $m=1$，}
\step{综上所述，实数 $m$ 的取值范围是 $\{m\mid m\leqslant-1\text{ 或 }m=1\}$.}
\solend
\fi

\item 对于任意实数 $x$，$[x]$ 表示不超过 $x$ 的最大整数，如 $[1.1]=1$，$[-1.2]=-2$，若
$A=\{y\mid y=[2x]+[3x],-1\leqslant x\leqslant2\}$，则 $A$ 中所有元素的和为\blankline[4.4em].

\ifanswer
\solbegin
% 原卷“法一：”“法二：”标号排红色，本稿按 高一53“情况1：”成例排 \stepm（加粗着色），见文件头
\stepm{法一：}{当 $-1\leqslant x<-\dfrac23$ 时，$[2x]=-2$，$[3x]=-3$，$f(x)=[2x]+[3x]=-5$，}
\step{当 $-\dfrac23\leqslant x<-\dfrac12$ 时，$[2x]=-2$，$[3x]=-2$，$f(x)=[2x]+[3x]=-4$；}
\step{当 $-\dfrac12\leqslant x<-\dfrac13$ 时，$[2x]=-1$，$[3x]=-2$，$f(x)=[2x]+[3x]=-3$；}
\step{当 $-\dfrac13\leqslant x<0$ 时，$[2x]=-1$，$[3x]=-1$，$f(x)=[2x]+[3x]=-2$；}
\step{当 $0\leqslant x<\dfrac13$ 时，$[2x]=0$，$[3x]=0$，$f(x)=[2x]+[3x]=0$；}
\step{当 $\dfrac13\leqslant x<\dfrac12$ 时，$[2x]=0$，$[3x]=1$，$f(x)=[2x]+[3x]=1$；}
\step{当 $\dfrac12\leqslant x<\dfrac23$ 时，$[2x]=1$，$[3x]=1$，$f(x)=[2x]+[3x]=2$；}
\step{当 $\dfrac23\leqslant x<1$ 时，$[2x]=1$，$[3x]=2$，$f(x)=[2x]+[3x]=3$；}
\step{当 $1\leqslant x<\dfrac43$ 时，$[2x]=2$，$[3x]=3$，$f(x)=[2x]+[3x]=5$；}
\step{当 $\dfrac43\leqslant x<\dfrac32$ 时，$[2x]=2$，$[3x]=4$，$f(x)=[2x]+[3x]=6$；}
\step{当 $\dfrac32\leqslant x<\dfrac53$ 时，$[2x]=3$，$[3x]=4$，$f(x)=[2x]+[3x]=7$；}
\step{当 $\dfrac53\leqslant x<2$ 时，$[2x]=3$，$[3x]=5$，$f(x)=[2x]+[3x]=8$；}
\step{当 $x=2$ 时，$[2x]=4$，$[3x]=6$，$f(x)=[2x]+[3x]=10$}
\step{故 $A$ 中所有元素的和为 $-5-4-3-2+0+1+2+3+5+6+7+8+10=28$.}
% 原卷“法二：”标号排红色，同上
\stepm{法二：}{设 $x=m+n$，其中 $m\in\Z$，$n\in[0,1)$，}
\step{若 $0\leqslant n<\dfrac13$，则 $0\leqslant2n<\dfrac23$，$0\leqslant3n<1$，}
\step{得 $[2x]=[2m+2n]=2m$，$[3x]=[3m+3n]=3m$，所以 $y=[2x]+[3x]=5m$；}
\step{若 $\dfrac13\leqslant n<\dfrac12$，则 $\dfrac23\leqslant2n<1$，$1\leqslant3n<\dfrac32$，}
\step{得 $[2x]=[2m+2n]=2m$，$[3x]=[3m+3n]=3m+1$，}
\step{所以 $y=[2x]+[3x]=5m+1$；}
\step{若 $\dfrac12\leqslant n<\dfrac23$，则 $1\leqslant2n<\dfrac43$，$\dfrac32\leqslant3n<2$，}
\step{得 $[2x]=[2m+2n]=2m+1$，$[3x]=[3m+3n]=3m+1$，}
\step{所以 $y=[2x]+[3x]=5m+2$；}
\step{若 $\dfrac23\leqslant n<1$，则 $\dfrac43\leqslant2n<2$，$2\leqslant3n<3$，}
\step{得 $[2x]=[2m+2n]=2m+1$，$[3x]=[3m+3n]=3m+2$，}
\step{所以 $y=[2x]+[3x]=5m+3$；}
\step{综上所述，当 $n\in[0,1)$ 时，所有函数值的和}
\step{为 $5m+5m+1+5m+2+5m+3=20m+6$，}
\step{又因为 $x\in[-1,2]$，此时 $m=-1,0,1$，且当 $x=2$ 时，$y=[2x]+[3x]=10$，}
\step{所有函数值的和为 $(-14)+6+26+10=28$，即 $A$ 中所有元素之和为 $28$.}
\solend
\fi

\end{enumerate}

\bigsec{二、选择题}

\begin{enumerate}
\setcounter{enumi}{12}

\item 下列写法中，正确的是\mbox{（\quad）}

\keepwithprev
A. $0\in\{\varnothing\}$\qquad B. $\varnothing\in\{\varnothing\}$

C. $\varnothing=\{\varnothing\}$\qquad D. $0\in\varnothing$

\ifanswer
\solbegin
\step{$0\notin\{\varnothing\}$，$\varnothing\in\{\varnothing\}$，故 $A$ 错误，$B$ 正确；}
\step{$\varnothing$ 是不含任何元素的集合，所以 $\varnothing\neq\{\varnothing\}$，$0\notin\varnothing$，故 $CD$ 错误.}
\step{故选 $B$.}
\solend
\fi

% 第 14 题答案“D”直接印在题干括号内“(  D  )”，且该题无【答案】/【解析】段；
% 按 高一37 第 13 题的成例括号留空、答案另提为 \ans{}（见文件头）
\item 若 $a<b<0$，则下列不等式中不能成立的是\mbox{（\quad）}

\keepwithprev
A. $\dfrac1a>\dfrac1b$\qquad B. $a-b>a$

C. $|a|>|b|$\qquad D. $a^3>b^3$

\ifanswer
\ans{$D$}
\fi

\item 已知 $x$、$y$、$z$ 为非零实数，代数式 $\dfrac x{|x|}+\dfrac y{|y|}+\dfrac z{|z|}+\dfrac{xyz}{|xyz|}$
的值所组成的集合是 $M$，则下列判断正确的是\mbox{（\quad）}

\keepwithprev
A. $-4\notin M$\qquad B. $2\in M$

C. $4\notin M$\qquad D. $0\in M$

\ifanswer
\solbegin
\step{当 $x$、$y$、$z$ 均为正数时，$xyz>0$，$\dfrac x{|x|}+\dfrac y{|y|}+\dfrac z{|z|}+\dfrac{xyz}{|xyz|}=1+1+1+1=4$，}
\step{当 $x$、$y$、$z$ 有 2 个正数，1 个负数时，不妨设 $x>0$，$y>0$，$z<0$，}
\step{则 $xyz<0$，$\dfrac x{|x|}+\dfrac y{|y|}+\dfrac z{|z|}+\dfrac{xyz}{|xyz|}=1+1-1-1=0$，}
\step{若 $x$、$y$、$z$ 有 1 个正数，2 个负数时，不妨设 $x>0$，$y<0$，$z<0$，}
\step{则 $xyz>0$，$\dfrac x{|x|}+\dfrac y{|y|}+\dfrac z{|z|}+\dfrac{xyz}{|xyz|}=1-1-1+1=0$，}
\step{当 $x$、$y$、$z$ 均为负数时，$xyz<0$，}
\step{$\dfrac x{|x|}+\dfrac y{|y|}+\dfrac z{|z|}+\dfrac{xyz}{|xyz|}=-1-1-1-1=-4$，}
\step{综上，$M=\{-4,0,4\}$，}
\step{所以 $-4\in M$，$A$ 错误；$2\notin M$，$B$ 错误；$4\in M$，$C$ 错误；}
\step{$0\in M$，$D$ 正确. 故选 $D$.}
\solend
\fi

\item 若 $xy>0$，且 $x+y=9$，则 $\dfrac1x+\dfrac4y$ 的最小值是\mbox{（\quad）}

\keepwithprev
A. $9$\qquad B. $\dfrac59$

C. $\dfrac49$\qquad D. $1$

\ifanswer
\solbegin
\step{$xy>0$，且 $x+y=9$，所以 $x>0$，$y>0$，}
\step{所以 $\dfrac1x+\dfrac4y=\dfrac19\left(\dfrac1x+\dfrac4y\right)(x+y)=\dfrac59+\dfrac19\left(\dfrac yx+\dfrac{4x}y\right)\geqslant\dfrac59+\dfrac29\cdot\sqrt{\dfrac yx\cdot\dfrac{4x}y}=\dfrac59+\dfrac49=1$，}
\step{当且仅当 $\dfrac yx=\dfrac{4x}y$，即 $x=3$，$y=6$ 时取等号，所以 $\dfrac1x+\dfrac4y$ 的最小值为 $1$.}
\step{故选 $D$.}
\solend
\fi

\end{enumerate}

\bigsec{三、解答题}

\begin{enumerate}
\setcounter{enumi}{16}

\item 求解下列不等式：

（1）$\dfrac{x-3}{x+1}\leqslant0$；

（2）$|2x-1|<3$.

\ifanswer
\solbegin
\stepm{(1)}{$\dfrac{x-3}{x+1}\leqslant0$，变形得 $(x-3)\cdot(x+1)\leqslant0$ 且 $x\neq-1$，}
\step{解得 $-1<x\leqslant3$，即不等式的解集为 $(-1,3]$.}
\stepm{(2)}{若 $|2x-1|<3$，则 $-3<2x-1<3$，解得 $-1<x<2$，即 $(-1,2)$.}
\solend
\fi

\ansspace[6]

\item 已知集合 $A=\{x\mid-2<x<15\}$,$B=\{x\mid m-6<x<2m-1,m\in\R\}$.

（1）当 $m=2$ 时，求 $A\cap B$；

（2）若 $A\cup B=A$，求实数 $m$ 的取值范围.

\ifanswer
\solbegin
\stepm{(1)}{当 $m=2$ 时，$A=\{x\mid-2<x<15\}$,$B=\{x\mid-4<x<3\}$，}
\step{此时 $A\cap B=\{x\mid-2<x<3\}$.}
\stepm{(2)}{由 $A\cup B=A$ 得 $B\subseteq A$.}
\step{当 $B=\varnothing$,$m-6\geqslant2m-1$，解得 $m\leqslant-5$；}
\step{当 $B\neq\varnothing$，结合 $B\subseteq A$ 得 $\begin{cases}m-6<2m-1\\ 2m-1\leqslant15\\ m-6\geqslant-2\end{cases}$，解得 $4\leqslant m\leqslant8$.}
\step{综上，$m\leqslant-5$ 或 $4\leqslant m\leqslant8$.}
\solend
\fi

\ansspace[6]

\item 已知 $a\in\R$，命题 $p$：关于 $x$ 的方程 $x^2-2x+a=0$ 无实根；命题 $q$：对任意 $x\in\R$，
不等式 $ax^2-ax+1>0$ 恒成立.

（1）若 $q$ 为真命题，求 $a$ 的取值范围；

（2）若 $p$、$q$ 中有且只有一个为真命题，求 $a$ 的取值范围.

\ifanswer
\solbegin
\stepm{(1)}{当 $q$ 为真命题时，}
\stepm{①}{当 $a=0$ 时，不等式化为 $1>0$，对任意 $x\in\R$ 恒成立，满足条件；}
\stepm{②}{当 $a\neq0$ 时，$ax^2-ax+1>0$ 为一元二次不等式，}
\step{若对任意 $x\in\R$ 恒成立，需满足 $\begin{cases}a>0\\ \Delta=(-a)^2-4a<0\end{cases}$，解得 $0<a<4$；}
\step{综上，$a$ 的取值范围为 $[0,4)$.}
\stepm{(2)}{$p$ 为真时，方程 $x^2-2x+a=0$ 无实根，}
\step{故判别式 $\Delta=4-4a<0$，解得 $a>1$，因此 $p$ 为假时，$a\leqslant1$；}
\step{由（1）得 $q$ 为假时 $a\geqslant4$ 或 $a<0$，}
\step{$p$、$q$ 有且只有一个真命题，分两类：}
\step{当 $p$ 假 $q$ 真时，满足 $\begin{cases}a\leqslant1\\ 0\leqslant a<4\end{cases}$，解得 $0\leqslant a\leqslant1$；}
\step{当 $p$ 真 $q$ 假时，满足 $\begin{cases}a>1\\ a\geqslant4\text{或}a<0\end{cases}$，解得 $a\geqslant4$.}
\step{综上，$a$ 的取值范围是 $[0,1]\cup[4,+\infty)$.}
\solend
\fi

\ansspace[8]

\item 某工厂生产 $x$($x>0$) 件产品的总成本 $C$（单位：元）与生产件数 $x$ 之间满足二次函数
$C(x)=\dfrac14x^2+50x+10000$.

（1）若要求生产总成本不超过 30000 元，求满足生产总成本要求的产品件数的取值范围；

（2）试求该工厂生产多少件产品时，可使每件产品的平均成本最小.

\ifanswer
\solbegin
\stepm{(1)}{若要求生产总成本不超过 30000 元，}
\step{则 $C(x)=\dfrac14x^2+50x+10000\leqslant30000$，}
\step{即 $x^2+200x-80000\leqslant0$，解得 $-400\leqslant x\leqslant200$，又 $x>0$，}
\step{所以满足生产总成本要求的产品件数的取值范围为 $(0,200]$，$x\in\N$；}
\stepm{(2)}{生产件数 $x$ 时每件的平均成本为 $\dfrac{\frac14x^2+50x+10000}x=\dfrac14x+50+\dfrac{10000}x$，}
\step{由基本不等式得 $\dfrac14x+50+\dfrac{10000}x\geqslant2\sqrt{\dfrac14x\cdot\dfrac{10000}x}+50=150$，}
\step{当且仅当 $\dfrac14x=\dfrac{10000}x$，即 $x=200$ 时取等号，}
\step{即该工厂生产 200 件产品时，每件产品的平均成本最小.}
\solend
\fi

\ansspace[8]

% 原卷题干“$(n\geqslant2,n\in N)$”后连印两个全角逗号“，，”，照录未改（见文件头“原卷笔误”①）
\item 已知有限集 $A=\{a_1,a_2,\cdots,a_n\}$($n\geqslant2$,$n\in\N$)，，如果 $A$ 中元素
$a_i$($i=1,2,3,\cdots,n$) 满足 $a_1+a_2+a_3+\cdots+a_n=a_1\times a_2\times\cdots\times a_n$，
就称 $A$ 为“和积平衡集”.

% 原卷 (2) 将“写方程”与“反证法证明”两项任务并作一句（“……表示），可尝试用反证法证明：……”），
% 照录未改（见文件头“原卷笔误”②）
（1）判断集合 $\{-1+\sqrt3,-1-\sqrt3\}$ 是否为“和积平衡集”；

（2）若 $a_1$、$a_2$ 是两个不同的正数，且 $\{a_1,a_2\}$ 是“和积平衡集”，写出以 $a_1$、$a_2$ 为根的
一个一元二次方程（系数可用 $a_1$、$a_2$ 表示），可尝试用反证法证明：$a_1$、$a_2$ 至少有一个大于 2；

（3）若 $a_i\geqslant1$，且 $a_i\in\N$，求所有符合条件的“和积平衡集”.

\ifanswer
\solbegin
\stepm{(1)}{对于集合 $\{-1+\sqrt3,-1-\sqrt3\}$，}
\step{因为 $(-1+\sqrt3)+(-1-\sqrt3)=-2$，$(-1+\sqrt3)\times(-1-\sqrt3)=-2$，}
\step{所以满足 $a_1+a_2=a_1\times a_2$，满足题意，}
\step{所以集合 $\{-1+\sqrt3,-1-\sqrt3\}$ 是和积平衡集.}
\stepm{(2)}{以 $a_1$、$a_2$ 为根的一个一元二次方程可为 $x^2-(a_1+a_2)x+a_1\cdot a_2=0$，}
\step{因为 $a_1$、$a_2$ 是两个不同的正数，且 $\{a_1,a_2\}$ 是和积平衡集，}
\step{所以 $a_1+a_2=a_1\cdot a_2$，}
\step{又 $a_1$ 和 $a_2$ 为 $x^2-(a_1+a_2)x+a_1\cdot a_2=0$ 的两个不同的正根，}
\step{所以 $\Delta=(a_1+a_2)^2-4a_1a_2=(a_1+a_2)^2-4(a_1+a_2)>0$，}
\step{又 $a_1+a_2>0$，则 $a_1+a_2>4$，}
\step{若 $a_1$、$a_2$ 都小于等于 2，则 $a_1+a_2<4$，矛盾，}
\step{所以 $a_1$、$a_2$ 至少有一个大于 2；}
\stepm{(3)}{设 $A$ 中的 $a_1<a_2<a_3<\cdots<a_n$，且 $a_i\in\N^*$，}
\step{由 $a_1+a_2+a_3+\cdots+a_n=a_1\times a_2\times\cdots\times a_n$，得 $n=2$ 时明显不满足，}
\step{所以 $n\geqslant3$，故 $a_1\times a_2\times\cdots\times a_n=a_1+a_2+a_3+\cdots+a_n\leqslant na_n$，}
\step{所以 $a_1\times a_2\times\cdots\times a_{n-1}\leqslant n$，}
\step{当 $n=3$ 时，$a_1\times a_2\leqslant3$，$a_i\in\N^*$，只能有 $a_1=1$，$a_2=2$，}
\step{由 $a_1+a_2+a_3=a_1\times a_2\times a_3$，得 $a_3=3$，此时的和积平衡集为 $\{1,2,3\}$，}
\step{当 $n\geqslant4$ 时，$a_1\times a_2\times\cdots\times a_{n-1}\geqslant(n-1)(n-2)\times\cdots\times2\times1\geqslant(n-1)(n-2)$，}
\step{又 $(n-1)(n-2)-n=n^2-4n+2=n(n-4)+2>0$，即 $(n-1)(n-2)>n$，}
\step{即 $a_1\times a_2\times\cdots\times a_{n-1}>n$，与 $a_1\times a_2\times\cdots\times a_{n-1}\leqslant n$ 矛盾，}
\step{所以 $n\geqslant4$ 不满足条件，}
\step{所以符合条件的和积平衡集有且只有 $n=3$，此时的和积平衡集为 $\{1,2,3\}$.}
\solend
\fi

\ansspace*[10]

\end{enumerate}

\end{document}
"
% 紧凑版：xelatex "\def\iscompact{1}% 高一63_吴淞中学_10月月考.tex
%   —— 高一上数学“2026-2027 年吴淞中学高一上 10 月月考”（讲评版）
% 试卷版：xelatex 高一63_吴淞中学_10月月考.tex
% 答案版：xelatex "\def\isanswer{1}\input{高一63_吴淞中学_10月月考.tex}"
% 紧凑版：xelatex "\def\iscompact{1}\input{高一63_吴淞中学_10月月考.tex}"
%
% 原始材料：微信公众号“魔都数学加油站”文章，作者破晓，连载“27学年高一63”；
% 原文标题“27学年高一63——2026-2027年上海市吴淞中学高一上10月月考”；
% 原文链接：https://mp.weixin.qq.com/s/TFQ4TLnt_l66t7t51es4OQ。
% 页面标注发布时间 2026-10-16 21:44（上海）。
% 正文为 10 张 PNG 页（`原图/高一63_吴淞中学_10月月考/01..10.png`）：
% 第 06、10 页 2560×3620，第 01、02、03、04、05、07、08、09 页 4762×6735
% （**同一篇各页分辨率不同**），均按页序读取。10 页全为试卷正文页（卷首标题在第 1 页页顶，
% 卷末解答题第 21 题解析止于第 10 页），无封面页、目录页、二维码或宣传页。
% 粉色斜排“魔都数学加油站”为卷面水印，与公众号页脚一并未录入。
%
% 卷首标题：照卷面印作“2026--2027 年吴淞中学高一上 10 月月考”（公众号标题中的“上海市”
% 卷面标题行没有，本稿照卷面），年份区间按全稿惯例排 en dash。原卷只有一行标题、无副标题。
%
% 篇幅与题号：原卷自第 3 题起编号（第 1、2 题不在本卷内）。填空题 3—12（10 题）、
% 选择题 13—16（4 题）、解答题 17—21（5 题），共 19 题。本稿沿用原节与原题号，
% 只改排为全稿统一的大节样式（原卷小节标题为“一、填空题”“二、选择题”“三、解答题”）。
%
% 讲评版：原卷每题下直接印【答案】或【解析】。**只印【答案】的 2 题**（第 4、7 题），
% 本稿只给 \ans{}；其余 17 题里 16 题都印【解析】（按全稿口径这类题不另提 \ans{}——第 13、15、16 题
% 的结论写在解析末句（“故选 B”等），亦照录在解析内；第 14 题则既无【答案】也无【解析】段。**另有 1 题答案印在题干括号内**——
% 第 14 题（“(  D  )”，且该题无【答案】/【解析】段），本稿按 高一37 第 13 题成例把括号留空、
% 答案另提为 \ans{$D$}。第 16 题题干括号为空但印有【解析】，不算括号内答案。
%
% 副标题：原卷未印副标题。本稿按“知识模块”口径标作“集合与常用逻辑用语、不等式”：
% 19 题中集合与常用逻辑用语 12 题（第 3、4、6、7、9、10、11、12、13、15、18、21 题），
% 不等式 7 题（含等式（韦达定理、恒等式）与不等式性质、基本不等式、二次不等式、恒成立：
% 第 5、8、14、16、17、19、20 题），两类均超过三分之一，故并标。口径同 高一59。
%
% 与原卷的排版差异：
%   ① 原卷每题下直接印【答案】/【解析】，本稿按全稿惯例将答案与解析置于答案版，试卷版保持干净；
%   ② 原卷未印副标题，本稿按“知识模块”口径补出；
%   ③ 原卷填空横线约四个字宽，本稿按全稿惯例统一排 \blankline[4.4em]；
%   ④ 原卷实数集排斜体 $R$（第 9、10、19 题题干与解析）或粗体 $\mathbf{R}$（第 18 题题干），
%      整数集排斜体 $Z$（第 12 题解析），自然数集排斜体 $N$（第 20 题解析、第 21 题题干
%      “$(n\geqslant2,n\in N)$”）或粗体 $\mathbf{N}$（第 21 题 (3) 题干）、斜体 $N^*$
%      （第 21 题 (3) 解析），本稿按全稿惯例统一改用 $\R$、$\Z$、$\N$、$\N^*$
%      （同 高一14 排版差异④：记号归一，粗体／斜体不作区分，$N$ 与 $N^*$ 的区别照原卷保留）；
%      空集记号原卷作 $\varnothing$，照录；
%   ⑤ 原卷在数学式内多处用半角逗号（如 $\{y\mid y=[2x]+[3x],-1\leqslant x\leqslant2\}$、
%      $m=-1,0,1$、$A=\{x\mid-2<x<15\}$,$B=\cdots$），均照录未强改；
%      句末句点按全稿惯例排西文句点；“和积平衡集”自定义名词照录；
%   ⑥ 第 12 题解析里的“法一：”“法二：”标号原卷排**红色**（全卷共 2 处），
%      同 高一53“情况1：”成例，本稿排 \stepm{法一：}{…}、\stepm{法二：}{…}
%      （加粗着色、不改黑色基调），冒号照原卷排全角；
%   ⑦ 第 17—21 题的子问原卷各自另起一行，本稿按全稿惯例用 enumerate 的 (1)(2)(3)；
%      解析中的子问标号一律用 \stepm{(1)}{…}；第 19 题解析中的①②分情形，按全稿惯例作
%      \stepm{①}{…}、\stepm{②}{…}；
%   ⑧ 第 14 题印在题干括号内的答案移入 \ans{}，见上；
%   ⑨ 每道解答题原卷未留作答空白，本稿按全稿惯例加 \ansspace；
%   ⑩ 选择题的作答括号原卷排半角“(　　)”（第 13、15、16 题括号内空白、第 14 题括号内印 D），
%      本稿按全稿惯例排 \mbox{（\quad）}。
%
% 原卷笔误/存疑（均照录未改，正文对应处另有 % 注）：
%   ① 第 21 题题干“$(n\geqslant2$,$n\in N)$”后连印两个全角逗号“，，”（衍字，
%      8 倍放大确认，非水印干扰），本稿照录；
%   ② 第 21 题 (2) 题干将“写出以 $a_1$、$a_2$ 为根的一个一元二次方程”与“用反证法证明：
%      $a_1$、$a_2$ 至少有一个大于 2”两项任务并作一句、之间仅以逗号相接（“……表示），
%      可尝试用反证法证明：……”），读来像抄并两小问，原卷如此，照录未改。
%
% 插图：无（全卷 10 页无任何图形、无“如图”字样）。
%
\documentclass[a4paper,11pt]{article}
\usepackage{common}

\begin{document}

\examtitlest[3]{2026--2027 年吴淞中学高一上 10 月月考}{集合与常用逻辑用语、不等式}

\bigsec{一、填空题}

\begin{enumerate}
\setcounter{enumi}{2}

\item 满足 $A\subseteq\{1,2,3,4\}$ 的集合 $A$ 共有\blankline[4.4em]个.

\ifanswer
\solbegin
\step{由题意，集合 $A$ 是集合 $\{1,2,3,4\}$ 的子集，共有 $2^4=16$ 个.}
\solend
\fi

\item 陈述句“$a>1$ 或 $b>3$”的否定形式为\blankline[4.4em].

\ifanswer
\ans{$a\leqslant1$ 且 $b\leqslant3$}
\fi

\item 已知方程 $3x^2+4x-4=0$ 的两个根为 $x_1$、$x_2$，则 $x_1^2+x_2^2=$\blankline[4.4em].

\ifanswer
\solbegin
\step{对于方程 $3x^2+4x-4=0$，因为 $\Delta=16-4\times3\times(-4)=64>0$，}
\step{所以方程 $3x^2+4x-4=0$ 的两个不相等的实根为 $x_1$、$x_2$，}
\step{由韦达定理可得 $x_1+x_2=-\dfrac43$，$x_1x_2=-\dfrac43$，}
\step{所以 $x_1^2+x_2^2=(x_1+x_2)^2-2x_1x_2=\dfrac{16}9+\dfrac83=\dfrac{40}9$.}
\solend
\fi

\item 设集合 $A=\{y\mid y=x^2+1\}$,$B=\{y\mid y=-x^2+2x+3\}$，则 $A\cap B=$\blankline[4.4em].

\ifanswer
\solbegin
\step{因为 $y=x^2+1\geqslant1$,$y=-x^2+2x+3=-(x-1)^2+4\leqslant4$，}
\step{所以 $A=\{y\mid y\geqslant1\}$,$B=\{y\mid y\leqslant4\}$，得 $A\cap B=\{y\mid1\leqslant y\leqslant4\}$.}
\solend
\fi

\item 若 $x\in\{1,x^2\}$，则 $x=$\blankline[4.4em].

\ifanswer
\ans{$0$}
\fi

\item 已知等式 $2x^2+3x+5=a(2x+1)(x+1)+c$，则 $a+c=$\blankline[4.4em].

\ifanswer
\solbegin
\step{因为 $2x^2+3x+5=a(2x+1)(x+1)+c$ 恒成立，}
\step{所以 $2x^2+3x+5=2ax^2+3ax+a+c$ 恒成立，则 $\begin{cases}2=2a\\ 3=3a\\ 5=a+c\end{cases}$，解得 $a+c=5$.}
\solend
\fi

\item 已知 $a\in\R$，设 $p:x<2$，$q:x<5a-3$，且 $p$ 是 $q$ 的必要非充分条件，则 $a$ 的取值范围是\blankline[4.4em].

\ifanswer
\solbegin
\step{设 $p:x<2$，$q:x<5a-3$，且 $p$ 是 $q$ 的必要非充分条件，}
\step{则 $\{x\mid x<5a-3\}\subset\{x\mid x<2\}$，所以 $5a-3<2$，即 $a<1$.}
\solend
\fi

\item 已知 $a\in\R$，集合 $A=\{x\mid ax^2+3x-1=0\}$，若集合 $A$ 中只有一个元素，则 $a$ 的取值范围是\blankline[4.4em].

\ifanswer
\solbegin
\step{当 $a=0$ 时，则有 $\{x\mid ax^2+3x-1=0\}=\{x\mid3x-1=0\}=\left\{\dfrac13\right\}$，符合题意；}
\step{当 $a\neq0$ 时，由题意得 $\Delta=9+4a=0$，解得 $a=-\dfrac94$. 符合题意；}
\step{综上所述，实数 $a$ 的取值集合为 $\left\{0,-\dfrac94\right\}$.}
\solend
\fi

\item 设集合 $A=\{x\mid x^2+2(m+1)x+m^2-1=0,m\in\R\}$，且 $A\subseteq\{0,-4\}$，则实数 $m$ 的取值范围是\blankline[4.4em].

\ifanswer
\solbegin
\step{由题意得 $A=\varnothing$ 或 $A=\{0\}$ 或 $A=\{-4\}$ 或 $A=\{0,-4\}$，}
\step{当 $A=\varnothing$ 时，方程 $x^2+2(m+1)x+m^2-1=0$ 无实根，}
\step{则 $\Delta=[2(m+1)]^2-4\times1\times(m^2-1)=8m+8<0$，解得 $m<-1$，}
\step{当 $A=\{0\}$ 时，方程 $x^2+2(m+1)x+m^2-1=0$ 有一个实根 $0$，}
\step{则 $\begin{cases}\Delta=8m+8=0\\ m^2-1=0\end{cases}$，解得 $m=-1$，}
\step{当 $A=\{-4\}$ 时，方程 $x^2+2(m+1)x+m^2-1=0$ 有一个实根 $-4$，}
\step{则 $\begin{cases}\Delta=8m+8=0\\ 16-8(m+1)+m^2-1=0\end{cases}$，无解，}
\step{当 $A=\{0,-4\}$ 时，方程 $x^2+2(m+1)x+m^2-1=0$ 有两个实根 $-4$ 和 $0$，}
\step{则 $\begin{cases}\Delta=8m+8>0\\ 16-8(m+1)+m^2-1=0\\ m^2-1=0\end{cases}$，解得 $m=1$，}
\step{综上所述，实数 $m$ 的取值范围是 $\{m\mid m\leqslant-1\text{ 或 }m=1\}$.}
\solend
\fi

\item 对于任意实数 $x$，$[x]$ 表示不超过 $x$ 的最大整数，如 $[1.1]=1$，$[-1.2]=-2$，若
$A=\{y\mid y=[2x]+[3x],-1\leqslant x\leqslant2\}$，则 $A$ 中所有元素的和为\blankline[4.4em].

\ifanswer
\solbegin
% 原卷“法一：”“法二：”标号排红色，本稿按 高一53“情况1：”成例排 \stepm（加粗着色），见文件头
\stepm{法一：}{当 $-1\leqslant x<-\dfrac23$ 时，$[2x]=-2$，$[3x]=-3$，$f(x)=[2x]+[3x]=-5$，}
\step{当 $-\dfrac23\leqslant x<-\dfrac12$ 时，$[2x]=-2$，$[3x]=-2$，$f(x)=[2x]+[3x]=-4$；}
\step{当 $-\dfrac12\leqslant x<-\dfrac13$ 时，$[2x]=-1$，$[3x]=-2$，$f(x)=[2x]+[3x]=-3$；}
\step{当 $-\dfrac13\leqslant x<0$ 时，$[2x]=-1$，$[3x]=-1$，$f(x)=[2x]+[3x]=-2$；}
\step{当 $0\leqslant x<\dfrac13$ 时，$[2x]=0$，$[3x]=0$，$f(x)=[2x]+[3x]=0$；}
\step{当 $\dfrac13\leqslant x<\dfrac12$ 时，$[2x]=0$，$[3x]=1$，$f(x)=[2x]+[3x]=1$；}
\step{当 $\dfrac12\leqslant x<\dfrac23$ 时，$[2x]=1$，$[3x]=1$，$f(x)=[2x]+[3x]=2$；}
\step{当 $\dfrac23\leqslant x<1$ 时，$[2x]=1$，$[3x]=2$，$f(x)=[2x]+[3x]=3$；}
\step{当 $1\leqslant x<\dfrac43$ 时，$[2x]=2$，$[3x]=3$，$f(x)=[2x]+[3x]=5$；}
\step{当 $\dfrac43\leqslant x<\dfrac32$ 时，$[2x]=2$，$[3x]=4$，$f(x)=[2x]+[3x]=6$；}
\step{当 $\dfrac32\leqslant x<\dfrac53$ 时，$[2x]=3$，$[3x]=4$，$f(x)=[2x]+[3x]=7$；}
\step{当 $\dfrac53\leqslant x<2$ 时，$[2x]=3$，$[3x]=5$，$f(x)=[2x]+[3x]=8$；}
\step{当 $x=2$ 时，$[2x]=4$，$[3x]=6$，$f(x)=[2x]+[3x]=10$}
\step{故 $A$ 中所有元素的和为 $-5-4-3-2+0+1+2+3+5+6+7+8+10=28$.}
% 原卷“法二：”标号排红色，同上
\stepm{法二：}{设 $x=m+n$，其中 $m\in\Z$，$n\in[0,1)$，}
\step{若 $0\leqslant n<\dfrac13$，则 $0\leqslant2n<\dfrac23$，$0\leqslant3n<1$，}
\step{得 $[2x]=[2m+2n]=2m$，$[3x]=[3m+3n]=3m$，所以 $y=[2x]+[3x]=5m$；}
\step{若 $\dfrac13\leqslant n<\dfrac12$，则 $\dfrac23\leqslant2n<1$，$1\leqslant3n<\dfrac32$，}
\step{得 $[2x]=[2m+2n]=2m$，$[3x]=[3m+3n]=3m+1$，}
\step{所以 $y=[2x]+[3x]=5m+1$；}
\step{若 $\dfrac12\leqslant n<\dfrac23$，则 $1\leqslant2n<\dfrac43$，$\dfrac32\leqslant3n<2$，}
\step{得 $[2x]=[2m+2n]=2m+1$，$[3x]=[3m+3n]=3m+1$，}
\step{所以 $y=[2x]+[3x]=5m+2$；}
\step{若 $\dfrac23\leqslant n<1$，则 $\dfrac43\leqslant2n<2$，$2\leqslant3n<3$，}
\step{得 $[2x]=[2m+2n]=2m+1$，$[3x]=[3m+3n]=3m+2$，}
\step{所以 $y=[2x]+[3x]=5m+3$；}
\step{综上所述，当 $n\in[0,1)$ 时，所有函数值的和}
\step{为 $5m+5m+1+5m+2+5m+3=20m+6$，}
\step{又因为 $x\in[-1,2]$，此时 $m=-1,0,1$，且当 $x=2$ 时，$y=[2x]+[3x]=10$，}
\step{所有函数值的和为 $(-14)+6+26+10=28$，即 $A$ 中所有元素之和为 $28$.}
\solend
\fi

\end{enumerate}

\bigsec{二、选择题}

\begin{enumerate}
\setcounter{enumi}{12}

\item 下列写法中，正确的是\mbox{（\quad）}

\keepwithprev
A. $0\in\{\varnothing\}$\qquad B. $\varnothing\in\{\varnothing\}$

C. $\varnothing=\{\varnothing\}$\qquad D. $0\in\varnothing$

\ifanswer
\solbegin
\step{$0\notin\{\varnothing\}$，$\varnothing\in\{\varnothing\}$，故 $A$ 错误，$B$ 正确；}
\step{$\varnothing$ 是不含任何元素的集合，所以 $\varnothing\neq\{\varnothing\}$，$0\notin\varnothing$，故 $CD$ 错误.}
\step{故选 $B$.}
\solend
\fi

% 第 14 题答案“D”直接印在题干括号内“(  D  )”，且该题无【答案】/【解析】段；
% 按 高一37 第 13 题的成例括号留空、答案另提为 \ans{}（见文件头）
\item 若 $a<b<0$，则下列不等式中不能成立的是\mbox{（\quad）}

\keepwithprev
A. $\dfrac1a>\dfrac1b$\qquad B. $a-b>a$

C. $|a|>|b|$\qquad D. $a^3>b^3$

\ifanswer
\ans{$D$}
\fi

\item 已知 $x$、$y$、$z$ 为非零实数，代数式 $\dfrac x{|x|}+\dfrac y{|y|}+\dfrac z{|z|}+\dfrac{xyz}{|xyz|}$
的值所组成的集合是 $M$，则下列判断正确的是\mbox{（\quad）}

\keepwithprev
A. $-4\notin M$\qquad B. $2\in M$

C. $4\notin M$\qquad D. $0\in M$

\ifanswer
\solbegin
\step{当 $x$、$y$、$z$ 均为正数时，$xyz>0$，$\dfrac x{|x|}+\dfrac y{|y|}+\dfrac z{|z|}+\dfrac{xyz}{|xyz|}=1+1+1+1=4$，}
\step{当 $x$、$y$、$z$ 有 2 个正数，1 个负数时，不妨设 $x>0$，$y>0$，$z<0$，}
\step{则 $xyz<0$，$\dfrac x{|x|}+\dfrac y{|y|}+\dfrac z{|z|}+\dfrac{xyz}{|xyz|}=1+1-1-1=0$，}
\step{若 $x$、$y$、$z$ 有 1 个正数，2 个负数时，不妨设 $x>0$，$y<0$，$z<0$，}
\step{则 $xyz>0$，$\dfrac x{|x|}+\dfrac y{|y|}+\dfrac z{|z|}+\dfrac{xyz}{|xyz|}=1-1-1+1=0$，}
\step{当 $x$、$y$、$z$ 均为负数时，$xyz<0$，}
\step{$\dfrac x{|x|}+\dfrac y{|y|}+\dfrac z{|z|}+\dfrac{xyz}{|xyz|}=-1-1-1-1=-4$，}
\step{综上，$M=\{-4,0,4\}$，}
\step{所以 $-4\in M$，$A$ 错误；$2\notin M$，$B$ 错误；$4\in M$，$C$ 错误；}
\step{$0\in M$，$D$ 正确. 故选 $D$.}
\solend
\fi

\item 若 $xy>0$，且 $x+y=9$，则 $\dfrac1x+\dfrac4y$ 的最小值是\mbox{（\quad）}

\keepwithprev
A. $9$\qquad B. $\dfrac59$

C. $\dfrac49$\qquad D. $1$

\ifanswer
\solbegin
\step{$xy>0$，且 $x+y=9$，所以 $x>0$，$y>0$，}
\step{所以 $\dfrac1x+\dfrac4y=\dfrac19\left(\dfrac1x+\dfrac4y\right)(x+y)=\dfrac59+\dfrac19\left(\dfrac yx+\dfrac{4x}y\right)\geqslant\dfrac59+\dfrac29\cdot\sqrt{\dfrac yx\cdot\dfrac{4x}y}=\dfrac59+\dfrac49=1$，}
\step{当且仅当 $\dfrac yx=\dfrac{4x}y$，即 $x=3$，$y=6$ 时取等号，所以 $\dfrac1x+\dfrac4y$ 的最小值为 $1$.}
\step{故选 $D$.}
\solend
\fi

\end{enumerate}

\bigsec{三、解答题}

\begin{enumerate}
\setcounter{enumi}{16}

\item 求解下列不等式：

（1）$\dfrac{x-3}{x+1}\leqslant0$；

（2）$|2x-1|<3$.

\ifanswer
\solbegin
\stepm{(1)}{$\dfrac{x-3}{x+1}\leqslant0$，变形得 $(x-3)\cdot(x+1)\leqslant0$ 且 $x\neq-1$，}
\step{解得 $-1<x\leqslant3$，即不等式的解集为 $(-1,3]$.}
\stepm{(2)}{若 $|2x-1|<3$，则 $-3<2x-1<3$，解得 $-1<x<2$，即 $(-1,2)$.}
\solend
\fi

\ansspace[6]

\item 已知集合 $A=\{x\mid-2<x<15\}$,$B=\{x\mid m-6<x<2m-1,m\in\R\}$.

（1）当 $m=2$ 时，求 $A\cap B$；

（2）若 $A\cup B=A$，求实数 $m$ 的取值范围.

\ifanswer
\solbegin
\stepm{(1)}{当 $m=2$ 时，$A=\{x\mid-2<x<15\}$,$B=\{x\mid-4<x<3\}$，}
\step{此时 $A\cap B=\{x\mid-2<x<3\}$.}
\stepm{(2)}{由 $A\cup B=A$ 得 $B\subseteq A$.}
\step{当 $B=\varnothing$,$m-6\geqslant2m-1$，解得 $m\leqslant-5$；}
\step{当 $B\neq\varnothing$，结合 $B\subseteq A$ 得 $\begin{cases}m-6<2m-1\\ 2m-1\leqslant15\\ m-6\geqslant-2\end{cases}$，解得 $4\leqslant m\leqslant8$.}
\step{综上，$m\leqslant-5$ 或 $4\leqslant m\leqslant8$.}
\solend
\fi

\ansspace[6]

\item 已知 $a\in\R$，命题 $p$：关于 $x$ 的方程 $x^2-2x+a=0$ 无实根；命题 $q$：对任意 $x\in\R$，
不等式 $ax^2-ax+1>0$ 恒成立.

（1）若 $q$ 为真命题，求 $a$ 的取值范围；

（2）若 $p$、$q$ 中有且只有一个为真命题，求 $a$ 的取值范围.

\ifanswer
\solbegin
\stepm{(1)}{当 $q$ 为真命题时，}
\stepm{①}{当 $a=0$ 时，不等式化为 $1>0$，对任意 $x\in\R$ 恒成立，满足条件；}
\stepm{②}{当 $a\neq0$ 时，$ax^2-ax+1>0$ 为一元二次不等式，}
\step{若对任意 $x\in\R$ 恒成立，需满足 $\begin{cases}a>0\\ \Delta=(-a)^2-4a<0\end{cases}$，解得 $0<a<4$；}
\step{综上，$a$ 的取值范围为 $[0,4)$.}
\stepm{(2)}{$p$ 为真时，方程 $x^2-2x+a=0$ 无实根，}
\step{故判别式 $\Delta=4-4a<0$，解得 $a>1$，因此 $p$ 为假时，$a\leqslant1$；}
\step{由（1）得 $q$ 为假时 $a\geqslant4$ 或 $a<0$，}
\step{$p$、$q$ 有且只有一个真命题，分两类：}
\step{当 $p$ 假 $q$ 真时，满足 $\begin{cases}a\leqslant1\\ 0\leqslant a<4\end{cases}$，解得 $0\leqslant a\leqslant1$；}
\step{当 $p$ 真 $q$ 假时，满足 $\begin{cases}a>1\\ a\geqslant4\text{或}a<0\end{cases}$，解得 $a\geqslant4$.}
\step{综上，$a$ 的取值范围是 $[0,1]\cup[4,+\infty)$.}
\solend
\fi

\ansspace[8]

\item 某工厂生产 $x$($x>0$) 件产品的总成本 $C$（单位：元）与生产件数 $x$ 之间满足二次函数
$C(x)=\dfrac14x^2+50x+10000$.

（1）若要求生产总成本不超过 30000 元，求满足生产总成本要求的产品件数的取值范围；

（2）试求该工厂生产多少件产品时，可使每件产品的平均成本最小.

\ifanswer
\solbegin
\stepm{(1)}{若要求生产总成本不超过 30000 元，}
\step{则 $C(x)=\dfrac14x^2+50x+10000\leqslant30000$，}
\step{即 $x^2+200x-80000\leqslant0$，解得 $-400\leqslant x\leqslant200$，又 $x>0$，}
\step{所以满足生产总成本要求的产品件数的取值范围为 $(0,200]$，$x\in\N$；}
\stepm{(2)}{生产件数 $x$ 时每件的平均成本为 $\dfrac{\frac14x^2+50x+10000}x=\dfrac14x+50+\dfrac{10000}x$，}
\step{由基本不等式得 $\dfrac14x+50+\dfrac{10000}x\geqslant2\sqrt{\dfrac14x\cdot\dfrac{10000}x}+50=150$，}
\step{当且仅当 $\dfrac14x=\dfrac{10000}x$，即 $x=200$ 时取等号，}
\step{即该工厂生产 200 件产品时，每件产品的平均成本最小.}
\solend
\fi

\ansspace[8]

% 原卷题干“$(n\geqslant2,n\in N)$”后连印两个全角逗号“，，”，照录未改（见文件头“原卷笔误”①）
\item 已知有限集 $A=\{a_1,a_2,\cdots,a_n\}$($n\geqslant2$,$n\in\N$)，，如果 $A$ 中元素
$a_i$($i=1,2,3,\cdots,n$) 满足 $a_1+a_2+a_3+\cdots+a_n=a_1\times a_2\times\cdots\times a_n$，
就称 $A$ 为“和积平衡集”.

% 原卷 (2) 将“写方程”与“反证法证明”两项任务并作一句（“……表示），可尝试用反证法证明：……”），
% 照录未改（见文件头“原卷笔误”②）
（1）判断集合 $\{-1+\sqrt3,-1-\sqrt3\}$ 是否为“和积平衡集”；

（2）若 $a_1$、$a_2$ 是两个不同的正数，且 $\{a_1,a_2\}$ 是“和积平衡集”，写出以 $a_1$、$a_2$ 为根的
一个一元二次方程（系数可用 $a_1$、$a_2$ 表示），可尝试用反证法证明：$a_1$、$a_2$ 至少有一个大于 2；

（3）若 $a_i\geqslant1$，且 $a_i\in\N$，求所有符合条件的“和积平衡集”.

\ifanswer
\solbegin
\stepm{(1)}{对于集合 $\{-1+\sqrt3,-1-\sqrt3\}$，}
\step{因为 $(-1+\sqrt3)+(-1-\sqrt3)=-2$，$(-1+\sqrt3)\times(-1-\sqrt3)=-2$，}
\step{所以满足 $a_1+a_2=a_1\times a_2$，满足题意，}
\step{所以集合 $\{-1+\sqrt3,-1-\sqrt3\}$ 是和积平衡集.}
\stepm{(2)}{以 $a_1$、$a_2$ 为根的一个一元二次方程可为 $x^2-(a_1+a_2)x+a_1\cdot a_2=0$，}
\step{因为 $a_1$、$a_2$ 是两个不同的正数，且 $\{a_1,a_2\}$ 是和积平衡集，}
\step{所以 $a_1+a_2=a_1\cdot a_2$，}
\step{又 $a_1$ 和 $a_2$ 为 $x^2-(a_1+a_2)x+a_1\cdot a_2=0$ 的两个不同的正根，}
\step{所以 $\Delta=(a_1+a_2)^2-4a_1a_2=(a_1+a_2)^2-4(a_1+a_2)>0$，}
\step{又 $a_1+a_2>0$，则 $a_1+a_2>4$，}
\step{若 $a_1$、$a_2$ 都小于等于 2，则 $a_1+a_2<4$，矛盾，}
\step{所以 $a_1$、$a_2$ 至少有一个大于 2；}
\stepm{(3)}{设 $A$ 中的 $a_1<a_2<a_3<\cdots<a_n$，且 $a_i\in\N^*$，}
\step{由 $a_1+a_2+a_3+\cdots+a_n=a_1\times a_2\times\cdots\times a_n$，得 $n=2$ 时明显不满足，}
\step{所以 $n\geqslant3$，故 $a_1\times a_2\times\cdots\times a_n=a_1+a_2+a_3+\cdots+a_n\leqslant na_n$，}
\step{所以 $a_1\times a_2\times\cdots\times a_{n-1}\leqslant n$，}
\step{当 $n=3$ 时，$a_1\times a_2\leqslant3$，$a_i\in\N^*$，只能有 $a_1=1$，$a_2=2$，}
\step{由 $a_1+a_2+a_3=a_1\times a_2\times a_3$，得 $a_3=3$，此时的和积平衡集为 $\{1,2,3\}$，}
\step{当 $n\geqslant4$ 时，$a_1\times a_2\times\cdots\times a_{n-1}\geqslant(n-1)(n-2)\times\cdots\times2\times1\geqslant(n-1)(n-2)$，}
\step{又 $(n-1)(n-2)-n=n^2-4n+2=n(n-4)+2>0$，即 $(n-1)(n-2)>n$，}
\step{即 $a_1\times a_2\times\cdots\times a_{n-1}>n$，与 $a_1\times a_2\times\cdots\times a_{n-1}\leqslant n$ 矛盾，}
\step{所以 $n\geqslant4$ 不满足条件，}
\step{所以符合条件的和积平衡集有且只有 $n=3$，此时的和积平衡集为 $\{1,2,3\}$.}
\solend
\fi

\ansspace*[10]

\end{enumerate}

\end{document}
"
%
% 原始材料：微信公众号“魔都数学加油站”文章，作者破晓，连载“27学年高一63”；
% 原文标题“27学年高一63——2026-2027年上海市吴淞中学高一上10月月考”；
% 原文链接：https://mp.weixin.qq.com/s/TFQ4TLnt_l66t7t51es4OQ。
% 页面标注发布时间 2026-10-16 21:44（上海）。
% 正文为 10 张 PNG 页（`原图/高一63_吴淞中学_10月月考/01..10.png`）：
% 第 06、10 页 2560×3620，第 01、02、03、04、05、07、08、09 页 4762×6735
% （**同一篇各页分辨率不同**），均按页序读取。10 页全为试卷正文页（卷首标题在第 1 页页顶，
% 卷末解答题第 21 题解析止于第 10 页），无封面页、目录页、二维码或宣传页。
% 粉色斜排“魔都数学加油站”为卷面水印，与公众号页脚一并未录入。
%
% 卷首标题：照卷面印作“2026--2027 年吴淞中学高一上 10 月月考”（公众号标题中的“上海市”
% 卷面标题行没有，本稿照卷面），年份区间按全稿惯例排 en dash。原卷只有一行标题、无副标题。
%
% 篇幅与题号：原卷自第 3 题起编号（第 1、2 题不在本卷内）。填空题 3—12（10 题）、
% 选择题 13—16（4 题）、解答题 17—21（5 题），共 19 题。本稿沿用原节与原题号，
% 只改排为全稿统一的大节样式（原卷小节标题为“一、填空题”“二、选择题”“三、解答题”）。
%
% 讲评版：原卷每题下直接印【答案】或【解析】。**只印【答案】的 2 题**（第 4、7 题），
% 本稿只给 \ans{}；其余 17 题里 16 题都印【解析】（按全稿口径这类题不另提 \ans{}——第 13、15、16 题
% 的结论写在解析末句（“故选 B”等），亦照录在解析内；第 14 题则既无【答案】也无【解析】段。**另有 1 题答案印在题干括号内**——
% 第 14 题（“(  D  )”，且该题无【答案】/【解析】段），本稿按 高一37 第 13 题成例把括号留空、
% 答案另提为 \ans{$D$}。第 16 题题干括号为空但印有【解析】，不算括号内答案。
%
% 副标题：原卷未印副标题。本稿按“知识模块”口径标作“集合与常用逻辑用语、不等式”：
% 19 题中集合与常用逻辑用语 12 题（第 3、4、6、7、9、10、11、12、13、15、18、21 题），
% 不等式 7 题（含等式（韦达定理、恒等式）与不等式性质、基本不等式、二次不等式、恒成立：
% 第 5、8、14、16、17、19、20 题），两类均超过三分之一，故并标。口径同 高一59。
%
% 与原卷的排版差异：
%   ① 原卷每题下直接印【答案】/【解析】，本稿按全稿惯例将答案与解析置于答案版，试卷版保持干净；
%   ② 原卷未印副标题，本稿按“知识模块”口径补出；
%   ③ 原卷填空横线约四个字宽，本稿按全稿惯例统一排 \blankline[4.4em]；
%   ④ 原卷实数集排斜体 $R$（第 9、10、19 题题干与解析）或粗体 $\mathbf{R}$（第 18 题题干），
%      整数集排斜体 $Z$（第 12 题解析），自然数集排斜体 $N$（第 20 题解析、第 21 题题干
%      “$(n\geqslant2,n\in N)$”）或粗体 $\mathbf{N}$（第 21 题 (3) 题干）、斜体 $N^*$
%      （第 21 题 (3) 解析），本稿按全稿惯例统一改用 $\R$、$\Z$、$\N$、$\N^*$
%      （同 高一14 排版差异④：记号归一，粗体／斜体不作区分，$N$ 与 $N^*$ 的区别照原卷保留）；
%      空集记号原卷作 $\varnothing$，照录；
%   ⑤ 原卷在数学式内多处用半角逗号（如 $\{y\mid y=[2x]+[3x],-1\leqslant x\leqslant2\}$、
%      $m=-1,0,1$、$A=\{x\mid-2<x<15\}$,$B=\cdots$），均照录未强改；
%      句末句点按全稿惯例排西文句点；“和积平衡集”自定义名词照录；
%   ⑥ 第 12 题解析里的“法一：”“法二：”标号原卷排**红色**（全卷共 2 处），
%      同 高一53“情况1：”成例，本稿排 \stepm{法一：}{…}、\stepm{法二：}{…}
%      （加粗着色、不改黑色基调），冒号照原卷排全角；
%   ⑦ 第 17—21 题的子问原卷各自另起一行，本稿按全稿惯例用 enumerate 的 (1)(2)(3)；
%      解析中的子问标号一律用 \stepm{(1)}{…}；第 19 题解析中的①②分情形，按全稿惯例作
%      \stepm{①}{…}、\stepm{②}{…}；
%   ⑧ 第 14 题印在题干括号内的答案移入 \ans{}，见上；
%   ⑨ 每道解答题原卷未留作答空白，本稿按全稿惯例加 \ansspace；
%   ⑩ 选择题的作答括号原卷排半角“(　　)”（第 13、15、16 题括号内空白、第 14 题括号内印 D），
%      本稿按全稿惯例排 \mbox{（\quad）}。
%
% 原卷笔误/存疑（均照录未改，正文对应处另有 % 注）：
%   ① 第 21 题题干“$(n\geqslant2$,$n\in N)$”后连印两个全角逗号“，，”（衍字，
%      8 倍放大确认，非水印干扰），本稿照录；
%   ② 第 21 题 (2) 题干将“写出以 $a_1$、$a_2$ 为根的一个一元二次方程”与“用反证法证明：
%      $a_1$、$a_2$ 至少有一个大于 2”两项任务并作一句、之间仅以逗号相接（“……表示），
%      可尝试用反证法证明：……”），读来像抄并两小问，原卷如此，照录未改。
%
% 插图：无（全卷 10 页无任何图形、无“如图”字样）。
%
\documentclass[a4paper,11pt]{article}
\usepackage{common}

\begin{document}

\examtitlest[3]{2026--2027 年吴淞中学高一上 10 月月考}{集合与常用逻辑用语、不等式}

\bigsec{一、填空题}

\begin{enumerate}
\setcounter{enumi}{2}

\item 满足 $A\subseteq\{1,2,3,4\}$ 的集合 $A$ 共有\blankline[4.4em]个.

\ifanswer
\solbegin
\step{由题意，集合 $A$ 是集合 $\{1,2,3,4\}$ 的子集，共有 $2^4=16$ 个.}
\solend
\fi

\item 陈述句“$a>1$ 或 $b>3$”的否定形式为\blankline[4.4em].

\ifanswer
\ans{$a\leqslant1$ 且 $b\leqslant3$}
\fi

\item 已知方程 $3x^2+4x-4=0$ 的两个根为 $x_1$、$x_2$，则 $x_1^2+x_2^2=$\blankline[4.4em].

\ifanswer
\solbegin
\step{对于方程 $3x^2+4x-4=0$，因为 $\Delta=16-4\times3\times(-4)=64>0$，}
\step{所以方程 $3x^2+4x-4=0$ 的两个不相等的实根为 $x_1$、$x_2$，}
\step{由韦达定理可得 $x_1+x_2=-\dfrac43$，$x_1x_2=-\dfrac43$，}
\step{所以 $x_1^2+x_2^2=(x_1+x_2)^2-2x_1x_2=\dfrac{16}9+\dfrac83=\dfrac{40}9$.}
\solend
\fi

\item 设集合 $A=\{y\mid y=x^2+1\}$,$B=\{y\mid y=-x^2+2x+3\}$，则 $A\cap B=$\blankline[4.4em].

\ifanswer
\solbegin
\step{因为 $y=x^2+1\geqslant1$,$y=-x^2+2x+3=-(x-1)^2+4\leqslant4$，}
\step{所以 $A=\{y\mid y\geqslant1\}$,$B=\{y\mid y\leqslant4\}$，得 $A\cap B=\{y\mid1\leqslant y\leqslant4\}$.}
\solend
\fi

\item 若 $x\in\{1,x^2\}$，则 $x=$\blankline[4.4em].

\ifanswer
\ans{$0$}
\fi

\item 已知等式 $2x^2+3x+5=a(2x+1)(x+1)+c$，则 $a+c=$\blankline[4.4em].

\ifanswer
\solbegin
\step{因为 $2x^2+3x+5=a(2x+1)(x+1)+c$ 恒成立，}
\step{所以 $2x^2+3x+5=2ax^2+3ax+a+c$ 恒成立，则 $\begin{cases}2=2a\\ 3=3a\\ 5=a+c\end{cases}$，解得 $a+c=5$.}
\solend
\fi

\item 已知 $a\in\R$，设 $p:x<2$，$q:x<5a-3$，且 $p$ 是 $q$ 的必要非充分条件，则 $a$ 的取值范围是\blankline[4.4em].

\ifanswer
\solbegin
\step{设 $p:x<2$，$q:x<5a-3$，且 $p$ 是 $q$ 的必要非充分条件，}
\step{则 $\{x\mid x<5a-3\}\subset\{x\mid x<2\}$，所以 $5a-3<2$，即 $a<1$.}
\solend
\fi

\item 已知 $a\in\R$，集合 $A=\{x\mid ax^2+3x-1=0\}$，若集合 $A$ 中只有一个元素，则 $a$ 的取值范围是\blankline[4.4em].

\ifanswer
\solbegin
\step{当 $a=0$ 时，则有 $\{x\mid ax^2+3x-1=0\}=\{x\mid3x-1=0\}=\left\{\dfrac13\right\}$，符合题意；}
\step{当 $a\neq0$ 时，由题意得 $\Delta=9+4a=0$，解得 $a=-\dfrac94$. 符合题意；}
\step{综上所述，实数 $a$ 的取值集合为 $\left\{0,-\dfrac94\right\}$.}
\solend
\fi

\item 设集合 $A=\{x\mid x^2+2(m+1)x+m^2-1=0,m\in\R\}$，且 $A\subseteq\{0,-4\}$，则实数 $m$ 的取值范围是\blankline[4.4em].

\ifanswer
\solbegin
\step{由题意得 $A=\varnothing$ 或 $A=\{0\}$ 或 $A=\{-4\}$ 或 $A=\{0,-4\}$，}
\step{当 $A=\varnothing$ 时，方程 $x^2+2(m+1)x+m^2-1=0$ 无实根，}
\step{则 $\Delta=[2(m+1)]^2-4\times1\times(m^2-1)=8m+8<0$，解得 $m<-1$，}
\step{当 $A=\{0\}$ 时，方程 $x^2+2(m+1)x+m^2-1=0$ 有一个实根 $0$，}
\step{则 $\begin{cases}\Delta=8m+8=0\\ m^2-1=0\end{cases}$，解得 $m=-1$，}
\step{当 $A=\{-4\}$ 时，方程 $x^2+2(m+1)x+m^2-1=0$ 有一个实根 $-4$，}
\step{则 $\begin{cases}\Delta=8m+8=0\\ 16-8(m+1)+m^2-1=0\end{cases}$，无解，}
\step{当 $A=\{0,-4\}$ 时，方程 $x^2+2(m+1)x+m^2-1=0$ 有两个实根 $-4$ 和 $0$，}
\step{则 $\begin{cases}\Delta=8m+8>0\\ 16-8(m+1)+m^2-1=0\\ m^2-1=0\end{cases}$，解得 $m=1$，}
\step{综上所述，实数 $m$ 的取值范围是 $\{m\mid m\leqslant-1\text{ 或 }m=1\}$.}
\solend
\fi

\item 对于任意实数 $x$，$[x]$ 表示不超过 $x$ 的最大整数，如 $[1.1]=1$，$[-1.2]=-2$，若
$A=\{y\mid y=[2x]+[3x],-1\leqslant x\leqslant2\}$，则 $A$ 中所有元素的和为\blankline[4.4em].

\ifanswer
\solbegin
% 原卷“法一：”“法二：”标号排红色，本稿按 高一53“情况1：”成例排 \stepm（加粗着色），见文件头
\stepm{法一：}{当 $-1\leqslant x<-\dfrac23$ 时，$[2x]=-2$，$[3x]=-3$，$f(x)=[2x]+[3x]=-5$，}
\step{当 $-\dfrac23\leqslant x<-\dfrac12$ 时，$[2x]=-2$，$[3x]=-2$，$f(x)=[2x]+[3x]=-4$；}
\step{当 $-\dfrac12\leqslant x<-\dfrac13$ 时，$[2x]=-1$，$[3x]=-2$，$f(x)=[2x]+[3x]=-3$；}
\step{当 $-\dfrac13\leqslant x<0$ 时，$[2x]=-1$，$[3x]=-1$，$f(x)=[2x]+[3x]=-2$；}
\step{当 $0\leqslant x<\dfrac13$ 时，$[2x]=0$，$[3x]=0$，$f(x)=[2x]+[3x]=0$；}
\step{当 $\dfrac13\leqslant x<\dfrac12$ 时，$[2x]=0$，$[3x]=1$，$f(x)=[2x]+[3x]=1$；}
\step{当 $\dfrac12\leqslant x<\dfrac23$ 时，$[2x]=1$，$[3x]=1$，$f(x)=[2x]+[3x]=2$；}
\step{当 $\dfrac23\leqslant x<1$ 时，$[2x]=1$，$[3x]=2$，$f(x)=[2x]+[3x]=3$；}
\step{当 $1\leqslant x<\dfrac43$ 时，$[2x]=2$，$[3x]=3$，$f(x)=[2x]+[3x]=5$；}
\step{当 $\dfrac43\leqslant x<\dfrac32$ 时，$[2x]=2$，$[3x]=4$，$f(x)=[2x]+[3x]=6$；}
\step{当 $\dfrac32\leqslant x<\dfrac53$ 时，$[2x]=3$，$[3x]=4$，$f(x)=[2x]+[3x]=7$；}
\step{当 $\dfrac53\leqslant x<2$ 时，$[2x]=3$，$[3x]=5$，$f(x)=[2x]+[3x]=8$；}
\step{当 $x=2$ 时，$[2x]=4$，$[3x]=6$，$f(x)=[2x]+[3x]=10$}
\step{故 $A$ 中所有元素的和为 $-5-4-3-2+0+1+2+3+5+6+7+8+10=28$.}
% 原卷“法二：”标号排红色，同上
\stepm{法二：}{设 $x=m+n$，其中 $m\in\Z$，$n\in[0,1)$，}
\step{若 $0\leqslant n<\dfrac13$，则 $0\leqslant2n<\dfrac23$，$0\leqslant3n<1$，}
\step{得 $[2x]=[2m+2n]=2m$，$[3x]=[3m+3n]=3m$，所以 $y=[2x]+[3x]=5m$；}
\step{若 $\dfrac13\leqslant n<\dfrac12$，则 $\dfrac23\leqslant2n<1$，$1\leqslant3n<\dfrac32$，}
\step{得 $[2x]=[2m+2n]=2m$，$[3x]=[3m+3n]=3m+1$，}
\step{所以 $y=[2x]+[3x]=5m+1$；}
\step{若 $\dfrac12\leqslant n<\dfrac23$，则 $1\leqslant2n<\dfrac43$，$\dfrac32\leqslant3n<2$，}
\step{得 $[2x]=[2m+2n]=2m+1$，$[3x]=[3m+3n]=3m+1$，}
\step{所以 $y=[2x]+[3x]=5m+2$；}
\step{若 $\dfrac23\leqslant n<1$，则 $\dfrac43\leqslant2n<2$，$2\leqslant3n<3$，}
\step{得 $[2x]=[2m+2n]=2m+1$，$[3x]=[3m+3n]=3m+2$，}
\step{所以 $y=[2x]+[3x]=5m+3$；}
\step{综上所述，当 $n\in[0,1)$ 时，所有函数值的和}
\step{为 $5m+5m+1+5m+2+5m+3=20m+6$，}
\step{又因为 $x\in[-1,2]$，此时 $m=-1,0,1$，且当 $x=2$ 时，$y=[2x]+[3x]=10$，}
\step{所有函数值的和为 $(-14)+6+26+10=28$，即 $A$ 中所有元素之和为 $28$.}
\solend
\fi

\end{enumerate}

\bigsec{二、选择题}

\begin{enumerate}
\setcounter{enumi}{12}

\item 下列写法中，正确的是\mbox{（\quad）}

\keepwithprev
A. $0\in\{\varnothing\}$\qquad B. $\varnothing\in\{\varnothing\}$

C. $\varnothing=\{\varnothing\}$\qquad D. $0\in\varnothing$

\ifanswer
\solbegin
\step{$0\notin\{\varnothing\}$，$\varnothing\in\{\varnothing\}$，故 $A$ 错误，$B$ 正确；}
\step{$\varnothing$ 是不含任何元素的集合，所以 $\varnothing\neq\{\varnothing\}$，$0\notin\varnothing$，故 $CD$ 错误.}
\step{故选 $B$.}
\solend
\fi

% 第 14 题答案“D”直接印在题干括号内“(  D  )”，且该题无【答案】/【解析】段；
% 按 高一37 第 13 题的成例括号留空、答案另提为 \ans{}（见文件头）
\item 若 $a<b<0$，则下列不等式中不能成立的是\mbox{（\quad）}

\keepwithprev
A. $\dfrac1a>\dfrac1b$\qquad B. $a-b>a$

C. $|a|>|b|$\qquad D. $a^3>b^3$

\ifanswer
\ans{$D$}
\fi

\item 已知 $x$、$y$、$z$ 为非零实数，代数式 $\dfrac x{|x|}+\dfrac y{|y|}+\dfrac z{|z|}+\dfrac{xyz}{|xyz|}$
的值所组成的集合是 $M$，则下列判断正确的是\mbox{（\quad）}

\keepwithprev
A. $-4\notin M$\qquad B. $2\in M$

C. $4\notin M$\qquad D. $0\in M$

\ifanswer
\solbegin
\step{当 $x$、$y$、$z$ 均为正数时，$xyz>0$，$\dfrac x{|x|}+\dfrac y{|y|}+\dfrac z{|z|}+\dfrac{xyz}{|xyz|}=1+1+1+1=4$，}
\step{当 $x$、$y$、$z$ 有 2 个正数，1 个负数时，不妨设 $x>0$，$y>0$，$z<0$，}
\step{则 $xyz<0$，$\dfrac x{|x|}+\dfrac y{|y|}+\dfrac z{|z|}+\dfrac{xyz}{|xyz|}=1+1-1-1=0$，}
\step{若 $x$、$y$、$z$ 有 1 个正数，2 个负数时，不妨设 $x>0$，$y<0$，$z<0$，}
\step{则 $xyz>0$，$\dfrac x{|x|}+\dfrac y{|y|}+\dfrac z{|z|}+\dfrac{xyz}{|xyz|}=1-1-1+1=0$，}
\step{当 $x$、$y$、$z$ 均为负数时，$xyz<0$，}
\step{$\dfrac x{|x|}+\dfrac y{|y|}+\dfrac z{|z|}+\dfrac{xyz}{|xyz|}=-1-1-1-1=-4$，}
\step{综上，$M=\{-4,0,4\}$，}
\step{所以 $-4\in M$，$A$ 错误；$2\notin M$，$B$ 错误；$4\in M$，$C$ 错误；}
\step{$0\in M$，$D$ 正确. 故选 $D$.}
\solend
\fi

\item 若 $xy>0$，且 $x+y=9$，则 $\dfrac1x+\dfrac4y$ 的最小值是\mbox{（\quad）}

\keepwithprev
A. $9$\qquad B. $\dfrac59$

C. $\dfrac49$\qquad D. $1$

\ifanswer
\solbegin
\step{$xy>0$，且 $x+y=9$，所以 $x>0$，$y>0$，}
\step{所以 $\dfrac1x+\dfrac4y=\dfrac19\left(\dfrac1x+\dfrac4y\right)(x+y)=\dfrac59+\dfrac19\left(\dfrac yx+\dfrac{4x}y\right)\geqslant\dfrac59+\dfrac29\cdot\sqrt{\dfrac yx\cdot\dfrac{4x}y}=\dfrac59+\dfrac49=1$，}
\step{当且仅当 $\dfrac yx=\dfrac{4x}y$，即 $x=3$，$y=6$ 时取等号，所以 $\dfrac1x+\dfrac4y$ 的最小值为 $1$.}
\step{故选 $D$.}
\solend
\fi

\end{enumerate}

\bigsec{三、解答题}

\begin{enumerate}
\setcounter{enumi}{16}

\item 求解下列不等式：

（1）$\dfrac{x-3}{x+1}\leqslant0$；

（2）$|2x-1|<3$.

\ifanswer
\solbegin
\stepm{(1)}{$\dfrac{x-3}{x+1}\leqslant0$，变形得 $(x-3)\cdot(x+1)\leqslant0$ 且 $x\neq-1$，}
\step{解得 $-1<x\leqslant3$，即不等式的解集为 $(-1,3]$.}
\stepm{(2)}{若 $|2x-1|<3$，则 $-3<2x-1<3$，解得 $-1<x<2$，即 $(-1,2)$.}
\solend
\fi

\ansspace[6]

\item 已知集合 $A=\{x\mid-2<x<15\}$,$B=\{x\mid m-6<x<2m-1,m\in\R\}$.

（1）当 $m=2$ 时，求 $A\cap B$；

（2）若 $A\cup B=A$，求实数 $m$ 的取值范围.

\ifanswer
\solbegin
\stepm{(1)}{当 $m=2$ 时，$A=\{x\mid-2<x<15\}$,$B=\{x\mid-4<x<3\}$，}
\step{此时 $A\cap B=\{x\mid-2<x<3\}$.}
\stepm{(2)}{由 $A\cup B=A$ 得 $B\subseteq A$.}
\step{当 $B=\varnothing$,$m-6\geqslant2m-1$，解得 $m\leqslant-5$；}
\step{当 $B\neq\varnothing$，结合 $B\subseteq A$ 得 $\begin{cases}m-6<2m-1\\ 2m-1\leqslant15\\ m-6\geqslant-2\end{cases}$，解得 $4\leqslant m\leqslant8$.}
\step{综上，$m\leqslant-5$ 或 $4\leqslant m\leqslant8$.}
\solend
\fi

\ansspace[6]

\item 已知 $a\in\R$，命题 $p$：关于 $x$ 的方程 $x^2-2x+a=0$ 无实根；命题 $q$：对任意 $x\in\R$，
不等式 $ax^2-ax+1>0$ 恒成立.

（1）若 $q$ 为真命题，求 $a$ 的取值范围；

（2）若 $p$、$q$ 中有且只有一个为真命题，求 $a$ 的取值范围.

\ifanswer
\solbegin
\stepm{(1)}{当 $q$ 为真命题时，}
\stepm{①}{当 $a=0$ 时，不等式化为 $1>0$，对任意 $x\in\R$ 恒成立，满足条件；}
\stepm{②}{当 $a\neq0$ 时，$ax^2-ax+1>0$ 为一元二次不等式，}
\step{若对任意 $x\in\R$ 恒成立，需满足 $\begin{cases}a>0\\ \Delta=(-a)^2-4a<0\end{cases}$，解得 $0<a<4$；}
\step{综上，$a$ 的取值范围为 $[0,4)$.}
\stepm{(2)}{$p$ 为真时，方程 $x^2-2x+a=0$ 无实根，}
\step{故判别式 $\Delta=4-4a<0$，解得 $a>1$，因此 $p$ 为假时，$a\leqslant1$；}
\step{由（1）得 $q$ 为假时 $a\geqslant4$ 或 $a<0$，}
\step{$p$、$q$ 有且只有一个真命题，分两类：}
\step{当 $p$ 假 $q$ 真时，满足 $\begin{cases}a\leqslant1\\ 0\leqslant a<4\end{cases}$，解得 $0\leqslant a\leqslant1$；}
\step{当 $p$ 真 $q$ 假时，满足 $\begin{cases}a>1\\ a\geqslant4\text{或}a<0\end{cases}$，解得 $a\geqslant4$.}
\step{综上，$a$ 的取值范围是 $[0,1]\cup[4,+\infty)$.}
\solend
\fi

\ansspace[8]

\item 某工厂生产 $x$($x>0$) 件产品的总成本 $C$（单位：元）与生产件数 $x$ 之间满足二次函数
$C(x)=\dfrac14x^2+50x+10000$.

（1）若要求生产总成本不超过 30000 元，求满足生产总成本要求的产品件数的取值范围；

（2）试求该工厂生产多少件产品时，可使每件产品的平均成本最小.

\ifanswer
\solbegin
\stepm{(1)}{若要求生产总成本不超过 30000 元，}
\step{则 $C(x)=\dfrac14x^2+50x+10000\leqslant30000$，}
\step{即 $x^2+200x-80000\leqslant0$，解得 $-400\leqslant x\leqslant200$，又 $x>0$，}
\step{所以满足生产总成本要求的产品件数的取值范围为 $(0,200]$，$x\in\N$；}
\stepm{(2)}{生产件数 $x$ 时每件的平均成本为 $\dfrac{\frac14x^2+50x+10000}x=\dfrac14x+50+\dfrac{10000}x$，}
\step{由基本不等式得 $\dfrac14x+50+\dfrac{10000}x\geqslant2\sqrt{\dfrac14x\cdot\dfrac{10000}x}+50=150$，}
\step{当且仅当 $\dfrac14x=\dfrac{10000}x$，即 $x=200$ 时取等号，}
\step{即该工厂生产 200 件产品时，每件产品的平均成本最小.}
\solend
\fi

\ansspace[8]

% 原卷题干“$(n\geqslant2,n\in N)$”后连印两个全角逗号“，，”，照录未改（见文件头“原卷笔误”①）
\item 已知有限集 $A=\{a_1,a_2,\cdots,a_n\}$($n\geqslant2$,$n\in\N$)，，如果 $A$ 中元素
$a_i$($i=1,2,3,\cdots,n$) 满足 $a_1+a_2+a_3+\cdots+a_n=a_1\times a_2\times\cdots\times a_n$，
就称 $A$ 为“和积平衡集”.

% 原卷 (2) 将“写方程”与“反证法证明”两项任务并作一句（“……表示），可尝试用反证法证明：……”），
% 照录未改（见文件头“原卷笔误”②）
（1）判断集合 $\{-1+\sqrt3,-1-\sqrt3\}$ 是否为“和积平衡集”；

（2）若 $a_1$、$a_2$ 是两个不同的正数，且 $\{a_1,a_2\}$ 是“和积平衡集”，写出以 $a_1$、$a_2$ 为根的
一个一元二次方程（系数可用 $a_1$、$a_2$ 表示），可尝试用反证法证明：$a_1$、$a_2$ 至少有一个大于 2；

（3）若 $a_i\geqslant1$，且 $a_i\in\N$，求所有符合条件的“和积平衡集”.

\ifanswer
\solbegin
\stepm{(1)}{对于集合 $\{-1+\sqrt3,-1-\sqrt3\}$，}
\step{因为 $(-1+\sqrt3)+(-1-\sqrt3)=-2$，$(-1+\sqrt3)\times(-1-\sqrt3)=-2$，}
\step{所以满足 $a_1+a_2=a_1\times a_2$，满足题意，}
\step{所以集合 $\{-1+\sqrt3,-1-\sqrt3\}$ 是和积平衡集.}
\stepm{(2)}{以 $a_1$、$a_2$ 为根的一个一元二次方程可为 $x^2-(a_1+a_2)x+a_1\cdot a_2=0$，}
\step{因为 $a_1$、$a_2$ 是两个不同的正数，且 $\{a_1,a_2\}$ 是和积平衡集，}
\step{所以 $a_1+a_2=a_1\cdot a_2$，}
\step{又 $a_1$ 和 $a_2$ 为 $x^2-(a_1+a_2)x+a_1\cdot a_2=0$ 的两个不同的正根，}
\step{所以 $\Delta=(a_1+a_2)^2-4a_1a_2=(a_1+a_2)^2-4(a_1+a_2)>0$，}
\step{又 $a_1+a_2>0$，则 $a_1+a_2>4$，}
\step{若 $a_1$、$a_2$ 都小于等于 2，则 $a_1+a_2<4$，矛盾，}
\step{所以 $a_1$、$a_2$ 至少有一个大于 2；}
\stepm{(3)}{设 $A$ 中的 $a_1<a_2<a_3<\cdots<a_n$，且 $a_i\in\N^*$，}
\step{由 $a_1+a_2+a_3+\cdots+a_n=a_1\times a_2\times\cdots\times a_n$，得 $n=2$ 时明显不满足，}
\step{所以 $n\geqslant3$，故 $a_1\times a_2\times\cdots\times a_n=a_1+a_2+a_3+\cdots+a_n\leqslant na_n$，}
\step{所以 $a_1\times a_2\times\cdots\times a_{n-1}\leqslant n$，}
\step{当 $n=3$ 时，$a_1\times a_2\leqslant3$，$a_i\in\N^*$，只能有 $a_1=1$，$a_2=2$，}
\step{由 $a_1+a_2+a_3=a_1\times a_2\times a_3$，得 $a_3=3$，此时的和积平衡集为 $\{1,2,3\}$，}
\step{当 $n\geqslant4$ 时，$a_1\times a_2\times\cdots\times a_{n-1}\geqslant(n-1)(n-2)\times\cdots\times2\times1\geqslant(n-1)(n-2)$，}
\step{又 $(n-1)(n-2)-n=n^2-4n+2=n(n-4)+2>0$，即 $(n-1)(n-2)>n$，}
\step{即 $a_1\times a_2\times\cdots\times a_{n-1}>n$，与 $a_1\times a_2\times\cdots\times a_{n-1}\leqslant n$ 矛盾，}
\step{所以 $n\geqslant4$ 不满足条件，}
\step{所以符合条件的和积平衡集有且只有 $n=3$，此时的和积平衡集为 $\{1,2,3\}$.}
\solend
\fi

\ansspace*[10]

\end{enumerate}

\end{document}
"
% 紧凑版：xelatex "\def\iscompact{1}% 高一63_吴淞中学_10月月考.tex
%   —— 高一上数学“2026-2027 年吴淞中学高一上 10 月月考”（讲评版）
% 试卷版：xelatex 高一63_吴淞中学_10月月考.tex
% 答案版：xelatex "\def\isanswer{1}% 高一63_吴淞中学_10月月考.tex
%   —— 高一上数学“2026-2027 年吴淞中学高一上 10 月月考”（讲评版）
% 试卷版：xelatex 高一63_吴淞中学_10月月考.tex
% 答案版：xelatex "\def\isanswer{1}\input{高一63_吴淞中学_10月月考.tex}"
% 紧凑版：xelatex "\def\iscompact{1}\input{高一63_吴淞中学_10月月考.tex}"
%
% 原始材料：微信公众号“魔都数学加油站”文章，作者破晓，连载“27学年高一63”；
% 原文标题“27学年高一63——2026-2027年上海市吴淞中学高一上10月月考”；
% 原文链接：https://mp.weixin.qq.com/s/TFQ4TLnt_l66t7t51es4OQ。
% 页面标注发布时间 2026-10-16 21:44（上海）。
% 正文为 10 张 PNG 页（`原图/高一63_吴淞中学_10月月考/01..10.png`）：
% 第 06、10 页 2560×3620，第 01、02、03、04、05、07、08、09 页 4762×6735
% （**同一篇各页分辨率不同**），均按页序读取。10 页全为试卷正文页（卷首标题在第 1 页页顶，
% 卷末解答题第 21 题解析止于第 10 页），无封面页、目录页、二维码或宣传页。
% 粉色斜排“魔都数学加油站”为卷面水印，与公众号页脚一并未录入。
%
% 卷首标题：照卷面印作“2026--2027 年吴淞中学高一上 10 月月考”（公众号标题中的“上海市”
% 卷面标题行没有，本稿照卷面），年份区间按全稿惯例排 en dash。原卷只有一行标题、无副标题。
%
% 篇幅与题号：原卷自第 3 题起编号（第 1、2 题不在本卷内）。填空题 3—12（10 题）、
% 选择题 13—16（4 题）、解答题 17—21（5 题），共 19 题。本稿沿用原节与原题号，
% 只改排为全稿统一的大节样式（原卷小节标题为“一、填空题”“二、选择题”“三、解答题”）。
%
% 讲评版：原卷每题下直接印【答案】或【解析】。**只印【答案】的 2 题**（第 4、7 题），
% 本稿只给 \ans{}；其余 17 题里 16 题都印【解析】（按全稿口径这类题不另提 \ans{}——第 13、15、16 题
% 的结论写在解析末句（“故选 B”等），亦照录在解析内；第 14 题则既无【答案】也无【解析】段。**另有 1 题答案印在题干括号内**——
% 第 14 题（“(  D  )”，且该题无【答案】/【解析】段），本稿按 高一37 第 13 题成例把括号留空、
% 答案另提为 \ans{$D$}。第 16 题题干括号为空但印有【解析】，不算括号内答案。
%
% 副标题：原卷未印副标题。本稿按“知识模块”口径标作“集合与常用逻辑用语、不等式”：
% 19 题中集合与常用逻辑用语 12 题（第 3、4、6、7、9、10、11、12、13、15、18、21 题），
% 不等式 7 题（含等式（韦达定理、恒等式）与不等式性质、基本不等式、二次不等式、恒成立：
% 第 5、8、14、16、17、19、20 题），两类均超过三分之一，故并标。口径同 高一59。
%
% 与原卷的排版差异：
%   ① 原卷每题下直接印【答案】/【解析】，本稿按全稿惯例将答案与解析置于答案版，试卷版保持干净；
%   ② 原卷未印副标题，本稿按“知识模块”口径补出；
%   ③ 原卷填空横线约四个字宽，本稿按全稿惯例统一排 \blankline[4.4em]；
%   ④ 原卷实数集排斜体 $R$（第 9、10、19 题题干与解析）或粗体 $\mathbf{R}$（第 18 题题干），
%      整数集排斜体 $Z$（第 12 题解析），自然数集排斜体 $N$（第 20 题解析、第 21 题题干
%      “$(n\geqslant2,n\in N)$”）或粗体 $\mathbf{N}$（第 21 题 (3) 题干）、斜体 $N^*$
%      （第 21 题 (3) 解析），本稿按全稿惯例统一改用 $\R$、$\Z$、$\N$、$\N^*$
%      （同 高一14 排版差异④：记号归一，粗体／斜体不作区分，$N$ 与 $N^*$ 的区别照原卷保留）；
%      空集记号原卷作 $\varnothing$，照录；
%   ⑤ 原卷在数学式内多处用半角逗号（如 $\{y\mid y=[2x]+[3x],-1\leqslant x\leqslant2\}$、
%      $m=-1,0,1$、$A=\{x\mid-2<x<15\}$,$B=\cdots$），均照录未强改；
%      句末句点按全稿惯例排西文句点；“和积平衡集”自定义名词照录；
%   ⑥ 第 12 题解析里的“法一：”“法二：”标号原卷排**红色**（全卷共 2 处），
%      同 高一53“情况1：”成例，本稿排 \stepm{法一：}{…}、\stepm{法二：}{…}
%      （加粗着色、不改黑色基调），冒号照原卷排全角；
%   ⑦ 第 17—21 题的子问原卷各自另起一行，本稿按全稿惯例用 enumerate 的 (1)(2)(3)；
%      解析中的子问标号一律用 \stepm{(1)}{…}；第 19 题解析中的①②分情形，按全稿惯例作
%      \stepm{①}{…}、\stepm{②}{…}；
%   ⑧ 第 14 题印在题干括号内的答案移入 \ans{}，见上；
%   ⑨ 每道解答题原卷未留作答空白，本稿按全稿惯例加 \ansspace；
%   ⑩ 选择题的作答括号原卷排半角“(　　)”（第 13、15、16 题括号内空白、第 14 题括号内印 D），
%      本稿按全稿惯例排 \mbox{（\quad）}。
%
% 原卷笔误/存疑（均照录未改，正文对应处另有 % 注）：
%   ① 第 21 题题干“$(n\geqslant2$,$n\in N)$”后连印两个全角逗号“，，”（衍字，
%      8 倍放大确认，非水印干扰），本稿照录；
%   ② 第 21 题 (2) 题干将“写出以 $a_1$、$a_2$ 为根的一个一元二次方程”与“用反证法证明：
%      $a_1$、$a_2$ 至少有一个大于 2”两项任务并作一句、之间仅以逗号相接（“……表示），
%      可尝试用反证法证明：……”），读来像抄并两小问，原卷如此，照录未改。
%
% 插图：无（全卷 10 页无任何图形、无“如图”字样）。
%
\documentclass[a4paper,11pt]{article}
\usepackage{common}

\begin{document}

\examtitlest[3]{2026--2027 年吴淞中学高一上 10 月月考}{集合与常用逻辑用语、不等式}

\bigsec{一、填空题}

\begin{enumerate}
\setcounter{enumi}{2}

\item 满足 $A\subseteq\{1,2,3,4\}$ 的集合 $A$ 共有\blankline[4.4em]个.

\ifanswer
\solbegin
\step{由题意，集合 $A$ 是集合 $\{1,2,3,4\}$ 的子集，共有 $2^4=16$ 个.}
\solend
\fi

\item 陈述句“$a>1$ 或 $b>3$”的否定形式为\blankline[4.4em].

\ifanswer
\ans{$a\leqslant1$ 且 $b\leqslant3$}
\fi

\item 已知方程 $3x^2+4x-4=0$ 的两个根为 $x_1$、$x_2$，则 $x_1^2+x_2^2=$\blankline[4.4em].

\ifanswer
\solbegin
\step{对于方程 $3x^2+4x-4=0$，因为 $\Delta=16-4\times3\times(-4)=64>0$，}
\step{所以方程 $3x^2+4x-4=0$ 的两个不相等的实根为 $x_1$、$x_2$，}
\step{由韦达定理可得 $x_1+x_2=-\dfrac43$，$x_1x_2=-\dfrac43$，}
\step{所以 $x_1^2+x_2^2=(x_1+x_2)^2-2x_1x_2=\dfrac{16}9+\dfrac83=\dfrac{40}9$.}
\solend
\fi

\item 设集合 $A=\{y\mid y=x^2+1\}$,$B=\{y\mid y=-x^2+2x+3\}$，则 $A\cap B=$\blankline[4.4em].

\ifanswer
\solbegin
\step{因为 $y=x^2+1\geqslant1$,$y=-x^2+2x+3=-(x-1)^2+4\leqslant4$，}
\step{所以 $A=\{y\mid y\geqslant1\}$,$B=\{y\mid y\leqslant4\}$，得 $A\cap B=\{y\mid1\leqslant y\leqslant4\}$.}
\solend
\fi

\item 若 $x\in\{1,x^2\}$，则 $x=$\blankline[4.4em].

\ifanswer
\ans{$0$}
\fi

\item 已知等式 $2x^2+3x+5=a(2x+1)(x+1)+c$，则 $a+c=$\blankline[4.4em].

\ifanswer
\solbegin
\step{因为 $2x^2+3x+5=a(2x+1)(x+1)+c$ 恒成立，}
\step{所以 $2x^2+3x+5=2ax^2+3ax+a+c$ 恒成立，则 $\begin{cases}2=2a\\ 3=3a\\ 5=a+c\end{cases}$，解得 $a+c=5$.}
\solend
\fi

\item 已知 $a\in\R$，设 $p:x<2$，$q:x<5a-3$，且 $p$ 是 $q$ 的必要非充分条件，则 $a$ 的取值范围是\blankline[4.4em].

\ifanswer
\solbegin
\step{设 $p:x<2$，$q:x<5a-3$，且 $p$ 是 $q$ 的必要非充分条件，}
\step{则 $\{x\mid x<5a-3\}\subset\{x\mid x<2\}$，所以 $5a-3<2$，即 $a<1$.}
\solend
\fi

\item 已知 $a\in\R$，集合 $A=\{x\mid ax^2+3x-1=0\}$，若集合 $A$ 中只有一个元素，则 $a$ 的取值范围是\blankline[4.4em].

\ifanswer
\solbegin
\step{当 $a=0$ 时，则有 $\{x\mid ax^2+3x-1=0\}=\{x\mid3x-1=0\}=\left\{\dfrac13\right\}$，符合题意；}
\step{当 $a\neq0$ 时，由题意得 $\Delta=9+4a=0$，解得 $a=-\dfrac94$. 符合题意；}
\step{综上所述，实数 $a$ 的取值集合为 $\left\{0,-\dfrac94\right\}$.}
\solend
\fi

\item 设集合 $A=\{x\mid x^2+2(m+1)x+m^2-1=0,m\in\R\}$，且 $A\subseteq\{0,-4\}$，则实数 $m$ 的取值范围是\blankline[4.4em].

\ifanswer
\solbegin
\step{由题意得 $A=\varnothing$ 或 $A=\{0\}$ 或 $A=\{-4\}$ 或 $A=\{0,-4\}$，}
\step{当 $A=\varnothing$ 时，方程 $x^2+2(m+1)x+m^2-1=0$ 无实根，}
\step{则 $\Delta=[2(m+1)]^2-4\times1\times(m^2-1)=8m+8<0$，解得 $m<-1$，}
\step{当 $A=\{0\}$ 时，方程 $x^2+2(m+1)x+m^2-1=0$ 有一个实根 $0$，}
\step{则 $\begin{cases}\Delta=8m+8=0\\ m^2-1=0\end{cases}$，解得 $m=-1$，}
\step{当 $A=\{-4\}$ 时，方程 $x^2+2(m+1)x+m^2-1=0$ 有一个实根 $-4$，}
\step{则 $\begin{cases}\Delta=8m+8=0\\ 16-8(m+1)+m^2-1=0\end{cases}$，无解，}
\step{当 $A=\{0,-4\}$ 时，方程 $x^2+2(m+1)x+m^2-1=0$ 有两个实根 $-4$ 和 $0$，}
\step{则 $\begin{cases}\Delta=8m+8>0\\ 16-8(m+1)+m^2-1=0\\ m^2-1=0\end{cases}$，解得 $m=1$，}
\step{综上所述，实数 $m$ 的取值范围是 $\{m\mid m\leqslant-1\text{ 或 }m=1\}$.}
\solend
\fi

\item 对于任意实数 $x$，$[x]$ 表示不超过 $x$ 的最大整数，如 $[1.1]=1$，$[-1.2]=-2$，若
$A=\{y\mid y=[2x]+[3x],-1\leqslant x\leqslant2\}$，则 $A$ 中所有元素的和为\blankline[4.4em].

\ifanswer
\solbegin
% 原卷“法一：”“法二：”标号排红色，本稿按 高一53“情况1：”成例排 \stepm（加粗着色），见文件头
\stepm{法一：}{当 $-1\leqslant x<-\dfrac23$ 时，$[2x]=-2$，$[3x]=-3$，$f(x)=[2x]+[3x]=-5$，}
\step{当 $-\dfrac23\leqslant x<-\dfrac12$ 时，$[2x]=-2$，$[3x]=-2$，$f(x)=[2x]+[3x]=-4$；}
\step{当 $-\dfrac12\leqslant x<-\dfrac13$ 时，$[2x]=-1$，$[3x]=-2$，$f(x)=[2x]+[3x]=-3$；}
\step{当 $-\dfrac13\leqslant x<0$ 时，$[2x]=-1$，$[3x]=-1$，$f(x)=[2x]+[3x]=-2$；}
\step{当 $0\leqslant x<\dfrac13$ 时，$[2x]=0$，$[3x]=0$，$f(x)=[2x]+[3x]=0$；}
\step{当 $\dfrac13\leqslant x<\dfrac12$ 时，$[2x]=0$，$[3x]=1$，$f(x)=[2x]+[3x]=1$；}
\step{当 $\dfrac12\leqslant x<\dfrac23$ 时，$[2x]=1$，$[3x]=1$，$f(x)=[2x]+[3x]=2$；}
\step{当 $\dfrac23\leqslant x<1$ 时，$[2x]=1$，$[3x]=2$，$f(x)=[2x]+[3x]=3$；}
\step{当 $1\leqslant x<\dfrac43$ 时，$[2x]=2$，$[3x]=3$，$f(x)=[2x]+[3x]=5$；}
\step{当 $\dfrac43\leqslant x<\dfrac32$ 时，$[2x]=2$，$[3x]=4$，$f(x)=[2x]+[3x]=6$；}
\step{当 $\dfrac32\leqslant x<\dfrac53$ 时，$[2x]=3$，$[3x]=4$，$f(x)=[2x]+[3x]=7$；}
\step{当 $\dfrac53\leqslant x<2$ 时，$[2x]=3$，$[3x]=5$，$f(x)=[2x]+[3x]=8$；}
\step{当 $x=2$ 时，$[2x]=4$，$[3x]=6$，$f(x)=[2x]+[3x]=10$}
\step{故 $A$ 中所有元素的和为 $-5-4-3-2+0+1+2+3+5+6+7+8+10=28$.}
% 原卷“法二：”标号排红色，同上
\stepm{法二：}{设 $x=m+n$，其中 $m\in\Z$，$n\in[0,1)$，}
\step{若 $0\leqslant n<\dfrac13$，则 $0\leqslant2n<\dfrac23$，$0\leqslant3n<1$，}
\step{得 $[2x]=[2m+2n]=2m$，$[3x]=[3m+3n]=3m$，所以 $y=[2x]+[3x]=5m$；}
\step{若 $\dfrac13\leqslant n<\dfrac12$，则 $\dfrac23\leqslant2n<1$，$1\leqslant3n<\dfrac32$，}
\step{得 $[2x]=[2m+2n]=2m$，$[3x]=[3m+3n]=3m+1$，}
\step{所以 $y=[2x]+[3x]=5m+1$；}
\step{若 $\dfrac12\leqslant n<\dfrac23$，则 $1\leqslant2n<\dfrac43$，$\dfrac32\leqslant3n<2$，}
\step{得 $[2x]=[2m+2n]=2m+1$，$[3x]=[3m+3n]=3m+1$，}
\step{所以 $y=[2x]+[3x]=5m+2$；}
\step{若 $\dfrac23\leqslant n<1$，则 $\dfrac43\leqslant2n<2$，$2\leqslant3n<3$，}
\step{得 $[2x]=[2m+2n]=2m+1$，$[3x]=[3m+3n]=3m+2$，}
\step{所以 $y=[2x]+[3x]=5m+3$；}
\step{综上所述，当 $n\in[0,1)$ 时，所有函数值的和}
\step{为 $5m+5m+1+5m+2+5m+3=20m+6$，}
\step{又因为 $x\in[-1,2]$，此时 $m=-1,0,1$，且当 $x=2$ 时，$y=[2x]+[3x]=10$，}
\step{所有函数值的和为 $(-14)+6+26+10=28$，即 $A$ 中所有元素之和为 $28$.}
\solend
\fi

\end{enumerate}

\bigsec{二、选择题}

\begin{enumerate}
\setcounter{enumi}{12}

\item 下列写法中，正确的是\mbox{（\quad）}

\keepwithprev
A. $0\in\{\varnothing\}$\qquad B. $\varnothing\in\{\varnothing\}$

C. $\varnothing=\{\varnothing\}$\qquad D. $0\in\varnothing$

\ifanswer
\solbegin
\step{$0\notin\{\varnothing\}$，$\varnothing\in\{\varnothing\}$，故 $A$ 错误，$B$ 正确；}
\step{$\varnothing$ 是不含任何元素的集合，所以 $\varnothing\neq\{\varnothing\}$，$0\notin\varnothing$，故 $CD$ 错误.}
\step{故选 $B$.}
\solend
\fi

% 第 14 题答案“D”直接印在题干括号内“(  D  )”，且该题无【答案】/【解析】段；
% 按 高一37 第 13 题的成例括号留空、答案另提为 \ans{}（见文件头）
\item 若 $a<b<0$，则下列不等式中不能成立的是\mbox{（\quad）}

\keepwithprev
A. $\dfrac1a>\dfrac1b$\qquad B. $a-b>a$

C. $|a|>|b|$\qquad D. $a^3>b^3$

\ifanswer
\ans{$D$}
\fi

\item 已知 $x$、$y$、$z$ 为非零实数，代数式 $\dfrac x{|x|}+\dfrac y{|y|}+\dfrac z{|z|}+\dfrac{xyz}{|xyz|}$
的值所组成的集合是 $M$，则下列判断正确的是\mbox{（\quad）}

\keepwithprev
A. $-4\notin M$\qquad B. $2\in M$

C. $4\notin M$\qquad D. $0\in M$

\ifanswer
\solbegin
\step{当 $x$、$y$、$z$ 均为正数时，$xyz>0$，$\dfrac x{|x|}+\dfrac y{|y|}+\dfrac z{|z|}+\dfrac{xyz}{|xyz|}=1+1+1+1=4$，}
\step{当 $x$、$y$、$z$ 有 2 个正数，1 个负数时，不妨设 $x>0$，$y>0$，$z<0$，}
\step{则 $xyz<0$，$\dfrac x{|x|}+\dfrac y{|y|}+\dfrac z{|z|}+\dfrac{xyz}{|xyz|}=1+1-1-1=0$，}
\step{若 $x$、$y$、$z$ 有 1 个正数，2 个负数时，不妨设 $x>0$，$y<0$，$z<0$，}
\step{则 $xyz>0$，$\dfrac x{|x|}+\dfrac y{|y|}+\dfrac z{|z|}+\dfrac{xyz}{|xyz|}=1-1-1+1=0$，}
\step{当 $x$、$y$、$z$ 均为负数时，$xyz<0$，}
\step{$\dfrac x{|x|}+\dfrac y{|y|}+\dfrac z{|z|}+\dfrac{xyz}{|xyz|}=-1-1-1-1=-4$，}
\step{综上，$M=\{-4,0,4\}$，}
\step{所以 $-4\in M$，$A$ 错误；$2\notin M$，$B$ 错误；$4\in M$，$C$ 错误；}
\step{$0\in M$，$D$ 正确. 故选 $D$.}
\solend
\fi

\item 若 $xy>0$，且 $x+y=9$，则 $\dfrac1x+\dfrac4y$ 的最小值是\mbox{（\quad）}

\keepwithprev
A. $9$\qquad B. $\dfrac59$

C. $\dfrac49$\qquad D. $1$

\ifanswer
\solbegin
\step{$xy>0$，且 $x+y=9$，所以 $x>0$，$y>0$，}
\step{所以 $\dfrac1x+\dfrac4y=\dfrac19\left(\dfrac1x+\dfrac4y\right)(x+y)=\dfrac59+\dfrac19\left(\dfrac yx+\dfrac{4x}y\right)\geqslant\dfrac59+\dfrac29\cdot\sqrt{\dfrac yx\cdot\dfrac{4x}y}=\dfrac59+\dfrac49=1$，}
\step{当且仅当 $\dfrac yx=\dfrac{4x}y$，即 $x=3$，$y=6$ 时取等号，所以 $\dfrac1x+\dfrac4y$ 的最小值为 $1$.}
\step{故选 $D$.}
\solend
\fi

\end{enumerate}

\bigsec{三、解答题}

\begin{enumerate}
\setcounter{enumi}{16}

\item 求解下列不等式：

（1）$\dfrac{x-3}{x+1}\leqslant0$；

（2）$|2x-1|<3$.

\ifanswer
\solbegin
\stepm{(1)}{$\dfrac{x-3}{x+1}\leqslant0$，变形得 $(x-3)\cdot(x+1)\leqslant0$ 且 $x\neq-1$，}
\step{解得 $-1<x\leqslant3$，即不等式的解集为 $(-1,3]$.}
\stepm{(2)}{若 $|2x-1|<3$，则 $-3<2x-1<3$，解得 $-1<x<2$，即 $(-1,2)$.}
\solend
\fi

\ansspace[6]

\item 已知集合 $A=\{x\mid-2<x<15\}$,$B=\{x\mid m-6<x<2m-1,m\in\R\}$.

（1）当 $m=2$ 时，求 $A\cap B$；

（2）若 $A\cup B=A$，求实数 $m$ 的取值范围.

\ifanswer
\solbegin
\stepm{(1)}{当 $m=2$ 时，$A=\{x\mid-2<x<15\}$,$B=\{x\mid-4<x<3\}$，}
\step{此时 $A\cap B=\{x\mid-2<x<3\}$.}
\stepm{(2)}{由 $A\cup B=A$ 得 $B\subseteq A$.}
\step{当 $B=\varnothing$,$m-6\geqslant2m-1$，解得 $m\leqslant-5$；}
\step{当 $B\neq\varnothing$，结合 $B\subseteq A$ 得 $\begin{cases}m-6<2m-1\\ 2m-1\leqslant15\\ m-6\geqslant-2\end{cases}$，解得 $4\leqslant m\leqslant8$.}
\step{综上，$m\leqslant-5$ 或 $4\leqslant m\leqslant8$.}
\solend
\fi

\ansspace[6]

\item 已知 $a\in\R$，命题 $p$：关于 $x$ 的方程 $x^2-2x+a=0$ 无实根；命题 $q$：对任意 $x\in\R$，
不等式 $ax^2-ax+1>0$ 恒成立.

（1）若 $q$ 为真命题，求 $a$ 的取值范围；

（2）若 $p$、$q$ 中有且只有一个为真命题，求 $a$ 的取值范围.

\ifanswer
\solbegin
\stepm{(1)}{当 $q$ 为真命题时，}
\stepm{①}{当 $a=0$ 时，不等式化为 $1>0$，对任意 $x\in\R$ 恒成立，满足条件；}
\stepm{②}{当 $a\neq0$ 时，$ax^2-ax+1>0$ 为一元二次不等式，}
\step{若对任意 $x\in\R$ 恒成立，需满足 $\begin{cases}a>0\\ \Delta=(-a)^2-4a<0\end{cases}$，解得 $0<a<4$；}
\step{综上，$a$ 的取值范围为 $[0,4)$.}
\stepm{(2)}{$p$ 为真时，方程 $x^2-2x+a=0$ 无实根，}
\step{故判别式 $\Delta=4-4a<0$，解得 $a>1$，因此 $p$ 为假时，$a\leqslant1$；}
\step{由（1）得 $q$ 为假时 $a\geqslant4$ 或 $a<0$，}
\step{$p$、$q$ 有且只有一个真命题，分两类：}
\step{当 $p$ 假 $q$ 真时，满足 $\begin{cases}a\leqslant1\\ 0\leqslant a<4\end{cases}$，解得 $0\leqslant a\leqslant1$；}
\step{当 $p$ 真 $q$ 假时，满足 $\begin{cases}a>1\\ a\geqslant4\text{或}a<0\end{cases}$，解得 $a\geqslant4$.}
\step{综上，$a$ 的取值范围是 $[0,1]\cup[4,+\infty)$.}
\solend
\fi

\ansspace[8]

\item 某工厂生产 $x$($x>0$) 件产品的总成本 $C$（单位：元）与生产件数 $x$ 之间满足二次函数
$C(x)=\dfrac14x^2+50x+10000$.

（1）若要求生产总成本不超过 30000 元，求满足生产总成本要求的产品件数的取值范围；

（2）试求该工厂生产多少件产品时，可使每件产品的平均成本最小.

\ifanswer
\solbegin
\stepm{(1)}{若要求生产总成本不超过 30000 元，}
\step{则 $C(x)=\dfrac14x^2+50x+10000\leqslant30000$，}
\step{即 $x^2+200x-80000\leqslant0$，解得 $-400\leqslant x\leqslant200$，又 $x>0$，}
\step{所以满足生产总成本要求的产品件数的取值范围为 $(0,200]$，$x\in\N$；}
\stepm{(2)}{生产件数 $x$ 时每件的平均成本为 $\dfrac{\frac14x^2+50x+10000}x=\dfrac14x+50+\dfrac{10000}x$，}
\step{由基本不等式得 $\dfrac14x+50+\dfrac{10000}x\geqslant2\sqrt{\dfrac14x\cdot\dfrac{10000}x}+50=150$，}
\step{当且仅当 $\dfrac14x=\dfrac{10000}x$，即 $x=200$ 时取等号，}
\step{即该工厂生产 200 件产品时，每件产品的平均成本最小.}
\solend
\fi

\ansspace[8]

% 原卷题干“$(n\geqslant2,n\in N)$”后连印两个全角逗号“，，”，照录未改（见文件头“原卷笔误”①）
\item 已知有限集 $A=\{a_1,a_2,\cdots,a_n\}$($n\geqslant2$,$n\in\N$)，，如果 $A$ 中元素
$a_i$($i=1,2,3,\cdots,n$) 满足 $a_1+a_2+a_3+\cdots+a_n=a_1\times a_2\times\cdots\times a_n$，
就称 $A$ 为“和积平衡集”.

% 原卷 (2) 将“写方程”与“反证法证明”两项任务并作一句（“……表示），可尝试用反证法证明：……”），
% 照录未改（见文件头“原卷笔误”②）
（1）判断集合 $\{-1+\sqrt3,-1-\sqrt3\}$ 是否为“和积平衡集”；

（2）若 $a_1$、$a_2$ 是两个不同的正数，且 $\{a_1,a_2\}$ 是“和积平衡集”，写出以 $a_1$、$a_2$ 为根的
一个一元二次方程（系数可用 $a_1$、$a_2$ 表示），可尝试用反证法证明：$a_1$、$a_2$ 至少有一个大于 2；

（3）若 $a_i\geqslant1$，且 $a_i\in\N$，求所有符合条件的“和积平衡集”.

\ifanswer
\solbegin
\stepm{(1)}{对于集合 $\{-1+\sqrt3,-1-\sqrt3\}$，}
\step{因为 $(-1+\sqrt3)+(-1-\sqrt3)=-2$，$(-1+\sqrt3)\times(-1-\sqrt3)=-2$，}
\step{所以满足 $a_1+a_2=a_1\times a_2$，满足题意，}
\step{所以集合 $\{-1+\sqrt3,-1-\sqrt3\}$ 是和积平衡集.}
\stepm{(2)}{以 $a_1$、$a_2$ 为根的一个一元二次方程可为 $x^2-(a_1+a_2)x+a_1\cdot a_2=0$，}
\step{因为 $a_1$、$a_2$ 是两个不同的正数，且 $\{a_1,a_2\}$ 是和积平衡集，}
\step{所以 $a_1+a_2=a_1\cdot a_2$，}
\step{又 $a_1$ 和 $a_2$ 为 $x^2-(a_1+a_2)x+a_1\cdot a_2=0$ 的两个不同的正根，}
\step{所以 $\Delta=(a_1+a_2)^2-4a_1a_2=(a_1+a_2)^2-4(a_1+a_2)>0$，}
\step{又 $a_1+a_2>0$，则 $a_1+a_2>4$，}
\step{若 $a_1$、$a_2$ 都小于等于 2，则 $a_1+a_2<4$，矛盾，}
\step{所以 $a_1$、$a_2$ 至少有一个大于 2；}
\stepm{(3)}{设 $A$ 中的 $a_1<a_2<a_3<\cdots<a_n$，且 $a_i\in\N^*$，}
\step{由 $a_1+a_2+a_3+\cdots+a_n=a_1\times a_2\times\cdots\times a_n$，得 $n=2$ 时明显不满足，}
\step{所以 $n\geqslant3$，故 $a_1\times a_2\times\cdots\times a_n=a_1+a_2+a_3+\cdots+a_n\leqslant na_n$，}
\step{所以 $a_1\times a_2\times\cdots\times a_{n-1}\leqslant n$，}
\step{当 $n=3$ 时，$a_1\times a_2\leqslant3$，$a_i\in\N^*$，只能有 $a_1=1$，$a_2=2$，}
\step{由 $a_1+a_2+a_3=a_1\times a_2\times a_3$，得 $a_3=3$，此时的和积平衡集为 $\{1,2,3\}$，}
\step{当 $n\geqslant4$ 时，$a_1\times a_2\times\cdots\times a_{n-1}\geqslant(n-1)(n-2)\times\cdots\times2\times1\geqslant(n-1)(n-2)$，}
\step{又 $(n-1)(n-2)-n=n^2-4n+2=n(n-4)+2>0$，即 $(n-1)(n-2)>n$，}
\step{即 $a_1\times a_2\times\cdots\times a_{n-1}>n$，与 $a_1\times a_2\times\cdots\times a_{n-1}\leqslant n$ 矛盾，}
\step{所以 $n\geqslant4$ 不满足条件，}
\step{所以符合条件的和积平衡集有且只有 $n=3$，此时的和积平衡集为 $\{1,2,3\}$.}
\solend
\fi

\ansspace*[10]

\end{enumerate}

\end{document}
"
% 紧凑版：xelatex "\def\iscompact{1}% 高一63_吴淞中学_10月月考.tex
%   —— 高一上数学“2026-2027 年吴淞中学高一上 10 月月考”（讲评版）
% 试卷版：xelatex 高一63_吴淞中学_10月月考.tex
% 答案版：xelatex "\def\isanswer{1}\input{高一63_吴淞中学_10月月考.tex}"
% 紧凑版：xelatex "\def\iscompact{1}\input{高一63_吴淞中学_10月月考.tex}"
%
% 原始材料：微信公众号“魔都数学加油站”文章，作者破晓，连载“27学年高一63”；
% 原文标题“27学年高一63——2026-2027年上海市吴淞中学高一上10月月考”；
% 原文链接：https://mp.weixin.qq.com/s/TFQ4TLnt_l66t7t51es4OQ。
% 页面标注发布时间 2026-10-16 21:44（上海）。
% 正文为 10 张 PNG 页（`原图/高一63_吴淞中学_10月月考/01..10.png`）：
% 第 06、10 页 2560×3620，第 01、02、03、04、05、07、08、09 页 4762×6735
% （**同一篇各页分辨率不同**），均按页序读取。10 页全为试卷正文页（卷首标题在第 1 页页顶，
% 卷末解答题第 21 题解析止于第 10 页），无封面页、目录页、二维码或宣传页。
% 粉色斜排“魔都数学加油站”为卷面水印，与公众号页脚一并未录入。
%
% 卷首标题：照卷面印作“2026--2027 年吴淞中学高一上 10 月月考”（公众号标题中的“上海市”
% 卷面标题行没有，本稿照卷面），年份区间按全稿惯例排 en dash。原卷只有一行标题、无副标题。
%
% 篇幅与题号：原卷自第 3 题起编号（第 1、2 题不在本卷内）。填空题 3—12（10 题）、
% 选择题 13—16（4 题）、解答题 17—21（5 题），共 19 题。本稿沿用原节与原题号，
% 只改排为全稿统一的大节样式（原卷小节标题为“一、填空题”“二、选择题”“三、解答题”）。
%
% 讲评版：原卷每题下直接印【答案】或【解析】。**只印【答案】的 2 题**（第 4、7 题），
% 本稿只给 \ans{}；其余 17 题里 16 题都印【解析】（按全稿口径这类题不另提 \ans{}——第 13、15、16 题
% 的结论写在解析末句（“故选 B”等），亦照录在解析内；第 14 题则既无【答案】也无【解析】段。**另有 1 题答案印在题干括号内**——
% 第 14 题（“(  D  )”，且该题无【答案】/【解析】段），本稿按 高一37 第 13 题成例把括号留空、
% 答案另提为 \ans{$D$}。第 16 题题干括号为空但印有【解析】，不算括号内答案。
%
% 副标题：原卷未印副标题。本稿按“知识模块”口径标作“集合与常用逻辑用语、不等式”：
% 19 题中集合与常用逻辑用语 12 题（第 3、4、6、7、9、10、11、12、13、15、18、21 题），
% 不等式 7 题（含等式（韦达定理、恒等式）与不等式性质、基本不等式、二次不等式、恒成立：
% 第 5、8、14、16、17、19、20 题），两类均超过三分之一，故并标。口径同 高一59。
%
% 与原卷的排版差异：
%   ① 原卷每题下直接印【答案】/【解析】，本稿按全稿惯例将答案与解析置于答案版，试卷版保持干净；
%   ② 原卷未印副标题，本稿按“知识模块”口径补出；
%   ③ 原卷填空横线约四个字宽，本稿按全稿惯例统一排 \blankline[4.4em]；
%   ④ 原卷实数集排斜体 $R$（第 9、10、19 题题干与解析）或粗体 $\mathbf{R}$（第 18 题题干），
%      整数集排斜体 $Z$（第 12 题解析），自然数集排斜体 $N$（第 20 题解析、第 21 题题干
%      “$(n\geqslant2,n\in N)$”）或粗体 $\mathbf{N}$（第 21 题 (3) 题干）、斜体 $N^*$
%      （第 21 题 (3) 解析），本稿按全稿惯例统一改用 $\R$、$\Z$、$\N$、$\N^*$
%      （同 高一14 排版差异④：记号归一，粗体／斜体不作区分，$N$ 与 $N^*$ 的区别照原卷保留）；
%      空集记号原卷作 $\varnothing$，照录；
%   ⑤ 原卷在数学式内多处用半角逗号（如 $\{y\mid y=[2x]+[3x],-1\leqslant x\leqslant2\}$、
%      $m=-1,0,1$、$A=\{x\mid-2<x<15\}$,$B=\cdots$），均照录未强改；
%      句末句点按全稿惯例排西文句点；“和积平衡集”自定义名词照录；
%   ⑥ 第 12 题解析里的“法一：”“法二：”标号原卷排**红色**（全卷共 2 处），
%      同 高一53“情况1：”成例，本稿排 \stepm{法一：}{…}、\stepm{法二：}{…}
%      （加粗着色、不改黑色基调），冒号照原卷排全角；
%   ⑦ 第 17—21 题的子问原卷各自另起一行，本稿按全稿惯例用 enumerate 的 (1)(2)(3)；
%      解析中的子问标号一律用 \stepm{(1)}{…}；第 19 题解析中的①②分情形，按全稿惯例作
%      \stepm{①}{…}、\stepm{②}{…}；
%   ⑧ 第 14 题印在题干括号内的答案移入 \ans{}，见上；
%   ⑨ 每道解答题原卷未留作答空白，本稿按全稿惯例加 \ansspace；
%   ⑩ 选择题的作答括号原卷排半角“(　　)”（第 13、15、16 题括号内空白、第 14 题括号内印 D），
%      本稿按全稿惯例排 \mbox{（\quad）}。
%
% 原卷笔误/存疑（均照录未改，正文对应处另有 % 注）：
%   ① 第 21 题题干“$(n\geqslant2$,$n\in N)$”后连印两个全角逗号“，，”（衍字，
%      8 倍放大确认，非水印干扰），本稿照录；
%   ② 第 21 题 (2) 题干将“写出以 $a_1$、$a_2$ 为根的一个一元二次方程”与“用反证法证明：
%      $a_1$、$a_2$ 至少有一个大于 2”两项任务并作一句、之间仅以逗号相接（“……表示），
%      可尝试用反证法证明：……”），读来像抄并两小问，原卷如此，照录未改。
%
% 插图：无（全卷 10 页无任何图形、无“如图”字样）。
%
\documentclass[a4paper,11pt]{article}
\usepackage{common}

\begin{document}

\examtitlest[3]{2026--2027 年吴淞中学高一上 10 月月考}{集合与常用逻辑用语、不等式}

\bigsec{一、填空题}

\begin{enumerate}
\setcounter{enumi}{2}

\item 满足 $A\subseteq\{1,2,3,4\}$ 的集合 $A$ 共有\blankline[4.4em]个.

\ifanswer
\solbegin
\step{由题意，集合 $A$ 是集合 $\{1,2,3,4\}$ 的子集，共有 $2^4=16$ 个.}
\solend
\fi

\item 陈述句“$a>1$ 或 $b>3$”的否定形式为\blankline[4.4em].

\ifanswer
\ans{$a\leqslant1$ 且 $b\leqslant3$}
\fi

\item 已知方程 $3x^2+4x-4=0$ 的两个根为 $x_1$、$x_2$，则 $x_1^2+x_2^2=$\blankline[4.4em].

\ifanswer
\solbegin
\step{对于方程 $3x^2+4x-4=0$，因为 $\Delta=16-4\times3\times(-4)=64>0$，}
\step{所以方程 $3x^2+4x-4=0$ 的两个不相等的实根为 $x_1$、$x_2$，}
\step{由韦达定理可得 $x_1+x_2=-\dfrac43$，$x_1x_2=-\dfrac43$，}
\step{所以 $x_1^2+x_2^2=(x_1+x_2)^2-2x_1x_2=\dfrac{16}9+\dfrac83=\dfrac{40}9$.}
\solend
\fi

\item 设集合 $A=\{y\mid y=x^2+1\}$,$B=\{y\mid y=-x^2+2x+3\}$，则 $A\cap B=$\blankline[4.4em].

\ifanswer
\solbegin
\step{因为 $y=x^2+1\geqslant1$,$y=-x^2+2x+3=-(x-1)^2+4\leqslant4$，}
\step{所以 $A=\{y\mid y\geqslant1\}$,$B=\{y\mid y\leqslant4\}$，得 $A\cap B=\{y\mid1\leqslant y\leqslant4\}$.}
\solend
\fi

\item 若 $x\in\{1,x^2\}$，则 $x=$\blankline[4.4em].

\ifanswer
\ans{$0$}
\fi

\item 已知等式 $2x^2+3x+5=a(2x+1)(x+1)+c$，则 $a+c=$\blankline[4.4em].

\ifanswer
\solbegin
\step{因为 $2x^2+3x+5=a(2x+1)(x+1)+c$ 恒成立，}
\step{所以 $2x^2+3x+5=2ax^2+3ax+a+c$ 恒成立，则 $\begin{cases}2=2a\\ 3=3a\\ 5=a+c\end{cases}$，解得 $a+c=5$.}
\solend
\fi

\item 已知 $a\in\R$，设 $p:x<2$，$q:x<5a-3$，且 $p$ 是 $q$ 的必要非充分条件，则 $a$ 的取值范围是\blankline[4.4em].

\ifanswer
\solbegin
\step{设 $p:x<2$，$q:x<5a-3$，且 $p$ 是 $q$ 的必要非充分条件，}
\step{则 $\{x\mid x<5a-3\}\subset\{x\mid x<2\}$，所以 $5a-3<2$，即 $a<1$.}
\solend
\fi

\item 已知 $a\in\R$，集合 $A=\{x\mid ax^2+3x-1=0\}$，若集合 $A$ 中只有一个元素，则 $a$ 的取值范围是\blankline[4.4em].

\ifanswer
\solbegin
\step{当 $a=0$ 时，则有 $\{x\mid ax^2+3x-1=0\}=\{x\mid3x-1=0\}=\left\{\dfrac13\right\}$，符合题意；}
\step{当 $a\neq0$ 时，由题意得 $\Delta=9+4a=0$，解得 $a=-\dfrac94$. 符合题意；}
\step{综上所述，实数 $a$ 的取值集合为 $\left\{0,-\dfrac94\right\}$.}
\solend
\fi

\item 设集合 $A=\{x\mid x^2+2(m+1)x+m^2-1=0,m\in\R\}$，且 $A\subseteq\{0,-4\}$，则实数 $m$ 的取值范围是\blankline[4.4em].

\ifanswer
\solbegin
\step{由题意得 $A=\varnothing$ 或 $A=\{0\}$ 或 $A=\{-4\}$ 或 $A=\{0,-4\}$，}
\step{当 $A=\varnothing$ 时，方程 $x^2+2(m+1)x+m^2-1=0$ 无实根，}
\step{则 $\Delta=[2(m+1)]^2-4\times1\times(m^2-1)=8m+8<0$，解得 $m<-1$，}
\step{当 $A=\{0\}$ 时，方程 $x^2+2(m+1)x+m^2-1=0$ 有一个实根 $0$，}
\step{则 $\begin{cases}\Delta=8m+8=0\\ m^2-1=0\end{cases}$，解得 $m=-1$，}
\step{当 $A=\{-4\}$ 时，方程 $x^2+2(m+1)x+m^2-1=0$ 有一个实根 $-4$，}
\step{则 $\begin{cases}\Delta=8m+8=0\\ 16-8(m+1)+m^2-1=0\end{cases}$，无解，}
\step{当 $A=\{0,-4\}$ 时，方程 $x^2+2(m+1)x+m^2-1=0$ 有两个实根 $-4$ 和 $0$，}
\step{则 $\begin{cases}\Delta=8m+8>0\\ 16-8(m+1)+m^2-1=0\\ m^2-1=0\end{cases}$，解得 $m=1$，}
\step{综上所述，实数 $m$ 的取值范围是 $\{m\mid m\leqslant-1\text{ 或 }m=1\}$.}
\solend
\fi

\item 对于任意实数 $x$，$[x]$ 表示不超过 $x$ 的最大整数，如 $[1.1]=1$，$[-1.2]=-2$，若
$A=\{y\mid y=[2x]+[3x],-1\leqslant x\leqslant2\}$，则 $A$ 中所有元素的和为\blankline[4.4em].

\ifanswer
\solbegin
% 原卷“法一：”“法二：”标号排红色，本稿按 高一53“情况1：”成例排 \stepm（加粗着色），见文件头
\stepm{法一：}{当 $-1\leqslant x<-\dfrac23$ 时，$[2x]=-2$，$[3x]=-3$，$f(x)=[2x]+[3x]=-5$，}
\step{当 $-\dfrac23\leqslant x<-\dfrac12$ 时，$[2x]=-2$，$[3x]=-2$，$f(x)=[2x]+[3x]=-4$；}
\step{当 $-\dfrac12\leqslant x<-\dfrac13$ 时，$[2x]=-1$，$[3x]=-2$，$f(x)=[2x]+[3x]=-3$；}
\step{当 $-\dfrac13\leqslant x<0$ 时，$[2x]=-1$，$[3x]=-1$，$f(x)=[2x]+[3x]=-2$；}
\step{当 $0\leqslant x<\dfrac13$ 时，$[2x]=0$，$[3x]=0$，$f(x)=[2x]+[3x]=0$；}
\step{当 $\dfrac13\leqslant x<\dfrac12$ 时，$[2x]=0$，$[3x]=1$，$f(x)=[2x]+[3x]=1$；}
\step{当 $\dfrac12\leqslant x<\dfrac23$ 时，$[2x]=1$，$[3x]=1$，$f(x)=[2x]+[3x]=2$；}
\step{当 $\dfrac23\leqslant x<1$ 时，$[2x]=1$，$[3x]=2$，$f(x)=[2x]+[3x]=3$；}
\step{当 $1\leqslant x<\dfrac43$ 时，$[2x]=2$，$[3x]=3$，$f(x)=[2x]+[3x]=5$；}
\step{当 $\dfrac43\leqslant x<\dfrac32$ 时，$[2x]=2$，$[3x]=4$，$f(x)=[2x]+[3x]=6$；}
\step{当 $\dfrac32\leqslant x<\dfrac53$ 时，$[2x]=3$，$[3x]=4$，$f(x)=[2x]+[3x]=7$；}
\step{当 $\dfrac53\leqslant x<2$ 时，$[2x]=3$，$[3x]=5$，$f(x)=[2x]+[3x]=8$；}
\step{当 $x=2$ 时，$[2x]=4$，$[3x]=6$，$f(x)=[2x]+[3x]=10$}
\step{故 $A$ 中所有元素的和为 $-5-4-3-2+0+1+2+3+5+6+7+8+10=28$.}
% 原卷“法二：”标号排红色，同上
\stepm{法二：}{设 $x=m+n$，其中 $m\in\Z$，$n\in[0,1)$，}
\step{若 $0\leqslant n<\dfrac13$，则 $0\leqslant2n<\dfrac23$，$0\leqslant3n<1$，}
\step{得 $[2x]=[2m+2n]=2m$，$[3x]=[3m+3n]=3m$，所以 $y=[2x]+[3x]=5m$；}
\step{若 $\dfrac13\leqslant n<\dfrac12$，则 $\dfrac23\leqslant2n<1$，$1\leqslant3n<\dfrac32$，}
\step{得 $[2x]=[2m+2n]=2m$，$[3x]=[3m+3n]=3m+1$，}
\step{所以 $y=[2x]+[3x]=5m+1$；}
\step{若 $\dfrac12\leqslant n<\dfrac23$，则 $1\leqslant2n<\dfrac43$，$\dfrac32\leqslant3n<2$，}
\step{得 $[2x]=[2m+2n]=2m+1$，$[3x]=[3m+3n]=3m+1$，}
\step{所以 $y=[2x]+[3x]=5m+2$；}
\step{若 $\dfrac23\leqslant n<1$，则 $\dfrac43\leqslant2n<2$，$2\leqslant3n<3$，}
\step{得 $[2x]=[2m+2n]=2m+1$，$[3x]=[3m+3n]=3m+2$，}
\step{所以 $y=[2x]+[3x]=5m+3$；}
\step{综上所述，当 $n\in[0,1)$ 时，所有函数值的和}
\step{为 $5m+5m+1+5m+2+5m+3=20m+6$，}
\step{又因为 $x\in[-1,2]$，此时 $m=-1,0,1$，且当 $x=2$ 时，$y=[2x]+[3x]=10$，}
\step{所有函数值的和为 $(-14)+6+26+10=28$，即 $A$ 中所有元素之和为 $28$.}
\solend
\fi

\end{enumerate}

\bigsec{二、选择题}

\begin{enumerate}
\setcounter{enumi}{12}

\item 下列写法中，正确的是\mbox{（\quad）}

\keepwithprev
A. $0\in\{\varnothing\}$\qquad B. $\varnothing\in\{\varnothing\}$

C. $\varnothing=\{\varnothing\}$\qquad D. $0\in\varnothing$

\ifanswer
\solbegin
\step{$0\notin\{\varnothing\}$，$\varnothing\in\{\varnothing\}$，故 $A$ 错误，$B$ 正确；}
\step{$\varnothing$ 是不含任何元素的集合，所以 $\varnothing\neq\{\varnothing\}$，$0\notin\varnothing$，故 $CD$ 错误.}
\step{故选 $B$.}
\solend
\fi

% 第 14 题答案“D”直接印在题干括号内“(  D  )”，且该题无【答案】/【解析】段；
% 按 高一37 第 13 题的成例括号留空、答案另提为 \ans{}（见文件头）
\item 若 $a<b<0$，则下列不等式中不能成立的是\mbox{（\quad）}

\keepwithprev
A. $\dfrac1a>\dfrac1b$\qquad B. $a-b>a$

C. $|a|>|b|$\qquad D. $a^3>b^3$

\ifanswer
\ans{$D$}
\fi

\item 已知 $x$、$y$、$z$ 为非零实数，代数式 $\dfrac x{|x|}+\dfrac y{|y|}+\dfrac z{|z|}+\dfrac{xyz}{|xyz|}$
的值所组成的集合是 $M$，则下列判断正确的是\mbox{（\quad）}

\keepwithprev
A. $-4\notin M$\qquad B. $2\in M$

C. $4\notin M$\qquad D. $0\in M$

\ifanswer
\solbegin
\step{当 $x$、$y$、$z$ 均为正数时，$xyz>0$，$\dfrac x{|x|}+\dfrac y{|y|}+\dfrac z{|z|}+\dfrac{xyz}{|xyz|}=1+1+1+1=4$，}
\step{当 $x$、$y$、$z$ 有 2 个正数，1 个负数时，不妨设 $x>0$，$y>0$，$z<0$，}
\step{则 $xyz<0$，$\dfrac x{|x|}+\dfrac y{|y|}+\dfrac z{|z|}+\dfrac{xyz}{|xyz|}=1+1-1-1=0$，}
\step{若 $x$、$y$、$z$ 有 1 个正数，2 个负数时，不妨设 $x>0$，$y<0$，$z<0$，}
\step{则 $xyz>0$，$\dfrac x{|x|}+\dfrac y{|y|}+\dfrac z{|z|}+\dfrac{xyz}{|xyz|}=1-1-1+1=0$，}
\step{当 $x$、$y$、$z$ 均为负数时，$xyz<0$，}
\step{$\dfrac x{|x|}+\dfrac y{|y|}+\dfrac z{|z|}+\dfrac{xyz}{|xyz|}=-1-1-1-1=-4$，}
\step{综上，$M=\{-4,0,4\}$，}
\step{所以 $-4\in M$，$A$ 错误；$2\notin M$，$B$ 错误；$4\in M$，$C$ 错误；}
\step{$0\in M$，$D$ 正确. 故选 $D$.}
\solend
\fi

\item 若 $xy>0$，且 $x+y=9$，则 $\dfrac1x+\dfrac4y$ 的最小值是\mbox{（\quad）}

\keepwithprev
A. $9$\qquad B. $\dfrac59$

C. $\dfrac49$\qquad D. $1$

\ifanswer
\solbegin
\step{$xy>0$，且 $x+y=9$，所以 $x>0$，$y>0$，}
\step{所以 $\dfrac1x+\dfrac4y=\dfrac19\left(\dfrac1x+\dfrac4y\right)(x+y)=\dfrac59+\dfrac19\left(\dfrac yx+\dfrac{4x}y\right)\geqslant\dfrac59+\dfrac29\cdot\sqrt{\dfrac yx\cdot\dfrac{4x}y}=\dfrac59+\dfrac49=1$，}
\step{当且仅当 $\dfrac yx=\dfrac{4x}y$，即 $x=3$，$y=6$ 时取等号，所以 $\dfrac1x+\dfrac4y$ 的最小值为 $1$.}
\step{故选 $D$.}
\solend
\fi

\end{enumerate}

\bigsec{三、解答题}

\begin{enumerate}
\setcounter{enumi}{16}

\item 求解下列不等式：

（1）$\dfrac{x-3}{x+1}\leqslant0$；

（2）$|2x-1|<3$.

\ifanswer
\solbegin
\stepm{(1)}{$\dfrac{x-3}{x+1}\leqslant0$，变形得 $(x-3)\cdot(x+1)\leqslant0$ 且 $x\neq-1$，}
\step{解得 $-1<x\leqslant3$，即不等式的解集为 $(-1,3]$.}
\stepm{(2)}{若 $|2x-1|<3$，则 $-3<2x-1<3$，解得 $-1<x<2$，即 $(-1,2)$.}
\solend
\fi

\ansspace[6]

\item 已知集合 $A=\{x\mid-2<x<15\}$,$B=\{x\mid m-6<x<2m-1,m\in\R\}$.

（1）当 $m=2$ 时，求 $A\cap B$；

（2）若 $A\cup B=A$，求实数 $m$ 的取值范围.

\ifanswer
\solbegin
\stepm{(1)}{当 $m=2$ 时，$A=\{x\mid-2<x<15\}$,$B=\{x\mid-4<x<3\}$，}
\step{此时 $A\cap B=\{x\mid-2<x<3\}$.}
\stepm{(2)}{由 $A\cup B=A$ 得 $B\subseteq A$.}
\step{当 $B=\varnothing$,$m-6\geqslant2m-1$，解得 $m\leqslant-5$；}
\step{当 $B\neq\varnothing$，结合 $B\subseteq A$ 得 $\begin{cases}m-6<2m-1\\ 2m-1\leqslant15\\ m-6\geqslant-2\end{cases}$，解得 $4\leqslant m\leqslant8$.}
\step{综上，$m\leqslant-5$ 或 $4\leqslant m\leqslant8$.}
\solend
\fi

\ansspace[6]

\item 已知 $a\in\R$，命题 $p$：关于 $x$ 的方程 $x^2-2x+a=0$ 无实根；命题 $q$：对任意 $x\in\R$，
不等式 $ax^2-ax+1>0$ 恒成立.

（1）若 $q$ 为真命题，求 $a$ 的取值范围；

（2）若 $p$、$q$ 中有且只有一个为真命题，求 $a$ 的取值范围.

\ifanswer
\solbegin
\stepm{(1)}{当 $q$ 为真命题时，}
\stepm{①}{当 $a=0$ 时，不等式化为 $1>0$，对任意 $x\in\R$ 恒成立，满足条件；}
\stepm{②}{当 $a\neq0$ 时，$ax^2-ax+1>0$ 为一元二次不等式，}
\step{若对任意 $x\in\R$ 恒成立，需满足 $\begin{cases}a>0\\ \Delta=(-a)^2-4a<0\end{cases}$，解得 $0<a<4$；}
\step{综上，$a$ 的取值范围为 $[0,4)$.}
\stepm{(2)}{$p$ 为真时，方程 $x^2-2x+a=0$ 无实根，}
\step{故判别式 $\Delta=4-4a<0$，解得 $a>1$，因此 $p$ 为假时，$a\leqslant1$；}
\step{由（1）得 $q$ 为假时 $a\geqslant4$ 或 $a<0$，}
\step{$p$、$q$ 有且只有一个真命题，分两类：}
\step{当 $p$ 假 $q$ 真时，满足 $\begin{cases}a\leqslant1\\ 0\leqslant a<4\end{cases}$，解得 $0\leqslant a\leqslant1$；}
\step{当 $p$ 真 $q$ 假时，满足 $\begin{cases}a>1\\ a\geqslant4\text{或}a<0\end{cases}$，解得 $a\geqslant4$.}
\step{综上，$a$ 的取值范围是 $[0,1]\cup[4,+\infty)$.}
\solend
\fi

\ansspace[8]

\item 某工厂生产 $x$($x>0$) 件产品的总成本 $C$（单位：元）与生产件数 $x$ 之间满足二次函数
$C(x)=\dfrac14x^2+50x+10000$.

（1）若要求生产总成本不超过 30000 元，求满足生产总成本要求的产品件数的取值范围；

（2）试求该工厂生产多少件产品时，可使每件产品的平均成本最小.

\ifanswer
\solbegin
\stepm{(1)}{若要求生产总成本不超过 30000 元，}
\step{则 $C(x)=\dfrac14x^2+50x+10000\leqslant30000$，}
\step{即 $x^2+200x-80000\leqslant0$，解得 $-400\leqslant x\leqslant200$，又 $x>0$，}
\step{所以满足生产总成本要求的产品件数的取值范围为 $(0,200]$，$x\in\N$；}
\stepm{(2)}{生产件数 $x$ 时每件的平均成本为 $\dfrac{\frac14x^2+50x+10000}x=\dfrac14x+50+\dfrac{10000}x$，}
\step{由基本不等式得 $\dfrac14x+50+\dfrac{10000}x\geqslant2\sqrt{\dfrac14x\cdot\dfrac{10000}x}+50=150$，}
\step{当且仅当 $\dfrac14x=\dfrac{10000}x$，即 $x=200$ 时取等号，}
\step{即该工厂生产 200 件产品时，每件产品的平均成本最小.}
\solend
\fi

\ansspace[8]

% 原卷题干“$(n\geqslant2,n\in N)$”后连印两个全角逗号“，，”，照录未改（见文件头“原卷笔误”①）
\item 已知有限集 $A=\{a_1,a_2,\cdots,a_n\}$($n\geqslant2$,$n\in\N$)，，如果 $A$ 中元素
$a_i$($i=1,2,3,\cdots,n$) 满足 $a_1+a_2+a_3+\cdots+a_n=a_1\times a_2\times\cdots\times a_n$，
就称 $A$ 为“和积平衡集”.

% 原卷 (2) 将“写方程”与“反证法证明”两项任务并作一句（“……表示），可尝试用反证法证明：……”），
% 照录未改（见文件头“原卷笔误”②）
（1）判断集合 $\{-1+\sqrt3,-1-\sqrt3\}$ 是否为“和积平衡集”；

（2）若 $a_1$、$a_2$ 是两个不同的正数，且 $\{a_1,a_2\}$ 是“和积平衡集”，写出以 $a_1$、$a_2$ 为根的
一个一元二次方程（系数可用 $a_1$、$a_2$ 表示），可尝试用反证法证明：$a_1$、$a_2$ 至少有一个大于 2；

（3）若 $a_i\geqslant1$，且 $a_i\in\N$，求所有符合条件的“和积平衡集”.

\ifanswer
\solbegin
\stepm{(1)}{对于集合 $\{-1+\sqrt3,-1-\sqrt3\}$，}
\step{因为 $(-1+\sqrt3)+(-1-\sqrt3)=-2$，$(-1+\sqrt3)\times(-1-\sqrt3)=-2$，}
\step{所以满足 $a_1+a_2=a_1\times a_2$，满足题意，}
\step{所以集合 $\{-1+\sqrt3,-1-\sqrt3\}$ 是和积平衡集.}
\stepm{(2)}{以 $a_1$、$a_2$ 为根的一个一元二次方程可为 $x^2-(a_1+a_2)x+a_1\cdot a_2=0$，}
\step{因为 $a_1$、$a_2$ 是两个不同的正数，且 $\{a_1,a_2\}$ 是和积平衡集，}
\step{所以 $a_1+a_2=a_1\cdot a_2$，}
\step{又 $a_1$ 和 $a_2$ 为 $x^2-(a_1+a_2)x+a_1\cdot a_2=0$ 的两个不同的正根，}
\step{所以 $\Delta=(a_1+a_2)^2-4a_1a_2=(a_1+a_2)^2-4(a_1+a_2)>0$，}
\step{又 $a_1+a_2>0$，则 $a_1+a_2>4$，}
\step{若 $a_1$、$a_2$ 都小于等于 2，则 $a_1+a_2<4$，矛盾，}
\step{所以 $a_1$、$a_2$ 至少有一个大于 2；}
\stepm{(3)}{设 $A$ 中的 $a_1<a_2<a_3<\cdots<a_n$，且 $a_i\in\N^*$，}
\step{由 $a_1+a_2+a_3+\cdots+a_n=a_1\times a_2\times\cdots\times a_n$，得 $n=2$ 时明显不满足，}
\step{所以 $n\geqslant3$，故 $a_1\times a_2\times\cdots\times a_n=a_1+a_2+a_3+\cdots+a_n\leqslant na_n$，}
\step{所以 $a_1\times a_2\times\cdots\times a_{n-1}\leqslant n$，}
\step{当 $n=3$ 时，$a_1\times a_2\leqslant3$，$a_i\in\N^*$，只能有 $a_1=1$，$a_2=2$，}
\step{由 $a_1+a_2+a_3=a_1\times a_2\times a_3$，得 $a_3=3$，此时的和积平衡集为 $\{1,2,3\}$，}
\step{当 $n\geqslant4$ 时，$a_1\times a_2\times\cdots\times a_{n-1}\geqslant(n-1)(n-2)\times\cdots\times2\times1\geqslant(n-1)(n-2)$，}
\step{又 $(n-1)(n-2)-n=n^2-4n+2=n(n-4)+2>0$，即 $(n-1)(n-2)>n$，}
\step{即 $a_1\times a_2\times\cdots\times a_{n-1}>n$，与 $a_1\times a_2\times\cdots\times a_{n-1}\leqslant n$ 矛盾，}
\step{所以 $n\geqslant4$ 不满足条件，}
\step{所以符合条件的和积平衡集有且只有 $n=3$，此时的和积平衡集为 $\{1,2,3\}$.}
\solend
\fi

\ansspace*[10]

\end{enumerate}

\end{document}
"
%
% 原始材料：微信公众号“魔都数学加油站”文章，作者破晓，连载“27学年高一63”；
% 原文标题“27学年高一63——2026-2027年上海市吴淞中学高一上10月月考”；
% 原文链接：https://mp.weixin.qq.com/s/TFQ4TLnt_l66t7t51es4OQ。
% 页面标注发布时间 2026-10-16 21:44（上海）。
% 正文为 10 张 PNG 页（`原图/高一63_吴淞中学_10月月考/01..10.png`）：
% 第 06、10 页 2560×3620，第 01、02、03、04、05、07、08、09 页 4762×6735
% （**同一篇各页分辨率不同**），均按页序读取。10 页全为试卷正文页（卷首标题在第 1 页页顶，
% 卷末解答题第 21 题解析止于第 10 页），无封面页、目录页、二维码或宣传页。
% 粉色斜排“魔都数学加油站”为卷面水印，与公众号页脚一并未录入。
%
% 卷首标题：照卷面印作“2026--2027 年吴淞中学高一上 10 月月考”（公众号标题中的“上海市”
% 卷面标题行没有，本稿照卷面），年份区间按全稿惯例排 en dash。原卷只有一行标题、无副标题。
%
% 篇幅与题号：原卷自第 3 题起编号（第 1、2 题不在本卷内）。填空题 3—12（10 题）、
% 选择题 13—16（4 题）、解答题 17—21（5 题），共 19 题。本稿沿用原节与原题号，
% 只改排为全稿统一的大节样式（原卷小节标题为“一、填空题”“二、选择题”“三、解答题”）。
%
% 讲评版：原卷每题下直接印【答案】或【解析】。**只印【答案】的 2 题**（第 4、7 题），
% 本稿只给 \ans{}；其余 17 题里 16 题都印【解析】（按全稿口径这类题不另提 \ans{}——第 13、15、16 题
% 的结论写在解析末句（“故选 B”等），亦照录在解析内；第 14 题则既无【答案】也无【解析】段。**另有 1 题答案印在题干括号内**——
% 第 14 题（“(  D  )”，且该题无【答案】/【解析】段），本稿按 高一37 第 13 题成例把括号留空、
% 答案另提为 \ans{$D$}。第 16 题题干括号为空但印有【解析】，不算括号内答案。
%
% 副标题：原卷未印副标题。本稿按“知识模块”口径标作“集合与常用逻辑用语、不等式”：
% 19 题中集合与常用逻辑用语 12 题（第 3、4、6、7、9、10、11、12、13、15、18、21 题），
% 不等式 7 题（含等式（韦达定理、恒等式）与不等式性质、基本不等式、二次不等式、恒成立：
% 第 5、8、14、16、17、19、20 题），两类均超过三分之一，故并标。口径同 高一59。
%
% 与原卷的排版差异：
%   ① 原卷每题下直接印【答案】/【解析】，本稿按全稿惯例将答案与解析置于答案版，试卷版保持干净；
%   ② 原卷未印副标题，本稿按“知识模块”口径补出；
%   ③ 原卷填空横线约四个字宽，本稿按全稿惯例统一排 \blankline[4.4em]；
%   ④ 原卷实数集排斜体 $R$（第 9、10、19 题题干与解析）或粗体 $\mathbf{R}$（第 18 题题干），
%      整数集排斜体 $Z$（第 12 题解析），自然数集排斜体 $N$（第 20 题解析、第 21 题题干
%      “$(n\geqslant2,n\in N)$”）或粗体 $\mathbf{N}$（第 21 题 (3) 题干）、斜体 $N^*$
%      （第 21 题 (3) 解析），本稿按全稿惯例统一改用 $\R$、$\Z$、$\N$、$\N^*$
%      （同 高一14 排版差异④：记号归一，粗体／斜体不作区分，$N$ 与 $N^*$ 的区别照原卷保留）；
%      空集记号原卷作 $\varnothing$，照录；
%   ⑤ 原卷在数学式内多处用半角逗号（如 $\{y\mid y=[2x]+[3x],-1\leqslant x\leqslant2\}$、
%      $m=-1,0,1$、$A=\{x\mid-2<x<15\}$,$B=\cdots$），均照录未强改；
%      句末句点按全稿惯例排西文句点；“和积平衡集”自定义名词照录；
%   ⑥ 第 12 题解析里的“法一：”“法二：”标号原卷排**红色**（全卷共 2 处），
%      同 高一53“情况1：”成例，本稿排 \stepm{法一：}{…}、\stepm{法二：}{…}
%      （加粗着色、不改黑色基调），冒号照原卷排全角；
%   ⑦ 第 17—21 题的子问原卷各自另起一行，本稿按全稿惯例用 enumerate 的 (1)(2)(3)；
%      解析中的子问标号一律用 \stepm{(1)}{…}；第 19 题解析中的①②分情形，按全稿惯例作
%      \stepm{①}{…}、\stepm{②}{…}；
%   ⑧ 第 14 题印在题干括号内的答案移入 \ans{}，见上；
%   ⑨ 每道解答题原卷未留作答空白，本稿按全稿惯例加 \ansspace；
%   ⑩ 选择题的作答括号原卷排半角“(　　)”（第 13、15、16 题括号内空白、第 14 题括号内印 D），
%      本稿按全稿惯例排 \mbox{（\quad）}。
%
% 原卷笔误/存疑（均照录未改，正文对应处另有 % 注）：
%   ① 第 21 题题干“$(n\geqslant2$,$n\in N)$”后连印两个全角逗号“，，”（衍字，
%      8 倍放大确认，非水印干扰），本稿照录；
%   ② 第 21 题 (2) 题干将“写出以 $a_1$、$a_2$ 为根的一个一元二次方程”与“用反证法证明：
%      $a_1$、$a_2$ 至少有一个大于 2”两项任务并作一句、之间仅以逗号相接（“……表示），
%      可尝试用反证法证明：……”），读来像抄并两小问，原卷如此，照录未改。
%
% 插图：无（全卷 10 页无任何图形、无“如图”字样）。
%
\documentclass[a4paper,11pt]{article}
\usepackage{common}

\begin{document}

\examtitlest[3]{2026--2027 年吴淞中学高一上 10 月月考}{集合与常用逻辑用语、不等式}

\bigsec{一、填空题}

\begin{enumerate}
\setcounter{enumi}{2}

\item 满足 $A\subseteq\{1,2,3,4\}$ 的集合 $A$ 共有\blankline[4.4em]个.

\ifanswer
\solbegin
\step{由题意，集合 $A$ 是集合 $\{1,2,3,4\}$ 的子集，共有 $2^4=16$ 个.}
\solend
\fi

\item 陈述句“$a>1$ 或 $b>3$”的否定形式为\blankline[4.4em].

\ifanswer
\ans{$a\leqslant1$ 且 $b\leqslant3$}
\fi

\item 已知方程 $3x^2+4x-4=0$ 的两个根为 $x_1$、$x_2$，则 $x_1^2+x_2^2=$\blankline[4.4em].

\ifanswer
\solbegin
\step{对于方程 $3x^2+4x-4=0$，因为 $\Delta=16-4\times3\times(-4)=64>0$，}
\step{所以方程 $3x^2+4x-4=0$ 的两个不相等的实根为 $x_1$、$x_2$，}
\step{由韦达定理可得 $x_1+x_2=-\dfrac43$，$x_1x_2=-\dfrac43$，}
\step{所以 $x_1^2+x_2^2=(x_1+x_2)^2-2x_1x_2=\dfrac{16}9+\dfrac83=\dfrac{40}9$.}
\solend
\fi

\item 设集合 $A=\{y\mid y=x^2+1\}$,$B=\{y\mid y=-x^2+2x+3\}$，则 $A\cap B=$\blankline[4.4em].

\ifanswer
\solbegin
\step{因为 $y=x^2+1\geqslant1$,$y=-x^2+2x+3=-(x-1)^2+4\leqslant4$，}
\step{所以 $A=\{y\mid y\geqslant1\}$,$B=\{y\mid y\leqslant4\}$，得 $A\cap B=\{y\mid1\leqslant y\leqslant4\}$.}
\solend
\fi

\item 若 $x\in\{1,x^2\}$，则 $x=$\blankline[4.4em].

\ifanswer
\ans{$0$}
\fi

\item 已知等式 $2x^2+3x+5=a(2x+1)(x+1)+c$，则 $a+c=$\blankline[4.4em].

\ifanswer
\solbegin
\step{因为 $2x^2+3x+5=a(2x+1)(x+1)+c$ 恒成立，}
\step{所以 $2x^2+3x+5=2ax^2+3ax+a+c$ 恒成立，则 $\begin{cases}2=2a\\ 3=3a\\ 5=a+c\end{cases}$，解得 $a+c=5$.}
\solend
\fi

\item 已知 $a\in\R$，设 $p:x<2$，$q:x<5a-3$，且 $p$ 是 $q$ 的必要非充分条件，则 $a$ 的取值范围是\blankline[4.4em].

\ifanswer
\solbegin
\step{设 $p:x<2$，$q:x<5a-3$，且 $p$ 是 $q$ 的必要非充分条件，}
\step{则 $\{x\mid x<5a-3\}\subset\{x\mid x<2\}$，所以 $5a-3<2$，即 $a<1$.}
\solend
\fi

\item 已知 $a\in\R$，集合 $A=\{x\mid ax^2+3x-1=0\}$，若集合 $A$ 中只有一个元素，则 $a$ 的取值范围是\blankline[4.4em].

\ifanswer
\solbegin
\step{当 $a=0$ 时，则有 $\{x\mid ax^2+3x-1=0\}=\{x\mid3x-1=0\}=\left\{\dfrac13\right\}$，符合题意；}
\step{当 $a\neq0$ 时，由题意得 $\Delta=9+4a=0$，解得 $a=-\dfrac94$. 符合题意；}
\step{综上所述，实数 $a$ 的取值集合为 $\left\{0,-\dfrac94\right\}$.}
\solend
\fi

\item 设集合 $A=\{x\mid x^2+2(m+1)x+m^2-1=0,m\in\R\}$，且 $A\subseteq\{0,-4\}$，则实数 $m$ 的取值范围是\blankline[4.4em].

\ifanswer
\solbegin
\step{由题意得 $A=\varnothing$ 或 $A=\{0\}$ 或 $A=\{-4\}$ 或 $A=\{0,-4\}$，}
\step{当 $A=\varnothing$ 时，方程 $x^2+2(m+1)x+m^2-1=0$ 无实根，}
\step{则 $\Delta=[2(m+1)]^2-4\times1\times(m^2-1)=8m+8<0$，解得 $m<-1$，}
\step{当 $A=\{0\}$ 时，方程 $x^2+2(m+1)x+m^2-1=0$ 有一个实根 $0$，}
\step{则 $\begin{cases}\Delta=8m+8=0\\ m^2-1=0\end{cases}$，解得 $m=-1$，}
\step{当 $A=\{-4\}$ 时，方程 $x^2+2(m+1)x+m^2-1=0$ 有一个实根 $-4$，}
\step{则 $\begin{cases}\Delta=8m+8=0\\ 16-8(m+1)+m^2-1=0\end{cases}$，无解，}
\step{当 $A=\{0,-4\}$ 时，方程 $x^2+2(m+1)x+m^2-1=0$ 有两个实根 $-4$ 和 $0$，}
\step{则 $\begin{cases}\Delta=8m+8>0\\ 16-8(m+1)+m^2-1=0\\ m^2-1=0\end{cases}$，解得 $m=1$，}
\step{综上所述，实数 $m$ 的取值范围是 $\{m\mid m\leqslant-1\text{ 或 }m=1\}$.}
\solend
\fi

\item 对于任意实数 $x$，$[x]$ 表示不超过 $x$ 的最大整数，如 $[1.1]=1$，$[-1.2]=-2$，若
$A=\{y\mid y=[2x]+[3x],-1\leqslant x\leqslant2\}$，则 $A$ 中所有元素的和为\blankline[4.4em].

\ifanswer
\solbegin
% 原卷“法一：”“法二：”标号排红色，本稿按 高一53“情况1：”成例排 \stepm（加粗着色），见文件头
\stepm{法一：}{当 $-1\leqslant x<-\dfrac23$ 时，$[2x]=-2$，$[3x]=-3$，$f(x)=[2x]+[3x]=-5$，}
\step{当 $-\dfrac23\leqslant x<-\dfrac12$ 时，$[2x]=-2$，$[3x]=-2$，$f(x)=[2x]+[3x]=-4$；}
\step{当 $-\dfrac12\leqslant x<-\dfrac13$ 时，$[2x]=-1$，$[3x]=-2$，$f(x)=[2x]+[3x]=-3$；}
\step{当 $-\dfrac13\leqslant x<0$ 时，$[2x]=-1$，$[3x]=-1$，$f(x)=[2x]+[3x]=-2$；}
\step{当 $0\leqslant x<\dfrac13$ 时，$[2x]=0$，$[3x]=0$，$f(x)=[2x]+[3x]=0$；}
\step{当 $\dfrac13\leqslant x<\dfrac12$ 时，$[2x]=0$，$[3x]=1$，$f(x)=[2x]+[3x]=1$；}
\step{当 $\dfrac12\leqslant x<\dfrac23$ 时，$[2x]=1$，$[3x]=1$，$f(x)=[2x]+[3x]=2$；}
\step{当 $\dfrac23\leqslant x<1$ 时，$[2x]=1$，$[3x]=2$，$f(x)=[2x]+[3x]=3$；}
\step{当 $1\leqslant x<\dfrac43$ 时，$[2x]=2$，$[3x]=3$，$f(x)=[2x]+[3x]=5$；}
\step{当 $\dfrac43\leqslant x<\dfrac32$ 时，$[2x]=2$，$[3x]=4$，$f(x)=[2x]+[3x]=6$；}
\step{当 $\dfrac32\leqslant x<\dfrac53$ 时，$[2x]=3$，$[3x]=4$，$f(x)=[2x]+[3x]=7$；}
\step{当 $\dfrac53\leqslant x<2$ 时，$[2x]=3$，$[3x]=5$，$f(x)=[2x]+[3x]=8$；}
\step{当 $x=2$ 时，$[2x]=4$，$[3x]=6$，$f(x)=[2x]+[3x]=10$}
\step{故 $A$ 中所有元素的和为 $-5-4-3-2+0+1+2+3+5+6+7+8+10=28$.}
% 原卷“法二：”标号排红色，同上
\stepm{法二：}{设 $x=m+n$，其中 $m\in\Z$，$n\in[0,1)$，}
\step{若 $0\leqslant n<\dfrac13$，则 $0\leqslant2n<\dfrac23$，$0\leqslant3n<1$，}
\step{得 $[2x]=[2m+2n]=2m$，$[3x]=[3m+3n]=3m$，所以 $y=[2x]+[3x]=5m$；}
\step{若 $\dfrac13\leqslant n<\dfrac12$，则 $\dfrac23\leqslant2n<1$，$1\leqslant3n<\dfrac32$，}
\step{得 $[2x]=[2m+2n]=2m$，$[3x]=[3m+3n]=3m+1$，}
\step{所以 $y=[2x]+[3x]=5m+1$；}
\step{若 $\dfrac12\leqslant n<\dfrac23$，则 $1\leqslant2n<\dfrac43$，$\dfrac32\leqslant3n<2$，}
\step{得 $[2x]=[2m+2n]=2m+1$，$[3x]=[3m+3n]=3m+1$，}
\step{所以 $y=[2x]+[3x]=5m+2$；}
\step{若 $\dfrac23\leqslant n<1$，则 $\dfrac43\leqslant2n<2$，$2\leqslant3n<3$，}
\step{得 $[2x]=[2m+2n]=2m+1$，$[3x]=[3m+3n]=3m+2$，}
\step{所以 $y=[2x]+[3x]=5m+3$；}
\step{综上所述，当 $n\in[0,1)$ 时，所有函数值的和}
\step{为 $5m+5m+1+5m+2+5m+3=20m+6$，}
\step{又因为 $x\in[-1,2]$，此时 $m=-1,0,1$，且当 $x=2$ 时，$y=[2x]+[3x]=10$，}
\step{所有函数值的和为 $(-14)+6+26+10=28$，即 $A$ 中所有元素之和为 $28$.}
\solend
\fi

\end{enumerate}

\bigsec{二、选择题}

\begin{enumerate}
\setcounter{enumi}{12}

\item 下列写法中，正确的是\mbox{（\quad）}

\keepwithprev
A. $0\in\{\varnothing\}$\qquad B. $\varnothing\in\{\varnothing\}$

C. $\varnothing=\{\varnothing\}$\qquad D. $0\in\varnothing$

\ifanswer
\solbegin
\step{$0\notin\{\varnothing\}$，$\varnothing\in\{\varnothing\}$，故 $A$ 错误，$B$ 正确；}
\step{$\varnothing$ 是不含任何元素的集合，所以 $\varnothing\neq\{\varnothing\}$，$0\notin\varnothing$，故 $CD$ 错误.}
\step{故选 $B$.}
\solend
\fi

% 第 14 题答案“D”直接印在题干括号内“(  D  )”，且该题无【答案】/【解析】段；
% 按 高一37 第 13 题的成例括号留空、答案另提为 \ans{}（见文件头）
\item 若 $a<b<0$，则下列不等式中不能成立的是\mbox{（\quad）}

\keepwithprev
A. $\dfrac1a>\dfrac1b$\qquad B. $a-b>a$

C. $|a|>|b|$\qquad D. $a^3>b^3$

\ifanswer
\ans{$D$}
\fi

\item 已知 $x$、$y$、$z$ 为非零实数，代数式 $\dfrac x{|x|}+\dfrac y{|y|}+\dfrac z{|z|}+\dfrac{xyz}{|xyz|}$
的值所组成的集合是 $M$，则下列判断正确的是\mbox{（\quad）}

\keepwithprev
A. $-4\notin M$\qquad B. $2\in M$

C. $4\notin M$\qquad D. $0\in M$

\ifanswer
\solbegin
\step{当 $x$、$y$、$z$ 均为正数时，$xyz>0$，$\dfrac x{|x|}+\dfrac y{|y|}+\dfrac z{|z|}+\dfrac{xyz}{|xyz|}=1+1+1+1=4$，}
\step{当 $x$、$y$、$z$ 有 2 个正数，1 个负数时，不妨设 $x>0$，$y>0$，$z<0$，}
\step{则 $xyz<0$，$\dfrac x{|x|}+\dfrac y{|y|}+\dfrac z{|z|}+\dfrac{xyz}{|xyz|}=1+1-1-1=0$，}
\step{若 $x$、$y$、$z$ 有 1 个正数，2 个负数时，不妨设 $x>0$，$y<0$，$z<0$，}
\step{则 $xyz>0$，$\dfrac x{|x|}+\dfrac y{|y|}+\dfrac z{|z|}+\dfrac{xyz}{|xyz|}=1-1-1+1=0$，}
\step{当 $x$、$y$、$z$ 均为负数时，$xyz<0$，}
\step{$\dfrac x{|x|}+\dfrac y{|y|}+\dfrac z{|z|}+\dfrac{xyz}{|xyz|}=-1-1-1-1=-4$，}
\step{综上，$M=\{-4,0,4\}$，}
\step{所以 $-4\in M$，$A$ 错误；$2\notin M$，$B$ 错误；$4\in M$，$C$ 错误；}
\step{$0\in M$，$D$ 正确. 故选 $D$.}
\solend
\fi

\item 若 $xy>0$，且 $x+y=9$，则 $\dfrac1x+\dfrac4y$ 的最小值是\mbox{（\quad）}

\keepwithprev
A. $9$\qquad B. $\dfrac59$

C. $\dfrac49$\qquad D. $1$

\ifanswer
\solbegin
\step{$xy>0$，且 $x+y=9$，所以 $x>0$，$y>0$，}
\step{所以 $\dfrac1x+\dfrac4y=\dfrac19\left(\dfrac1x+\dfrac4y\right)(x+y)=\dfrac59+\dfrac19\left(\dfrac yx+\dfrac{4x}y\right)\geqslant\dfrac59+\dfrac29\cdot\sqrt{\dfrac yx\cdot\dfrac{4x}y}=\dfrac59+\dfrac49=1$，}
\step{当且仅当 $\dfrac yx=\dfrac{4x}y$，即 $x=3$，$y=6$ 时取等号，所以 $\dfrac1x+\dfrac4y$ 的最小值为 $1$.}
\step{故选 $D$.}
\solend
\fi

\end{enumerate}

\bigsec{三、解答题}

\begin{enumerate}
\setcounter{enumi}{16}

\item 求解下列不等式：

（1）$\dfrac{x-3}{x+1}\leqslant0$；

（2）$|2x-1|<3$.

\ifanswer
\solbegin
\stepm{(1)}{$\dfrac{x-3}{x+1}\leqslant0$，变形得 $(x-3)\cdot(x+1)\leqslant0$ 且 $x\neq-1$，}
\step{解得 $-1<x\leqslant3$，即不等式的解集为 $(-1,3]$.}
\stepm{(2)}{若 $|2x-1|<3$，则 $-3<2x-1<3$，解得 $-1<x<2$，即 $(-1,2)$.}
\solend
\fi

\ansspace[6]

\item 已知集合 $A=\{x\mid-2<x<15\}$,$B=\{x\mid m-6<x<2m-1,m\in\R\}$.

（1）当 $m=2$ 时，求 $A\cap B$；

（2）若 $A\cup B=A$，求实数 $m$ 的取值范围.

\ifanswer
\solbegin
\stepm{(1)}{当 $m=2$ 时，$A=\{x\mid-2<x<15\}$,$B=\{x\mid-4<x<3\}$，}
\step{此时 $A\cap B=\{x\mid-2<x<3\}$.}
\stepm{(2)}{由 $A\cup B=A$ 得 $B\subseteq A$.}
\step{当 $B=\varnothing$,$m-6\geqslant2m-1$，解得 $m\leqslant-5$；}
\step{当 $B\neq\varnothing$，结合 $B\subseteq A$ 得 $\begin{cases}m-6<2m-1\\ 2m-1\leqslant15\\ m-6\geqslant-2\end{cases}$，解得 $4\leqslant m\leqslant8$.}
\step{综上，$m\leqslant-5$ 或 $4\leqslant m\leqslant8$.}
\solend
\fi

\ansspace[6]

\item 已知 $a\in\R$，命题 $p$：关于 $x$ 的方程 $x^2-2x+a=0$ 无实根；命题 $q$：对任意 $x\in\R$，
不等式 $ax^2-ax+1>0$ 恒成立.

（1）若 $q$ 为真命题，求 $a$ 的取值范围；

（2）若 $p$、$q$ 中有且只有一个为真命题，求 $a$ 的取值范围.

\ifanswer
\solbegin
\stepm{(1)}{当 $q$ 为真命题时，}
\stepm{①}{当 $a=0$ 时，不等式化为 $1>0$，对任意 $x\in\R$ 恒成立，满足条件；}
\stepm{②}{当 $a\neq0$ 时，$ax^2-ax+1>0$ 为一元二次不等式，}
\step{若对任意 $x\in\R$ 恒成立，需满足 $\begin{cases}a>0\\ \Delta=(-a)^2-4a<0\end{cases}$，解得 $0<a<4$；}
\step{综上，$a$ 的取值范围为 $[0,4)$.}
\stepm{(2)}{$p$ 为真时，方程 $x^2-2x+a=0$ 无实根，}
\step{故判别式 $\Delta=4-4a<0$，解得 $a>1$，因此 $p$ 为假时，$a\leqslant1$；}
\step{由（1）得 $q$ 为假时 $a\geqslant4$ 或 $a<0$，}
\step{$p$、$q$ 有且只有一个真命题，分两类：}
\step{当 $p$ 假 $q$ 真时，满足 $\begin{cases}a\leqslant1\\ 0\leqslant a<4\end{cases}$，解得 $0\leqslant a\leqslant1$；}
\step{当 $p$ 真 $q$ 假时，满足 $\begin{cases}a>1\\ a\geqslant4\text{或}a<0\end{cases}$，解得 $a\geqslant4$.}
\step{综上，$a$ 的取值范围是 $[0,1]\cup[4,+\infty)$.}
\solend
\fi

\ansspace[8]

\item 某工厂生产 $x$($x>0$) 件产品的总成本 $C$（单位：元）与生产件数 $x$ 之间满足二次函数
$C(x)=\dfrac14x^2+50x+10000$.

（1）若要求生产总成本不超过 30000 元，求满足生产总成本要求的产品件数的取值范围；

（2）试求该工厂生产多少件产品时，可使每件产品的平均成本最小.

\ifanswer
\solbegin
\stepm{(1)}{若要求生产总成本不超过 30000 元，}
\step{则 $C(x)=\dfrac14x^2+50x+10000\leqslant30000$，}
\step{即 $x^2+200x-80000\leqslant0$，解得 $-400\leqslant x\leqslant200$，又 $x>0$，}
\step{所以满足生产总成本要求的产品件数的取值范围为 $(0,200]$，$x\in\N$；}
\stepm{(2)}{生产件数 $x$ 时每件的平均成本为 $\dfrac{\frac14x^2+50x+10000}x=\dfrac14x+50+\dfrac{10000}x$，}
\step{由基本不等式得 $\dfrac14x+50+\dfrac{10000}x\geqslant2\sqrt{\dfrac14x\cdot\dfrac{10000}x}+50=150$，}
\step{当且仅当 $\dfrac14x=\dfrac{10000}x$，即 $x=200$ 时取等号，}
\step{即该工厂生产 200 件产品时，每件产品的平均成本最小.}
\solend
\fi

\ansspace[8]

% 原卷题干“$(n\geqslant2,n\in N)$”后连印两个全角逗号“，，”，照录未改（见文件头“原卷笔误”①）
\item 已知有限集 $A=\{a_1,a_2,\cdots,a_n\}$($n\geqslant2$,$n\in\N$)，，如果 $A$ 中元素
$a_i$($i=1,2,3,\cdots,n$) 满足 $a_1+a_2+a_3+\cdots+a_n=a_1\times a_2\times\cdots\times a_n$，
就称 $A$ 为“和积平衡集”.

% 原卷 (2) 将“写方程”与“反证法证明”两项任务并作一句（“……表示），可尝试用反证法证明：……”），
% 照录未改（见文件头“原卷笔误”②）
（1）判断集合 $\{-1+\sqrt3,-1-\sqrt3\}$ 是否为“和积平衡集”；

（2）若 $a_1$、$a_2$ 是两个不同的正数，且 $\{a_1,a_2\}$ 是“和积平衡集”，写出以 $a_1$、$a_2$ 为根的
一个一元二次方程（系数可用 $a_1$、$a_2$ 表示），可尝试用反证法证明：$a_1$、$a_2$ 至少有一个大于 2；

（3）若 $a_i\geqslant1$，且 $a_i\in\N$，求所有符合条件的“和积平衡集”.

\ifanswer
\solbegin
\stepm{(1)}{对于集合 $\{-1+\sqrt3,-1-\sqrt3\}$，}
\step{因为 $(-1+\sqrt3)+(-1-\sqrt3)=-2$，$(-1+\sqrt3)\times(-1-\sqrt3)=-2$，}
\step{所以满足 $a_1+a_2=a_1\times a_2$，满足题意，}
\step{所以集合 $\{-1+\sqrt3,-1-\sqrt3\}$ 是和积平衡集.}
\stepm{(2)}{以 $a_1$、$a_2$ 为根的一个一元二次方程可为 $x^2-(a_1+a_2)x+a_1\cdot a_2=0$，}
\step{因为 $a_1$、$a_2$ 是两个不同的正数，且 $\{a_1,a_2\}$ 是和积平衡集，}
\step{所以 $a_1+a_2=a_1\cdot a_2$，}
\step{又 $a_1$ 和 $a_2$ 为 $x^2-(a_1+a_2)x+a_1\cdot a_2=0$ 的两个不同的正根，}
\step{所以 $\Delta=(a_1+a_2)^2-4a_1a_2=(a_1+a_2)^2-4(a_1+a_2)>0$，}
\step{又 $a_1+a_2>0$，则 $a_1+a_2>4$，}
\step{若 $a_1$、$a_2$ 都小于等于 2，则 $a_1+a_2<4$，矛盾，}
\step{所以 $a_1$、$a_2$ 至少有一个大于 2；}
\stepm{(3)}{设 $A$ 中的 $a_1<a_2<a_3<\cdots<a_n$，且 $a_i\in\N^*$，}
\step{由 $a_1+a_2+a_3+\cdots+a_n=a_1\times a_2\times\cdots\times a_n$，得 $n=2$ 时明显不满足，}
\step{所以 $n\geqslant3$，故 $a_1\times a_2\times\cdots\times a_n=a_1+a_2+a_3+\cdots+a_n\leqslant na_n$，}
\step{所以 $a_1\times a_2\times\cdots\times a_{n-1}\leqslant n$，}
\step{当 $n=3$ 时，$a_1\times a_2\leqslant3$，$a_i\in\N^*$，只能有 $a_1=1$，$a_2=2$，}
\step{由 $a_1+a_2+a_3=a_1\times a_2\times a_3$，得 $a_3=3$，此时的和积平衡集为 $\{1,2,3\}$，}
\step{当 $n\geqslant4$ 时，$a_1\times a_2\times\cdots\times a_{n-1}\geqslant(n-1)(n-2)\times\cdots\times2\times1\geqslant(n-1)(n-2)$，}
\step{又 $(n-1)(n-2)-n=n^2-4n+2=n(n-4)+2>0$，即 $(n-1)(n-2)>n$，}
\step{即 $a_1\times a_2\times\cdots\times a_{n-1}>n$，与 $a_1\times a_2\times\cdots\times a_{n-1}\leqslant n$ 矛盾，}
\step{所以 $n\geqslant4$ 不满足条件，}
\step{所以符合条件的和积平衡集有且只有 $n=3$，此时的和积平衡集为 $\{1,2,3\}$.}
\solend
\fi

\ansspace*[10]

\end{enumerate}

\end{document}
"
%
% 原始材料：微信公众号“魔都数学加油站”文章，作者破晓，连载“27学年高一63”；
% 原文标题“27学年高一63——2026-2027年上海市吴淞中学高一上10月月考”；
% 原文链接：https://mp.weixin.qq.com/s/TFQ4TLnt_l66t7t51es4OQ。
% 页面标注发布时间 2026-10-16 21:44（上海）。
% 正文为 10 张 PNG 页（`原图/高一63_吴淞中学_10月月考/01..10.png`）：
% 第 06、10 页 2560×3620，第 01、02、03、04、05、07、08、09 页 4762×6735
% （**同一篇各页分辨率不同**），均按页序读取。10 页全为试卷正文页（卷首标题在第 1 页页顶，
% 卷末解答题第 21 题解析止于第 10 页），无封面页、目录页、二维码或宣传页。
% 粉色斜排“魔都数学加油站”为卷面水印，与公众号页脚一并未录入。
%
% 卷首标题：照卷面印作“2026--2027 年吴淞中学高一上 10 月月考”（公众号标题中的“上海市”
% 卷面标题行没有，本稿照卷面），年份区间按全稿惯例排 en dash。原卷只有一行标题、无副标题。
%
% 篇幅与题号：原卷自第 3 题起编号（第 1、2 题不在本卷内）。填空题 3—12（10 题）、
% 选择题 13—16（4 题）、解答题 17—21（5 题），共 19 题。本稿沿用原节与原题号，
% 只改排为全稿统一的大节样式（原卷小节标题为“一、填空题”“二、选择题”“三、解答题”）。
%
% 讲评版：原卷每题下直接印【答案】或【解析】。**只印【答案】的 2 题**（第 4、7 题），
% 本稿只给 \ans{}；其余 17 题里 16 题都印【解析】（按全稿口径这类题不另提 \ans{}——第 13、15、16 题
% 的结论写在解析末句（“故选 B”等），亦照录在解析内；第 14 题则既无【答案】也无【解析】段。**另有 1 题答案印在题干括号内**——
% 第 14 题（“(  D  )”，且该题无【答案】/【解析】段），本稿按 高一37 第 13 题成例把括号留空、
% 答案另提为 \ans{$D$}。第 16 题题干括号为空但印有【解析】，不算括号内答案。
%
% 副标题：原卷未印副标题。本稿按“知识模块”口径标作“集合与常用逻辑用语、不等式”：
% 19 题中集合与常用逻辑用语 12 题（第 3、4、6、7、9、10、11、12、13、15、18、21 题），
% 不等式 7 题（含等式（韦达定理、恒等式）与不等式性质、基本不等式、二次不等式、恒成立：
% 第 5、8、14、16、17、19、20 题），两类均超过三分之一，故并标。口径同 高一59。
%
% 与原卷的排版差异：
%   ① 原卷每题下直接印【答案】/【解析】，本稿按全稿惯例将答案与解析置于答案版，试卷版保持干净；
%   ② 原卷未印副标题，本稿按“知识模块”口径补出；
%   ③ 原卷填空横线约四个字宽，本稿按全稿惯例统一排 \blankline[4.4em]；
%   ④ 原卷实数集排斜体 $R$（第 9、10、19 题题干与解析）或粗体 $\mathbf{R}$（第 18 题题干），
%      整数集排斜体 $Z$（第 12 题解析），自然数集排斜体 $N$（第 20 题解析、第 21 题题干
%      “$(n\geqslant2,n\in N)$”）或粗体 $\mathbf{N}$（第 21 题 (3) 题干）、斜体 $N^*$
%      （第 21 题 (3) 解析），本稿按全稿惯例统一改用 $\R$、$\Z$、$\N$、$\N^*$
%      （同 高一14 排版差异④：记号归一，粗体／斜体不作区分，$N$ 与 $N^*$ 的区别照原卷保留）；
%      空集记号原卷作 $\varnothing$，照录；
%   ⑤ 原卷在数学式内多处用半角逗号（如 $\{y\mid y=[2x]+[3x],-1\leqslant x\leqslant2\}$、
%      $m=-1,0,1$、$A=\{x\mid-2<x<15\}$,$B=\cdots$），均照录未强改；
%      句末句点按全稿惯例排西文句点；“和积平衡集”自定义名词照录；
%   ⑥ 第 12 题解析里的“法一：”“法二：”标号原卷排**红色**（全卷共 2 处），
%      同 高一53“情况1：”成例，本稿排 \stepm{法一：}{…}、\stepm{法二：}{…}
%      （加粗着色、不改黑色基调），冒号照原卷排全角；
%   ⑦ 第 17—21 题的子问原卷各自另起一行，本稿按全稿惯例用 enumerate 的 (1)(2)(3)；
%      解析中的子问标号一律用 \stepm{(1)}{…}；第 19 题解析中的①②分情形，按全稿惯例作
%      \stepm{①}{…}、\stepm{②}{…}；
%   ⑧ 第 14 题印在题干括号内的答案移入 \ans{}，见上；
%   ⑨ 每道解答题原卷未留作答空白，本稿按全稿惯例加 \ansspace；
%   ⑩ 选择题的作答括号原卷排半角“(　　)”（第 13、15、16 题括号内空白、第 14 题括号内印 D），
%      本稿按全稿惯例排 \mbox{（\quad）}。
%
% 原卷笔误/存疑（均照录未改，正文对应处另有 % 注）：
%   ① 第 21 题题干“$(n\geqslant2$,$n\in N)$”后连印两个全角逗号“，，”（衍字，
%      8 倍放大确认，非水印干扰），本稿照录；
%   ② 第 21 题 (2) 题干将“写出以 $a_1$、$a_2$ 为根的一个一元二次方程”与“用反证法证明：
%      $a_1$、$a_2$ 至少有一个大于 2”两项任务并作一句、之间仅以逗号相接（“……表示），
%      可尝试用反证法证明：……”），读来像抄并两小问，原卷如此，照录未改。
%
% 插图：无（全卷 10 页无任何图形、无“如图”字样）。
%
\documentclass[a4paper,11pt]{article}
\usepackage{common}

\begin{document}

\examtitlest[3]{2026--2027 年吴淞中学高一上 10 月月考}{集合与常用逻辑用语、不等式}

\bigsec{一、填空题}

\begin{enumerate}
\setcounter{enumi}{2}

\item 满足 $A\subseteq\{1,2,3,4\}$ 的集合 $A$ 共有\blankline[4.4em]个.

\ifanswer
\solbegin
\step{由题意，集合 $A$ 是集合 $\{1,2,3,4\}$ 的子集，共有 $2^4=16$ 个.}
\solend
\fi

\item 陈述句“$a>1$ 或 $b>3$”的否定形式为\blankline[4.4em].

\ifanswer
\ans{$a\leqslant1$ 且 $b\leqslant3$}
\fi

\item 已知方程 $3x^2+4x-4=0$ 的两个根为 $x_1$、$x_2$，则 $x_1^2+x_2^2=$\blankline[4.4em].

\ifanswer
\solbegin
\step{对于方程 $3x^2+4x-4=0$，因为 $\Delta=16-4\times3\times(-4)=64>0$，}
\step{所以方程 $3x^2+4x-4=0$ 的两个不相等的实根为 $x_1$、$x_2$，}
\step{由韦达定理可得 $x_1+x_2=-\dfrac43$，$x_1x_2=-\dfrac43$，}
\step{所以 $x_1^2+x_2^2=(x_1+x_2)^2-2x_1x_2=\dfrac{16}9+\dfrac83=\dfrac{40}9$.}
\solend
\fi

\item 设集合 $A=\{y\mid y=x^2+1\}$,$B=\{y\mid y=-x^2+2x+3\}$，则 $A\cap B=$\blankline[4.4em].

\ifanswer
\solbegin
\step{因为 $y=x^2+1\geqslant1$,$y=-x^2+2x+3=-(x-1)^2+4\leqslant4$，}
\step{所以 $A=\{y\mid y\geqslant1\}$,$B=\{y\mid y\leqslant4\}$，得 $A\cap B=\{y\mid1\leqslant y\leqslant4\}$.}
\solend
\fi

\item 若 $x\in\{1,x^2\}$，则 $x=$\blankline[4.4em].

\ifanswer
\ans{$0$}
\fi

\item 已知等式 $2x^2+3x+5=a(2x+1)(x+1)+c$，则 $a+c=$\blankline[4.4em].

\ifanswer
\solbegin
\step{因为 $2x^2+3x+5=a(2x+1)(x+1)+c$ 恒成立，}
\step{所以 $2x^2+3x+5=2ax^2+3ax+a+c$ 恒成立，则 $\begin{cases}2=2a\\ 3=3a\\ 5=a+c\end{cases}$，解得 $a+c=5$.}
\solend
\fi

\item 已知 $a\in\R$，设 $p:x<2$，$q:x<5a-3$，且 $p$ 是 $q$ 的必要非充分条件，则 $a$ 的取值范围是\blankline[4.4em].

\ifanswer
\solbegin
\step{设 $p:x<2$，$q:x<5a-3$，且 $p$ 是 $q$ 的必要非充分条件，}
\step{则 $\{x\mid x<5a-3\}\subset\{x\mid x<2\}$，所以 $5a-3<2$，即 $a<1$.}
\solend
\fi

\item 已知 $a\in\R$，集合 $A=\{x\mid ax^2+3x-1=0\}$，若集合 $A$ 中只有一个元素，则 $a$ 的取值范围是\blankline[4.4em].

\ifanswer
\solbegin
\step{当 $a=0$ 时，则有 $\{x\mid ax^2+3x-1=0\}=\{x\mid3x-1=0\}=\left\{\dfrac13\right\}$，符合题意；}
\step{当 $a\neq0$ 时，由题意得 $\Delta=9+4a=0$，解得 $a=-\dfrac94$. 符合题意；}
\step{综上所述，实数 $a$ 的取值集合为 $\left\{0,-\dfrac94\right\}$.}
\solend
\fi

\item 设集合 $A=\{x\mid x^2+2(m+1)x+m^2-1=0,m\in\R\}$，且 $A\subseteq\{0,-4\}$，则实数 $m$ 的取值范围是\blankline[4.4em].

\ifanswer
\solbegin
\step{由题意得 $A=\varnothing$ 或 $A=\{0\}$ 或 $A=\{-4\}$ 或 $A=\{0,-4\}$，}
\step{当 $A=\varnothing$ 时，方程 $x^2+2(m+1)x+m^2-1=0$ 无实根，}
\step{则 $\Delta=[2(m+1)]^2-4\times1\times(m^2-1)=8m+8<0$，解得 $m<-1$，}
\step{当 $A=\{0\}$ 时，方程 $x^2+2(m+1)x+m^2-1=0$ 有一个实根 $0$，}
\step{则 $\begin{cases}\Delta=8m+8=0\\ m^2-1=0\end{cases}$，解得 $m=-1$，}
\step{当 $A=\{-4\}$ 时，方程 $x^2+2(m+1)x+m^2-1=0$ 有一个实根 $-4$，}
\step{则 $\begin{cases}\Delta=8m+8=0\\ 16-8(m+1)+m^2-1=0\end{cases}$，无解，}
\step{当 $A=\{0,-4\}$ 时，方程 $x^2+2(m+1)x+m^2-1=0$ 有两个实根 $-4$ 和 $0$，}
\step{则 $\begin{cases}\Delta=8m+8>0\\ 16-8(m+1)+m^2-1=0\\ m^2-1=0\end{cases}$，解得 $m=1$，}
\step{综上所述，实数 $m$ 的取值范围是 $\{m\mid m\leqslant-1\text{ 或 }m=1\}$.}
\solend
\fi

\item 对于任意实数 $x$，$[x]$ 表示不超过 $x$ 的最大整数，如 $[1.1]=1$，$[-1.2]=-2$，若
$A=\{y\mid y=[2x]+[3x],-1\leqslant x\leqslant2\}$，则 $A$ 中所有元素的和为\blankline[4.4em].

\ifanswer
\solbegin
% 原卷“法一：”“法二：”标号排红色，本稿按 高一53“情况1：”成例排 \stepm（加粗着色），见文件头
\stepm{法一：}{当 $-1\leqslant x<-\dfrac23$ 时，$[2x]=-2$，$[3x]=-3$，$f(x)=[2x]+[3x]=-5$，}
\step{当 $-\dfrac23\leqslant x<-\dfrac12$ 时，$[2x]=-2$，$[3x]=-2$，$f(x)=[2x]+[3x]=-4$；}
\step{当 $-\dfrac12\leqslant x<-\dfrac13$ 时，$[2x]=-1$，$[3x]=-2$，$f(x)=[2x]+[3x]=-3$；}
\step{当 $-\dfrac13\leqslant x<0$ 时，$[2x]=-1$，$[3x]=-1$，$f(x)=[2x]+[3x]=-2$；}
\step{当 $0\leqslant x<\dfrac13$ 时，$[2x]=0$，$[3x]=0$，$f(x)=[2x]+[3x]=0$；}
\step{当 $\dfrac13\leqslant x<\dfrac12$ 时，$[2x]=0$，$[3x]=1$，$f(x)=[2x]+[3x]=1$；}
\step{当 $\dfrac12\leqslant x<\dfrac23$ 时，$[2x]=1$，$[3x]=1$，$f(x)=[2x]+[3x]=2$；}
\step{当 $\dfrac23\leqslant x<1$ 时，$[2x]=1$，$[3x]=2$，$f(x)=[2x]+[3x]=3$；}
\step{当 $1\leqslant x<\dfrac43$ 时，$[2x]=2$，$[3x]=3$，$f(x)=[2x]+[3x]=5$；}
\step{当 $\dfrac43\leqslant x<\dfrac32$ 时，$[2x]=2$，$[3x]=4$，$f(x)=[2x]+[3x]=6$；}
\step{当 $\dfrac32\leqslant x<\dfrac53$ 时，$[2x]=3$，$[3x]=4$，$f(x)=[2x]+[3x]=7$；}
\step{当 $\dfrac53\leqslant x<2$ 时，$[2x]=3$，$[3x]=5$，$f(x)=[2x]+[3x]=8$；}
\step{当 $x=2$ 时，$[2x]=4$，$[3x]=6$，$f(x)=[2x]+[3x]=10$}
\step{故 $A$ 中所有元素的和为 $-5-4-3-2+0+1+2+3+5+6+7+8+10=28$.}
% 原卷“法二：”标号排红色，同上
\stepm{法二：}{设 $x=m+n$，其中 $m\in\Z$，$n\in[0,1)$，}
\step{若 $0\leqslant n<\dfrac13$，则 $0\leqslant2n<\dfrac23$，$0\leqslant3n<1$，}
\step{得 $[2x]=[2m+2n]=2m$，$[3x]=[3m+3n]=3m$，所以 $y=[2x]+[3x]=5m$；}
\step{若 $\dfrac13\leqslant n<\dfrac12$，则 $\dfrac23\leqslant2n<1$，$1\leqslant3n<\dfrac32$，}
\step{得 $[2x]=[2m+2n]=2m$，$[3x]=[3m+3n]=3m+1$，}
\step{所以 $y=[2x]+[3x]=5m+1$；}
\step{若 $\dfrac12\leqslant n<\dfrac23$，则 $1\leqslant2n<\dfrac43$，$\dfrac32\leqslant3n<2$，}
\step{得 $[2x]=[2m+2n]=2m+1$，$[3x]=[3m+3n]=3m+1$，}
\step{所以 $y=[2x]+[3x]=5m+2$；}
\step{若 $\dfrac23\leqslant n<1$，则 $\dfrac43\leqslant2n<2$，$2\leqslant3n<3$，}
\step{得 $[2x]=[2m+2n]=2m+1$，$[3x]=[3m+3n]=3m+2$，}
\step{所以 $y=[2x]+[3x]=5m+3$；}
\step{综上所述，当 $n\in[0,1)$ 时，所有函数值的和}
\step{为 $5m+5m+1+5m+2+5m+3=20m+6$，}
\step{又因为 $x\in[-1,2]$，此时 $m=-1,0,1$，且当 $x=2$ 时，$y=[2x]+[3x]=10$，}
\step{所有函数值的和为 $(-14)+6+26+10=28$，即 $A$ 中所有元素之和为 $28$.}
\solend
\fi

\end{enumerate}

\bigsec{二、选择题}

\begin{enumerate}
\setcounter{enumi}{12}

\item 下列写法中，正确的是\mbox{（\quad）}

\keepwithprev
A. $0\in\{\varnothing\}$\qquad B. $\varnothing\in\{\varnothing\}$

C. $\varnothing=\{\varnothing\}$\qquad D. $0\in\varnothing$

\ifanswer
\solbegin
\step{$0\notin\{\varnothing\}$，$\varnothing\in\{\varnothing\}$，故 $A$ 错误，$B$ 正确；}
\step{$\varnothing$ 是不含任何元素的集合，所以 $\varnothing\neq\{\varnothing\}$，$0\notin\varnothing$，故 $CD$ 错误.}
\step{故选 $B$.}
\solend
\fi

% 第 14 题答案“D”直接印在题干括号内“(  D  )”，且该题无【答案】/【解析】段；
% 按 高一37 第 13 题的成例括号留空、答案另提为 \ans{}（见文件头）
\item 若 $a<b<0$，则下列不等式中不能成立的是\mbox{（\quad）}

\keepwithprev
A. $\dfrac1a>\dfrac1b$\qquad B. $a-b>a$

C. $|a|>|b|$\qquad D. $a^3>b^3$

\ifanswer
\ans{$D$}
\fi

\item 已知 $x$、$y$、$z$ 为非零实数，代数式 $\dfrac x{|x|}+\dfrac y{|y|}+\dfrac z{|z|}+\dfrac{xyz}{|xyz|}$
的值所组成的集合是 $M$，则下列判断正确的是\mbox{（\quad）}

\keepwithprev
A. $-4\notin M$\qquad B. $2\in M$

C. $4\notin M$\qquad D. $0\in M$

\ifanswer
\solbegin
\step{当 $x$、$y$、$z$ 均为正数时，$xyz>0$，$\dfrac x{|x|}+\dfrac y{|y|}+\dfrac z{|z|}+\dfrac{xyz}{|xyz|}=1+1+1+1=4$，}
\step{当 $x$、$y$、$z$ 有 2 个正数，1 个负数时，不妨设 $x>0$，$y>0$，$z<0$，}
\step{则 $xyz<0$，$\dfrac x{|x|}+\dfrac y{|y|}+\dfrac z{|z|}+\dfrac{xyz}{|xyz|}=1+1-1-1=0$，}
\step{若 $x$、$y$、$z$ 有 1 个正数，2 个负数时，不妨设 $x>0$，$y<0$，$z<0$，}
\step{则 $xyz>0$，$\dfrac x{|x|}+\dfrac y{|y|}+\dfrac z{|z|}+\dfrac{xyz}{|xyz|}=1-1-1+1=0$，}
\step{当 $x$、$y$、$z$ 均为负数时，$xyz<0$，}
\step{$\dfrac x{|x|}+\dfrac y{|y|}+\dfrac z{|z|}+\dfrac{xyz}{|xyz|}=-1-1-1-1=-4$，}
\step{综上，$M=\{-4,0,4\}$，}
\step{所以 $-4\in M$，$A$ 错误；$2\notin M$，$B$ 错误；$4\in M$，$C$ 错误；}
\step{$0\in M$，$D$ 正确. 故选 $D$.}
\solend
\fi

\item 若 $xy>0$，且 $x+y=9$，则 $\dfrac1x+\dfrac4y$ 的最小值是\mbox{（\quad）}

\keepwithprev
A. $9$\qquad B. $\dfrac59$

C. $\dfrac49$\qquad D. $1$

\ifanswer
\solbegin
\step{$xy>0$，且 $x+y=9$，所以 $x>0$，$y>0$，}
\step{所以 $\dfrac1x+\dfrac4y=\dfrac19\left(\dfrac1x+\dfrac4y\right)(x+y)=\dfrac59+\dfrac19\left(\dfrac yx+\dfrac{4x}y\right)\geqslant\dfrac59+\dfrac29\cdot\sqrt{\dfrac yx\cdot\dfrac{4x}y}=\dfrac59+\dfrac49=1$，}
\step{当且仅当 $\dfrac yx=\dfrac{4x}y$，即 $x=3$，$y=6$ 时取等号，所以 $\dfrac1x+\dfrac4y$ 的最小值为 $1$.}
\step{故选 $D$.}
\solend
\fi

\end{enumerate}

\bigsec{三、解答题}

\begin{enumerate}
\setcounter{enumi}{16}

\item 求解下列不等式：

（1）$\dfrac{x-3}{x+1}\leqslant0$；

（2）$|2x-1|<3$.

\ifanswer
\solbegin
\stepm{(1)}{$\dfrac{x-3}{x+1}\leqslant0$，变形得 $(x-3)\cdot(x+1)\leqslant0$ 且 $x\neq-1$，}
\step{解得 $-1<x\leqslant3$，即不等式的解集为 $(-1,3]$.}
\stepm{(2)}{若 $|2x-1|<3$，则 $-3<2x-1<3$，解得 $-1<x<2$，即 $(-1,2)$.}
\solend
\fi

\ansspace[6]

\item 已知集合 $A=\{x\mid-2<x<15\}$,$B=\{x\mid m-6<x<2m-1,m\in\R\}$.

（1）当 $m=2$ 时，求 $A\cap B$；

（2）若 $A\cup B=A$，求实数 $m$ 的取值范围.

\ifanswer
\solbegin
\stepm{(1)}{当 $m=2$ 时，$A=\{x\mid-2<x<15\}$,$B=\{x\mid-4<x<3\}$，}
\step{此时 $A\cap B=\{x\mid-2<x<3\}$.}
\stepm{(2)}{由 $A\cup B=A$ 得 $B\subseteq A$.}
\step{当 $B=\varnothing$,$m-6\geqslant2m-1$，解得 $m\leqslant-5$；}
\step{当 $B\neq\varnothing$，结合 $B\subseteq A$ 得 $\begin{cases}m-6<2m-1\\ 2m-1\leqslant15\\ m-6\geqslant-2\end{cases}$，解得 $4\leqslant m\leqslant8$.}
\step{综上，$m\leqslant-5$ 或 $4\leqslant m\leqslant8$.}
\solend
\fi

\ansspace[6]

\item 已知 $a\in\R$，命题 $p$：关于 $x$ 的方程 $x^2-2x+a=0$ 无实根；命题 $q$：对任意 $x\in\R$，
不等式 $ax^2-ax+1>0$ 恒成立.

（1）若 $q$ 为真命题，求 $a$ 的取值范围；

（2）若 $p$、$q$ 中有且只有一个为真命题，求 $a$ 的取值范围.

\ifanswer
\solbegin
\stepm{(1)}{当 $q$ 为真命题时，}
\stepm{①}{当 $a=0$ 时，不等式化为 $1>0$，对任意 $x\in\R$ 恒成立，满足条件；}
\stepm{②}{当 $a\neq0$ 时，$ax^2-ax+1>0$ 为一元二次不等式，}
\step{若对任意 $x\in\R$ 恒成立，需满足 $\begin{cases}a>0\\ \Delta=(-a)^2-4a<0\end{cases}$，解得 $0<a<4$；}
\step{综上，$a$ 的取值范围为 $[0,4)$.}
\stepm{(2)}{$p$ 为真时，方程 $x^2-2x+a=0$ 无实根，}
\step{故判别式 $\Delta=4-4a<0$，解得 $a>1$，因此 $p$ 为假时，$a\leqslant1$；}
\step{由（1）得 $q$ 为假时 $a\geqslant4$ 或 $a<0$，}
\step{$p$、$q$ 有且只有一个真命题，分两类：}
\step{当 $p$ 假 $q$ 真时，满足 $\begin{cases}a\leqslant1\\ 0\leqslant a<4\end{cases}$，解得 $0\leqslant a\leqslant1$；}
\step{当 $p$ 真 $q$ 假时，满足 $\begin{cases}a>1\\ a\geqslant4\text{或}a<0\end{cases}$，解得 $a\geqslant4$.}
\step{综上，$a$ 的取值范围是 $[0,1]\cup[4,+\infty)$.}
\solend
\fi

\ansspace[8]

\item 某工厂生产 $x$($x>0$) 件产品的总成本 $C$（单位：元）与生产件数 $x$ 之间满足二次函数
$C(x)=\dfrac14x^2+50x+10000$.

（1）若要求生产总成本不超过 30000 元，求满足生产总成本要求的产品件数的取值范围；

（2）试求该工厂生产多少件产品时，可使每件产品的平均成本最小.

\ifanswer
\solbegin
\stepm{(1)}{若要求生产总成本不超过 30000 元，}
\step{则 $C(x)=\dfrac14x^2+50x+10000\leqslant30000$，}
\step{即 $x^2+200x-80000\leqslant0$，解得 $-400\leqslant x\leqslant200$，又 $x>0$，}
\step{所以满足生产总成本要求的产品件数的取值范围为 $(0,200]$，$x\in\N$；}
\stepm{(2)}{生产件数 $x$ 时每件的平均成本为 $\dfrac{\frac14x^2+50x+10000}x=\dfrac14x+50+\dfrac{10000}x$，}
\step{由基本不等式得 $\dfrac14x+50+\dfrac{10000}x\geqslant2\sqrt{\dfrac14x\cdot\dfrac{10000}x}+50=150$，}
\step{当且仅当 $\dfrac14x=\dfrac{10000}x$，即 $x=200$ 时取等号，}
\step{即该工厂生产 200 件产品时，每件产品的平均成本最小.}
\solend
\fi

\ansspace[8]

% 原卷题干“$(n\geqslant2,n\in N)$”后连印两个全角逗号“，，”，照录未改（见文件头“原卷笔误”①）
\item 已知有限集 $A=\{a_1,a_2,\cdots,a_n\}$($n\geqslant2$,$n\in\N$)，，如果 $A$ 中元素
$a_i$($i=1,2,3,\cdots,n$) 满足 $a_1+a_2+a_3+\cdots+a_n=a_1\times a_2\times\cdots\times a_n$，
就称 $A$ 为“和积平衡集”.

% 原卷 (2) 将“写方程”与“反证法证明”两项任务并作一句（“……表示），可尝试用反证法证明：……”），
% 照录未改（见文件头“原卷笔误”②）
（1）判断集合 $\{-1+\sqrt3,-1-\sqrt3\}$ 是否为“和积平衡集”；

（2）若 $a_1$、$a_2$ 是两个不同的正数，且 $\{a_1,a_2\}$ 是“和积平衡集”，写出以 $a_1$、$a_2$ 为根的
一个一元二次方程（系数可用 $a_1$、$a_2$ 表示），可尝试用反证法证明：$a_1$、$a_2$ 至少有一个大于 2；

（3）若 $a_i\geqslant1$，且 $a_i\in\N$，求所有符合条件的“和积平衡集”.

\ifanswer
\solbegin
\stepm{(1)}{对于集合 $\{-1+\sqrt3,-1-\sqrt3\}$，}
\step{因为 $(-1+\sqrt3)+(-1-\sqrt3)=-2$，$(-1+\sqrt3)\times(-1-\sqrt3)=-2$，}
\step{所以满足 $a_1+a_2=a_1\times a_2$，满足题意，}
\step{所以集合 $\{-1+\sqrt3,-1-\sqrt3\}$ 是和积平衡集.}
\stepm{(2)}{以 $a_1$、$a_2$ 为根的一个一元二次方程可为 $x^2-(a_1+a_2)x+a_1\cdot a_2=0$，}
\step{因为 $a_1$、$a_2$ 是两个不同的正数，且 $\{a_1,a_2\}$ 是和积平衡集，}
\step{所以 $a_1+a_2=a_1\cdot a_2$，}
\step{又 $a_1$ 和 $a_2$ 为 $x^2-(a_1+a_2)x+a_1\cdot a_2=0$ 的两个不同的正根，}
\step{所以 $\Delta=(a_1+a_2)^2-4a_1a_2=(a_1+a_2)^2-4(a_1+a_2)>0$，}
\step{又 $a_1+a_2>0$，则 $a_1+a_2>4$，}
\step{若 $a_1$、$a_2$ 都小于等于 2，则 $a_1+a_2<4$，矛盾，}
\step{所以 $a_1$、$a_2$ 至少有一个大于 2；}
\stepm{(3)}{设 $A$ 中的 $a_1<a_2<a_3<\cdots<a_n$，且 $a_i\in\N^*$，}
\step{由 $a_1+a_2+a_3+\cdots+a_n=a_1\times a_2\times\cdots\times a_n$，得 $n=2$ 时明显不满足，}
\step{所以 $n\geqslant3$，故 $a_1\times a_2\times\cdots\times a_n=a_1+a_2+a_3+\cdots+a_n\leqslant na_n$，}
\step{所以 $a_1\times a_2\times\cdots\times a_{n-1}\leqslant n$，}
\step{当 $n=3$ 时，$a_1\times a_2\leqslant3$，$a_i\in\N^*$，只能有 $a_1=1$，$a_2=2$，}
\step{由 $a_1+a_2+a_3=a_1\times a_2\times a_3$，得 $a_3=3$，此时的和积平衡集为 $\{1,2,3\}$，}
\step{当 $n\geqslant4$ 时，$a_1\times a_2\times\cdots\times a_{n-1}\geqslant(n-1)(n-2)\times\cdots\times2\times1\geqslant(n-1)(n-2)$，}
\step{又 $(n-1)(n-2)-n=n^2-4n+2=n(n-4)+2>0$，即 $(n-1)(n-2)>n$，}
\step{即 $a_1\times a_2\times\cdots\times a_{n-1}>n$，与 $a_1\times a_2\times\cdots\times a_{n-1}\leqslant n$ 矛盾，}
\step{所以 $n\geqslant4$ 不满足条件，}
\step{所以符合条件的和积平衡集有且只有 $n=3$，此时的和积平衡集为 $\{1,2,3\}$.}
\solend
\fi

\ansspace*[10]

\end{enumerate}

\end{document}
"
% 紧凑版：xelatex "\def\iscompact{1}% 高一63_吴淞中学_10月月考.tex
%   —— 高一上数学“2026-2027 年吴淞中学高一上 10 月月考”（讲评版）
% 试卷版：xelatex 高一63_吴淞中学_10月月考.tex
% 答案版：xelatex "\def\isanswer{1}% 高一63_吴淞中学_10月月考.tex
%   —— 高一上数学“2026-2027 年吴淞中学高一上 10 月月考”（讲评版）
% 试卷版：xelatex 高一63_吴淞中学_10月月考.tex
% 答案版：xelatex "\def\isanswer{1}% 高一63_吴淞中学_10月月考.tex
%   —— 高一上数学“2026-2027 年吴淞中学高一上 10 月月考”（讲评版）
% 试卷版：xelatex 高一63_吴淞中学_10月月考.tex
% 答案版：xelatex "\def\isanswer{1}\input{高一63_吴淞中学_10月月考.tex}"
% 紧凑版：xelatex "\def\iscompact{1}\input{高一63_吴淞中学_10月月考.tex}"
%
% 原始材料：微信公众号“魔都数学加油站”文章，作者破晓，连载“27学年高一63”；
% 原文标题“27学年高一63——2026-2027年上海市吴淞中学高一上10月月考”；
% 原文链接：https://mp.weixin.qq.com/s/TFQ4TLnt_l66t7t51es4OQ。
% 页面标注发布时间 2026-10-16 21:44（上海）。
% 正文为 10 张 PNG 页（`原图/高一63_吴淞中学_10月月考/01..10.png`）：
% 第 06、10 页 2560×3620，第 01、02、03、04、05、07、08、09 页 4762×6735
% （**同一篇各页分辨率不同**），均按页序读取。10 页全为试卷正文页（卷首标题在第 1 页页顶，
% 卷末解答题第 21 题解析止于第 10 页），无封面页、目录页、二维码或宣传页。
% 粉色斜排“魔都数学加油站”为卷面水印，与公众号页脚一并未录入。
%
% 卷首标题：照卷面印作“2026--2027 年吴淞中学高一上 10 月月考”（公众号标题中的“上海市”
% 卷面标题行没有，本稿照卷面），年份区间按全稿惯例排 en dash。原卷只有一行标题、无副标题。
%
% 篇幅与题号：原卷自第 3 题起编号（第 1、2 题不在本卷内）。填空题 3—12（10 题）、
% 选择题 13—16（4 题）、解答题 17—21（5 题），共 19 题。本稿沿用原节与原题号，
% 只改排为全稿统一的大节样式（原卷小节标题为“一、填空题”“二、选择题”“三、解答题”）。
%
% 讲评版：原卷每题下直接印【答案】或【解析】。**只印【答案】的 2 题**（第 4、7 题），
% 本稿只给 \ans{}；其余 17 题里 16 题都印【解析】（按全稿口径这类题不另提 \ans{}——第 13、15、16 题
% 的结论写在解析末句（“故选 B”等），亦照录在解析内；第 14 题则既无【答案】也无【解析】段。**另有 1 题答案印在题干括号内**——
% 第 14 题（“(  D  )”，且该题无【答案】/【解析】段），本稿按 高一37 第 13 题成例把括号留空、
% 答案另提为 \ans{$D$}。第 16 题题干括号为空但印有【解析】，不算括号内答案。
%
% 副标题：原卷未印副标题。本稿按“知识模块”口径标作“集合与常用逻辑用语、不等式”：
% 19 题中集合与常用逻辑用语 12 题（第 3、4、6、7、9、10、11、12、13、15、18、21 题），
% 不等式 7 题（含等式（韦达定理、恒等式）与不等式性质、基本不等式、二次不等式、恒成立：
% 第 5、8、14、16、17、19、20 题），两类均超过三分之一，故并标。口径同 高一59。
%
% 与原卷的排版差异：
%   ① 原卷每题下直接印【答案】/【解析】，本稿按全稿惯例将答案与解析置于答案版，试卷版保持干净；
%   ② 原卷未印副标题，本稿按“知识模块”口径补出；
%   ③ 原卷填空横线约四个字宽，本稿按全稿惯例统一排 \blankline[4.4em]；
%   ④ 原卷实数集排斜体 $R$（第 9、10、19 题题干与解析）或粗体 $\mathbf{R}$（第 18 题题干），
%      整数集排斜体 $Z$（第 12 题解析），自然数集排斜体 $N$（第 20 题解析、第 21 题题干
%      “$(n\geqslant2,n\in N)$”）或粗体 $\mathbf{N}$（第 21 题 (3) 题干）、斜体 $N^*$
%      （第 21 题 (3) 解析），本稿按全稿惯例统一改用 $\R$、$\Z$、$\N$、$\N^*$
%      （同 高一14 排版差异④：记号归一，粗体／斜体不作区分，$N$ 与 $N^*$ 的区别照原卷保留）；
%      空集记号原卷作 $\varnothing$，照录；
%   ⑤ 原卷在数学式内多处用半角逗号（如 $\{y\mid y=[2x]+[3x],-1\leqslant x\leqslant2\}$、
%      $m=-1,0,1$、$A=\{x\mid-2<x<15\}$,$B=\cdots$），均照录未强改；
%      句末句点按全稿惯例排西文句点；“和积平衡集”自定义名词照录；
%   ⑥ 第 12 题解析里的“法一：”“法二：”标号原卷排**红色**（全卷共 2 处），
%      同 高一53“情况1：”成例，本稿排 \stepm{法一：}{…}、\stepm{法二：}{…}
%      （加粗着色、不改黑色基调），冒号照原卷排全角；
%   ⑦ 第 17—21 题的子问原卷各自另起一行，本稿按全稿惯例用 enumerate 的 (1)(2)(3)；
%      解析中的子问标号一律用 \stepm{(1)}{…}；第 19 题解析中的①②分情形，按全稿惯例作
%      \stepm{①}{…}、\stepm{②}{…}；
%   ⑧ 第 14 题印在题干括号内的答案移入 \ans{}，见上；
%   ⑨ 每道解答题原卷未留作答空白，本稿按全稿惯例加 \ansspace；
%   ⑩ 选择题的作答括号原卷排半角“(　　)”（第 13、15、16 题括号内空白、第 14 题括号内印 D），
%      本稿按全稿惯例排 \mbox{（\quad）}。
%
% 原卷笔误/存疑（均照录未改，正文对应处另有 % 注）：
%   ① 第 21 题题干“$(n\geqslant2$,$n\in N)$”后连印两个全角逗号“，，”（衍字，
%      8 倍放大确认，非水印干扰），本稿照录；
%   ② 第 21 题 (2) 题干将“写出以 $a_1$、$a_2$ 为根的一个一元二次方程”与“用反证法证明：
%      $a_1$、$a_2$ 至少有一个大于 2”两项任务并作一句、之间仅以逗号相接（“……表示），
%      可尝试用反证法证明：……”），读来像抄并两小问，原卷如此，照录未改。
%
% 插图：无（全卷 10 页无任何图形、无“如图”字样）。
%
\documentclass[a4paper,11pt]{article}
\usepackage{common}

\begin{document}

\examtitlest[3]{2026--2027 年吴淞中学高一上 10 月月考}{集合与常用逻辑用语、不等式}

\bigsec{一、填空题}

\begin{enumerate}
\setcounter{enumi}{2}

\item 满足 $A\subseteq\{1,2,3,4\}$ 的集合 $A$ 共有\blankline[4.4em]个.

\ifanswer
\solbegin
\step{由题意，集合 $A$ 是集合 $\{1,2,3,4\}$ 的子集，共有 $2^4=16$ 个.}
\solend
\fi

\item 陈述句“$a>1$ 或 $b>3$”的否定形式为\blankline[4.4em].

\ifanswer
\ans{$a\leqslant1$ 且 $b\leqslant3$}
\fi

\item 已知方程 $3x^2+4x-4=0$ 的两个根为 $x_1$、$x_2$，则 $x_1^2+x_2^2=$\blankline[4.4em].

\ifanswer
\solbegin
\step{对于方程 $3x^2+4x-4=0$，因为 $\Delta=16-4\times3\times(-4)=64>0$，}
\step{所以方程 $3x^2+4x-4=0$ 的两个不相等的实根为 $x_1$、$x_2$，}
\step{由韦达定理可得 $x_1+x_2=-\dfrac43$，$x_1x_2=-\dfrac43$，}
\step{所以 $x_1^2+x_2^2=(x_1+x_2)^2-2x_1x_2=\dfrac{16}9+\dfrac83=\dfrac{40}9$.}
\solend
\fi

\item 设集合 $A=\{y\mid y=x^2+1\}$,$B=\{y\mid y=-x^2+2x+3\}$，则 $A\cap B=$\blankline[4.4em].

\ifanswer
\solbegin
\step{因为 $y=x^2+1\geqslant1$,$y=-x^2+2x+3=-(x-1)^2+4\leqslant4$，}
\step{所以 $A=\{y\mid y\geqslant1\}$,$B=\{y\mid y\leqslant4\}$，得 $A\cap B=\{y\mid1\leqslant y\leqslant4\}$.}
\solend
\fi

\item 若 $x\in\{1,x^2\}$，则 $x=$\blankline[4.4em].

\ifanswer
\ans{$0$}
\fi

\item 已知等式 $2x^2+3x+5=a(2x+1)(x+1)+c$，则 $a+c=$\blankline[4.4em].

\ifanswer
\solbegin
\step{因为 $2x^2+3x+5=a(2x+1)(x+1)+c$ 恒成立，}
\step{所以 $2x^2+3x+5=2ax^2+3ax+a+c$ 恒成立，则 $\begin{cases}2=2a\\ 3=3a\\ 5=a+c\end{cases}$，解得 $a+c=5$.}
\solend
\fi

\item 已知 $a\in\R$，设 $p:x<2$，$q:x<5a-3$，且 $p$ 是 $q$ 的必要非充分条件，则 $a$ 的取值范围是\blankline[4.4em].

\ifanswer
\solbegin
\step{设 $p:x<2$，$q:x<5a-3$，且 $p$ 是 $q$ 的必要非充分条件，}
\step{则 $\{x\mid x<5a-3\}\subset\{x\mid x<2\}$，所以 $5a-3<2$，即 $a<1$.}
\solend
\fi

\item 已知 $a\in\R$，集合 $A=\{x\mid ax^2+3x-1=0\}$，若集合 $A$ 中只有一个元素，则 $a$ 的取值范围是\blankline[4.4em].

\ifanswer
\solbegin
\step{当 $a=0$ 时，则有 $\{x\mid ax^2+3x-1=0\}=\{x\mid3x-1=0\}=\left\{\dfrac13\right\}$，符合题意；}
\step{当 $a\neq0$ 时，由题意得 $\Delta=9+4a=0$，解得 $a=-\dfrac94$. 符合题意；}
\step{综上所述，实数 $a$ 的取值集合为 $\left\{0,-\dfrac94\right\}$.}
\solend
\fi

\item 设集合 $A=\{x\mid x^2+2(m+1)x+m^2-1=0,m\in\R\}$，且 $A\subseteq\{0,-4\}$，则实数 $m$ 的取值范围是\blankline[4.4em].

\ifanswer
\solbegin
\step{由题意得 $A=\varnothing$ 或 $A=\{0\}$ 或 $A=\{-4\}$ 或 $A=\{0,-4\}$，}
\step{当 $A=\varnothing$ 时，方程 $x^2+2(m+1)x+m^2-1=0$ 无实根，}
\step{则 $\Delta=[2(m+1)]^2-4\times1\times(m^2-1)=8m+8<0$，解得 $m<-1$，}
\step{当 $A=\{0\}$ 时，方程 $x^2+2(m+1)x+m^2-1=0$ 有一个实根 $0$，}
\step{则 $\begin{cases}\Delta=8m+8=0\\ m^2-1=0\end{cases}$，解得 $m=-1$，}
\step{当 $A=\{-4\}$ 时，方程 $x^2+2(m+1)x+m^2-1=0$ 有一个实根 $-4$，}
\step{则 $\begin{cases}\Delta=8m+8=0\\ 16-8(m+1)+m^2-1=0\end{cases}$，无解，}
\step{当 $A=\{0,-4\}$ 时，方程 $x^2+2(m+1)x+m^2-1=0$ 有两个实根 $-4$ 和 $0$，}
\step{则 $\begin{cases}\Delta=8m+8>0\\ 16-8(m+1)+m^2-1=0\\ m^2-1=0\end{cases}$，解得 $m=1$，}
\step{综上所述，实数 $m$ 的取值范围是 $\{m\mid m\leqslant-1\text{ 或 }m=1\}$.}
\solend
\fi

\item 对于任意实数 $x$，$[x]$ 表示不超过 $x$ 的最大整数，如 $[1.1]=1$，$[-1.2]=-2$，若
$A=\{y\mid y=[2x]+[3x],-1\leqslant x\leqslant2\}$，则 $A$ 中所有元素的和为\blankline[4.4em].

\ifanswer
\solbegin
% 原卷“法一：”“法二：”标号排红色，本稿按 高一53“情况1：”成例排 \stepm（加粗着色），见文件头
\stepm{法一：}{当 $-1\leqslant x<-\dfrac23$ 时，$[2x]=-2$，$[3x]=-3$，$f(x)=[2x]+[3x]=-5$，}
\step{当 $-\dfrac23\leqslant x<-\dfrac12$ 时，$[2x]=-2$，$[3x]=-2$，$f(x)=[2x]+[3x]=-4$；}
\step{当 $-\dfrac12\leqslant x<-\dfrac13$ 时，$[2x]=-1$，$[3x]=-2$，$f(x)=[2x]+[3x]=-3$；}
\step{当 $-\dfrac13\leqslant x<0$ 时，$[2x]=-1$，$[3x]=-1$，$f(x)=[2x]+[3x]=-2$；}
\step{当 $0\leqslant x<\dfrac13$ 时，$[2x]=0$，$[3x]=0$，$f(x)=[2x]+[3x]=0$；}
\step{当 $\dfrac13\leqslant x<\dfrac12$ 时，$[2x]=0$，$[3x]=1$，$f(x)=[2x]+[3x]=1$；}
\step{当 $\dfrac12\leqslant x<\dfrac23$ 时，$[2x]=1$，$[3x]=1$，$f(x)=[2x]+[3x]=2$；}
\step{当 $\dfrac23\leqslant x<1$ 时，$[2x]=1$，$[3x]=2$，$f(x)=[2x]+[3x]=3$；}
\step{当 $1\leqslant x<\dfrac43$ 时，$[2x]=2$，$[3x]=3$，$f(x)=[2x]+[3x]=5$；}
\step{当 $\dfrac43\leqslant x<\dfrac32$ 时，$[2x]=2$，$[3x]=4$，$f(x)=[2x]+[3x]=6$；}
\step{当 $\dfrac32\leqslant x<\dfrac53$ 时，$[2x]=3$，$[3x]=4$，$f(x)=[2x]+[3x]=7$；}
\step{当 $\dfrac53\leqslant x<2$ 时，$[2x]=3$，$[3x]=5$，$f(x)=[2x]+[3x]=8$；}
\step{当 $x=2$ 时，$[2x]=4$，$[3x]=6$，$f(x)=[2x]+[3x]=10$}
\step{故 $A$ 中所有元素的和为 $-5-4-3-2+0+1+2+3+5+6+7+8+10=28$.}
% 原卷“法二：”标号排红色，同上
\stepm{法二：}{设 $x=m+n$，其中 $m\in\Z$，$n\in[0,1)$，}
\step{若 $0\leqslant n<\dfrac13$，则 $0\leqslant2n<\dfrac23$，$0\leqslant3n<1$，}
\step{得 $[2x]=[2m+2n]=2m$，$[3x]=[3m+3n]=3m$，所以 $y=[2x]+[3x]=5m$；}
\step{若 $\dfrac13\leqslant n<\dfrac12$，则 $\dfrac23\leqslant2n<1$，$1\leqslant3n<\dfrac32$，}
\step{得 $[2x]=[2m+2n]=2m$，$[3x]=[3m+3n]=3m+1$，}
\step{所以 $y=[2x]+[3x]=5m+1$；}
\step{若 $\dfrac12\leqslant n<\dfrac23$，则 $1\leqslant2n<\dfrac43$，$\dfrac32\leqslant3n<2$，}
\step{得 $[2x]=[2m+2n]=2m+1$，$[3x]=[3m+3n]=3m+1$，}
\step{所以 $y=[2x]+[3x]=5m+2$；}
\step{若 $\dfrac23\leqslant n<1$，则 $\dfrac43\leqslant2n<2$，$2\leqslant3n<3$，}
\step{得 $[2x]=[2m+2n]=2m+1$，$[3x]=[3m+3n]=3m+2$，}
\step{所以 $y=[2x]+[3x]=5m+3$；}
\step{综上所述，当 $n\in[0,1)$ 时，所有函数值的和}
\step{为 $5m+5m+1+5m+2+5m+3=20m+6$，}
\step{又因为 $x\in[-1,2]$，此时 $m=-1,0,1$，且当 $x=2$ 时，$y=[2x]+[3x]=10$，}
\step{所有函数值的和为 $(-14)+6+26+10=28$，即 $A$ 中所有元素之和为 $28$.}
\solend
\fi

\end{enumerate}

\bigsec{二、选择题}

\begin{enumerate}
\setcounter{enumi}{12}

\item 下列写法中，正确的是\mbox{（\quad）}

\keepwithprev
A. $0\in\{\varnothing\}$\qquad B. $\varnothing\in\{\varnothing\}$

C. $\varnothing=\{\varnothing\}$\qquad D. $0\in\varnothing$

\ifanswer
\solbegin
\step{$0\notin\{\varnothing\}$，$\varnothing\in\{\varnothing\}$，故 $A$ 错误，$B$ 正确；}
\step{$\varnothing$ 是不含任何元素的集合，所以 $\varnothing\neq\{\varnothing\}$，$0\notin\varnothing$，故 $CD$ 错误.}
\step{故选 $B$.}
\solend
\fi

% 第 14 题答案“D”直接印在题干括号内“(  D  )”，且该题无【答案】/【解析】段；
% 按 高一37 第 13 题的成例括号留空、答案另提为 \ans{}（见文件头）
\item 若 $a<b<0$，则下列不等式中不能成立的是\mbox{（\quad）}

\keepwithprev
A. $\dfrac1a>\dfrac1b$\qquad B. $a-b>a$

C. $|a|>|b|$\qquad D. $a^3>b^3$

\ifanswer
\ans{$D$}
\fi

\item 已知 $x$、$y$、$z$ 为非零实数，代数式 $\dfrac x{|x|}+\dfrac y{|y|}+\dfrac z{|z|}+\dfrac{xyz}{|xyz|}$
的值所组成的集合是 $M$，则下列判断正确的是\mbox{（\quad）}

\keepwithprev
A. $-4\notin M$\qquad B. $2\in M$

C. $4\notin M$\qquad D. $0\in M$

\ifanswer
\solbegin
\step{当 $x$、$y$、$z$ 均为正数时，$xyz>0$，$\dfrac x{|x|}+\dfrac y{|y|}+\dfrac z{|z|}+\dfrac{xyz}{|xyz|}=1+1+1+1=4$，}
\step{当 $x$、$y$、$z$ 有 2 个正数，1 个负数时，不妨设 $x>0$，$y>0$，$z<0$，}
\step{则 $xyz<0$，$\dfrac x{|x|}+\dfrac y{|y|}+\dfrac z{|z|}+\dfrac{xyz}{|xyz|}=1+1-1-1=0$，}
\step{若 $x$、$y$、$z$ 有 1 个正数，2 个负数时，不妨设 $x>0$，$y<0$，$z<0$，}
\step{则 $xyz>0$，$\dfrac x{|x|}+\dfrac y{|y|}+\dfrac z{|z|}+\dfrac{xyz}{|xyz|}=1-1-1+1=0$，}
\step{当 $x$、$y$、$z$ 均为负数时，$xyz<0$，}
\step{$\dfrac x{|x|}+\dfrac y{|y|}+\dfrac z{|z|}+\dfrac{xyz}{|xyz|}=-1-1-1-1=-4$，}
\step{综上，$M=\{-4,0,4\}$，}
\step{所以 $-4\in M$，$A$ 错误；$2\notin M$，$B$ 错误；$4\in M$，$C$ 错误；}
\step{$0\in M$，$D$ 正确. 故选 $D$.}
\solend
\fi

\item 若 $xy>0$，且 $x+y=9$，则 $\dfrac1x+\dfrac4y$ 的最小值是\mbox{（\quad）}

\keepwithprev
A. $9$\qquad B. $\dfrac59$

C. $\dfrac49$\qquad D. $1$

\ifanswer
\solbegin
\step{$xy>0$，且 $x+y=9$，所以 $x>0$，$y>0$，}
\step{所以 $\dfrac1x+\dfrac4y=\dfrac19\left(\dfrac1x+\dfrac4y\right)(x+y)=\dfrac59+\dfrac19\left(\dfrac yx+\dfrac{4x}y\right)\geqslant\dfrac59+\dfrac29\cdot\sqrt{\dfrac yx\cdot\dfrac{4x}y}=\dfrac59+\dfrac49=1$，}
\step{当且仅当 $\dfrac yx=\dfrac{4x}y$，即 $x=3$，$y=6$ 时取等号，所以 $\dfrac1x+\dfrac4y$ 的最小值为 $1$.}
\step{故选 $D$.}
\solend
\fi

\end{enumerate}

\bigsec{三、解答题}

\begin{enumerate}
\setcounter{enumi}{16}

\item 求解下列不等式：

（1）$\dfrac{x-3}{x+1}\leqslant0$；

（2）$|2x-1|<3$.

\ifanswer
\solbegin
\stepm{(1)}{$\dfrac{x-3}{x+1}\leqslant0$，变形得 $(x-3)\cdot(x+1)\leqslant0$ 且 $x\neq-1$，}
\step{解得 $-1<x\leqslant3$，即不等式的解集为 $(-1,3]$.}
\stepm{(2)}{若 $|2x-1|<3$，则 $-3<2x-1<3$，解得 $-1<x<2$，即 $(-1,2)$.}
\solend
\fi

\ansspace[6]

\item 已知集合 $A=\{x\mid-2<x<15\}$,$B=\{x\mid m-6<x<2m-1,m\in\R\}$.

（1）当 $m=2$ 时，求 $A\cap B$；

（2）若 $A\cup B=A$，求实数 $m$ 的取值范围.

\ifanswer
\solbegin
\stepm{(1)}{当 $m=2$ 时，$A=\{x\mid-2<x<15\}$,$B=\{x\mid-4<x<3\}$，}
\step{此时 $A\cap B=\{x\mid-2<x<3\}$.}
\stepm{(2)}{由 $A\cup B=A$ 得 $B\subseteq A$.}
\step{当 $B=\varnothing$,$m-6\geqslant2m-1$，解得 $m\leqslant-5$；}
\step{当 $B\neq\varnothing$，结合 $B\subseteq A$ 得 $\begin{cases}m-6<2m-1\\ 2m-1\leqslant15\\ m-6\geqslant-2\end{cases}$，解得 $4\leqslant m\leqslant8$.}
\step{综上，$m\leqslant-5$ 或 $4\leqslant m\leqslant8$.}
\solend
\fi

\ansspace[6]

\item 已知 $a\in\R$，命题 $p$：关于 $x$ 的方程 $x^2-2x+a=0$ 无实根；命题 $q$：对任意 $x\in\R$，
不等式 $ax^2-ax+1>0$ 恒成立.

（1）若 $q$ 为真命题，求 $a$ 的取值范围；

（2）若 $p$、$q$ 中有且只有一个为真命题，求 $a$ 的取值范围.

\ifanswer
\solbegin
\stepm{(1)}{当 $q$ 为真命题时，}
\stepm{①}{当 $a=0$ 时，不等式化为 $1>0$，对任意 $x\in\R$ 恒成立，满足条件；}
\stepm{②}{当 $a\neq0$ 时，$ax^2-ax+1>0$ 为一元二次不等式，}
\step{若对任意 $x\in\R$ 恒成立，需满足 $\begin{cases}a>0\\ \Delta=(-a)^2-4a<0\end{cases}$，解得 $0<a<4$；}
\step{综上，$a$ 的取值范围为 $[0,4)$.}
\stepm{(2)}{$p$ 为真时，方程 $x^2-2x+a=0$ 无实根，}
\step{故判别式 $\Delta=4-4a<0$，解得 $a>1$，因此 $p$ 为假时，$a\leqslant1$；}
\step{由（1）得 $q$ 为假时 $a\geqslant4$ 或 $a<0$，}
\step{$p$、$q$ 有且只有一个真命题，分两类：}
\step{当 $p$ 假 $q$ 真时，满足 $\begin{cases}a\leqslant1\\ 0\leqslant a<4\end{cases}$，解得 $0\leqslant a\leqslant1$；}
\step{当 $p$ 真 $q$ 假时，满足 $\begin{cases}a>1\\ a\geqslant4\text{或}a<0\end{cases}$，解得 $a\geqslant4$.}
\step{综上，$a$ 的取值范围是 $[0,1]\cup[4,+\infty)$.}
\solend
\fi

\ansspace[8]

\item 某工厂生产 $x$($x>0$) 件产品的总成本 $C$（单位：元）与生产件数 $x$ 之间满足二次函数
$C(x)=\dfrac14x^2+50x+10000$.

（1）若要求生产总成本不超过 30000 元，求满足生产总成本要求的产品件数的取值范围；

（2）试求该工厂生产多少件产品时，可使每件产品的平均成本最小.

\ifanswer
\solbegin
\stepm{(1)}{若要求生产总成本不超过 30000 元，}
\step{则 $C(x)=\dfrac14x^2+50x+10000\leqslant30000$，}
\step{即 $x^2+200x-80000\leqslant0$，解得 $-400\leqslant x\leqslant200$，又 $x>0$，}
\step{所以满足生产总成本要求的产品件数的取值范围为 $(0,200]$，$x\in\N$；}
\stepm{(2)}{生产件数 $x$ 时每件的平均成本为 $\dfrac{\frac14x^2+50x+10000}x=\dfrac14x+50+\dfrac{10000}x$，}
\step{由基本不等式得 $\dfrac14x+50+\dfrac{10000}x\geqslant2\sqrt{\dfrac14x\cdot\dfrac{10000}x}+50=150$，}
\step{当且仅当 $\dfrac14x=\dfrac{10000}x$，即 $x=200$ 时取等号，}
\step{即该工厂生产 200 件产品时，每件产品的平均成本最小.}
\solend
\fi

\ansspace[8]

% 原卷题干“$(n\geqslant2,n\in N)$”后连印两个全角逗号“，，”，照录未改（见文件头“原卷笔误”①）
\item 已知有限集 $A=\{a_1,a_2,\cdots,a_n\}$($n\geqslant2$,$n\in\N$)，，如果 $A$ 中元素
$a_i$($i=1,2,3,\cdots,n$) 满足 $a_1+a_2+a_3+\cdots+a_n=a_1\times a_2\times\cdots\times a_n$，
就称 $A$ 为“和积平衡集”.

% 原卷 (2) 将“写方程”与“反证法证明”两项任务并作一句（“……表示），可尝试用反证法证明：……”），
% 照录未改（见文件头“原卷笔误”②）
（1）判断集合 $\{-1+\sqrt3,-1-\sqrt3\}$ 是否为“和积平衡集”；

（2）若 $a_1$、$a_2$ 是两个不同的正数，且 $\{a_1,a_2\}$ 是“和积平衡集”，写出以 $a_1$、$a_2$ 为根的
一个一元二次方程（系数可用 $a_1$、$a_2$ 表示），可尝试用反证法证明：$a_1$、$a_2$ 至少有一个大于 2；

（3）若 $a_i\geqslant1$，且 $a_i\in\N$，求所有符合条件的“和积平衡集”.

\ifanswer
\solbegin
\stepm{(1)}{对于集合 $\{-1+\sqrt3,-1-\sqrt3\}$，}
\step{因为 $(-1+\sqrt3)+(-1-\sqrt3)=-2$，$(-1+\sqrt3)\times(-1-\sqrt3)=-2$，}
\step{所以满足 $a_1+a_2=a_1\times a_2$，满足题意，}
\step{所以集合 $\{-1+\sqrt3,-1-\sqrt3\}$ 是和积平衡集.}
\stepm{(2)}{以 $a_1$、$a_2$ 为根的一个一元二次方程可为 $x^2-(a_1+a_2)x+a_1\cdot a_2=0$，}
\step{因为 $a_1$、$a_2$ 是两个不同的正数，且 $\{a_1,a_2\}$ 是和积平衡集，}
\step{所以 $a_1+a_2=a_1\cdot a_2$，}
\step{又 $a_1$ 和 $a_2$ 为 $x^2-(a_1+a_2)x+a_1\cdot a_2=0$ 的两个不同的正根，}
\step{所以 $\Delta=(a_1+a_2)^2-4a_1a_2=(a_1+a_2)^2-4(a_1+a_2)>0$，}
\step{又 $a_1+a_2>0$，则 $a_1+a_2>4$，}
\step{若 $a_1$、$a_2$ 都小于等于 2，则 $a_1+a_2<4$，矛盾，}
\step{所以 $a_1$、$a_2$ 至少有一个大于 2；}
\stepm{(3)}{设 $A$ 中的 $a_1<a_2<a_3<\cdots<a_n$，且 $a_i\in\N^*$，}
\step{由 $a_1+a_2+a_3+\cdots+a_n=a_1\times a_2\times\cdots\times a_n$，得 $n=2$ 时明显不满足，}
\step{所以 $n\geqslant3$，故 $a_1\times a_2\times\cdots\times a_n=a_1+a_2+a_3+\cdots+a_n\leqslant na_n$，}
\step{所以 $a_1\times a_2\times\cdots\times a_{n-1}\leqslant n$，}
\step{当 $n=3$ 时，$a_1\times a_2\leqslant3$，$a_i\in\N^*$，只能有 $a_1=1$，$a_2=2$，}
\step{由 $a_1+a_2+a_3=a_1\times a_2\times a_3$，得 $a_3=3$，此时的和积平衡集为 $\{1,2,3\}$，}
\step{当 $n\geqslant4$ 时，$a_1\times a_2\times\cdots\times a_{n-1}\geqslant(n-1)(n-2)\times\cdots\times2\times1\geqslant(n-1)(n-2)$，}
\step{又 $(n-1)(n-2)-n=n^2-4n+2=n(n-4)+2>0$，即 $(n-1)(n-2)>n$，}
\step{即 $a_1\times a_2\times\cdots\times a_{n-1}>n$，与 $a_1\times a_2\times\cdots\times a_{n-1}\leqslant n$ 矛盾，}
\step{所以 $n\geqslant4$ 不满足条件，}
\step{所以符合条件的和积平衡集有且只有 $n=3$，此时的和积平衡集为 $\{1,2,3\}$.}
\solend
\fi

\ansspace*[10]

\end{enumerate}

\end{document}
"
% 紧凑版：xelatex "\def\iscompact{1}% 高一63_吴淞中学_10月月考.tex
%   —— 高一上数学“2026-2027 年吴淞中学高一上 10 月月考”（讲评版）
% 试卷版：xelatex 高一63_吴淞中学_10月月考.tex
% 答案版：xelatex "\def\isanswer{1}\input{高一63_吴淞中学_10月月考.tex}"
% 紧凑版：xelatex "\def\iscompact{1}\input{高一63_吴淞中学_10月月考.tex}"
%
% 原始材料：微信公众号“魔都数学加油站”文章，作者破晓，连载“27学年高一63”；
% 原文标题“27学年高一63——2026-2027年上海市吴淞中学高一上10月月考”；
% 原文链接：https://mp.weixin.qq.com/s/TFQ4TLnt_l66t7t51es4OQ。
% 页面标注发布时间 2026-10-16 21:44（上海）。
% 正文为 10 张 PNG 页（`原图/高一63_吴淞中学_10月月考/01..10.png`）：
% 第 06、10 页 2560×3620，第 01、02、03、04、05、07、08、09 页 4762×6735
% （**同一篇各页分辨率不同**），均按页序读取。10 页全为试卷正文页（卷首标题在第 1 页页顶，
% 卷末解答题第 21 题解析止于第 10 页），无封面页、目录页、二维码或宣传页。
% 粉色斜排“魔都数学加油站”为卷面水印，与公众号页脚一并未录入。
%
% 卷首标题：照卷面印作“2026--2027 年吴淞中学高一上 10 月月考”（公众号标题中的“上海市”
% 卷面标题行没有，本稿照卷面），年份区间按全稿惯例排 en dash。原卷只有一行标题、无副标题。
%
% 篇幅与题号：原卷自第 3 题起编号（第 1、2 题不在本卷内）。填空题 3—12（10 题）、
% 选择题 13—16（4 题）、解答题 17—21（5 题），共 19 题。本稿沿用原节与原题号，
% 只改排为全稿统一的大节样式（原卷小节标题为“一、填空题”“二、选择题”“三、解答题”）。
%
% 讲评版：原卷每题下直接印【答案】或【解析】。**只印【答案】的 2 题**（第 4、7 题），
% 本稿只给 \ans{}；其余 17 题里 16 题都印【解析】（按全稿口径这类题不另提 \ans{}——第 13、15、16 题
% 的结论写在解析末句（“故选 B”等），亦照录在解析内；第 14 题则既无【答案】也无【解析】段。**另有 1 题答案印在题干括号内**——
% 第 14 题（“(  D  )”，且该题无【答案】/【解析】段），本稿按 高一37 第 13 题成例把括号留空、
% 答案另提为 \ans{$D$}。第 16 题题干括号为空但印有【解析】，不算括号内答案。
%
% 副标题：原卷未印副标题。本稿按“知识模块”口径标作“集合与常用逻辑用语、不等式”：
% 19 题中集合与常用逻辑用语 12 题（第 3、4、6、7、9、10、11、12、13、15、18、21 题），
% 不等式 7 题（含等式（韦达定理、恒等式）与不等式性质、基本不等式、二次不等式、恒成立：
% 第 5、8、14、16、17、19、20 题），两类均超过三分之一，故并标。口径同 高一59。
%
% 与原卷的排版差异：
%   ① 原卷每题下直接印【答案】/【解析】，本稿按全稿惯例将答案与解析置于答案版，试卷版保持干净；
%   ② 原卷未印副标题，本稿按“知识模块”口径补出；
%   ③ 原卷填空横线约四个字宽，本稿按全稿惯例统一排 \blankline[4.4em]；
%   ④ 原卷实数集排斜体 $R$（第 9、10、19 题题干与解析）或粗体 $\mathbf{R}$（第 18 题题干），
%      整数集排斜体 $Z$（第 12 题解析），自然数集排斜体 $N$（第 20 题解析、第 21 题题干
%      “$(n\geqslant2,n\in N)$”）或粗体 $\mathbf{N}$（第 21 题 (3) 题干）、斜体 $N^*$
%      （第 21 题 (3) 解析），本稿按全稿惯例统一改用 $\R$、$\Z$、$\N$、$\N^*$
%      （同 高一14 排版差异④：记号归一，粗体／斜体不作区分，$N$ 与 $N^*$ 的区别照原卷保留）；
%      空集记号原卷作 $\varnothing$，照录；
%   ⑤ 原卷在数学式内多处用半角逗号（如 $\{y\mid y=[2x]+[3x],-1\leqslant x\leqslant2\}$、
%      $m=-1,0,1$、$A=\{x\mid-2<x<15\}$,$B=\cdots$），均照录未强改；
%      句末句点按全稿惯例排西文句点；“和积平衡集”自定义名词照录；
%   ⑥ 第 12 题解析里的“法一：”“法二：”标号原卷排**红色**（全卷共 2 处），
%      同 高一53“情况1：”成例，本稿排 \stepm{法一：}{…}、\stepm{法二：}{…}
%      （加粗着色、不改黑色基调），冒号照原卷排全角；
%   ⑦ 第 17—21 题的子问原卷各自另起一行，本稿按全稿惯例用 enumerate 的 (1)(2)(3)；
%      解析中的子问标号一律用 \stepm{(1)}{…}；第 19 题解析中的①②分情形，按全稿惯例作
%      \stepm{①}{…}、\stepm{②}{…}；
%   ⑧ 第 14 题印在题干括号内的答案移入 \ans{}，见上；
%   ⑨ 每道解答题原卷未留作答空白，本稿按全稿惯例加 \ansspace；
%   ⑩ 选择题的作答括号原卷排半角“(　　)”（第 13、15、16 题括号内空白、第 14 题括号内印 D），
%      本稿按全稿惯例排 \mbox{（\quad）}。
%
% 原卷笔误/存疑（均照录未改，正文对应处另有 % 注）：
%   ① 第 21 题题干“$(n\geqslant2$,$n\in N)$”后连印两个全角逗号“，，”（衍字，
%      8 倍放大确认，非水印干扰），本稿照录；
%   ② 第 21 题 (2) 题干将“写出以 $a_1$、$a_2$ 为根的一个一元二次方程”与“用反证法证明：
%      $a_1$、$a_2$ 至少有一个大于 2”两项任务并作一句、之间仅以逗号相接（“……表示），
%      可尝试用反证法证明：……”），读来像抄并两小问，原卷如此，照录未改。
%
% 插图：无（全卷 10 页无任何图形、无“如图”字样）。
%
\documentclass[a4paper,11pt]{article}
\usepackage{common}

\begin{document}

\examtitlest[3]{2026--2027 年吴淞中学高一上 10 月月考}{集合与常用逻辑用语、不等式}

\bigsec{一、填空题}

\begin{enumerate}
\setcounter{enumi}{2}

\item 满足 $A\subseteq\{1,2,3,4\}$ 的集合 $A$ 共有\blankline[4.4em]个.

\ifanswer
\solbegin
\step{由题意，集合 $A$ 是集合 $\{1,2,3,4\}$ 的子集，共有 $2^4=16$ 个.}
\solend
\fi

\item 陈述句“$a>1$ 或 $b>3$”的否定形式为\blankline[4.4em].

\ifanswer
\ans{$a\leqslant1$ 且 $b\leqslant3$}
\fi

\item 已知方程 $3x^2+4x-4=0$ 的两个根为 $x_1$、$x_2$，则 $x_1^2+x_2^2=$\blankline[4.4em].

\ifanswer
\solbegin
\step{对于方程 $3x^2+4x-4=0$，因为 $\Delta=16-4\times3\times(-4)=64>0$，}
\step{所以方程 $3x^2+4x-4=0$ 的两个不相等的实根为 $x_1$、$x_2$，}
\step{由韦达定理可得 $x_1+x_2=-\dfrac43$，$x_1x_2=-\dfrac43$，}
\step{所以 $x_1^2+x_2^2=(x_1+x_2)^2-2x_1x_2=\dfrac{16}9+\dfrac83=\dfrac{40}9$.}
\solend
\fi

\item 设集合 $A=\{y\mid y=x^2+1\}$,$B=\{y\mid y=-x^2+2x+3\}$，则 $A\cap B=$\blankline[4.4em].

\ifanswer
\solbegin
\step{因为 $y=x^2+1\geqslant1$,$y=-x^2+2x+3=-(x-1)^2+4\leqslant4$，}
\step{所以 $A=\{y\mid y\geqslant1\}$,$B=\{y\mid y\leqslant4\}$，得 $A\cap B=\{y\mid1\leqslant y\leqslant4\}$.}
\solend
\fi

\item 若 $x\in\{1,x^2\}$，则 $x=$\blankline[4.4em].

\ifanswer
\ans{$0$}
\fi

\item 已知等式 $2x^2+3x+5=a(2x+1)(x+1)+c$，则 $a+c=$\blankline[4.4em].

\ifanswer
\solbegin
\step{因为 $2x^2+3x+5=a(2x+1)(x+1)+c$ 恒成立，}
\step{所以 $2x^2+3x+5=2ax^2+3ax+a+c$ 恒成立，则 $\begin{cases}2=2a\\ 3=3a\\ 5=a+c\end{cases}$，解得 $a+c=5$.}
\solend
\fi

\item 已知 $a\in\R$，设 $p:x<2$，$q:x<5a-3$，且 $p$ 是 $q$ 的必要非充分条件，则 $a$ 的取值范围是\blankline[4.4em].

\ifanswer
\solbegin
\step{设 $p:x<2$，$q:x<5a-3$，且 $p$ 是 $q$ 的必要非充分条件，}
\step{则 $\{x\mid x<5a-3\}\subset\{x\mid x<2\}$，所以 $5a-3<2$，即 $a<1$.}
\solend
\fi

\item 已知 $a\in\R$，集合 $A=\{x\mid ax^2+3x-1=0\}$，若集合 $A$ 中只有一个元素，则 $a$ 的取值范围是\blankline[4.4em].

\ifanswer
\solbegin
\step{当 $a=0$ 时，则有 $\{x\mid ax^2+3x-1=0\}=\{x\mid3x-1=0\}=\left\{\dfrac13\right\}$，符合题意；}
\step{当 $a\neq0$ 时，由题意得 $\Delta=9+4a=0$，解得 $a=-\dfrac94$. 符合题意；}
\step{综上所述，实数 $a$ 的取值集合为 $\left\{0,-\dfrac94\right\}$.}
\solend
\fi

\item 设集合 $A=\{x\mid x^2+2(m+1)x+m^2-1=0,m\in\R\}$，且 $A\subseteq\{0,-4\}$，则实数 $m$ 的取值范围是\blankline[4.4em].

\ifanswer
\solbegin
\step{由题意得 $A=\varnothing$ 或 $A=\{0\}$ 或 $A=\{-4\}$ 或 $A=\{0,-4\}$，}
\step{当 $A=\varnothing$ 时，方程 $x^2+2(m+1)x+m^2-1=0$ 无实根，}
\step{则 $\Delta=[2(m+1)]^2-4\times1\times(m^2-1)=8m+8<0$，解得 $m<-1$，}
\step{当 $A=\{0\}$ 时，方程 $x^2+2(m+1)x+m^2-1=0$ 有一个实根 $0$，}
\step{则 $\begin{cases}\Delta=8m+8=0\\ m^2-1=0\end{cases}$，解得 $m=-1$，}
\step{当 $A=\{-4\}$ 时，方程 $x^2+2(m+1)x+m^2-1=0$ 有一个实根 $-4$，}
\step{则 $\begin{cases}\Delta=8m+8=0\\ 16-8(m+1)+m^2-1=0\end{cases}$，无解，}
\step{当 $A=\{0,-4\}$ 时，方程 $x^2+2(m+1)x+m^2-1=0$ 有两个实根 $-4$ 和 $0$，}
\step{则 $\begin{cases}\Delta=8m+8>0\\ 16-8(m+1)+m^2-1=0\\ m^2-1=0\end{cases}$，解得 $m=1$，}
\step{综上所述，实数 $m$ 的取值范围是 $\{m\mid m\leqslant-1\text{ 或 }m=1\}$.}
\solend
\fi

\item 对于任意实数 $x$，$[x]$ 表示不超过 $x$ 的最大整数，如 $[1.1]=1$，$[-1.2]=-2$，若
$A=\{y\mid y=[2x]+[3x],-1\leqslant x\leqslant2\}$，则 $A$ 中所有元素的和为\blankline[4.4em].

\ifanswer
\solbegin
% 原卷“法一：”“法二：”标号排红色，本稿按 高一53“情况1：”成例排 \stepm（加粗着色），见文件头
\stepm{法一：}{当 $-1\leqslant x<-\dfrac23$ 时，$[2x]=-2$，$[3x]=-3$，$f(x)=[2x]+[3x]=-5$，}
\step{当 $-\dfrac23\leqslant x<-\dfrac12$ 时，$[2x]=-2$，$[3x]=-2$，$f(x)=[2x]+[3x]=-4$；}
\step{当 $-\dfrac12\leqslant x<-\dfrac13$ 时，$[2x]=-1$，$[3x]=-2$，$f(x)=[2x]+[3x]=-3$；}
\step{当 $-\dfrac13\leqslant x<0$ 时，$[2x]=-1$，$[3x]=-1$，$f(x)=[2x]+[3x]=-2$；}
\step{当 $0\leqslant x<\dfrac13$ 时，$[2x]=0$，$[3x]=0$，$f(x)=[2x]+[3x]=0$；}
\step{当 $\dfrac13\leqslant x<\dfrac12$ 时，$[2x]=0$，$[3x]=1$，$f(x)=[2x]+[3x]=1$；}
\step{当 $\dfrac12\leqslant x<\dfrac23$ 时，$[2x]=1$，$[3x]=1$，$f(x)=[2x]+[3x]=2$；}
\step{当 $\dfrac23\leqslant x<1$ 时，$[2x]=1$，$[3x]=2$，$f(x)=[2x]+[3x]=3$；}
\step{当 $1\leqslant x<\dfrac43$ 时，$[2x]=2$，$[3x]=3$，$f(x)=[2x]+[3x]=5$；}
\step{当 $\dfrac43\leqslant x<\dfrac32$ 时，$[2x]=2$，$[3x]=4$，$f(x)=[2x]+[3x]=6$；}
\step{当 $\dfrac32\leqslant x<\dfrac53$ 时，$[2x]=3$，$[3x]=4$，$f(x)=[2x]+[3x]=7$；}
\step{当 $\dfrac53\leqslant x<2$ 时，$[2x]=3$，$[3x]=5$，$f(x)=[2x]+[3x]=8$；}
\step{当 $x=2$ 时，$[2x]=4$，$[3x]=6$，$f(x)=[2x]+[3x]=10$}
\step{故 $A$ 中所有元素的和为 $-5-4-3-2+0+1+2+3+5+6+7+8+10=28$.}
% 原卷“法二：”标号排红色，同上
\stepm{法二：}{设 $x=m+n$，其中 $m\in\Z$，$n\in[0,1)$，}
\step{若 $0\leqslant n<\dfrac13$，则 $0\leqslant2n<\dfrac23$，$0\leqslant3n<1$，}
\step{得 $[2x]=[2m+2n]=2m$，$[3x]=[3m+3n]=3m$，所以 $y=[2x]+[3x]=5m$；}
\step{若 $\dfrac13\leqslant n<\dfrac12$，则 $\dfrac23\leqslant2n<1$，$1\leqslant3n<\dfrac32$，}
\step{得 $[2x]=[2m+2n]=2m$，$[3x]=[3m+3n]=3m+1$，}
\step{所以 $y=[2x]+[3x]=5m+1$；}
\step{若 $\dfrac12\leqslant n<\dfrac23$，则 $1\leqslant2n<\dfrac43$，$\dfrac32\leqslant3n<2$，}
\step{得 $[2x]=[2m+2n]=2m+1$，$[3x]=[3m+3n]=3m+1$，}
\step{所以 $y=[2x]+[3x]=5m+2$；}
\step{若 $\dfrac23\leqslant n<1$，则 $\dfrac43\leqslant2n<2$，$2\leqslant3n<3$，}
\step{得 $[2x]=[2m+2n]=2m+1$，$[3x]=[3m+3n]=3m+2$，}
\step{所以 $y=[2x]+[3x]=5m+3$；}
\step{综上所述，当 $n\in[0,1)$ 时，所有函数值的和}
\step{为 $5m+5m+1+5m+2+5m+3=20m+6$，}
\step{又因为 $x\in[-1,2]$，此时 $m=-1,0,1$，且当 $x=2$ 时，$y=[2x]+[3x]=10$，}
\step{所有函数值的和为 $(-14)+6+26+10=28$，即 $A$ 中所有元素之和为 $28$.}
\solend
\fi

\end{enumerate}

\bigsec{二、选择题}

\begin{enumerate}
\setcounter{enumi}{12}

\item 下列写法中，正确的是\mbox{（\quad）}

\keepwithprev
A. $0\in\{\varnothing\}$\qquad B. $\varnothing\in\{\varnothing\}$

C. $\varnothing=\{\varnothing\}$\qquad D. $0\in\varnothing$

\ifanswer
\solbegin
\step{$0\notin\{\varnothing\}$，$\varnothing\in\{\varnothing\}$，故 $A$ 错误，$B$ 正确；}
\step{$\varnothing$ 是不含任何元素的集合，所以 $\varnothing\neq\{\varnothing\}$，$0\notin\varnothing$，故 $CD$ 错误.}
\step{故选 $B$.}
\solend
\fi

% 第 14 题答案“D”直接印在题干括号内“(  D  )”，且该题无【答案】/【解析】段；
% 按 高一37 第 13 题的成例括号留空、答案另提为 \ans{}（见文件头）
\item 若 $a<b<0$，则下列不等式中不能成立的是\mbox{（\quad）}

\keepwithprev
A. $\dfrac1a>\dfrac1b$\qquad B. $a-b>a$

C. $|a|>|b|$\qquad D. $a^3>b^3$

\ifanswer
\ans{$D$}
\fi

\item 已知 $x$、$y$、$z$ 为非零实数，代数式 $\dfrac x{|x|}+\dfrac y{|y|}+\dfrac z{|z|}+\dfrac{xyz}{|xyz|}$
的值所组成的集合是 $M$，则下列判断正确的是\mbox{（\quad）}

\keepwithprev
A. $-4\notin M$\qquad B. $2\in M$

C. $4\notin M$\qquad D. $0\in M$

\ifanswer
\solbegin
\step{当 $x$、$y$、$z$ 均为正数时，$xyz>0$，$\dfrac x{|x|}+\dfrac y{|y|}+\dfrac z{|z|}+\dfrac{xyz}{|xyz|}=1+1+1+1=4$，}
\step{当 $x$、$y$、$z$ 有 2 个正数，1 个负数时，不妨设 $x>0$，$y>0$，$z<0$，}
\step{则 $xyz<0$，$\dfrac x{|x|}+\dfrac y{|y|}+\dfrac z{|z|}+\dfrac{xyz}{|xyz|}=1+1-1-1=0$，}
\step{若 $x$、$y$、$z$ 有 1 个正数，2 个负数时，不妨设 $x>0$，$y<0$，$z<0$，}
\step{则 $xyz>0$，$\dfrac x{|x|}+\dfrac y{|y|}+\dfrac z{|z|}+\dfrac{xyz}{|xyz|}=1-1-1+1=0$，}
\step{当 $x$、$y$、$z$ 均为负数时，$xyz<0$，}
\step{$\dfrac x{|x|}+\dfrac y{|y|}+\dfrac z{|z|}+\dfrac{xyz}{|xyz|}=-1-1-1-1=-4$，}
\step{综上，$M=\{-4,0,4\}$，}
\step{所以 $-4\in M$，$A$ 错误；$2\notin M$，$B$ 错误；$4\in M$，$C$ 错误；}
\step{$0\in M$，$D$ 正确. 故选 $D$.}
\solend
\fi

\item 若 $xy>0$，且 $x+y=9$，则 $\dfrac1x+\dfrac4y$ 的最小值是\mbox{（\quad）}

\keepwithprev
A. $9$\qquad B. $\dfrac59$

C. $\dfrac49$\qquad D. $1$

\ifanswer
\solbegin
\step{$xy>0$，且 $x+y=9$，所以 $x>0$，$y>0$，}
\step{所以 $\dfrac1x+\dfrac4y=\dfrac19\left(\dfrac1x+\dfrac4y\right)(x+y)=\dfrac59+\dfrac19\left(\dfrac yx+\dfrac{4x}y\right)\geqslant\dfrac59+\dfrac29\cdot\sqrt{\dfrac yx\cdot\dfrac{4x}y}=\dfrac59+\dfrac49=1$，}
\step{当且仅当 $\dfrac yx=\dfrac{4x}y$，即 $x=3$，$y=6$ 时取等号，所以 $\dfrac1x+\dfrac4y$ 的最小值为 $1$.}
\step{故选 $D$.}
\solend
\fi

\end{enumerate}

\bigsec{三、解答题}

\begin{enumerate}
\setcounter{enumi}{16}

\item 求解下列不等式：

（1）$\dfrac{x-3}{x+1}\leqslant0$；

（2）$|2x-1|<3$.

\ifanswer
\solbegin
\stepm{(1)}{$\dfrac{x-3}{x+1}\leqslant0$，变形得 $(x-3)\cdot(x+1)\leqslant0$ 且 $x\neq-1$，}
\step{解得 $-1<x\leqslant3$，即不等式的解集为 $(-1,3]$.}
\stepm{(2)}{若 $|2x-1|<3$，则 $-3<2x-1<3$，解得 $-1<x<2$，即 $(-1,2)$.}
\solend
\fi

\ansspace[6]

\item 已知集合 $A=\{x\mid-2<x<15\}$,$B=\{x\mid m-6<x<2m-1,m\in\R\}$.

（1）当 $m=2$ 时，求 $A\cap B$；

（2）若 $A\cup B=A$，求实数 $m$ 的取值范围.

\ifanswer
\solbegin
\stepm{(1)}{当 $m=2$ 时，$A=\{x\mid-2<x<15\}$,$B=\{x\mid-4<x<3\}$，}
\step{此时 $A\cap B=\{x\mid-2<x<3\}$.}
\stepm{(2)}{由 $A\cup B=A$ 得 $B\subseteq A$.}
\step{当 $B=\varnothing$,$m-6\geqslant2m-1$，解得 $m\leqslant-5$；}
\step{当 $B\neq\varnothing$，结合 $B\subseteq A$ 得 $\begin{cases}m-6<2m-1\\ 2m-1\leqslant15\\ m-6\geqslant-2\end{cases}$，解得 $4\leqslant m\leqslant8$.}
\step{综上，$m\leqslant-5$ 或 $4\leqslant m\leqslant8$.}
\solend
\fi

\ansspace[6]

\item 已知 $a\in\R$，命题 $p$：关于 $x$ 的方程 $x^2-2x+a=0$ 无实根；命题 $q$：对任意 $x\in\R$，
不等式 $ax^2-ax+1>0$ 恒成立.

（1）若 $q$ 为真命题，求 $a$ 的取值范围；

（2）若 $p$、$q$ 中有且只有一个为真命题，求 $a$ 的取值范围.

\ifanswer
\solbegin
\stepm{(1)}{当 $q$ 为真命题时，}
\stepm{①}{当 $a=0$ 时，不等式化为 $1>0$，对任意 $x\in\R$ 恒成立，满足条件；}
\stepm{②}{当 $a\neq0$ 时，$ax^2-ax+1>0$ 为一元二次不等式，}
\step{若对任意 $x\in\R$ 恒成立，需满足 $\begin{cases}a>0\\ \Delta=(-a)^2-4a<0\end{cases}$，解得 $0<a<4$；}
\step{综上，$a$ 的取值范围为 $[0,4)$.}
\stepm{(2)}{$p$ 为真时，方程 $x^2-2x+a=0$ 无实根，}
\step{故判别式 $\Delta=4-4a<0$，解得 $a>1$，因此 $p$ 为假时，$a\leqslant1$；}
\step{由（1）得 $q$ 为假时 $a\geqslant4$ 或 $a<0$，}
\step{$p$、$q$ 有且只有一个真命题，分两类：}
\step{当 $p$ 假 $q$ 真时，满足 $\begin{cases}a\leqslant1\\ 0\leqslant a<4\end{cases}$，解得 $0\leqslant a\leqslant1$；}
\step{当 $p$ 真 $q$ 假时，满足 $\begin{cases}a>1\\ a\geqslant4\text{或}a<0\end{cases}$，解得 $a\geqslant4$.}
\step{综上，$a$ 的取值范围是 $[0,1]\cup[4,+\infty)$.}
\solend
\fi

\ansspace[8]

\item 某工厂生产 $x$($x>0$) 件产品的总成本 $C$（单位：元）与生产件数 $x$ 之间满足二次函数
$C(x)=\dfrac14x^2+50x+10000$.

（1）若要求生产总成本不超过 30000 元，求满足生产总成本要求的产品件数的取值范围；

（2）试求该工厂生产多少件产品时，可使每件产品的平均成本最小.

\ifanswer
\solbegin
\stepm{(1)}{若要求生产总成本不超过 30000 元，}
\step{则 $C(x)=\dfrac14x^2+50x+10000\leqslant30000$，}
\step{即 $x^2+200x-80000\leqslant0$，解得 $-400\leqslant x\leqslant200$，又 $x>0$，}
\step{所以满足生产总成本要求的产品件数的取值范围为 $(0,200]$，$x\in\N$；}
\stepm{(2)}{生产件数 $x$ 时每件的平均成本为 $\dfrac{\frac14x^2+50x+10000}x=\dfrac14x+50+\dfrac{10000}x$，}
\step{由基本不等式得 $\dfrac14x+50+\dfrac{10000}x\geqslant2\sqrt{\dfrac14x\cdot\dfrac{10000}x}+50=150$，}
\step{当且仅当 $\dfrac14x=\dfrac{10000}x$，即 $x=200$ 时取等号，}
\step{即该工厂生产 200 件产品时，每件产品的平均成本最小.}
\solend
\fi

\ansspace[8]

% 原卷题干“$(n\geqslant2,n\in N)$”后连印两个全角逗号“，，”，照录未改（见文件头“原卷笔误”①）
\item 已知有限集 $A=\{a_1,a_2,\cdots,a_n\}$($n\geqslant2$,$n\in\N$)，，如果 $A$ 中元素
$a_i$($i=1,2,3,\cdots,n$) 满足 $a_1+a_2+a_3+\cdots+a_n=a_1\times a_2\times\cdots\times a_n$，
就称 $A$ 为“和积平衡集”.

% 原卷 (2) 将“写方程”与“反证法证明”两项任务并作一句（“……表示），可尝试用反证法证明：……”），
% 照录未改（见文件头“原卷笔误”②）
（1）判断集合 $\{-1+\sqrt3,-1-\sqrt3\}$ 是否为“和积平衡集”；

（2）若 $a_1$、$a_2$ 是两个不同的正数，且 $\{a_1,a_2\}$ 是“和积平衡集”，写出以 $a_1$、$a_2$ 为根的
一个一元二次方程（系数可用 $a_1$、$a_2$ 表示），可尝试用反证法证明：$a_1$、$a_2$ 至少有一个大于 2；

（3）若 $a_i\geqslant1$，且 $a_i\in\N$，求所有符合条件的“和积平衡集”.

\ifanswer
\solbegin
\stepm{(1)}{对于集合 $\{-1+\sqrt3,-1-\sqrt3\}$，}
\step{因为 $(-1+\sqrt3)+(-1-\sqrt3)=-2$，$(-1+\sqrt3)\times(-1-\sqrt3)=-2$，}
\step{所以满足 $a_1+a_2=a_1\times a_2$，满足题意，}
\step{所以集合 $\{-1+\sqrt3,-1-\sqrt3\}$ 是和积平衡集.}
\stepm{(2)}{以 $a_1$、$a_2$ 为根的一个一元二次方程可为 $x^2-(a_1+a_2)x+a_1\cdot a_2=0$，}
\step{因为 $a_1$、$a_2$ 是两个不同的正数，且 $\{a_1,a_2\}$ 是和积平衡集，}
\step{所以 $a_1+a_2=a_1\cdot a_2$，}
\step{又 $a_1$ 和 $a_2$ 为 $x^2-(a_1+a_2)x+a_1\cdot a_2=0$ 的两个不同的正根，}
\step{所以 $\Delta=(a_1+a_2)^2-4a_1a_2=(a_1+a_2)^2-4(a_1+a_2)>0$，}
\step{又 $a_1+a_2>0$，则 $a_1+a_2>4$，}
\step{若 $a_1$、$a_2$ 都小于等于 2，则 $a_1+a_2<4$，矛盾，}
\step{所以 $a_1$、$a_2$ 至少有一个大于 2；}
\stepm{(3)}{设 $A$ 中的 $a_1<a_2<a_3<\cdots<a_n$，且 $a_i\in\N^*$，}
\step{由 $a_1+a_2+a_3+\cdots+a_n=a_1\times a_2\times\cdots\times a_n$，得 $n=2$ 时明显不满足，}
\step{所以 $n\geqslant3$，故 $a_1\times a_2\times\cdots\times a_n=a_1+a_2+a_3+\cdots+a_n\leqslant na_n$，}
\step{所以 $a_1\times a_2\times\cdots\times a_{n-1}\leqslant n$，}
\step{当 $n=3$ 时，$a_1\times a_2\leqslant3$，$a_i\in\N^*$，只能有 $a_1=1$，$a_2=2$，}
\step{由 $a_1+a_2+a_3=a_1\times a_2\times a_3$，得 $a_3=3$，此时的和积平衡集为 $\{1,2,3\}$，}
\step{当 $n\geqslant4$ 时，$a_1\times a_2\times\cdots\times a_{n-1}\geqslant(n-1)(n-2)\times\cdots\times2\times1\geqslant(n-1)(n-2)$，}
\step{又 $(n-1)(n-2)-n=n^2-4n+2=n(n-4)+2>0$，即 $(n-1)(n-2)>n$，}
\step{即 $a_1\times a_2\times\cdots\times a_{n-1}>n$，与 $a_1\times a_2\times\cdots\times a_{n-1}\leqslant n$ 矛盾，}
\step{所以 $n\geqslant4$ 不满足条件，}
\step{所以符合条件的和积平衡集有且只有 $n=3$，此时的和积平衡集为 $\{1,2,3\}$.}
\solend
\fi

\ansspace*[10]

\end{enumerate}

\end{document}
"
%
% 原始材料：微信公众号“魔都数学加油站”文章，作者破晓，连载“27学年高一63”；
% 原文标题“27学年高一63——2026-2027年上海市吴淞中学高一上10月月考”；
% 原文链接：https://mp.weixin.qq.com/s/TFQ4TLnt_l66t7t51es4OQ。
% 页面标注发布时间 2026-10-16 21:44（上海）。
% 正文为 10 张 PNG 页（`原图/高一63_吴淞中学_10月月考/01..10.png`）：
% 第 06、10 页 2560×3620，第 01、02、03、04、05、07、08、09 页 4762×6735
% （**同一篇各页分辨率不同**），均按页序读取。10 页全为试卷正文页（卷首标题在第 1 页页顶，
% 卷末解答题第 21 题解析止于第 10 页），无封面页、目录页、二维码或宣传页。
% 粉色斜排“魔都数学加油站”为卷面水印，与公众号页脚一并未录入。
%
% 卷首标题：照卷面印作“2026--2027 年吴淞中学高一上 10 月月考”（公众号标题中的“上海市”
% 卷面标题行没有，本稿照卷面），年份区间按全稿惯例排 en dash。原卷只有一行标题、无副标题。
%
% 篇幅与题号：原卷自第 3 题起编号（第 1、2 题不在本卷内）。填空题 3—12（10 题）、
% 选择题 13—16（4 题）、解答题 17—21（5 题），共 19 题。本稿沿用原节与原题号，
% 只改排为全稿统一的大节样式（原卷小节标题为“一、填空题”“二、选择题”“三、解答题”）。
%
% 讲评版：原卷每题下直接印【答案】或【解析】。**只印【答案】的 2 题**（第 4、7 题），
% 本稿只给 \ans{}；其余 17 题里 16 题都印【解析】（按全稿口径这类题不另提 \ans{}——第 13、15、16 题
% 的结论写在解析末句（“故选 B”等），亦照录在解析内；第 14 题则既无【答案】也无【解析】段。**另有 1 题答案印在题干括号内**——
% 第 14 题（“(  D  )”，且该题无【答案】/【解析】段），本稿按 高一37 第 13 题成例把括号留空、
% 答案另提为 \ans{$D$}。第 16 题题干括号为空但印有【解析】，不算括号内答案。
%
% 副标题：原卷未印副标题。本稿按“知识模块”口径标作“集合与常用逻辑用语、不等式”：
% 19 题中集合与常用逻辑用语 12 题（第 3、4、6、7、9、10、11、12、13、15、18、21 题），
% 不等式 7 题（含等式（韦达定理、恒等式）与不等式性质、基本不等式、二次不等式、恒成立：
% 第 5、8、14、16、17、19、20 题），两类均超过三分之一，故并标。口径同 高一59。
%
% 与原卷的排版差异：
%   ① 原卷每题下直接印【答案】/【解析】，本稿按全稿惯例将答案与解析置于答案版，试卷版保持干净；
%   ② 原卷未印副标题，本稿按“知识模块”口径补出；
%   ③ 原卷填空横线约四个字宽，本稿按全稿惯例统一排 \blankline[4.4em]；
%   ④ 原卷实数集排斜体 $R$（第 9、10、19 题题干与解析）或粗体 $\mathbf{R}$（第 18 题题干），
%      整数集排斜体 $Z$（第 12 题解析），自然数集排斜体 $N$（第 20 题解析、第 21 题题干
%      “$(n\geqslant2,n\in N)$”）或粗体 $\mathbf{N}$（第 21 题 (3) 题干）、斜体 $N^*$
%      （第 21 题 (3) 解析），本稿按全稿惯例统一改用 $\R$、$\Z$、$\N$、$\N^*$
%      （同 高一14 排版差异④：记号归一，粗体／斜体不作区分，$N$ 与 $N^*$ 的区别照原卷保留）；
%      空集记号原卷作 $\varnothing$，照录；
%   ⑤ 原卷在数学式内多处用半角逗号（如 $\{y\mid y=[2x]+[3x],-1\leqslant x\leqslant2\}$、
%      $m=-1,0,1$、$A=\{x\mid-2<x<15\}$,$B=\cdots$），均照录未强改；
%      句末句点按全稿惯例排西文句点；“和积平衡集”自定义名词照录；
%   ⑥ 第 12 题解析里的“法一：”“法二：”标号原卷排**红色**（全卷共 2 处），
%      同 高一53“情况1：”成例，本稿排 \stepm{法一：}{…}、\stepm{法二：}{…}
%      （加粗着色、不改黑色基调），冒号照原卷排全角；
%   ⑦ 第 17—21 题的子问原卷各自另起一行，本稿按全稿惯例用 enumerate 的 (1)(2)(3)；
%      解析中的子问标号一律用 \stepm{(1)}{…}；第 19 题解析中的①②分情形，按全稿惯例作
%      \stepm{①}{…}、\stepm{②}{…}；
%   ⑧ 第 14 题印在题干括号内的答案移入 \ans{}，见上；
%   ⑨ 每道解答题原卷未留作答空白，本稿按全稿惯例加 \ansspace；
%   ⑩ 选择题的作答括号原卷排半角“(　　)”（第 13、15、16 题括号内空白、第 14 题括号内印 D），
%      本稿按全稿惯例排 \mbox{（\quad）}。
%
% 原卷笔误/存疑（均照录未改，正文对应处另有 % 注）：
%   ① 第 21 题题干“$(n\geqslant2$,$n\in N)$”后连印两个全角逗号“，，”（衍字，
%      8 倍放大确认，非水印干扰），本稿照录；
%   ② 第 21 题 (2) 题干将“写出以 $a_1$、$a_2$ 为根的一个一元二次方程”与“用反证法证明：
%      $a_1$、$a_2$ 至少有一个大于 2”两项任务并作一句、之间仅以逗号相接（“……表示），
%      可尝试用反证法证明：……”），读来像抄并两小问，原卷如此，照录未改。
%
% 插图：无（全卷 10 页无任何图形、无“如图”字样）。
%
\documentclass[a4paper,11pt]{article}
\usepackage{common}

\begin{document}

\examtitlest[3]{2026--2027 年吴淞中学高一上 10 月月考}{集合与常用逻辑用语、不等式}

\bigsec{一、填空题}

\begin{enumerate}
\setcounter{enumi}{2}

\item 满足 $A\subseteq\{1,2,3,4\}$ 的集合 $A$ 共有\blankline[4.4em]个.

\ifanswer
\solbegin
\step{由题意，集合 $A$ 是集合 $\{1,2,3,4\}$ 的子集，共有 $2^4=16$ 个.}
\solend
\fi

\item 陈述句“$a>1$ 或 $b>3$”的否定形式为\blankline[4.4em].

\ifanswer
\ans{$a\leqslant1$ 且 $b\leqslant3$}
\fi

\item 已知方程 $3x^2+4x-4=0$ 的两个根为 $x_1$、$x_2$，则 $x_1^2+x_2^2=$\blankline[4.4em].

\ifanswer
\solbegin
\step{对于方程 $3x^2+4x-4=0$，因为 $\Delta=16-4\times3\times(-4)=64>0$，}
\step{所以方程 $3x^2+4x-4=0$ 的两个不相等的实根为 $x_1$、$x_2$，}
\step{由韦达定理可得 $x_1+x_2=-\dfrac43$，$x_1x_2=-\dfrac43$，}
\step{所以 $x_1^2+x_2^2=(x_1+x_2)^2-2x_1x_2=\dfrac{16}9+\dfrac83=\dfrac{40}9$.}
\solend
\fi

\item 设集合 $A=\{y\mid y=x^2+1\}$,$B=\{y\mid y=-x^2+2x+3\}$，则 $A\cap B=$\blankline[4.4em].

\ifanswer
\solbegin
\step{因为 $y=x^2+1\geqslant1$,$y=-x^2+2x+3=-(x-1)^2+4\leqslant4$，}
\step{所以 $A=\{y\mid y\geqslant1\}$,$B=\{y\mid y\leqslant4\}$，得 $A\cap B=\{y\mid1\leqslant y\leqslant4\}$.}
\solend
\fi

\item 若 $x\in\{1,x^2\}$，则 $x=$\blankline[4.4em].

\ifanswer
\ans{$0$}
\fi

\item 已知等式 $2x^2+3x+5=a(2x+1)(x+1)+c$，则 $a+c=$\blankline[4.4em].

\ifanswer
\solbegin
\step{因为 $2x^2+3x+5=a(2x+1)(x+1)+c$ 恒成立，}
\step{所以 $2x^2+3x+5=2ax^2+3ax+a+c$ 恒成立，则 $\begin{cases}2=2a\\ 3=3a\\ 5=a+c\end{cases}$，解得 $a+c=5$.}
\solend
\fi

\item 已知 $a\in\R$，设 $p:x<2$，$q:x<5a-3$，且 $p$ 是 $q$ 的必要非充分条件，则 $a$ 的取值范围是\blankline[4.4em].

\ifanswer
\solbegin
\step{设 $p:x<2$，$q:x<5a-3$，且 $p$ 是 $q$ 的必要非充分条件，}
\step{则 $\{x\mid x<5a-3\}\subset\{x\mid x<2\}$，所以 $5a-3<2$，即 $a<1$.}
\solend
\fi

\item 已知 $a\in\R$，集合 $A=\{x\mid ax^2+3x-1=0\}$，若集合 $A$ 中只有一个元素，则 $a$ 的取值范围是\blankline[4.4em].

\ifanswer
\solbegin
\step{当 $a=0$ 时，则有 $\{x\mid ax^2+3x-1=0\}=\{x\mid3x-1=0\}=\left\{\dfrac13\right\}$，符合题意；}
\step{当 $a\neq0$ 时，由题意得 $\Delta=9+4a=0$，解得 $a=-\dfrac94$. 符合题意；}
\step{综上所述，实数 $a$ 的取值集合为 $\left\{0,-\dfrac94\right\}$.}
\solend
\fi

\item 设集合 $A=\{x\mid x^2+2(m+1)x+m^2-1=0,m\in\R\}$，且 $A\subseteq\{0,-4\}$，则实数 $m$ 的取值范围是\blankline[4.4em].

\ifanswer
\solbegin
\step{由题意得 $A=\varnothing$ 或 $A=\{0\}$ 或 $A=\{-4\}$ 或 $A=\{0,-4\}$，}
\step{当 $A=\varnothing$ 时，方程 $x^2+2(m+1)x+m^2-1=0$ 无实根，}
\step{则 $\Delta=[2(m+1)]^2-4\times1\times(m^2-1)=8m+8<0$，解得 $m<-1$，}
\step{当 $A=\{0\}$ 时，方程 $x^2+2(m+1)x+m^2-1=0$ 有一个实根 $0$，}
\step{则 $\begin{cases}\Delta=8m+8=0\\ m^2-1=0\end{cases}$，解得 $m=-1$，}
\step{当 $A=\{-4\}$ 时，方程 $x^2+2(m+1)x+m^2-1=0$ 有一个实根 $-4$，}
\step{则 $\begin{cases}\Delta=8m+8=0\\ 16-8(m+1)+m^2-1=0\end{cases}$，无解，}
\step{当 $A=\{0,-4\}$ 时，方程 $x^2+2(m+1)x+m^2-1=0$ 有两个实根 $-4$ 和 $0$，}
\step{则 $\begin{cases}\Delta=8m+8>0\\ 16-8(m+1)+m^2-1=0\\ m^2-1=0\end{cases}$，解得 $m=1$，}
\step{综上所述，实数 $m$ 的取值范围是 $\{m\mid m\leqslant-1\text{ 或 }m=1\}$.}
\solend
\fi

\item 对于任意实数 $x$，$[x]$ 表示不超过 $x$ 的最大整数，如 $[1.1]=1$，$[-1.2]=-2$，若
$A=\{y\mid y=[2x]+[3x],-1\leqslant x\leqslant2\}$，则 $A$ 中所有元素的和为\blankline[4.4em].

\ifanswer
\solbegin
% 原卷“法一：”“法二：”标号排红色，本稿按 高一53“情况1：”成例排 \stepm（加粗着色），见文件头
\stepm{法一：}{当 $-1\leqslant x<-\dfrac23$ 时，$[2x]=-2$，$[3x]=-3$，$f(x)=[2x]+[3x]=-5$，}
\step{当 $-\dfrac23\leqslant x<-\dfrac12$ 时，$[2x]=-2$，$[3x]=-2$，$f(x)=[2x]+[3x]=-4$；}
\step{当 $-\dfrac12\leqslant x<-\dfrac13$ 时，$[2x]=-1$，$[3x]=-2$，$f(x)=[2x]+[3x]=-3$；}
\step{当 $-\dfrac13\leqslant x<0$ 时，$[2x]=-1$，$[3x]=-1$，$f(x)=[2x]+[3x]=-2$；}
\step{当 $0\leqslant x<\dfrac13$ 时，$[2x]=0$，$[3x]=0$，$f(x)=[2x]+[3x]=0$；}
\step{当 $\dfrac13\leqslant x<\dfrac12$ 时，$[2x]=0$，$[3x]=1$，$f(x)=[2x]+[3x]=1$；}
\step{当 $\dfrac12\leqslant x<\dfrac23$ 时，$[2x]=1$，$[3x]=1$，$f(x)=[2x]+[3x]=2$；}
\step{当 $\dfrac23\leqslant x<1$ 时，$[2x]=1$，$[3x]=2$，$f(x)=[2x]+[3x]=3$；}
\step{当 $1\leqslant x<\dfrac43$ 时，$[2x]=2$，$[3x]=3$，$f(x)=[2x]+[3x]=5$；}
\step{当 $\dfrac43\leqslant x<\dfrac32$ 时，$[2x]=2$，$[3x]=4$，$f(x)=[2x]+[3x]=6$；}
\step{当 $\dfrac32\leqslant x<\dfrac53$ 时，$[2x]=3$，$[3x]=4$，$f(x)=[2x]+[3x]=7$；}
\step{当 $\dfrac53\leqslant x<2$ 时，$[2x]=3$，$[3x]=5$，$f(x)=[2x]+[3x]=8$；}
\step{当 $x=2$ 时，$[2x]=4$，$[3x]=6$，$f(x)=[2x]+[3x]=10$}
\step{故 $A$ 中所有元素的和为 $-5-4-3-2+0+1+2+3+5+6+7+8+10=28$.}
% 原卷“法二：”标号排红色，同上
\stepm{法二：}{设 $x=m+n$，其中 $m\in\Z$，$n\in[0,1)$，}
\step{若 $0\leqslant n<\dfrac13$，则 $0\leqslant2n<\dfrac23$，$0\leqslant3n<1$，}
\step{得 $[2x]=[2m+2n]=2m$，$[3x]=[3m+3n]=3m$，所以 $y=[2x]+[3x]=5m$；}
\step{若 $\dfrac13\leqslant n<\dfrac12$，则 $\dfrac23\leqslant2n<1$，$1\leqslant3n<\dfrac32$，}
\step{得 $[2x]=[2m+2n]=2m$，$[3x]=[3m+3n]=3m+1$，}
\step{所以 $y=[2x]+[3x]=5m+1$；}
\step{若 $\dfrac12\leqslant n<\dfrac23$，则 $1\leqslant2n<\dfrac43$，$\dfrac32\leqslant3n<2$，}
\step{得 $[2x]=[2m+2n]=2m+1$，$[3x]=[3m+3n]=3m+1$，}
\step{所以 $y=[2x]+[3x]=5m+2$；}
\step{若 $\dfrac23\leqslant n<1$，则 $\dfrac43\leqslant2n<2$，$2\leqslant3n<3$，}
\step{得 $[2x]=[2m+2n]=2m+1$，$[3x]=[3m+3n]=3m+2$，}
\step{所以 $y=[2x]+[3x]=5m+3$；}
\step{综上所述，当 $n\in[0,1)$ 时，所有函数值的和}
\step{为 $5m+5m+1+5m+2+5m+3=20m+6$，}
\step{又因为 $x\in[-1,2]$，此时 $m=-1,0,1$，且当 $x=2$ 时，$y=[2x]+[3x]=10$，}
\step{所有函数值的和为 $(-14)+6+26+10=28$，即 $A$ 中所有元素之和为 $28$.}
\solend
\fi

\end{enumerate}

\bigsec{二、选择题}

\begin{enumerate}
\setcounter{enumi}{12}

\item 下列写法中，正确的是\mbox{（\quad）}

\keepwithprev
A. $0\in\{\varnothing\}$\qquad B. $\varnothing\in\{\varnothing\}$

C. $\varnothing=\{\varnothing\}$\qquad D. $0\in\varnothing$

\ifanswer
\solbegin
\step{$0\notin\{\varnothing\}$，$\varnothing\in\{\varnothing\}$，故 $A$ 错误，$B$ 正确；}
\step{$\varnothing$ 是不含任何元素的集合，所以 $\varnothing\neq\{\varnothing\}$，$0\notin\varnothing$，故 $CD$ 错误.}
\step{故选 $B$.}
\solend
\fi

% 第 14 题答案“D”直接印在题干括号内“(  D  )”，且该题无【答案】/【解析】段；
% 按 高一37 第 13 题的成例括号留空、答案另提为 \ans{}（见文件头）
\item 若 $a<b<0$，则下列不等式中不能成立的是\mbox{（\quad）}

\keepwithprev
A. $\dfrac1a>\dfrac1b$\qquad B. $a-b>a$

C. $|a|>|b|$\qquad D. $a^3>b^3$

\ifanswer
\ans{$D$}
\fi

\item 已知 $x$、$y$、$z$ 为非零实数，代数式 $\dfrac x{|x|}+\dfrac y{|y|}+\dfrac z{|z|}+\dfrac{xyz}{|xyz|}$
的值所组成的集合是 $M$，则下列判断正确的是\mbox{（\quad）}

\keepwithprev
A. $-4\notin M$\qquad B. $2\in M$

C. $4\notin M$\qquad D. $0\in M$

\ifanswer
\solbegin
\step{当 $x$、$y$、$z$ 均为正数时，$xyz>0$，$\dfrac x{|x|}+\dfrac y{|y|}+\dfrac z{|z|}+\dfrac{xyz}{|xyz|}=1+1+1+1=4$，}
\step{当 $x$、$y$、$z$ 有 2 个正数，1 个负数时，不妨设 $x>0$，$y>0$，$z<0$，}
\step{则 $xyz<0$，$\dfrac x{|x|}+\dfrac y{|y|}+\dfrac z{|z|}+\dfrac{xyz}{|xyz|}=1+1-1-1=0$，}
\step{若 $x$、$y$、$z$ 有 1 个正数，2 个负数时，不妨设 $x>0$，$y<0$，$z<0$，}
\step{则 $xyz>0$，$\dfrac x{|x|}+\dfrac y{|y|}+\dfrac z{|z|}+\dfrac{xyz}{|xyz|}=1-1-1+1=0$，}
\step{当 $x$、$y$、$z$ 均为负数时，$xyz<0$，}
\step{$\dfrac x{|x|}+\dfrac y{|y|}+\dfrac z{|z|}+\dfrac{xyz}{|xyz|}=-1-1-1-1=-4$，}
\step{综上，$M=\{-4,0,4\}$，}
\step{所以 $-4\in M$，$A$ 错误；$2\notin M$，$B$ 错误；$4\in M$，$C$ 错误；}
\step{$0\in M$，$D$ 正确. 故选 $D$.}
\solend
\fi

\item 若 $xy>0$，且 $x+y=9$，则 $\dfrac1x+\dfrac4y$ 的最小值是\mbox{（\quad）}

\keepwithprev
A. $9$\qquad B. $\dfrac59$

C. $\dfrac49$\qquad D. $1$

\ifanswer
\solbegin
\step{$xy>0$，且 $x+y=9$，所以 $x>0$，$y>0$，}
\step{所以 $\dfrac1x+\dfrac4y=\dfrac19\left(\dfrac1x+\dfrac4y\right)(x+y)=\dfrac59+\dfrac19\left(\dfrac yx+\dfrac{4x}y\right)\geqslant\dfrac59+\dfrac29\cdot\sqrt{\dfrac yx\cdot\dfrac{4x}y}=\dfrac59+\dfrac49=1$，}
\step{当且仅当 $\dfrac yx=\dfrac{4x}y$，即 $x=3$，$y=6$ 时取等号，所以 $\dfrac1x+\dfrac4y$ 的最小值为 $1$.}
\step{故选 $D$.}
\solend
\fi

\end{enumerate}

\bigsec{三、解答题}

\begin{enumerate}
\setcounter{enumi}{16}

\item 求解下列不等式：

（1）$\dfrac{x-3}{x+1}\leqslant0$；

（2）$|2x-1|<3$.

\ifanswer
\solbegin
\stepm{(1)}{$\dfrac{x-3}{x+1}\leqslant0$，变形得 $(x-3)\cdot(x+1)\leqslant0$ 且 $x\neq-1$，}
\step{解得 $-1<x\leqslant3$，即不等式的解集为 $(-1,3]$.}
\stepm{(2)}{若 $|2x-1|<3$，则 $-3<2x-1<3$，解得 $-1<x<2$，即 $(-1,2)$.}
\solend
\fi

\ansspace[6]

\item 已知集合 $A=\{x\mid-2<x<15\}$,$B=\{x\mid m-6<x<2m-1,m\in\R\}$.

（1）当 $m=2$ 时，求 $A\cap B$；

（2）若 $A\cup B=A$，求实数 $m$ 的取值范围.

\ifanswer
\solbegin
\stepm{(1)}{当 $m=2$ 时，$A=\{x\mid-2<x<15\}$,$B=\{x\mid-4<x<3\}$，}
\step{此时 $A\cap B=\{x\mid-2<x<3\}$.}
\stepm{(2)}{由 $A\cup B=A$ 得 $B\subseteq A$.}
\step{当 $B=\varnothing$,$m-6\geqslant2m-1$，解得 $m\leqslant-5$；}
\step{当 $B\neq\varnothing$，结合 $B\subseteq A$ 得 $\begin{cases}m-6<2m-1\\ 2m-1\leqslant15\\ m-6\geqslant-2\end{cases}$，解得 $4\leqslant m\leqslant8$.}
\step{综上，$m\leqslant-5$ 或 $4\leqslant m\leqslant8$.}
\solend
\fi

\ansspace[6]

\item 已知 $a\in\R$，命题 $p$：关于 $x$ 的方程 $x^2-2x+a=0$ 无实根；命题 $q$：对任意 $x\in\R$，
不等式 $ax^2-ax+1>0$ 恒成立.

（1）若 $q$ 为真命题，求 $a$ 的取值范围；

（2）若 $p$、$q$ 中有且只有一个为真命题，求 $a$ 的取值范围.

\ifanswer
\solbegin
\stepm{(1)}{当 $q$ 为真命题时，}
\stepm{①}{当 $a=0$ 时，不等式化为 $1>0$，对任意 $x\in\R$ 恒成立，满足条件；}
\stepm{②}{当 $a\neq0$ 时，$ax^2-ax+1>0$ 为一元二次不等式，}
\step{若对任意 $x\in\R$ 恒成立，需满足 $\begin{cases}a>0\\ \Delta=(-a)^2-4a<0\end{cases}$，解得 $0<a<4$；}
\step{综上，$a$ 的取值范围为 $[0,4)$.}
\stepm{(2)}{$p$ 为真时，方程 $x^2-2x+a=0$ 无实根，}
\step{故判别式 $\Delta=4-4a<0$，解得 $a>1$，因此 $p$ 为假时，$a\leqslant1$；}
\step{由（1）得 $q$ 为假时 $a\geqslant4$ 或 $a<0$，}
\step{$p$、$q$ 有且只有一个真命题，分两类：}
\step{当 $p$ 假 $q$ 真时，满足 $\begin{cases}a\leqslant1\\ 0\leqslant a<4\end{cases}$，解得 $0\leqslant a\leqslant1$；}
\step{当 $p$ 真 $q$ 假时，满足 $\begin{cases}a>1\\ a\geqslant4\text{或}a<0\end{cases}$，解得 $a\geqslant4$.}
\step{综上，$a$ 的取值范围是 $[0,1]\cup[4,+\infty)$.}
\solend
\fi

\ansspace[8]

\item 某工厂生产 $x$($x>0$) 件产品的总成本 $C$（单位：元）与生产件数 $x$ 之间满足二次函数
$C(x)=\dfrac14x^2+50x+10000$.

（1）若要求生产总成本不超过 30000 元，求满足生产总成本要求的产品件数的取值范围；

（2）试求该工厂生产多少件产品时，可使每件产品的平均成本最小.

\ifanswer
\solbegin
\stepm{(1)}{若要求生产总成本不超过 30000 元，}
\step{则 $C(x)=\dfrac14x^2+50x+10000\leqslant30000$，}
\step{即 $x^2+200x-80000\leqslant0$，解得 $-400\leqslant x\leqslant200$，又 $x>0$，}
\step{所以满足生产总成本要求的产品件数的取值范围为 $(0,200]$，$x\in\N$；}
\stepm{(2)}{生产件数 $x$ 时每件的平均成本为 $\dfrac{\frac14x^2+50x+10000}x=\dfrac14x+50+\dfrac{10000}x$，}
\step{由基本不等式得 $\dfrac14x+50+\dfrac{10000}x\geqslant2\sqrt{\dfrac14x\cdot\dfrac{10000}x}+50=150$，}
\step{当且仅当 $\dfrac14x=\dfrac{10000}x$，即 $x=200$ 时取等号，}
\step{即该工厂生产 200 件产品时，每件产品的平均成本最小.}
\solend
\fi

\ansspace[8]

% 原卷题干“$(n\geqslant2,n\in N)$”后连印两个全角逗号“，，”，照录未改（见文件头“原卷笔误”①）
\item 已知有限集 $A=\{a_1,a_2,\cdots,a_n\}$($n\geqslant2$,$n\in\N$)，，如果 $A$ 中元素
$a_i$($i=1,2,3,\cdots,n$) 满足 $a_1+a_2+a_3+\cdots+a_n=a_1\times a_2\times\cdots\times a_n$，
就称 $A$ 为“和积平衡集”.

% 原卷 (2) 将“写方程”与“反证法证明”两项任务并作一句（“……表示），可尝试用反证法证明：……”），
% 照录未改（见文件头“原卷笔误”②）
（1）判断集合 $\{-1+\sqrt3,-1-\sqrt3\}$ 是否为“和积平衡集”；

（2）若 $a_1$、$a_2$ 是两个不同的正数，且 $\{a_1,a_2\}$ 是“和积平衡集”，写出以 $a_1$、$a_2$ 为根的
一个一元二次方程（系数可用 $a_1$、$a_2$ 表示），可尝试用反证法证明：$a_1$、$a_2$ 至少有一个大于 2；

（3）若 $a_i\geqslant1$，且 $a_i\in\N$，求所有符合条件的“和积平衡集”.

\ifanswer
\solbegin
\stepm{(1)}{对于集合 $\{-1+\sqrt3,-1-\sqrt3\}$，}
\step{因为 $(-1+\sqrt3)+(-1-\sqrt3)=-2$，$(-1+\sqrt3)\times(-1-\sqrt3)=-2$，}
\step{所以满足 $a_1+a_2=a_1\times a_2$，满足题意，}
\step{所以集合 $\{-1+\sqrt3,-1-\sqrt3\}$ 是和积平衡集.}
\stepm{(2)}{以 $a_1$、$a_2$ 为根的一个一元二次方程可为 $x^2-(a_1+a_2)x+a_1\cdot a_2=0$，}
\step{因为 $a_1$、$a_2$ 是两个不同的正数，且 $\{a_1,a_2\}$ 是和积平衡集，}
\step{所以 $a_1+a_2=a_1\cdot a_2$，}
\step{又 $a_1$ 和 $a_2$ 为 $x^2-(a_1+a_2)x+a_1\cdot a_2=0$ 的两个不同的正根，}
\step{所以 $\Delta=(a_1+a_2)^2-4a_1a_2=(a_1+a_2)^2-4(a_1+a_2)>0$，}
\step{又 $a_1+a_2>0$，则 $a_1+a_2>4$，}
\step{若 $a_1$、$a_2$ 都小于等于 2，则 $a_1+a_2<4$，矛盾，}
\step{所以 $a_1$、$a_2$ 至少有一个大于 2；}
\stepm{(3)}{设 $A$ 中的 $a_1<a_2<a_3<\cdots<a_n$，且 $a_i\in\N^*$，}
\step{由 $a_1+a_2+a_3+\cdots+a_n=a_1\times a_2\times\cdots\times a_n$，得 $n=2$ 时明显不满足，}
\step{所以 $n\geqslant3$，故 $a_1\times a_2\times\cdots\times a_n=a_1+a_2+a_3+\cdots+a_n\leqslant na_n$，}
\step{所以 $a_1\times a_2\times\cdots\times a_{n-1}\leqslant n$，}
\step{当 $n=3$ 时，$a_1\times a_2\leqslant3$，$a_i\in\N^*$，只能有 $a_1=1$，$a_2=2$，}
\step{由 $a_1+a_2+a_3=a_1\times a_2\times a_3$，得 $a_3=3$，此时的和积平衡集为 $\{1,2,3\}$，}
\step{当 $n\geqslant4$ 时，$a_1\times a_2\times\cdots\times a_{n-1}\geqslant(n-1)(n-2)\times\cdots\times2\times1\geqslant(n-1)(n-2)$，}
\step{又 $(n-1)(n-2)-n=n^2-4n+2=n(n-4)+2>0$，即 $(n-1)(n-2)>n$，}
\step{即 $a_1\times a_2\times\cdots\times a_{n-1}>n$，与 $a_1\times a_2\times\cdots\times a_{n-1}\leqslant n$ 矛盾，}
\step{所以 $n\geqslant4$ 不满足条件，}
\step{所以符合条件的和积平衡集有且只有 $n=3$，此时的和积平衡集为 $\{1,2,3\}$.}
\solend
\fi

\ansspace*[10]

\end{enumerate}

\end{document}
"
% 紧凑版：xelatex "\def\iscompact{1}% 高一63_吴淞中学_10月月考.tex
%   —— 高一上数学“2026-2027 年吴淞中学高一上 10 月月考”（讲评版）
% 试卷版：xelatex 高一63_吴淞中学_10月月考.tex
% 答案版：xelatex "\def\isanswer{1}% 高一63_吴淞中学_10月月考.tex
%   —— 高一上数学“2026-2027 年吴淞中学高一上 10 月月考”（讲评版）
% 试卷版：xelatex 高一63_吴淞中学_10月月考.tex
% 答案版：xelatex "\def\isanswer{1}\input{高一63_吴淞中学_10月月考.tex}"
% 紧凑版：xelatex "\def\iscompact{1}\input{高一63_吴淞中学_10月月考.tex}"
%
% 原始材料：微信公众号“魔都数学加油站”文章，作者破晓，连载“27学年高一63”；
% 原文标题“27学年高一63——2026-2027年上海市吴淞中学高一上10月月考”；
% 原文链接：https://mp.weixin.qq.com/s/TFQ4TLnt_l66t7t51es4OQ。
% 页面标注发布时间 2026-10-16 21:44（上海）。
% 正文为 10 张 PNG 页（`原图/高一63_吴淞中学_10月月考/01..10.png`）：
% 第 06、10 页 2560×3620，第 01、02、03、04、05、07、08、09 页 4762×6735
% （**同一篇各页分辨率不同**），均按页序读取。10 页全为试卷正文页（卷首标题在第 1 页页顶，
% 卷末解答题第 21 题解析止于第 10 页），无封面页、目录页、二维码或宣传页。
% 粉色斜排“魔都数学加油站”为卷面水印，与公众号页脚一并未录入。
%
% 卷首标题：照卷面印作“2026--2027 年吴淞中学高一上 10 月月考”（公众号标题中的“上海市”
% 卷面标题行没有，本稿照卷面），年份区间按全稿惯例排 en dash。原卷只有一行标题、无副标题。
%
% 篇幅与题号：原卷自第 3 题起编号（第 1、2 题不在本卷内）。填空题 3—12（10 题）、
% 选择题 13—16（4 题）、解答题 17—21（5 题），共 19 题。本稿沿用原节与原题号，
% 只改排为全稿统一的大节样式（原卷小节标题为“一、填空题”“二、选择题”“三、解答题”）。
%
% 讲评版：原卷每题下直接印【答案】或【解析】。**只印【答案】的 2 题**（第 4、7 题），
% 本稿只给 \ans{}；其余 17 题里 16 题都印【解析】（按全稿口径这类题不另提 \ans{}——第 13、15、16 题
% 的结论写在解析末句（“故选 B”等），亦照录在解析内；第 14 题则既无【答案】也无【解析】段。**另有 1 题答案印在题干括号内**——
% 第 14 题（“(  D  )”，且该题无【答案】/【解析】段），本稿按 高一37 第 13 题成例把括号留空、
% 答案另提为 \ans{$D$}。第 16 题题干括号为空但印有【解析】，不算括号内答案。
%
% 副标题：原卷未印副标题。本稿按“知识模块”口径标作“集合与常用逻辑用语、不等式”：
% 19 题中集合与常用逻辑用语 12 题（第 3、4、6、7、9、10、11、12、13、15、18、21 题），
% 不等式 7 题（含等式（韦达定理、恒等式）与不等式性质、基本不等式、二次不等式、恒成立：
% 第 5、8、14、16、17、19、20 题），两类均超过三分之一，故并标。口径同 高一59。
%
% 与原卷的排版差异：
%   ① 原卷每题下直接印【答案】/【解析】，本稿按全稿惯例将答案与解析置于答案版，试卷版保持干净；
%   ② 原卷未印副标题，本稿按“知识模块”口径补出；
%   ③ 原卷填空横线约四个字宽，本稿按全稿惯例统一排 \blankline[4.4em]；
%   ④ 原卷实数集排斜体 $R$（第 9、10、19 题题干与解析）或粗体 $\mathbf{R}$（第 18 题题干），
%      整数集排斜体 $Z$（第 12 题解析），自然数集排斜体 $N$（第 20 题解析、第 21 题题干
%      “$(n\geqslant2,n\in N)$”）或粗体 $\mathbf{N}$（第 21 题 (3) 题干）、斜体 $N^*$
%      （第 21 题 (3) 解析），本稿按全稿惯例统一改用 $\R$、$\Z$、$\N$、$\N^*$
%      （同 高一14 排版差异④：记号归一，粗体／斜体不作区分，$N$ 与 $N^*$ 的区别照原卷保留）；
%      空集记号原卷作 $\varnothing$，照录；
%   ⑤ 原卷在数学式内多处用半角逗号（如 $\{y\mid y=[2x]+[3x],-1\leqslant x\leqslant2\}$、
%      $m=-1,0,1$、$A=\{x\mid-2<x<15\}$,$B=\cdots$），均照录未强改；
%      句末句点按全稿惯例排西文句点；“和积平衡集”自定义名词照录；
%   ⑥ 第 12 题解析里的“法一：”“法二：”标号原卷排**红色**（全卷共 2 处），
%      同 高一53“情况1：”成例，本稿排 \stepm{法一：}{…}、\stepm{法二：}{…}
%      （加粗着色、不改黑色基调），冒号照原卷排全角；
%   ⑦ 第 17—21 题的子问原卷各自另起一行，本稿按全稿惯例用 enumerate 的 (1)(2)(3)；
%      解析中的子问标号一律用 \stepm{(1)}{…}；第 19 题解析中的①②分情形，按全稿惯例作
%      \stepm{①}{…}、\stepm{②}{…}；
%   ⑧ 第 14 题印在题干括号内的答案移入 \ans{}，见上；
%   ⑨ 每道解答题原卷未留作答空白，本稿按全稿惯例加 \ansspace；
%   ⑩ 选择题的作答括号原卷排半角“(　　)”（第 13、15、16 题括号内空白、第 14 题括号内印 D），
%      本稿按全稿惯例排 \mbox{（\quad）}。
%
% 原卷笔误/存疑（均照录未改，正文对应处另有 % 注）：
%   ① 第 21 题题干“$(n\geqslant2$,$n\in N)$”后连印两个全角逗号“，，”（衍字，
%      8 倍放大确认，非水印干扰），本稿照录；
%   ② 第 21 题 (2) 题干将“写出以 $a_1$、$a_2$ 为根的一个一元二次方程”与“用反证法证明：
%      $a_1$、$a_2$ 至少有一个大于 2”两项任务并作一句、之间仅以逗号相接（“……表示），
%      可尝试用反证法证明：……”），读来像抄并两小问，原卷如此，照录未改。
%
% 插图：无（全卷 10 页无任何图形、无“如图”字样）。
%
\documentclass[a4paper,11pt]{article}
\usepackage{common}

\begin{document}

\examtitlest[3]{2026--2027 年吴淞中学高一上 10 月月考}{集合与常用逻辑用语、不等式}

\bigsec{一、填空题}

\begin{enumerate}
\setcounter{enumi}{2}

\item 满足 $A\subseteq\{1,2,3,4\}$ 的集合 $A$ 共有\blankline[4.4em]个.

\ifanswer
\solbegin
\step{由题意，集合 $A$ 是集合 $\{1,2,3,4\}$ 的子集，共有 $2^4=16$ 个.}
\solend
\fi

\item 陈述句“$a>1$ 或 $b>3$”的否定形式为\blankline[4.4em].

\ifanswer
\ans{$a\leqslant1$ 且 $b\leqslant3$}
\fi

\item 已知方程 $3x^2+4x-4=0$ 的两个根为 $x_1$、$x_2$，则 $x_1^2+x_2^2=$\blankline[4.4em].

\ifanswer
\solbegin
\step{对于方程 $3x^2+4x-4=0$，因为 $\Delta=16-4\times3\times(-4)=64>0$，}
\step{所以方程 $3x^2+4x-4=0$ 的两个不相等的实根为 $x_1$、$x_2$，}
\step{由韦达定理可得 $x_1+x_2=-\dfrac43$，$x_1x_2=-\dfrac43$，}
\step{所以 $x_1^2+x_2^2=(x_1+x_2)^2-2x_1x_2=\dfrac{16}9+\dfrac83=\dfrac{40}9$.}
\solend
\fi

\item 设集合 $A=\{y\mid y=x^2+1\}$,$B=\{y\mid y=-x^2+2x+3\}$，则 $A\cap B=$\blankline[4.4em].

\ifanswer
\solbegin
\step{因为 $y=x^2+1\geqslant1$,$y=-x^2+2x+3=-(x-1)^2+4\leqslant4$，}
\step{所以 $A=\{y\mid y\geqslant1\}$,$B=\{y\mid y\leqslant4\}$，得 $A\cap B=\{y\mid1\leqslant y\leqslant4\}$.}
\solend
\fi

\item 若 $x\in\{1,x^2\}$，则 $x=$\blankline[4.4em].

\ifanswer
\ans{$0$}
\fi

\item 已知等式 $2x^2+3x+5=a(2x+1)(x+1)+c$，则 $a+c=$\blankline[4.4em].

\ifanswer
\solbegin
\step{因为 $2x^2+3x+5=a(2x+1)(x+1)+c$ 恒成立，}
\step{所以 $2x^2+3x+5=2ax^2+3ax+a+c$ 恒成立，则 $\begin{cases}2=2a\\ 3=3a\\ 5=a+c\end{cases}$，解得 $a+c=5$.}
\solend
\fi

\item 已知 $a\in\R$，设 $p:x<2$，$q:x<5a-3$，且 $p$ 是 $q$ 的必要非充分条件，则 $a$ 的取值范围是\blankline[4.4em].

\ifanswer
\solbegin
\step{设 $p:x<2$，$q:x<5a-3$，且 $p$ 是 $q$ 的必要非充分条件，}
\step{则 $\{x\mid x<5a-3\}\subset\{x\mid x<2\}$，所以 $5a-3<2$，即 $a<1$.}
\solend
\fi

\item 已知 $a\in\R$，集合 $A=\{x\mid ax^2+3x-1=0\}$，若集合 $A$ 中只有一个元素，则 $a$ 的取值范围是\blankline[4.4em].

\ifanswer
\solbegin
\step{当 $a=0$ 时，则有 $\{x\mid ax^2+3x-1=0\}=\{x\mid3x-1=0\}=\left\{\dfrac13\right\}$，符合题意；}
\step{当 $a\neq0$ 时，由题意得 $\Delta=9+4a=0$，解得 $a=-\dfrac94$. 符合题意；}
\step{综上所述，实数 $a$ 的取值集合为 $\left\{0,-\dfrac94\right\}$.}
\solend
\fi

\item 设集合 $A=\{x\mid x^2+2(m+1)x+m^2-1=0,m\in\R\}$，且 $A\subseteq\{0,-4\}$，则实数 $m$ 的取值范围是\blankline[4.4em].

\ifanswer
\solbegin
\step{由题意得 $A=\varnothing$ 或 $A=\{0\}$ 或 $A=\{-4\}$ 或 $A=\{0,-4\}$，}
\step{当 $A=\varnothing$ 时，方程 $x^2+2(m+1)x+m^2-1=0$ 无实根，}
\step{则 $\Delta=[2(m+1)]^2-4\times1\times(m^2-1)=8m+8<0$，解得 $m<-1$，}
\step{当 $A=\{0\}$ 时，方程 $x^2+2(m+1)x+m^2-1=0$ 有一个实根 $0$，}
\step{则 $\begin{cases}\Delta=8m+8=0\\ m^2-1=0\end{cases}$，解得 $m=-1$，}
\step{当 $A=\{-4\}$ 时，方程 $x^2+2(m+1)x+m^2-1=0$ 有一个实根 $-4$，}
\step{则 $\begin{cases}\Delta=8m+8=0\\ 16-8(m+1)+m^2-1=0\end{cases}$，无解，}
\step{当 $A=\{0,-4\}$ 时，方程 $x^2+2(m+1)x+m^2-1=0$ 有两个实根 $-4$ 和 $0$，}
\step{则 $\begin{cases}\Delta=8m+8>0\\ 16-8(m+1)+m^2-1=0\\ m^2-1=0\end{cases}$，解得 $m=1$，}
\step{综上所述，实数 $m$ 的取值范围是 $\{m\mid m\leqslant-1\text{ 或 }m=1\}$.}
\solend
\fi

\item 对于任意实数 $x$，$[x]$ 表示不超过 $x$ 的最大整数，如 $[1.1]=1$，$[-1.2]=-2$，若
$A=\{y\mid y=[2x]+[3x],-1\leqslant x\leqslant2\}$，则 $A$ 中所有元素的和为\blankline[4.4em].

\ifanswer
\solbegin
% 原卷“法一：”“法二：”标号排红色，本稿按 高一53“情况1：”成例排 \stepm（加粗着色），见文件头
\stepm{法一：}{当 $-1\leqslant x<-\dfrac23$ 时，$[2x]=-2$，$[3x]=-3$，$f(x)=[2x]+[3x]=-5$，}
\step{当 $-\dfrac23\leqslant x<-\dfrac12$ 时，$[2x]=-2$，$[3x]=-2$，$f(x)=[2x]+[3x]=-4$；}
\step{当 $-\dfrac12\leqslant x<-\dfrac13$ 时，$[2x]=-1$，$[3x]=-2$，$f(x)=[2x]+[3x]=-3$；}
\step{当 $-\dfrac13\leqslant x<0$ 时，$[2x]=-1$，$[3x]=-1$，$f(x)=[2x]+[3x]=-2$；}
\step{当 $0\leqslant x<\dfrac13$ 时，$[2x]=0$，$[3x]=0$，$f(x)=[2x]+[3x]=0$；}
\step{当 $\dfrac13\leqslant x<\dfrac12$ 时，$[2x]=0$，$[3x]=1$，$f(x)=[2x]+[3x]=1$；}
\step{当 $\dfrac12\leqslant x<\dfrac23$ 时，$[2x]=1$，$[3x]=1$，$f(x)=[2x]+[3x]=2$；}
\step{当 $\dfrac23\leqslant x<1$ 时，$[2x]=1$，$[3x]=2$，$f(x)=[2x]+[3x]=3$；}
\step{当 $1\leqslant x<\dfrac43$ 时，$[2x]=2$，$[3x]=3$，$f(x)=[2x]+[3x]=5$；}
\step{当 $\dfrac43\leqslant x<\dfrac32$ 时，$[2x]=2$，$[3x]=4$，$f(x)=[2x]+[3x]=6$；}
\step{当 $\dfrac32\leqslant x<\dfrac53$ 时，$[2x]=3$，$[3x]=4$，$f(x)=[2x]+[3x]=7$；}
\step{当 $\dfrac53\leqslant x<2$ 时，$[2x]=3$，$[3x]=5$，$f(x)=[2x]+[3x]=8$；}
\step{当 $x=2$ 时，$[2x]=4$，$[3x]=6$，$f(x)=[2x]+[3x]=10$}
\step{故 $A$ 中所有元素的和为 $-5-4-3-2+0+1+2+3+5+6+7+8+10=28$.}
% 原卷“法二：”标号排红色，同上
\stepm{法二：}{设 $x=m+n$，其中 $m\in\Z$，$n\in[0,1)$，}
\step{若 $0\leqslant n<\dfrac13$，则 $0\leqslant2n<\dfrac23$，$0\leqslant3n<1$，}
\step{得 $[2x]=[2m+2n]=2m$，$[3x]=[3m+3n]=3m$，所以 $y=[2x]+[3x]=5m$；}
\step{若 $\dfrac13\leqslant n<\dfrac12$，则 $\dfrac23\leqslant2n<1$，$1\leqslant3n<\dfrac32$，}
\step{得 $[2x]=[2m+2n]=2m$，$[3x]=[3m+3n]=3m+1$，}
\step{所以 $y=[2x]+[3x]=5m+1$；}
\step{若 $\dfrac12\leqslant n<\dfrac23$，则 $1\leqslant2n<\dfrac43$，$\dfrac32\leqslant3n<2$，}
\step{得 $[2x]=[2m+2n]=2m+1$，$[3x]=[3m+3n]=3m+1$，}
\step{所以 $y=[2x]+[3x]=5m+2$；}
\step{若 $\dfrac23\leqslant n<1$，则 $\dfrac43\leqslant2n<2$，$2\leqslant3n<3$，}
\step{得 $[2x]=[2m+2n]=2m+1$，$[3x]=[3m+3n]=3m+2$，}
\step{所以 $y=[2x]+[3x]=5m+3$；}
\step{综上所述，当 $n\in[0,1)$ 时，所有函数值的和}
\step{为 $5m+5m+1+5m+2+5m+3=20m+6$，}
\step{又因为 $x\in[-1,2]$，此时 $m=-1,0,1$，且当 $x=2$ 时，$y=[2x]+[3x]=10$，}
\step{所有函数值的和为 $(-14)+6+26+10=28$，即 $A$ 中所有元素之和为 $28$.}
\solend
\fi

\end{enumerate}

\bigsec{二、选择题}

\begin{enumerate}
\setcounter{enumi}{12}

\item 下列写法中，正确的是\mbox{（\quad）}

\keepwithprev
A. $0\in\{\varnothing\}$\qquad B. $\varnothing\in\{\varnothing\}$

C. $\varnothing=\{\varnothing\}$\qquad D. $0\in\varnothing$

\ifanswer
\solbegin
\step{$0\notin\{\varnothing\}$，$\varnothing\in\{\varnothing\}$，故 $A$ 错误，$B$ 正确；}
\step{$\varnothing$ 是不含任何元素的集合，所以 $\varnothing\neq\{\varnothing\}$，$0\notin\varnothing$，故 $CD$ 错误.}
\step{故选 $B$.}
\solend
\fi

% 第 14 题答案“D”直接印在题干括号内“(  D  )”，且该题无【答案】/【解析】段；
% 按 高一37 第 13 题的成例括号留空、答案另提为 \ans{}（见文件头）
\item 若 $a<b<0$，则下列不等式中不能成立的是\mbox{（\quad）}

\keepwithprev
A. $\dfrac1a>\dfrac1b$\qquad B. $a-b>a$

C. $|a|>|b|$\qquad D. $a^3>b^3$

\ifanswer
\ans{$D$}
\fi

\item 已知 $x$、$y$、$z$ 为非零实数，代数式 $\dfrac x{|x|}+\dfrac y{|y|}+\dfrac z{|z|}+\dfrac{xyz}{|xyz|}$
的值所组成的集合是 $M$，则下列判断正确的是\mbox{（\quad）}

\keepwithprev
A. $-4\notin M$\qquad B. $2\in M$

C. $4\notin M$\qquad D. $0\in M$

\ifanswer
\solbegin
\step{当 $x$、$y$、$z$ 均为正数时，$xyz>0$，$\dfrac x{|x|}+\dfrac y{|y|}+\dfrac z{|z|}+\dfrac{xyz}{|xyz|}=1+1+1+1=4$，}
\step{当 $x$、$y$、$z$ 有 2 个正数，1 个负数时，不妨设 $x>0$，$y>0$，$z<0$，}
\step{则 $xyz<0$，$\dfrac x{|x|}+\dfrac y{|y|}+\dfrac z{|z|}+\dfrac{xyz}{|xyz|}=1+1-1-1=0$，}
\step{若 $x$、$y$、$z$ 有 1 个正数，2 个负数时，不妨设 $x>0$，$y<0$，$z<0$，}
\step{则 $xyz>0$，$\dfrac x{|x|}+\dfrac y{|y|}+\dfrac z{|z|}+\dfrac{xyz}{|xyz|}=1-1-1+1=0$，}
\step{当 $x$、$y$、$z$ 均为负数时，$xyz<0$，}
\step{$\dfrac x{|x|}+\dfrac y{|y|}+\dfrac z{|z|}+\dfrac{xyz}{|xyz|}=-1-1-1-1=-4$，}
\step{综上，$M=\{-4,0,4\}$，}
\step{所以 $-4\in M$，$A$ 错误；$2\notin M$，$B$ 错误；$4\in M$，$C$ 错误；}
\step{$0\in M$，$D$ 正确. 故选 $D$.}
\solend
\fi

\item 若 $xy>0$，且 $x+y=9$，则 $\dfrac1x+\dfrac4y$ 的最小值是\mbox{（\quad）}

\keepwithprev
A. $9$\qquad B. $\dfrac59$

C. $\dfrac49$\qquad D. $1$

\ifanswer
\solbegin
\step{$xy>0$，且 $x+y=9$，所以 $x>0$，$y>0$，}
\step{所以 $\dfrac1x+\dfrac4y=\dfrac19\left(\dfrac1x+\dfrac4y\right)(x+y)=\dfrac59+\dfrac19\left(\dfrac yx+\dfrac{4x}y\right)\geqslant\dfrac59+\dfrac29\cdot\sqrt{\dfrac yx\cdot\dfrac{4x}y}=\dfrac59+\dfrac49=1$，}
\step{当且仅当 $\dfrac yx=\dfrac{4x}y$，即 $x=3$，$y=6$ 时取等号，所以 $\dfrac1x+\dfrac4y$ 的最小值为 $1$.}
\step{故选 $D$.}
\solend
\fi

\end{enumerate}

\bigsec{三、解答题}

\begin{enumerate}
\setcounter{enumi}{16}

\item 求解下列不等式：

（1）$\dfrac{x-3}{x+1}\leqslant0$；

（2）$|2x-1|<3$.

\ifanswer
\solbegin
\stepm{(1)}{$\dfrac{x-3}{x+1}\leqslant0$，变形得 $(x-3)\cdot(x+1)\leqslant0$ 且 $x\neq-1$，}
\step{解得 $-1<x\leqslant3$，即不等式的解集为 $(-1,3]$.}
\stepm{(2)}{若 $|2x-1|<3$，则 $-3<2x-1<3$，解得 $-1<x<2$，即 $(-1,2)$.}
\solend
\fi

\ansspace[6]

\item 已知集合 $A=\{x\mid-2<x<15\}$,$B=\{x\mid m-6<x<2m-1,m\in\R\}$.

（1）当 $m=2$ 时，求 $A\cap B$；

（2）若 $A\cup B=A$，求实数 $m$ 的取值范围.

\ifanswer
\solbegin
\stepm{(1)}{当 $m=2$ 时，$A=\{x\mid-2<x<15\}$,$B=\{x\mid-4<x<3\}$，}
\step{此时 $A\cap B=\{x\mid-2<x<3\}$.}
\stepm{(2)}{由 $A\cup B=A$ 得 $B\subseteq A$.}
\step{当 $B=\varnothing$,$m-6\geqslant2m-1$，解得 $m\leqslant-5$；}
\step{当 $B\neq\varnothing$，结合 $B\subseteq A$ 得 $\begin{cases}m-6<2m-1\\ 2m-1\leqslant15\\ m-6\geqslant-2\end{cases}$，解得 $4\leqslant m\leqslant8$.}
\step{综上，$m\leqslant-5$ 或 $4\leqslant m\leqslant8$.}
\solend
\fi

\ansspace[6]

\item 已知 $a\in\R$，命题 $p$：关于 $x$ 的方程 $x^2-2x+a=0$ 无实根；命题 $q$：对任意 $x\in\R$，
不等式 $ax^2-ax+1>0$ 恒成立.

（1）若 $q$ 为真命题，求 $a$ 的取值范围；

（2）若 $p$、$q$ 中有且只有一个为真命题，求 $a$ 的取值范围.

\ifanswer
\solbegin
\stepm{(1)}{当 $q$ 为真命题时，}
\stepm{①}{当 $a=0$ 时，不等式化为 $1>0$，对任意 $x\in\R$ 恒成立，满足条件；}
\stepm{②}{当 $a\neq0$ 时，$ax^2-ax+1>0$ 为一元二次不等式，}
\step{若对任意 $x\in\R$ 恒成立，需满足 $\begin{cases}a>0\\ \Delta=(-a)^2-4a<0\end{cases}$，解得 $0<a<4$；}
\step{综上，$a$ 的取值范围为 $[0,4)$.}
\stepm{(2)}{$p$ 为真时，方程 $x^2-2x+a=0$ 无实根，}
\step{故判别式 $\Delta=4-4a<0$，解得 $a>1$，因此 $p$ 为假时，$a\leqslant1$；}
\step{由（1）得 $q$ 为假时 $a\geqslant4$ 或 $a<0$，}
\step{$p$、$q$ 有且只有一个真命题，分两类：}
\step{当 $p$ 假 $q$ 真时，满足 $\begin{cases}a\leqslant1\\ 0\leqslant a<4\end{cases}$，解得 $0\leqslant a\leqslant1$；}
\step{当 $p$ 真 $q$ 假时，满足 $\begin{cases}a>1\\ a\geqslant4\text{或}a<0\end{cases}$，解得 $a\geqslant4$.}
\step{综上，$a$ 的取值范围是 $[0,1]\cup[4,+\infty)$.}
\solend
\fi

\ansspace[8]

\item 某工厂生产 $x$($x>0$) 件产品的总成本 $C$（单位：元）与生产件数 $x$ 之间满足二次函数
$C(x)=\dfrac14x^2+50x+10000$.

（1）若要求生产总成本不超过 30000 元，求满足生产总成本要求的产品件数的取值范围；

（2）试求该工厂生产多少件产品时，可使每件产品的平均成本最小.

\ifanswer
\solbegin
\stepm{(1)}{若要求生产总成本不超过 30000 元，}
\step{则 $C(x)=\dfrac14x^2+50x+10000\leqslant30000$，}
\step{即 $x^2+200x-80000\leqslant0$，解得 $-400\leqslant x\leqslant200$，又 $x>0$，}
\step{所以满足生产总成本要求的产品件数的取值范围为 $(0,200]$，$x\in\N$；}
\stepm{(2)}{生产件数 $x$ 时每件的平均成本为 $\dfrac{\frac14x^2+50x+10000}x=\dfrac14x+50+\dfrac{10000}x$，}
\step{由基本不等式得 $\dfrac14x+50+\dfrac{10000}x\geqslant2\sqrt{\dfrac14x\cdot\dfrac{10000}x}+50=150$，}
\step{当且仅当 $\dfrac14x=\dfrac{10000}x$，即 $x=200$ 时取等号，}
\step{即该工厂生产 200 件产品时，每件产品的平均成本最小.}
\solend
\fi

\ansspace[8]

% 原卷题干“$(n\geqslant2,n\in N)$”后连印两个全角逗号“，，”，照录未改（见文件头“原卷笔误”①）
\item 已知有限集 $A=\{a_1,a_2,\cdots,a_n\}$($n\geqslant2$,$n\in\N$)，，如果 $A$ 中元素
$a_i$($i=1,2,3,\cdots,n$) 满足 $a_1+a_2+a_3+\cdots+a_n=a_1\times a_2\times\cdots\times a_n$，
就称 $A$ 为“和积平衡集”.

% 原卷 (2) 将“写方程”与“反证法证明”两项任务并作一句（“……表示），可尝试用反证法证明：……”），
% 照录未改（见文件头“原卷笔误”②）
（1）判断集合 $\{-1+\sqrt3,-1-\sqrt3\}$ 是否为“和积平衡集”；

（2）若 $a_1$、$a_2$ 是两个不同的正数，且 $\{a_1,a_2\}$ 是“和积平衡集”，写出以 $a_1$、$a_2$ 为根的
一个一元二次方程（系数可用 $a_1$、$a_2$ 表示），可尝试用反证法证明：$a_1$、$a_2$ 至少有一个大于 2；

（3）若 $a_i\geqslant1$，且 $a_i\in\N$，求所有符合条件的“和积平衡集”.

\ifanswer
\solbegin
\stepm{(1)}{对于集合 $\{-1+\sqrt3,-1-\sqrt3\}$，}
\step{因为 $(-1+\sqrt3)+(-1-\sqrt3)=-2$，$(-1+\sqrt3)\times(-1-\sqrt3)=-2$，}
\step{所以满足 $a_1+a_2=a_1\times a_2$，满足题意，}
\step{所以集合 $\{-1+\sqrt3,-1-\sqrt3\}$ 是和积平衡集.}
\stepm{(2)}{以 $a_1$、$a_2$ 为根的一个一元二次方程可为 $x^2-(a_1+a_2)x+a_1\cdot a_2=0$，}
\step{因为 $a_1$、$a_2$ 是两个不同的正数，且 $\{a_1,a_2\}$ 是和积平衡集，}
\step{所以 $a_1+a_2=a_1\cdot a_2$，}
\step{又 $a_1$ 和 $a_2$ 为 $x^2-(a_1+a_2)x+a_1\cdot a_2=0$ 的两个不同的正根，}
\step{所以 $\Delta=(a_1+a_2)^2-4a_1a_2=(a_1+a_2)^2-4(a_1+a_2)>0$，}
\step{又 $a_1+a_2>0$，则 $a_1+a_2>4$，}
\step{若 $a_1$、$a_2$ 都小于等于 2，则 $a_1+a_2<4$，矛盾，}
\step{所以 $a_1$、$a_2$ 至少有一个大于 2；}
\stepm{(3)}{设 $A$ 中的 $a_1<a_2<a_3<\cdots<a_n$，且 $a_i\in\N^*$，}
\step{由 $a_1+a_2+a_3+\cdots+a_n=a_1\times a_2\times\cdots\times a_n$，得 $n=2$ 时明显不满足，}
\step{所以 $n\geqslant3$，故 $a_1\times a_2\times\cdots\times a_n=a_1+a_2+a_3+\cdots+a_n\leqslant na_n$，}
\step{所以 $a_1\times a_2\times\cdots\times a_{n-1}\leqslant n$，}
\step{当 $n=3$ 时，$a_1\times a_2\leqslant3$，$a_i\in\N^*$，只能有 $a_1=1$，$a_2=2$，}
\step{由 $a_1+a_2+a_3=a_1\times a_2\times a_3$，得 $a_3=3$，此时的和积平衡集为 $\{1,2,3\}$，}
\step{当 $n\geqslant4$ 时，$a_1\times a_2\times\cdots\times a_{n-1}\geqslant(n-1)(n-2)\times\cdots\times2\times1\geqslant(n-1)(n-2)$，}
\step{又 $(n-1)(n-2)-n=n^2-4n+2=n(n-4)+2>0$，即 $(n-1)(n-2)>n$，}
\step{即 $a_1\times a_2\times\cdots\times a_{n-1}>n$，与 $a_1\times a_2\times\cdots\times a_{n-1}\leqslant n$ 矛盾，}
\step{所以 $n\geqslant4$ 不满足条件，}
\step{所以符合条件的和积平衡集有且只有 $n=3$，此时的和积平衡集为 $\{1,2,3\}$.}
\solend
\fi

\ansspace*[10]

\end{enumerate}

\end{document}
"
% 紧凑版：xelatex "\def\iscompact{1}% 高一63_吴淞中学_10月月考.tex
%   —— 高一上数学“2026-2027 年吴淞中学高一上 10 月月考”（讲评版）
% 试卷版：xelatex 高一63_吴淞中学_10月月考.tex
% 答案版：xelatex "\def\isanswer{1}\input{高一63_吴淞中学_10月月考.tex}"
% 紧凑版：xelatex "\def\iscompact{1}\input{高一63_吴淞中学_10月月考.tex}"
%
% 原始材料：微信公众号“魔都数学加油站”文章，作者破晓，连载“27学年高一63”；
% 原文标题“27学年高一63——2026-2027年上海市吴淞中学高一上10月月考”；
% 原文链接：https://mp.weixin.qq.com/s/TFQ4TLnt_l66t7t51es4OQ。
% 页面标注发布时间 2026-10-16 21:44（上海）。
% 正文为 10 张 PNG 页（`原图/高一63_吴淞中学_10月月考/01..10.png`）：
% 第 06、10 页 2560×3620，第 01、02、03、04、05、07、08、09 页 4762×6735
% （**同一篇各页分辨率不同**），均按页序读取。10 页全为试卷正文页（卷首标题在第 1 页页顶，
% 卷末解答题第 21 题解析止于第 10 页），无封面页、目录页、二维码或宣传页。
% 粉色斜排“魔都数学加油站”为卷面水印，与公众号页脚一并未录入。
%
% 卷首标题：照卷面印作“2026--2027 年吴淞中学高一上 10 月月考”（公众号标题中的“上海市”
% 卷面标题行没有，本稿照卷面），年份区间按全稿惯例排 en dash。原卷只有一行标题、无副标题。
%
% 篇幅与题号：原卷自第 3 题起编号（第 1、2 题不在本卷内）。填空题 3—12（10 题）、
% 选择题 13—16（4 题）、解答题 17—21（5 题），共 19 题。本稿沿用原节与原题号，
% 只改排为全稿统一的大节样式（原卷小节标题为“一、填空题”“二、选择题”“三、解答题”）。
%
% 讲评版：原卷每题下直接印【答案】或【解析】。**只印【答案】的 2 题**（第 4、7 题），
% 本稿只给 \ans{}；其余 17 题里 16 题都印【解析】（按全稿口径这类题不另提 \ans{}——第 13、15、16 题
% 的结论写在解析末句（“故选 B”等），亦照录在解析内；第 14 题则既无【答案】也无【解析】段。**另有 1 题答案印在题干括号内**——
% 第 14 题（“(  D  )”，且该题无【答案】/【解析】段），本稿按 高一37 第 13 题成例把括号留空、
% 答案另提为 \ans{$D$}。第 16 题题干括号为空但印有【解析】，不算括号内答案。
%
% 副标题：原卷未印副标题。本稿按“知识模块”口径标作“集合与常用逻辑用语、不等式”：
% 19 题中集合与常用逻辑用语 12 题（第 3、4、6、7、9、10、11、12、13、15、18、21 题），
% 不等式 7 题（含等式（韦达定理、恒等式）与不等式性质、基本不等式、二次不等式、恒成立：
% 第 5、8、14、16、17、19、20 题），两类均超过三分之一，故并标。口径同 高一59。
%
% 与原卷的排版差异：
%   ① 原卷每题下直接印【答案】/【解析】，本稿按全稿惯例将答案与解析置于答案版，试卷版保持干净；
%   ② 原卷未印副标题，本稿按“知识模块”口径补出；
%   ③ 原卷填空横线约四个字宽，本稿按全稿惯例统一排 \blankline[4.4em]；
%   ④ 原卷实数集排斜体 $R$（第 9、10、19 题题干与解析）或粗体 $\mathbf{R}$（第 18 题题干），
%      整数集排斜体 $Z$（第 12 题解析），自然数集排斜体 $N$（第 20 题解析、第 21 题题干
%      “$(n\geqslant2,n\in N)$”）或粗体 $\mathbf{N}$（第 21 题 (3) 题干）、斜体 $N^*$
%      （第 21 题 (3) 解析），本稿按全稿惯例统一改用 $\R$、$\Z$、$\N$、$\N^*$
%      （同 高一14 排版差异④：记号归一，粗体／斜体不作区分，$N$ 与 $N^*$ 的区别照原卷保留）；
%      空集记号原卷作 $\varnothing$，照录；
%   ⑤ 原卷在数学式内多处用半角逗号（如 $\{y\mid y=[2x]+[3x],-1\leqslant x\leqslant2\}$、
%      $m=-1,0,1$、$A=\{x\mid-2<x<15\}$,$B=\cdots$），均照录未强改；
%      句末句点按全稿惯例排西文句点；“和积平衡集”自定义名词照录；
%   ⑥ 第 12 题解析里的“法一：”“法二：”标号原卷排**红色**（全卷共 2 处），
%      同 高一53“情况1：”成例，本稿排 \stepm{法一：}{…}、\stepm{法二：}{…}
%      （加粗着色、不改黑色基调），冒号照原卷排全角；
%   ⑦ 第 17—21 题的子问原卷各自另起一行，本稿按全稿惯例用 enumerate 的 (1)(2)(3)；
%      解析中的子问标号一律用 \stepm{(1)}{…}；第 19 题解析中的①②分情形，按全稿惯例作
%      \stepm{①}{…}、\stepm{②}{…}；
%   ⑧ 第 14 题印在题干括号内的答案移入 \ans{}，见上；
%   ⑨ 每道解答题原卷未留作答空白，本稿按全稿惯例加 \ansspace；
%   ⑩ 选择题的作答括号原卷排半角“(　　)”（第 13、15、16 题括号内空白、第 14 题括号内印 D），
%      本稿按全稿惯例排 \mbox{（\quad）}。
%
% 原卷笔误/存疑（均照录未改，正文对应处另有 % 注）：
%   ① 第 21 题题干“$(n\geqslant2$,$n\in N)$”后连印两个全角逗号“，，”（衍字，
%      8 倍放大确认，非水印干扰），本稿照录；
%   ② 第 21 题 (2) 题干将“写出以 $a_1$、$a_2$ 为根的一个一元二次方程”与“用反证法证明：
%      $a_1$、$a_2$ 至少有一个大于 2”两项任务并作一句、之间仅以逗号相接（“……表示），
%      可尝试用反证法证明：……”），读来像抄并两小问，原卷如此，照录未改。
%
% 插图：无（全卷 10 页无任何图形、无“如图”字样）。
%
\documentclass[a4paper,11pt]{article}
\usepackage{common}

\begin{document}

\examtitlest[3]{2026--2027 年吴淞中学高一上 10 月月考}{集合与常用逻辑用语、不等式}

\bigsec{一、填空题}

\begin{enumerate}
\setcounter{enumi}{2}

\item 满足 $A\subseteq\{1,2,3,4\}$ 的集合 $A$ 共有\blankline[4.4em]个.

\ifanswer
\solbegin
\step{由题意，集合 $A$ 是集合 $\{1,2,3,4\}$ 的子集，共有 $2^4=16$ 个.}
\solend
\fi

\item 陈述句“$a>1$ 或 $b>3$”的否定形式为\blankline[4.4em].

\ifanswer
\ans{$a\leqslant1$ 且 $b\leqslant3$}
\fi

\item 已知方程 $3x^2+4x-4=0$ 的两个根为 $x_1$、$x_2$，则 $x_1^2+x_2^2=$\blankline[4.4em].

\ifanswer
\solbegin
\step{对于方程 $3x^2+4x-4=0$，因为 $\Delta=16-4\times3\times(-4)=64>0$，}
\step{所以方程 $3x^2+4x-4=0$ 的两个不相等的实根为 $x_1$、$x_2$，}
\step{由韦达定理可得 $x_1+x_2=-\dfrac43$，$x_1x_2=-\dfrac43$，}
\step{所以 $x_1^2+x_2^2=(x_1+x_2)^2-2x_1x_2=\dfrac{16}9+\dfrac83=\dfrac{40}9$.}
\solend
\fi

\item 设集合 $A=\{y\mid y=x^2+1\}$,$B=\{y\mid y=-x^2+2x+3\}$，则 $A\cap B=$\blankline[4.4em].

\ifanswer
\solbegin
\step{因为 $y=x^2+1\geqslant1$,$y=-x^2+2x+3=-(x-1)^2+4\leqslant4$，}
\step{所以 $A=\{y\mid y\geqslant1\}$,$B=\{y\mid y\leqslant4\}$，得 $A\cap B=\{y\mid1\leqslant y\leqslant4\}$.}
\solend
\fi

\item 若 $x\in\{1,x^2\}$，则 $x=$\blankline[4.4em].

\ifanswer
\ans{$0$}
\fi

\item 已知等式 $2x^2+3x+5=a(2x+1)(x+1)+c$，则 $a+c=$\blankline[4.4em].

\ifanswer
\solbegin
\step{因为 $2x^2+3x+5=a(2x+1)(x+1)+c$ 恒成立，}
\step{所以 $2x^2+3x+5=2ax^2+3ax+a+c$ 恒成立，则 $\begin{cases}2=2a\\ 3=3a\\ 5=a+c\end{cases}$，解得 $a+c=5$.}
\solend
\fi

\item 已知 $a\in\R$，设 $p:x<2$，$q:x<5a-3$，且 $p$ 是 $q$ 的必要非充分条件，则 $a$ 的取值范围是\blankline[4.4em].

\ifanswer
\solbegin
\step{设 $p:x<2$，$q:x<5a-3$，且 $p$ 是 $q$ 的必要非充分条件，}
\step{则 $\{x\mid x<5a-3\}\subset\{x\mid x<2\}$，所以 $5a-3<2$，即 $a<1$.}
\solend
\fi

\item 已知 $a\in\R$，集合 $A=\{x\mid ax^2+3x-1=0\}$，若集合 $A$ 中只有一个元素，则 $a$ 的取值范围是\blankline[4.4em].

\ifanswer
\solbegin
\step{当 $a=0$ 时，则有 $\{x\mid ax^2+3x-1=0\}=\{x\mid3x-1=0\}=\left\{\dfrac13\right\}$，符合题意；}
\step{当 $a\neq0$ 时，由题意得 $\Delta=9+4a=0$，解得 $a=-\dfrac94$. 符合题意；}
\step{综上所述，实数 $a$ 的取值集合为 $\left\{0,-\dfrac94\right\}$.}
\solend
\fi

\item 设集合 $A=\{x\mid x^2+2(m+1)x+m^2-1=0,m\in\R\}$，且 $A\subseteq\{0,-4\}$，则实数 $m$ 的取值范围是\blankline[4.4em].

\ifanswer
\solbegin
\step{由题意得 $A=\varnothing$ 或 $A=\{0\}$ 或 $A=\{-4\}$ 或 $A=\{0,-4\}$，}
\step{当 $A=\varnothing$ 时，方程 $x^2+2(m+1)x+m^2-1=0$ 无实根，}
\step{则 $\Delta=[2(m+1)]^2-4\times1\times(m^2-1)=8m+8<0$，解得 $m<-1$，}
\step{当 $A=\{0\}$ 时，方程 $x^2+2(m+1)x+m^2-1=0$ 有一个实根 $0$，}
\step{则 $\begin{cases}\Delta=8m+8=0\\ m^2-1=0\end{cases}$，解得 $m=-1$，}
\step{当 $A=\{-4\}$ 时，方程 $x^2+2(m+1)x+m^2-1=0$ 有一个实根 $-4$，}
\step{则 $\begin{cases}\Delta=8m+8=0\\ 16-8(m+1)+m^2-1=0\end{cases}$，无解，}
\step{当 $A=\{0,-4\}$ 时，方程 $x^2+2(m+1)x+m^2-1=0$ 有两个实根 $-4$ 和 $0$，}
\step{则 $\begin{cases}\Delta=8m+8>0\\ 16-8(m+1)+m^2-1=0\\ m^2-1=0\end{cases}$，解得 $m=1$，}
\step{综上所述，实数 $m$ 的取值范围是 $\{m\mid m\leqslant-1\text{ 或 }m=1\}$.}
\solend
\fi

\item 对于任意实数 $x$，$[x]$ 表示不超过 $x$ 的最大整数，如 $[1.1]=1$，$[-1.2]=-2$，若
$A=\{y\mid y=[2x]+[3x],-1\leqslant x\leqslant2\}$，则 $A$ 中所有元素的和为\blankline[4.4em].

\ifanswer
\solbegin
% 原卷“法一：”“法二：”标号排红色，本稿按 高一53“情况1：”成例排 \stepm（加粗着色），见文件头
\stepm{法一：}{当 $-1\leqslant x<-\dfrac23$ 时，$[2x]=-2$，$[3x]=-3$，$f(x)=[2x]+[3x]=-5$，}
\step{当 $-\dfrac23\leqslant x<-\dfrac12$ 时，$[2x]=-2$，$[3x]=-2$，$f(x)=[2x]+[3x]=-4$；}
\step{当 $-\dfrac12\leqslant x<-\dfrac13$ 时，$[2x]=-1$，$[3x]=-2$，$f(x)=[2x]+[3x]=-3$；}
\step{当 $-\dfrac13\leqslant x<0$ 时，$[2x]=-1$，$[3x]=-1$，$f(x)=[2x]+[3x]=-2$；}
\step{当 $0\leqslant x<\dfrac13$ 时，$[2x]=0$，$[3x]=0$，$f(x)=[2x]+[3x]=0$；}
\step{当 $\dfrac13\leqslant x<\dfrac12$ 时，$[2x]=0$，$[3x]=1$，$f(x)=[2x]+[3x]=1$；}
\step{当 $\dfrac12\leqslant x<\dfrac23$ 时，$[2x]=1$，$[3x]=1$，$f(x)=[2x]+[3x]=2$；}
\step{当 $\dfrac23\leqslant x<1$ 时，$[2x]=1$，$[3x]=2$，$f(x)=[2x]+[3x]=3$；}
\step{当 $1\leqslant x<\dfrac43$ 时，$[2x]=2$，$[3x]=3$，$f(x)=[2x]+[3x]=5$；}
\step{当 $\dfrac43\leqslant x<\dfrac32$ 时，$[2x]=2$，$[3x]=4$，$f(x)=[2x]+[3x]=6$；}
\step{当 $\dfrac32\leqslant x<\dfrac53$ 时，$[2x]=3$，$[3x]=4$，$f(x)=[2x]+[3x]=7$；}
\step{当 $\dfrac53\leqslant x<2$ 时，$[2x]=3$，$[3x]=5$，$f(x)=[2x]+[3x]=8$；}
\step{当 $x=2$ 时，$[2x]=4$，$[3x]=6$，$f(x)=[2x]+[3x]=10$}
\step{故 $A$ 中所有元素的和为 $-5-4-3-2+0+1+2+3+5+6+7+8+10=28$.}
% 原卷“法二：”标号排红色，同上
\stepm{法二：}{设 $x=m+n$，其中 $m\in\Z$，$n\in[0,1)$，}
\step{若 $0\leqslant n<\dfrac13$，则 $0\leqslant2n<\dfrac23$，$0\leqslant3n<1$，}
\step{得 $[2x]=[2m+2n]=2m$，$[3x]=[3m+3n]=3m$，所以 $y=[2x]+[3x]=5m$；}
\step{若 $\dfrac13\leqslant n<\dfrac12$，则 $\dfrac23\leqslant2n<1$，$1\leqslant3n<\dfrac32$，}
\step{得 $[2x]=[2m+2n]=2m$，$[3x]=[3m+3n]=3m+1$，}
\step{所以 $y=[2x]+[3x]=5m+1$；}
\step{若 $\dfrac12\leqslant n<\dfrac23$，则 $1\leqslant2n<\dfrac43$，$\dfrac32\leqslant3n<2$，}
\step{得 $[2x]=[2m+2n]=2m+1$，$[3x]=[3m+3n]=3m+1$，}
\step{所以 $y=[2x]+[3x]=5m+2$；}
\step{若 $\dfrac23\leqslant n<1$，则 $\dfrac43\leqslant2n<2$，$2\leqslant3n<3$，}
\step{得 $[2x]=[2m+2n]=2m+1$，$[3x]=[3m+3n]=3m+2$，}
\step{所以 $y=[2x]+[3x]=5m+3$；}
\step{综上所述，当 $n\in[0,1)$ 时，所有函数值的和}
\step{为 $5m+5m+1+5m+2+5m+3=20m+6$，}
\step{又因为 $x\in[-1,2]$，此时 $m=-1,0,1$，且当 $x=2$ 时，$y=[2x]+[3x]=10$，}
\step{所有函数值的和为 $(-14)+6+26+10=28$，即 $A$ 中所有元素之和为 $28$.}
\solend
\fi

\end{enumerate}

\bigsec{二、选择题}

\begin{enumerate}
\setcounter{enumi}{12}

\item 下列写法中，正确的是\mbox{（\quad）}

\keepwithprev
A. $0\in\{\varnothing\}$\qquad B. $\varnothing\in\{\varnothing\}$

C. $\varnothing=\{\varnothing\}$\qquad D. $0\in\varnothing$

\ifanswer
\solbegin
\step{$0\notin\{\varnothing\}$，$\varnothing\in\{\varnothing\}$，故 $A$ 错误，$B$ 正确；}
\step{$\varnothing$ 是不含任何元素的集合，所以 $\varnothing\neq\{\varnothing\}$，$0\notin\varnothing$，故 $CD$ 错误.}
\step{故选 $B$.}
\solend
\fi

% 第 14 题答案“D”直接印在题干括号内“(  D  )”，且该题无【答案】/【解析】段；
% 按 高一37 第 13 题的成例括号留空、答案另提为 \ans{}（见文件头）
\item 若 $a<b<0$，则下列不等式中不能成立的是\mbox{（\quad）}

\keepwithprev
A. $\dfrac1a>\dfrac1b$\qquad B. $a-b>a$

C. $|a|>|b|$\qquad D. $a^3>b^3$

\ifanswer
\ans{$D$}
\fi

\item 已知 $x$、$y$、$z$ 为非零实数，代数式 $\dfrac x{|x|}+\dfrac y{|y|}+\dfrac z{|z|}+\dfrac{xyz}{|xyz|}$
的值所组成的集合是 $M$，则下列判断正确的是\mbox{（\quad）}

\keepwithprev
A. $-4\notin M$\qquad B. $2\in M$

C. $4\notin M$\qquad D. $0\in M$

\ifanswer
\solbegin
\step{当 $x$、$y$、$z$ 均为正数时，$xyz>0$，$\dfrac x{|x|}+\dfrac y{|y|}+\dfrac z{|z|}+\dfrac{xyz}{|xyz|}=1+1+1+1=4$，}
\step{当 $x$、$y$、$z$ 有 2 个正数，1 个负数时，不妨设 $x>0$，$y>0$，$z<0$，}
\step{则 $xyz<0$，$\dfrac x{|x|}+\dfrac y{|y|}+\dfrac z{|z|}+\dfrac{xyz}{|xyz|}=1+1-1-1=0$，}
\step{若 $x$、$y$、$z$ 有 1 个正数，2 个负数时，不妨设 $x>0$，$y<0$，$z<0$，}
\step{则 $xyz>0$，$\dfrac x{|x|}+\dfrac y{|y|}+\dfrac z{|z|}+\dfrac{xyz}{|xyz|}=1-1-1+1=0$，}
\step{当 $x$、$y$、$z$ 均为负数时，$xyz<0$，}
\step{$\dfrac x{|x|}+\dfrac y{|y|}+\dfrac z{|z|}+\dfrac{xyz}{|xyz|}=-1-1-1-1=-4$，}
\step{综上，$M=\{-4,0,4\}$，}
\step{所以 $-4\in M$，$A$ 错误；$2\notin M$，$B$ 错误；$4\in M$，$C$ 错误；}
\step{$0\in M$，$D$ 正确. 故选 $D$.}
\solend
\fi

\item 若 $xy>0$，且 $x+y=9$，则 $\dfrac1x+\dfrac4y$ 的最小值是\mbox{（\quad）}

\keepwithprev
A. $9$\qquad B. $\dfrac59$

C. $\dfrac49$\qquad D. $1$

\ifanswer
\solbegin
\step{$xy>0$，且 $x+y=9$，所以 $x>0$，$y>0$，}
\step{所以 $\dfrac1x+\dfrac4y=\dfrac19\left(\dfrac1x+\dfrac4y\right)(x+y)=\dfrac59+\dfrac19\left(\dfrac yx+\dfrac{4x}y\right)\geqslant\dfrac59+\dfrac29\cdot\sqrt{\dfrac yx\cdot\dfrac{4x}y}=\dfrac59+\dfrac49=1$，}
\step{当且仅当 $\dfrac yx=\dfrac{4x}y$，即 $x=3$，$y=6$ 时取等号，所以 $\dfrac1x+\dfrac4y$ 的最小值为 $1$.}
\step{故选 $D$.}
\solend
\fi

\end{enumerate}

\bigsec{三、解答题}

\begin{enumerate}
\setcounter{enumi}{16}

\item 求解下列不等式：

（1）$\dfrac{x-3}{x+1}\leqslant0$；

（2）$|2x-1|<3$.

\ifanswer
\solbegin
\stepm{(1)}{$\dfrac{x-3}{x+1}\leqslant0$，变形得 $(x-3)\cdot(x+1)\leqslant0$ 且 $x\neq-1$，}
\step{解得 $-1<x\leqslant3$，即不等式的解集为 $(-1,3]$.}
\stepm{(2)}{若 $|2x-1|<3$，则 $-3<2x-1<3$，解得 $-1<x<2$，即 $(-1,2)$.}
\solend
\fi

\ansspace[6]

\item 已知集合 $A=\{x\mid-2<x<15\}$,$B=\{x\mid m-6<x<2m-1,m\in\R\}$.

（1）当 $m=2$ 时，求 $A\cap B$；

（2）若 $A\cup B=A$，求实数 $m$ 的取值范围.

\ifanswer
\solbegin
\stepm{(1)}{当 $m=2$ 时，$A=\{x\mid-2<x<15\}$,$B=\{x\mid-4<x<3\}$，}
\step{此时 $A\cap B=\{x\mid-2<x<3\}$.}
\stepm{(2)}{由 $A\cup B=A$ 得 $B\subseteq A$.}
\step{当 $B=\varnothing$,$m-6\geqslant2m-1$，解得 $m\leqslant-5$；}
\step{当 $B\neq\varnothing$，结合 $B\subseteq A$ 得 $\begin{cases}m-6<2m-1\\ 2m-1\leqslant15\\ m-6\geqslant-2\end{cases}$，解得 $4\leqslant m\leqslant8$.}
\step{综上，$m\leqslant-5$ 或 $4\leqslant m\leqslant8$.}
\solend
\fi

\ansspace[6]

\item 已知 $a\in\R$，命题 $p$：关于 $x$ 的方程 $x^2-2x+a=0$ 无实根；命题 $q$：对任意 $x\in\R$，
不等式 $ax^2-ax+1>0$ 恒成立.

（1）若 $q$ 为真命题，求 $a$ 的取值范围；

（2）若 $p$、$q$ 中有且只有一个为真命题，求 $a$ 的取值范围.

\ifanswer
\solbegin
\stepm{(1)}{当 $q$ 为真命题时，}
\stepm{①}{当 $a=0$ 时，不等式化为 $1>0$，对任意 $x\in\R$ 恒成立，满足条件；}
\stepm{②}{当 $a\neq0$ 时，$ax^2-ax+1>0$ 为一元二次不等式，}
\step{若对任意 $x\in\R$ 恒成立，需满足 $\begin{cases}a>0\\ \Delta=(-a)^2-4a<0\end{cases}$，解得 $0<a<4$；}
\step{综上，$a$ 的取值范围为 $[0,4)$.}
\stepm{(2)}{$p$ 为真时，方程 $x^2-2x+a=0$ 无实根，}
\step{故判别式 $\Delta=4-4a<0$，解得 $a>1$，因此 $p$ 为假时，$a\leqslant1$；}
\step{由（1）得 $q$ 为假时 $a\geqslant4$ 或 $a<0$，}
\step{$p$、$q$ 有且只有一个真命题，分两类：}
\step{当 $p$ 假 $q$ 真时，满足 $\begin{cases}a\leqslant1\\ 0\leqslant a<4\end{cases}$，解得 $0\leqslant a\leqslant1$；}
\step{当 $p$ 真 $q$ 假时，满足 $\begin{cases}a>1\\ a\geqslant4\text{或}a<0\end{cases}$，解得 $a\geqslant4$.}
\step{综上，$a$ 的取值范围是 $[0,1]\cup[4,+\infty)$.}
\solend
\fi

\ansspace[8]

\item 某工厂生产 $x$($x>0$) 件产品的总成本 $C$（单位：元）与生产件数 $x$ 之间满足二次函数
$C(x)=\dfrac14x^2+50x+10000$.

（1）若要求生产总成本不超过 30000 元，求满足生产总成本要求的产品件数的取值范围；

（2）试求该工厂生产多少件产品时，可使每件产品的平均成本最小.

\ifanswer
\solbegin
\stepm{(1)}{若要求生产总成本不超过 30000 元，}
\step{则 $C(x)=\dfrac14x^2+50x+10000\leqslant30000$，}
\step{即 $x^2+200x-80000\leqslant0$，解得 $-400\leqslant x\leqslant200$，又 $x>0$，}
\step{所以满足生产总成本要求的产品件数的取值范围为 $(0,200]$，$x\in\N$；}
\stepm{(2)}{生产件数 $x$ 时每件的平均成本为 $\dfrac{\frac14x^2+50x+10000}x=\dfrac14x+50+\dfrac{10000}x$，}
\step{由基本不等式得 $\dfrac14x+50+\dfrac{10000}x\geqslant2\sqrt{\dfrac14x\cdot\dfrac{10000}x}+50=150$，}
\step{当且仅当 $\dfrac14x=\dfrac{10000}x$，即 $x=200$ 时取等号，}
\step{即该工厂生产 200 件产品时，每件产品的平均成本最小.}
\solend
\fi

\ansspace[8]

% 原卷题干“$(n\geqslant2,n\in N)$”后连印两个全角逗号“，，”，照录未改（见文件头“原卷笔误”①）
\item 已知有限集 $A=\{a_1,a_2,\cdots,a_n\}$($n\geqslant2$,$n\in\N$)，，如果 $A$ 中元素
$a_i$($i=1,2,3,\cdots,n$) 满足 $a_1+a_2+a_3+\cdots+a_n=a_1\times a_2\times\cdots\times a_n$，
就称 $A$ 为“和积平衡集”.

% 原卷 (2) 将“写方程”与“反证法证明”两项任务并作一句（“……表示），可尝试用反证法证明：……”），
% 照录未改（见文件头“原卷笔误”②）
（1）判断集合 $\{-1+\sqrt3,-1-\sqrt3\}$ 是否为“和积平衡集”；

（2）若 $a_1$、$a_2$ 是两个不同的正数，且 $\{a_1,a_2\}$ 是“和积平衡集”，写出以 $a_1$、$a_2$ 为根的
一个一元二次方程（系数可用 $a_1$、$a_2$ 表示），可尝试用反证法证明：$a_1$、$a_2$ 至少有一个大于 2；

（3）若 $a_i\geqslant1$，且 $a_i\in\N$，求所有符合条件的“和积平衡集”.

\ifanswer
\solbegin
\stepm{(1)}{对于集合 $\{-1+\sqrt3,-1-\sqrt3\}$，}
\step{因为 $(-1+\sqrt3)+(-1-\sqrt3)=-2$，$(-1+\sqrt3)\times(-1-\sqrt3)=-2$，}
\step{所以满足 $a_1+a_2=a_1\times a_2$，满足题意，}
\step{所以集合 $\{-1+\sqrt3,-1-\sqrt3\}$ 是和积平衡集.}
\stepm{(2)}{以 $a_1$、$a_2$ 为根的一个一元二次方程可为 $x^2-(a_1+a_2)x+a_1\cdot a_2=0$，}
\step{因为 $a_1$、$a_2$ 是两个不同的正数，且 $\{a_1,a_2\}$ 是和积平衡集，}
\step{所以 $a_1+a_2=a_1\cdot a_2$，}
\step{又 $a_1$ 和 $a_2$ 为 $x^2-(a_1+a_2)x+a_1\cdot a_2=0$ 的两个不同的正根，}
\step{所以 $\Delta=(a_1+a_2)^2-4a_1a_2=(a_1+a_2)^2-4(a_1+a_2)>0$，}
\step{又 $a_1+a_2>0$，则 $a_1+a_2>4$，}
\step{若 $a_1$、$a_2$ 都小于等于 2，则 $a_1+a_2<4$，矛盾，}
\step{所以 $a_1$、$a_2$ 至少有一个大于 2；}
\stepm{(3)}{设 $A$ 中的 $a_1<a_2<a_3<\cdots<a_n$，且 $a_i\in\N^*$，}
\step{由 $a_1+a_2+a_3+\cdots+a_n=a_1\times a_2\times\cdots\times a_n$，得 $n=2$ 时明显不满足，}
\step{所以 $n\geqslant3$，故 $a_1\times a_2\times\cdots\times a_n=a_1+a_2+a_3+\cdots+a_n\leqslant na_n$，}
\step{所以 $a_1\times a_2\times\cdots\times a_{n-1}\leqslant n$，}
\step{当 $n=3$ 时，$a_1\times a_2\leqslant3$，$a_i\in\N^*$，只能有 $a_1=1$，$a_2=2$，}
\step{由 $a_1+a_2+a_3=a_1\times a_2\times a_3$，得 $a_3=3$，此时的和积平衡集为 $\{1,2,3\}$，}
\step{当 $n\geqslant4$ 时，$a_1\times a_2\times\cdots\times a_{n-1}\geqslant(n-1)(n-2)\times\cdots\times2\times1\geqslant(n-1)(n-2)$，}
\step{又 $(n-1)(n-2)-n=n^2-4n+2=n(n-4)+2>0$，即 $(n-1)(n-2)>n$，}
\step{即 $a_1\times a_2\times\cdots\times a_{n-1}>n$，与 $a_1\times a_2\times\cdots\times a_{n-1}\leqslant n$ 矛盾，}
\step{所以 $n\geqslant4$ 不满足条件，}
\step{所以符合条件的和积平衡集有且只有 $n=3$，此时的和积平衡集为 $\{1,2,3\}$.}
\solend
\fi

\ansspace*[10]

\end{enumerate}

\end{document}
"
%
% 原始材料：微信公众号“魔都数学加油站”文章，作者破晓，连载“27学年高一63”；
% 原文标题“27学年高一63——2026-2027年上海市吴淞中学高一上10月月考”；
% 原文链接：https://mp.weixin.qq.com/s/TFQ4TLnt_l66t7t51es4OQ。
% 页面标注发布时间 2026-10-16 21:44（上海）。
% 正文为 10 张 PNG 页（`原图/高一63_吴淞中学_10月月考/01..10.png`）：
% 第 06、10 页 2560×3620，第 01、02、03、04、05、07、08、09 页 4762×6735
% （**同一篇各页分辨率不同**），均按页序读取。10 页全为试卷正文页（卷首标题在第 1 页页顶，
% 卷末解答题第 21 题解析止于第 10 页），无封面页、目录页、二维码或宣传页。
% 粉色斜排“魔都数学加油站”为卷面水印，与公众号页脚一并未录入。
%
% 卷首标题：照卷面印作“2026--2027 年吴淞中学高一上 10 月月考”（公众号标题中的“上海市”
% 卷面标题行没有，本稿照卷面），年份区间按全稿惯例排 en dash。原卷只有一行标题、无副标题。
%
% 篇幅与题号：原卷自第 3 题起编号（第 1、2 题不在本卷内）。填空题 3—12（10 题）、
% 选择题 13—16（4 题）、解答题 17—21（5 题），共 19 题。本稿沿用原节与原题号，
% 只改排为全稿统一的大节样式（原卷小节标题为“一、填空题”“二、选择题”“三、解答题”）。
%
% 讲评版：原卷每题下直接印【答案】或【解析】。**只印【答案】的 2 题**（第 4、7 题），
% 本稿只给 \ans{}；其余 17 题里 16 题都印【解析】（按全稿口径这类题不另提 \ans{}——第 13、15、16 题
% 的结论写在解析末句（“故选 B”等），亦照录在解析内；第 14 题则既无【答案】也无【解析】段。**另有 1 题答案印在题干括号内**——
% 第 14 题（“(  D  )”，且该题无【答案】/【解析】段），本稿按 高一37 第 13 题成例把括号留空、
% 答案另提为 \ans{$D$}。第 16 题题干括号为空但印有【解析】，不算括号内答案。
%
% 副标题：原卷未印副标题。本稿按“知识模块”口径标作“集合与常用逻辑用语、不等式”：
% 19 题中集合与常用逻辑用语 12 题（第 3、4、6、7、9、10、11、12、13、15、18、21 题），
% 不等式 7 题（含等式（韦达定理、恒等式）与不等式性质、基本不等式、二次不等式、恒成立：
% 第 5、8、14、16、17、19、20 题），两类均超过三分之一，故并标。口径同 高一59。
%
% 与原卷的排版差异：
%   ① 原卷每题下直接印【答案】/【解析】，本稿按全稿惯例将答案与解析置于答案版，试卷版保持干净；
%   ② 原卷未印副标题，本稿按“知识模块”口径补出；
%   ③ 原卷填空横线约四个字宽，本稿按全稿惯例统一排 \blankline[4.4em]；
%   ④ 原卷实数集排斜体 $R$（第 9、10、19 题题干与解析）或粗体 $\mathbf{R}$（第 18 题题干），
%      整数集排斜体 $Z$（第 12 题解析），自然数集排斜体 $N$（第 20 题解析、第 21 题题干
%      “$(n\geqslant2,n\in N)$”）或粗体 $\mathbf{N}$（第 21 题 (3) 题干）、斜体 $N^*$
%      （第 21 题 (3) 解析），本稿按全稿惯例统一改用 $\R$、$\Z$、$\N$、$\N^*$
%      （同 高一14 排版差异④：记号归一，粗体／斜体不作区分，$N$ 与 $N^*$ 的区别照原卷保留）；
%      空集记号原卷作 $\varnothing$，照录；
%   ⑤ 原卷在数学式内多处用半角逗号（如 $\{y\mid y=[2x]+[3x],-1\leqslant x\leqslant2\}$、
%      $m=-1,0,1$、$A=\{x\mid-2<x<15\}$,$B=\cdots$），均照录未强改；
%      句末句点按全稿惯例排西文句点；“和积平衡集”自定义名词照录；
%   ⑥ 第 12 题解析里的“法一：”“法二：”标号原卷排**红色**（全卷共 2 处），
%      同 高一53“情况1：”成例，本稿排 \stepm{法一：}{…}、\stepm{法二：}{…}
%      （加粗着色、不改黑色基调），冒号照原卷排全角；
%   ⑦ 第 17—21 题的子问原卷各自另起一行，本稿按全稿惯例用 enumerate 的 (1)(2)(3)；
%      解析中的子问标号一律用 \stepm{(1)}{…}；第 19 题解析中的①②分情形，按全稿惯例作
%      \stepm{①}{…}、\stepm{②}{…}；
%   ⑧ 第 14 题印在题干括号内的答案移入 \ans{}，见上；
%   ⑨ 每道解答题原卷未留作答空白，本稿按全稿惯例加 \ansspace；
%   ⑩ 选择题的作答括号原卷排半角“(　　)”（第 13、15、16 题括号内空白、第 14 题括号内印 D），
%      本稿按全稿惯例排 \mbox{（\quad）}。
%
% 原卷笔误/存疑（均照录未改，正文对应处另有 % 注）：
%   ① 第 21 题题干“$(n\geqslant2$,$n\in N)$”后连印两个全角逗号“，，”（衍字，
%      8 倍放大确认，非水印干扰），本稿照录；
%   ② 第 21 题 (2) 题干将“写出以 $a_1$、$a_2$ 为根的一个一元二次方程”与“用反证法证明：
%      $a_1$、$a_2$ 至少有一个大于 2”两项任务并作一句、之间仅以逗号相接（“……表示），
%      可尝试用反证法证明：……”），读来像抄并两小问，原卷如此，照录未改。
%
% 插图：无（全卷 10 页无任何图形、无“如图”字样）。
%
\documentclass[a4paper,11pt]{article}
\usepackage{common}

\begin{document}

\examtitlest[3]{2026--2027 年吴淞中学高一上 10 月月考}{集合与常用逻辑用语、不等式}

\bigsec{一、填空题}

\begin{enumerate}
\setcounter{enumi}{2}

\item 满足 $A\subseteq\{1,2,3,4\}$ 的集合 $A$ 共有\blankline[4.4em]个.

\ifanswer
\solbegin
\step{由题意，集合 $A$ 是集合 $\{1,2,3,4\}$ 的子集，共有 $2^4=16$ 个.}
\solend
\fi

\item 陈述句“$a>1$ 或 $b>3$”的否定形式为\blankline[4.4em].

\ifanswer
\ans{$a\leqslant1$ 且 $b\leqslant3$}
\fi

\item 已知方程 $3x^2+4x-4=0$ 的两个根为 $x_1$、$x_2$，则 $x_1^2+x_2^2=$\blankline[4.4em].

\ifanswer
\solbegin
\step{对于方程 $3x^2+4x-4=0$，因为 $\Delta=16-4\times3\times(-4)=64>0$，}
\step{所以方程 $3x^2+4x-4=0$ 的两个不相等的实根为 $x_1$、$x_2$，}
\step{由韦达定理可得 $x_1+x_2=-\dfrac43$，$x_1x_2=-\dfrac43$，}
\step{所以 $x_1^2+x_2^2=(x_1+x_2)^2-2x_1x_2=\dfrac{16}9+\dfrac83=\dfrac{40}9$.}
\solend
\fi

\item 设集合 $A=\{y\mid y=x^2+1\}$,$B=\{y\mid y=-x^2+2x+3\}$，则 $A\cap B=$\blankline[4.4em].

\ifanswer
\solbegin
\step{因为 $y=x^2+1\geqslant1$,$y=-x^2+2x+3=-(x-1)^2+4\leqslant4$，}
\step{所以 $A=\{y\mid y\geqslant1\}$,$B=\{y\mid y\leqslant4\}$，得 $A\cap B=\{y\mid1\leqslant y\leqslant4\}$.}
\solend
\fi

\item 若 $x\in\{1,x^2\}$，则 $x=$\blankline[4.4em].

\ifanswer
\ans{$0$}
\fi

\item 已知等式 $2x^2+3x+5=a(2x+1)(x+1)+c$，则 $a+c=$\blankline[4.4em].

\ifanswer
\solbegin
\step{因为 $2x^2+3x+5=a(2x+1)(x+1)+c$ 恒成立，}
\step{所以 $2x^2+3x+5=2ax^2+3ax+a+c$ 恒成立，则 $\begin{cases}2=2a\\ 3=3a\\ 5=a+c\end{cases}$，解得 $a+c=5$.}
\solend
\fi

\item 已知 $a\in\R$，设 $p:x<2$，$q:x<5a-3$，且 $p$ 是 $q$ 的必要非充分条件，则 $a$ 的取值范围是\blankline[4.4em].

\ifanswer
\solbegin
\step{设 $p:x<2$，$q:x<5a-3$，且 $p$ 是 $q$ 的必要非充分条件，}
\step{则 $\{x\mid x<5a-3\}\subset\{x\mid x<2\}$，所以 $5a-3<2$，即 $a<1$.}
\solend
\fi

\item 已知 $a\in\R$，集合 $A=\{x\mid ax^2+3x-1=0\}$，若集合 $A$ 中只有一个元素，则 $a$ 的取值范围是\blankline[4.4em].

\ifanswer
\solbegin
\step{当 $a=0$ 时，则有 $\{x\mid ax^2+3x-1=0\}=\{x\mid3x-1=0\}=\left\{\dfrac13\right\}$，符合题意；}
\step{当 $a\neq0$ 时，由题意得 $\Delta=9+4a=0$，解得 $a=-\dfrac94$. 符合题意；}
\step{综上所述，实数 $a$ 的取值集合为 $\left\{0,-\dfrac94\right\}$.}
\solend
\fi

\item 设集合 $A=\{x\mid x^2+2(m+1)x+m^2-1=0,m\in\R\}$，且 $A\subseteq\{0,-4\}$，则实数 $m$ 的取值范围是\blankline[4.4em].

\ifanswer
\solbegin
\step{由题意得 $A=\varnothing$ 或 $A=\{0\}$ 或 $A=\{-4\}$ 或 $A=\{0,-4\}$，}
\step{当 $A=\varnothing$ 时，方程 $x^2+2(m+1)x+m^2-1=0$ 无实根，}
\step{则 $\Delta=[2(m+1)]^2-4\times1\times(m^2-1)=8m+8<0$，解得 $m<-1$，}
\step{当 $A=\{0\}$ 时，方程 $x^2+2(m+1)x+m^2-1=0$ 有一个实根 $0$，}
\step{则 $\begin{cases}\Delta=8m+8=0\\ m^2-1=0\end{cases}$，解得 $m=-1$，}
\step{当 $A=\{-4\}$ 时，方程 $x^2+2(m+1)x+m^2-1=0$ 有一个实根 $-4$，}
\step{则 $\begin{cases}\Delta=8m+8=0\\ 16-8(m+1)+m^2-1=0\end{cases}$，无解，}
\step{当 $A=\{0,-4\}$ 时，方程 $x^2+2(m+1)x+m^2-1=0$ 有两个实根 $-4$ 和 $0$，}
\step{则 $\begin{cases}\Delta=8m+8>0\\ 16-8(m+1)+m^2-1=0\\ m^2-1=0\end{cases}$，解得 $m=1$，}
\step{综上所述，实数 $m$ 的取值范围是 $\{m\mid m\leqslant-1\text{ 或 }m=1\}$.}
\solend
\fi

\item 对于任意实数 $x$，$[x]$ 表示不超过 $x$ 的最大整数，如 $[1.1]=1$，$[-1.2]=-2$，若
$A=\{y\mid y=[2x]+[3x],-1\leqslant x\leqslant2\}$，则 $A$ 中所有元素的和为\blankline[4.4em].

\ifanswer
\solbegin
% 原卷“法一：”“法二：”标号排红色，本稿按 高一53“情况1：”成例排 \stepm（加粗着色），见文件头
\stepm{法一：}{当 $-1\leqslant x<-\dfrac23$ 时，$[2x]=-2$，$[3x]=-3$，$f(x)=[2x]+[3x]=-5$，}
\step{当 $-\dfrac23\leqslant x<-\dfrac12$ 时，$[2x]=-2$，$[3x]=-2$，$f(x)=[2x]+[3x]=-4$；}
\step{当 $-\dfrac12\leqslant x<-\dfrac13$ 时，$[2x]=-1$，$[3x]=-2$，$f(x)=[2x]+[3x]=-3$；}
\step{当 $-\dfrac13\leqslant x<0$ 时，$[2x]=-1$，$[3x]=-1$，$f(x)=[2x]+[3x]=-2$；}
\step{当 $0\leqslant x<\dfrac13$ 时，$[2x]=0$，$[3x]=0$，$f(x)=[2x]+[3x]=0$；}
\step{当 $\dfrac13\leqslant x<\dfrac12$ 时，$[2x]=0$，$[3x]=1$，$f(x)=[2x]+[3x]=1$；}
\step{当 $\dfrac12\leqslant x<\dfrac23$ 时，$[2x]=1$，$[3x]=1$，$f(x)=[2x]+[3x]=2$；}
\step{当 $\dfrac23\leqslant x<1$ 时，$[2x]=1$，$[3x]=2$，$f(x)=[2x]+[3x]=3$；}
\step{当 $1\leqslant x<\dfrac43$ 时，$[2x]=2$，$[3x]=3$，$f(x)=[2x]+[3x]=5$；}
\step{当 $\dfrac43\leqslant x<\dfrac32$ 时，$[2x]=2$，$[3x]=4$，$f(x)=[2x]+[3x]=6$；}
\step{当 $\dfrac32\leqslant x<\dfrac53$ 时，$[2x]=3$，$[3x]=4$，$f(x)=[2x]+[3x]=7$；}
\step{当 $\dfrac53\leqslant x<2$ 时，$[2x]=3$，$[3x]=5$，$f(x)=[2x]+[3x]=8$；}
\step{当 $x=2$ 时，$[2x]=4$，$[3x]=6$，$f(x)=[2x]+[3x]=10$}
\step{故 $A$ 中所有元素的和为 $-5-4-3-2+0+1+2+3+5+6+7+8+10=28$.}
% 原卷“法二：”标号排红色，同上
\stepm{法二：}{设 $x=m+n$，其中 $m\in\Z$，$n\in[0,1)$，}
\step{若 $0\leqslant n<\dfrac13$，则 $0\leqslant2n<\dfrac23$，$0\leqslant3n<1$，}
\step{得 $[2x]=[2m+2n]=2m$，$[3x]=[3m+3n]=3m$，所以 $y=[2x]+[3x]=5m$；}
\step{若 $\dfrac13\leqslant n<\dfrac12$，则 $\dfrac23\leqslant2n<1$，$1\leqslant3n<\dfrac32$，}
\step{得 $[2x]=[2m+2n]=2m$，$[3x]=[3m+3n]=3m+1$，}
\step{所以 $y=[2x]+[3x]=5m+1$；}
\step{若 $\dfrac12\leqslant n<\dfrac23$，则 $1\leqslant2n<\dfrac43$，$\dfrac32\leqslant3n<2$，}
\step{得 $[2x]=[2m+2n]=2m+1$，$[3x]=[3m+3n]=3m+1$，}
\step{所以 $y=[2x]+[3x]=5m+2$；}
\step{若 $\dfrac23\leqslant n<1$，则 $\dfrac43\leqslant2n<2$，$2\leqslant3n<3$，}
\step{得 $[2x]=[2m+2n]=2m+1$，$[3x]=[3m+3n]=3m+2$，}
\step{所以 $y=[2x]+[3x]=5m+3$；}
\step{综上所述，当 $n\in[0,1)$ 时，所有函数值的和}
\step{为 $5m+5m+1+5m+2+5m+3=20m+6$，}
\step{又因为 $x\in[-1,2]$，此时 $m=-1,0,1$，且当 $x=2$ 时，$y=[2x]+[3x]=10$，}
\step{所有函数值的和为 $(-14)+6+26+10=28$，即 $A$ 中所有元素之和为 $28$.}
\solend
\fi

\end{enumerate}

\bigsec{二、选择题}

\begin{enumerate}
\setcounter{enumi}{12}

\item 下列写法中，正确的是\mbox{（\quad）}

\keepwithprev
A. $0\in\{\varnothing\}$\qquad B. $\varnothing\in\{\varnothing\}$

C. $\varnothing=\{\varnothing\}$\qquad D. $0\in\varnothing$

\ifanswer
\solbegin
\step{$0\notin\{\varnothing\}$，$\varnothing\in\{\varnothing\}$，故 $A$ 错误，$B$ 正确；}
\step{$\varnothing$ 是不含任何元素的集合，所以 $\varnothing\neq\{\varnothing\}$，$0\notin\varnothing$，故 $CD$ 错误.}
\step{故选 $B$.}
\solend
\fi

% 第 14 题答案“D”直接印在题干括号内“(  D  )”，且该题无【答案】/【解析】段；
% 按 高一37 第 13 题的成例括号留空、答案另提为 \ans{}（见文件头）
\item 若 $a<b<0$，则下列不等式中不能成立的是\mbox{（\quad）}

\keepwithprev
A. $\dfrac1a>\dfrac1b$\qquad B. $a-b>a$

C. $|a|>|b|$\qquad D. $a^3>b^3$

\ifanswer
\ans{$D$}
\fi

\item 已知 $x$、$y$、$z$ 为非零实数，代数式 $\dfrac x{|x|}+\dfrac y{|y|}+\dfrac z{|z|}+\dfrac{xyz}{|xyz|}$
的值所组成的集合是 $M$，则下列判断正确的是\mbox{（\quad）}

\keepwithprev
A. $-4\notin M$\qquad B. $2\in M$

C. $4\notin M$\qquad D. $0\in M$

\ifanswer
\solbegin
\step{当 $x$、$y$、$z$ 均为正数时，$xyz>0$，$\dfrac x{|x|}+\dfrac y{|y|}+\dfrac z{|z|}+\dfrac{xyz}{|xyz|}=1+1+1+1=4$，}
\step{当 $x$、$y$、$z$ 有 2 个正数，1 个负数时，不妨设 $x>0$，$y>0$，$z<0$，}
\step{则 $xyz<0$，$\dfrac x{|x|}+\dfrac y{|y|}+\dfrac z{|z|}+\dfrac{xyz}{|xyz|}=1+1-1-1=0$，}
\step{若 $x$、$y$、$z$ 有 1 个正数，2 个负数时，不妨设 $x>0$，$y<0$，$z<0$，}
\step{则 $xyz>0$，$\dfrac x{|x|}+\dfrac y{|y|}+\dfrac z{|z|}+\dfrac{xyz}{|xyz|}=1-1-1+1=0$，}
\step{当 $x$、$y$、$z$ 均为负数时，$xyz<0$，}
\step{$\dfrac x{|x|}+\dfrac y{|y|}+\dfrac z{|z|}+\dfrac{xyz}{|xyz|}=-1-1-1-1=-4$，}
\step{综上，$M=\{-4,0,4\}$，}
\step{所以 $-4\in M$，$A$ 错误；$2\notin M$，$B$ 错误；$4\in M$，$C$ 错误；}
\step{$0\in M$，$D$ 正确. 故选 $D$.}
\solend
\fi

\item 若 $xy>0$，且 $x+y=9$，则 $\dfrac1x+\dfrac4y$ 的最小值是\mbox{（\quad）}

\keepwithprev
A. $9$\qquad B. $\dfrac59$

C. $\dfrac49$\qquad D. $1$

\ifanswer
\solbegin
\step{$xy>0$，且 $x+y=9$，所以 $x>0$，$y>0$，}
\step{所以 $\dfrac1x+\dfrac4y=\dfrac19\left(\dfrac1x+\dfrac4y\right)(x+y)=\dfrac59+\dfrac19\left(\dfrac yx+\dfrac{4x}y\right)\geqslant\dfrac59+\dfrac29\cdot\sqrt{\dfrac yx\cdot\dfrac{4x}y}=\dfrac59+\dfrac49=1$，}
\step{当且仅当 $\dfrac yx=\dfrac{4x}y$，即 $x=3$，$y=6$ 时取等号，所以 $\dfrac1x+\dfrac4y$ 的最小值为 $1$.}
\step{故选 $D$.}
\solend
\fi

\end{enumerate}

\bigsec{三、解答题}

\begin{enumerate}
\setcounter{enumi}{16}

\item 求解下列不等式：

（1）$\dfrac{x-3}{x+1}\leqslant0$；

（2）$|2x-1|<3$.

\ifanswer
\solbegin
\stepm{(1)}{$\dfrac{x-3}{x+1}\leqslant0$，变形得 $(x-3)\cdot(x+1)\leqslant0$ 且 $x\neq-1$，}
\step{解得 $-1<x\leqslant3$，即不等式的解集为 $(-1,3]$.}
\stepm{(2)}{若 $|2x-1|<3$，则 $-3<2x-1<3$，解得 $-1<x<2$，即 $(-1,2)$.}
\solend
\fi

\ansspace[6]

\item 已知集合 $A=\{x\mid-2<x<15\}$,$B=\{x\mid m-6<x<2m-1,m\in\R\}$.

（1）当 $m=2$ 时，求 $A\cap B$；

（2）若 $A\cup B=A$，求实数 $m$ 的取值范围.

\ifanswer
\solbegin
\stepm{(1)}{当 $m=2$ 时，$A=\{x\mid-2<x<15\}$,$B=\{x\mid-4<x<3\}$，}
\step{此时 $A\cap B=\{x\mid-2<x<3\}$.}
\stepm{(2)}{由 $A\cup B=A$ 得 $B\subseteq A$.}
\step{当 $B=\varnothing$,$m-6\geqslant2m-1$，解得 $m\leqslant-5$；}
\step{当 $B\neq\varnothing$，结合 $B\subseteq A$ 得 $\begin{cases}m-6<2m-1\\ 2m-1\leqslant15\\ m-6\geqslant-2\end{cases}$，解得 $4\leqslant m\leqslant8$.}
\step{综上，$m\leqslant-5$ 或 $4\leqslant m\leqslant8$.}
\solend
\fi

\ansspace[6]

\item 已知 $a\in\R$，命题 $p$：关于 $x$ 的方程 $x^2-2x+a=0$ 无实根；命题 $q$：对任意 $x\in\R$，
不等式 $ax^2-ax+1>0$ 恒成立.

（1）若 $q$ 为真命题，求 $a$ 的取值范围；

（2）若 $p$、$q$ 中有且只有一个为真命题，求 $a$ 的取值范围.

\ifanswer
\solbegin
\stepm{(1)}{当 $q$ 为真命题时，}
\stepm{①}{当 $a=0$ 时，不等式化为 $1>0$，对任意 $x\in\R$ 恒成立，满足条件；}
\stepm{②}{当 $a\neq0$ 时，$ax^2-ax+1>0$ 为一元二次不等式，}
\step{若对任意 $x\in\R$ 恒成立，需满足 $\begin{cases}a>0\\ \Delta=(-a)^2-4a<0\end{cases}$，解得 $0<a<4$；}
\step{综上，$a$ 的取值范围为 $[0,4)$.}
\stepm{(2)}{$p$ 为真时，方程 $x^2-2x+a=0$ 无实根，}
\step{故判别式 $\Delta=4-4a<0$，解得 $a>1$，因此 $p$ 为假时，$a\leqslant1$；}
\step{由（1）得 $q$ 为假时 $a\geqslant4$ 或 $a<0$，}
\step{$p$、$q$ 有且只有一个真命题，分两类：}
\step{当 $p$ 假 $q$ 真时，满足 $\begin{cases}a\leqslant1\\ 0\leqslant a<4\end{cases}$，解得 $0\leqslant a\leqslant1$；}
\step{当 $p$ 真 $q$ 假时，满足 $\begin{cases}a>1\\ a\geqslant4\text{或}a<0\end{cases}$，解得 $a\geqslant4$.}
\step{综上，$a$ 的取值范围是 $[0,1]\cup[4,+\infty)$.}
\solend
\fi

\ansspace[8]

\item 某工厂生产 $x$($x>0$) 件产品的总成本 $C$（单位：元）与生产件数 $x$ 之间满足二次函数
$C(x)=\dfrac14x^2+50x+10000$.

（1）若要求生产总成本不超过 30000 元，求满足生产总成本要求的产品件数的取值范围；

（2）试求该工厂生产多少件产品时，可使每件产品的平均成本最小.

\ifanswer
\solbegin
\stepm{(1)}{若要求生产总成本不超过 30000 元，}
\step{则 $C(x)=\dfrac14x^2+50x+10000\leqslant30000$，}
\step{即 $x^2+200x-80000\leqslant0$，解得 $-400\leqslant x\leqslant200$，又 $x>0$，}
\step{所以满足生产总成本要求的产品件数的取值范围为 $(0,200]$，$x\in\N$；}
\stepm{(2)}{生产件数 $x$ 时每件的平均成本为 $\dfrac{\frac14x^2+50x+10000}x=\dfrac14x+50+\dfrac{10000}x$，}
\step{由基本不等式得 $\dfrac14x+50+\dfrac{10000}x\geqslant2\sqrt{\dfrac14x\cdot\dfrac{10000}x}+50=150$，}
\step{当且仅当 $\dfrac14x=\dfrac{10000}x$，即 $x=200$ 时取等号，}
\step{即该工厂生产 200 件产品时，每件产品的平均成本最小.}
\solend
\fi

\ansspace[8]

% 原卷题干“$(n\geqslant2,n\in N)$”后连印两个全角逗号“，，”，照录未改（见文件头“原卷笔误”①）
\item 已知有限集 $A=\{a_1,a_2,\cdots,a_n\}$($n\geqslant2$,$n\in\N$)，，如果 $A$ 中元素
$a_i$($i=1,2,3,\cdots,n$) 满足 $a_1+a_2+a_3+\cdots+a_n=a_1\times a_2\times\cdots\times a_n$，
就称 $A$ 为“和积平衡集”.

% 原卷 (2) 将“写方程”与“反证法证明”两项任务并作一句（“……表示），可尝试用反证法证明：……”），
% 照录未改（见文件头“原卷笔误”②）
（1）判断集合 $\{-1+\sqrt3,-1-\sqrt3\}$ 是否为“和积平衡集”；

（2）若 $a_1$、$a_2$ 是两个不同的正数，且 $\{a_1,a_2\}$ 是“和积平衡集”，写出以 $a_1$、$a_2$ 为根的
一个一元二次方程（系数可用 $a_1$、$a_2$ 表示），可尝试用反证法证明：$a_1$、$a_2$ 至少有一个大于 2；

（3）若 $a_i\geqslant1$，且 $a_i\in\N$，求所有符合条件的“和积平衡集”.

\ifanswer
\solbegin
\stepm{(1)}{对于集合 $\{-1+\sqrt3,-1-\sqrt3\}$，}
\step{因为 $(-1+\sqrt3)+(-1-\sqrt3)=-2$，$(-1+\sqrt3)\times(-1-\sqrt3)=-2$，}
\step{所以满足 $a_1+a_2=a_1\times a_2$，满足题意，}
\step{所以集合 $\{-1+\sqrt3,-1-\sqrt3\}$ 是和积平衡集.}
\stepm{(2)}{以 $a_1$、$a_2$ 为根的一个一元二次方程可为 $x^2-(a_1+a_2)x+a_1\cdot a_2=0$，}
\step{因为 $a_1$、$a_2$ 是两个不同的正数，且 $\{a_1,a_2\}$ 是和积平衡集，}
\step{所以 $a_1+a_2=a_1\cdot a_2$，}
\step{又 $a_1$ 和 $a_2$ 为 $x^2-(a_1+a_2)x+a_1\cdot a_2=0$ 的两个不同的正根，}
\step{所以 $\Delta=(a_1+a_2)^2-4a_1a_2=(a_1+a_2)^2-4(a_1+a_2)>0$，}
\step{又 $a_1+a_2>0$，则 $a_1+a_2>4$，}
\step{若 $a_1$、$a_2$ 都小于等于 2，则 $a_1+a_2<4$，矛盾，}
\step{所以 $a_1$、$a_2$ 至少有一个大于 2；}
\stepm{(3)}{设 $A$ 中的 $a_1<a_2<a_3<\cdots<a_n$，且 $a_i\in\N^*$，}
\step{由 $a_1+a_2+a_3+\cdots+a_n=a_1\times a_2\times\cdots\times a_n$，得 $n=2$ 时明显不满足，}
\step{所以 $n\geqslant3$，故 $a_1\times a_2\times\cdots\times a_n=a_1+a_2+a_3+\cdots+a_n\leqslant na_n$，}
\step{所以 $a_1\times a_2\times\cdots\times a_{n-1}\leqslant n$，}
\step{当 $n=3$ 时，$a_1\times a_2\leqslant3$，$a_i\in\N^*$，只能有 $a_1=1$，$a_2=2$，}
\step{由 $a_1+a_2+a_3=a_1\times a_2\times a_3$，得 $a_3=3$，此时的和积平衡集为 $\{1,2,3\}$，}
\step{当 $n\geqslant4$ 时，$a_1\times a_2\times\cdots\times a_{n-1}\geqslant(n-1)(n-2)\times\cdots\times2\times1\geqslant(n-1)(n-2)$，}
\step{又 $(n-1)(n-2)-n=n^2-4n+2=n(n-4)+2>0$，即 $(n-1)(n-2)>n$，}
\step{即 $a_1\times a_2\times\cdots\times a_{n-1}>n$，与 $a_1\times a_2\times\cdots\times a_{n-1}\leqslant n$ 矛盾，}
\step{所以 $n\geqslant4$ 不满足条件，}
\step{所以符合条件的和积平衡集有且只有 $n=3$，此时的和积平衡集为 $\{1,2,3\}$.}
\solend
\fi

\ansspace*[10]

\end{enumerate}

\end{document}
"
%
% 原始材料：微信公众号“魔都数学加油站”文章，作者破晓，连载“27学年高一63”；
% 原文标题“27学年高一63——2026-2027年上海市吴淞中学高一上10月月考”；
% 原文链接：https://mp.weixin.qq.com/s/TFQ4TLnt_l66t7t51es4OQ。
% 页面标注发布时间 2026-10-16 21:44（上海）。
% 正文为 10 张 PNG 页（`原图/高一63_吴淞中学_10月月考/01..10.png`）：
% 第 06、10 页 2560×3620，第 01、02、03、04、05、07、08、09 页 4762×6735
% （**同一篇各页分辨率不同**），均按页序读取。10 页全为试卷正文页（卷首标题在第 1 页页顶，
% 卷末解答题第 21 题解析止于第 10 页），无封面页、目录页、二维码或宣传页。
% 粉色斜排“魔都数学加油站”为卷面水印，与公众号页脚一并未录入。
%
% 卷首标题：照卷面印作“2026--2027 年吴淞中学高一上 10 月月考”（公众号标题中的“上海市”
% 卷面标题行没有，本稿照卷面），年份区间按全稿惯例排 en dash。原卷只有一行标题、无副标题。
%
% 篇幅与题号：原卷自第 3 题起编号（第 1、2 题不在本卷内）。填空题 3—12（10 题）、
% 选择题 13—16（4 题）、解答题 17—21（5 题），共 19 题。本稿沿用原节与原题号，
% 只改排为全稿统一的大节样式（原卷小节标题为“一、填空题”“二、选择题”“三、解答题”）。
%
% 讲评版：原卷每题下直接印【答案】或【解析】。**只印【答案】的 2 题**（第 4、7 题），
% 本稿只给 \ans{}；其余 17 题里 16 题都印【解析】（按全稿口径这类题不另提 \ans{}——第 13、15、16 题
% 的结论写在解析末句（“故选 B”等），亦照录在解析内；第 14 题则既无【答案】也无【解析】段。**另有 1 题答案印在题干括号内**——
% 第 14 题（“(  D  )”，且该题无【答案】/【解析】段），本稿按 高一37 第 13 题成例把括号留空、
% 答案另提为 \ans{$D$}。第 16 题题干括号为空但印有【解析】，不算括号内答案。
%
% 副标题：原卷未印副标题。本稿按“知识模块”口径标作“集合与常用逻辑用语、不等式”：
% 19 题中集合与常用逻辑用语 12 题（第 3、4、6、7、9、10、11、12、13、15、18、21 题），
% 不等式 7 题（含等式（韦达定理、恒等式）与不等式性质、基本不等式、二次不等式、恒成立：
% 第 5、8、14、16、17、19、20 题），两类均超过三分之一，故并标。口径同 高一59。
%
% 与原卷的排版差异：
%   ① 原卷每题下直接印【答案】/【解析】，本稿按全稿惯例将答案与解析置于答案版，试卷版保持干净；
%   ② 原卷未印副标题，本稿按“知识模块”口径补出；
%   ③ 原卷填空横线约四个字宽，本稿按全稿惯例统一排 \blankline[4.4em]；
%   ④ 原卷实数集排斜体 $R$（第 9、10、19 题题干与解析）或粗体 $\mathbf{R}$（第 18 题题干），
%      整数集排斜体 $Z$（第 12 题解析），自然数集排斜体 $N$（第 20 题解析、第 21 题题干
%      “$(n\geqslant2,n\in N)$”）或粗体 $\mathbf{N}$（第 21 题 (3) 题干）、斜体 $N^*$
%      （第 21 题 (3) 解析），本稿按全稿惯例统一改用 $\R$、$\Z$、$\N$、$\N^*$
%      （同 高一14 排版差异④：记号归一，粗体／斜体不作区分，$N$ 与 $N^*$ 的区别照原卷保留）；
%      空集记号原卷作 $\varnothing$，照录；
%   ⑤ 原卷在数学式内多处用半角逗号（如 $\{y\mid y=[2x]+[3x],-1\leqslant x\leqslant2\}$、
%      $m=-1,0,1$、$A=\{x\mid-2<x<15\}$,$B=\cdots$），均照录未强改；
%      句末句点按全稿惯例排西文句点；“和积平衡集”自定义名词照录；
%   ⑥ 第 12 题解析里的“法一：”“法二：”标号原卷排**红色**（全卷共 2 处），
%      同 高一53“情况1：”成例，本稿排 \stepm{法一：}{…}、\stepm{法二：}{…}
%      （加粗着色、不改黑色基调），冒号照原卷排全角；
%   ⑦ 第 17—21 题的子问原卷各自另起一行，本稿按全稿惯例用 enumerate 的 (1)(2)(3)；
%      解析中的子问标号一律用 \stepm{(1)}{…}；第 19 题解析中的①②分情形，按全稿惯例作
%      \stepm{①}{…}、\stepm{②}{…}；
%   ⑧ 第 14 题印在题干括号内的答案移入 \ans{}，见上；
%   ⑨ 每道解答题原卷未留作答空白，本稿按全稿惯例加 \ansspace；
%   ⑩ 选择题的作答括号原卷排半角“(　　)”（第 13、15、16 题括号内空白、第 14 题括号内印 D），
%      本稿按全稿惯例排 \mbox{（\quad）}。
%
% 原卷笔误/存疑（均照录未改，正文对应处另有 % 注）：
%   ① 第 21 题题干“$(n\geqslant2$,$n\in N)$”后连印两个全角逗号“，，”（衍字，
%      8 倍放大确认，非水印干扰），本稿照录；
%   ② 第 21 题 (2) 题干将“写出以 $a_1$、$a_2$ 为根的一个一元二次方程”与“用反证法证明：
%      $a_1$、$a_2$ 至少有一个大于 2”两项任务并作一句、之间仅以逗号相接（“……表示），
%      可尝试用反证法证明：……”），读来像抄并两小问，原卷如此，照录未改。
%
% 插图：无（全卷 10 页无任何图形、无“如图”字样）。
%
\documentclass[a4paper,11pt]{article}
\usepackage{common}

\begin{document}

\examtitlest[3]{2026--2027 年吴淞中学高一上 10 月月考}{集合与常用逻辑用语、不等式}

\bigsec{一、填空题}

\begin{enumerate}
\setcounter{enumi}{2}

\item 满足 $A\subseteq\{1,2,3,4\}$ 的集合 $A$ 共有\blankline[4.4em]个.

\ifanswer
\solbegin
\step{由题意，集合 $A$ 是集合 $\{1,2,3,4\}$ 的子集，共有 $2^4=16$ 个.}
\solend
\fi

\item 陈述句“$a>1$ 或 $b>3$”的否定形式为\blankline[4.4em].

\ifanswer
\ans{$a\leqslant1$ 且 $b\leqslant3$}
\fi

\item 已知方程 $3x^2+4x-4=0$ 的两个根为 $x_1$、$x_2$，则 $x_1^2+x_2^2=$\blankline[4.4em].

\ifanswer
\solbegin
\step{对于方程 $3x^2+4x-4=0$，因为 $\Delta=16-4\times3\times(-4)=64>0$，}
\step{所以方程 $3x^2+4x-4=0$ 的两个不相等的实根为 $x_1$、$x_2$，}
\step{由韦达定理可得 $x_1+x_2=-\dfrac43$，$x_1x_2=-\dfrac43$，}
\step{所以 $x_1^2+x_2^2=(x_1+x_2)^2-2x_1x_2=\dfrac{16}9+\dfrac83=\dfrac{40}9$.}
\solend
\fi

\item 设集合 $A=\{y\mid y=x^2+1\}$,$B=\{y\mid y=-x^2+2x+3\}$，则 $A\cap B=$\blankline[4.4em].

\ifanswer
\solbegin
\step{因为 $y=x^2+1\geqslant1$,$y=-x^2+2x+3=-(x-1)^2+4\leqslant4$，}
\step{所以 $A=\{y\mid y\geqslant1\}$,$B=\{y\mid y\leqslant4\}$，得 $A\cap B=\{y\mid1\leqslant y\leqslant4\}$.}
\solend
\fi

\item 若 $x\in\{1,x^2\}$，则 $x=$\blankline[4.4em].

\ifanswer
\ans{$0$}
\fi

\item 已知等式 $2x^2+3x+5=a(2x+1)(x+1)+c$，则 $a+c=$\blankline[4.4em].

\ifanswer
\solbegin
\step{因为 $2x^2+3x+5=a(2x+1)(x+1)+c$ 恒成立，}
\step{所以 $2x^2+3x+5=2ax^2+3ax+a+c$ 恒成立，则 $\begin{cases}2=2a\\ 3=3a\\ 5=a+c\end{cases}$，解得 $a+c=5$.}
\solend
\fi

\item 已知 $a\in\R$，设 $p:x<2$，$q:x<5a-3$，且 $p$ 是 $q$ 的必要非充分条件，则 $a$ 的取值范围是\blankline[4.4em].

\ifanswer
\solbegin
\step{设 $p:x<2$，$q:x<5a-3$，且 $p$ 是 $q$ 的必要非充分条件，}
\step{则 $\{x\mid x<5a-3\}\subset\{x\mid x<2\}$，所以 $5a-3<2$，即 $a<1$.}
\solend
\fi

\item 已知 $a\in\R$，集合 $A=\{x\mid ax^2+3x-1=0\}$，若集合 $A$ 中只有一个元素，则 $a$ 的取值范围是\blankline[4.4em].

\ifanswer
\solbegin
\step{当 $a=0$ 时，则有 $\{x\mid ax^2+3x-1=0\}=\{x\mid3x-1=0\}=\left\{\dfrac13\right\}$，符合题意；}
\step{当 $a\neq0$ 时，由题意得 $\Delta=9+4a=0$，解得 $a=-\dfrac94$. 符合题意；}
\step{综上所述，实数 $a$ 的取值集合为 $\left\{0,-\dfrac94\right\}$.}
\solend
\fi

\item 设集合 $A=\{x\mid x^2+2(m+1)x+m^2-1=0,m\in\R\}$，且 $A\subseteq\{0,-4\}$，则实数 $m$ 的取值范围是\blankline[4.4em].

\ifanswer
\solbegin
\step{由题意得 $A=\varnothing$ 或 $A=\{0\}$ 或 $A=\{-4\}$ 或 $A=\{0,-4\}$，}
\step{当 $A=\varnothing$ 时，方程 $x^2+2(m+1)x+m^2-1=0$ 无实根，}
\step{则 $\Delta=[2(m+1)]^2-4\times1\times(m^2-1)=8m+8<0$，解得 $m<-1$，}
\step{当 $A=\{0\}$ 时，方程 $x^2+2(m+1)x+m^2-1=0$ 有一个实根 $0$，}
\step{则 $\begin{cases}\Delta=8m+8=0\\ m^2-1=0\end{cases}$，解得 $m=-1$，}
\step{当 $A=\{-4\}$ 时，方程 $x^2+2(m+1)x+m^2-1=0$ 有一个实根 $-4$，}
\step{则 $\begin{cases}\Delta=8m+8=0\\ 16-8(m+1)+m^2-1=0\end{cases}$，无解，}
\step{当 $A=\{0,-4\}$ 时，方程 $x^2+2(m+1)x+m^2-1=0$ 有两个实根 $-4$ 和 $0$，}
\step{则 $\begin{cases}\Delta=8m+8>0\\ 16-8(m+1)+m^2-1=0\\ m^2-1=0\end{cases}$，解得 $m=1$，}
\step{综上所述，实数 $m$ 的取值范围是 $\{m\mid m\leqslant-1\text{ 或 }m=1\}$.}
\solend
\fi

\item 对于任意实数 $x$，$[x]$ 表示不超过 $x$ 的最大整数，如 $[1.1]=1$，$[-1.2]=-2$，若
$A=\{y\mid y=[2x]+[3x],-1\leqslant x\leqslant2\}$，则 $A$ 中所有元素的和为\blankline[4.4em].

\ifanswer
\solbegin
% 原卷“法一：”“法二：”标号排红色，本稿按 高一53“情况1：”成例排 \stepm（加粗着色），见文件头
\stepm{法一：}{当 $-1\leqslant x<-\dfrac23$ 时，$[2x]=-2$，$[3x]=-3$，$f(x)=[2x]+[3x]=-5$，}
\step{当 $-\dfrac23\leqslant x<-\dfrac12$ 时，$[2x]=-2$，$[3x]=-2$，$f(x)=[2x]+[3x]=-4$；}
\step{当 $-\dfrac12\leqslant x<-\dfrac13$ 时，$[2x]=-1$，$[3x]=-2$，$f(x)=[2x]+[3x]=-3$；}
\step{当 $-\dfrac13\leqslant x<0$ 时，$[2x]=-1$，$[3x]=-1$，$f(x)=[2x]+[3x]=-2$；}
\step{当 $0\leqslant x<\dfrac13$ 时，$[2x]=0$，$[3x]=0$，$f(x)=[2x]+[3x]=0$；}
\step{当 $\dfrac13\leqslant x<\dfrac12$ 时，$[2x]=0$，$[3x]=1$，$f(x)=[2x]+[3x]=1$；}
\step{当 $\dfrac12\leqslant x<\dfrac23$ 时，$[2x]=1$，$[3x]=1$，$f(x)=[2x]+[3x]=2$；}
\step{当 $\dfrac23\leqslant x<1$ 时，$[2x]=1$，$[3x]=2$，$f(x)=[2x]+[3x]=3$；}
\step{当 $1\leqslant x<\dfrac43$ 时，$[2x]=2$，$[3x]=3$，$f(x)=[2x]+[3x]=5$；}
\step{当 $\dfrac43\leqslant x<\dfrac32$ 时，$[2x]=2$，$[3x]=4$，$f(x)=[2x]+[3x]=6$；}
\step{当 $\dfrac32\leqslant x<\dfrac53$ 时，$[2x]=3$，$[3x]=4$，$f(x)=[2x]+[3x]=7$；}
\step{当 $\dfrac53\leqslant x<2$ 时，$[2x]=3$，$[3x]=5$，$f(x)=[2x]+[3x]=8$；}
\step{当 $x=2$ 时，$[2x]=4$，$[3x]=6$，$f(x)=[2x]+[3x]=10$}
\step{故 $A$ 中所有元素的和为 $-5-4-3-2+0+1+2+3+5+6+7+8+10=28$.}
% 原卷“法二：”标号排红色，同上
\stepm{法二：}{设 $x=m+n$，其中 $m\in\Z$，$n\in[0,1)$，}
\step{若 $0\leqslant n<\dfrac13$，则 $0\leqslant2n<\dfrac23$，$0\leqslant3n<1$，}
\step{得 $[2x]=[2m+2n]=2m$，$[3x]=[3m+3n]=3m$，所以 $y=[2x]+[3x]=5m$；}
\step{若 $\dfrac13\leqslant n<\dfrac12$，则 $\dfrac23\leqslant2n<1$，$1\leqslant3n<\dfrac32$，}
\step{得 $[2x]=[2m+2n]=2m$，$[3x]=[3m+3n]=3m+1$，}
\step{所以 $y=[2x]+[3x]=5m+1$；}
\step{若 $\dfrac12\leqslant n<\dfrac23$，则 $1\leqslant2n<\dfrac43$，$\dfrac32\leqslant3n<2$，}
\step{得 $[2x]=[2m+2n]=2m+1$，$[3x]=[3m+3n]=3m+1$，}
\step{所以 $y=[2x]+[3x]=5m+2$；}
\step{若 $\dfrac23\leqslant n<1$，则 $\dfrac43\leqslant2n<2$，$2\leqslant3n<3$，}
\step{得 $[2x]=[2m+2n]=2m+1$，$[3x]=[3m+3n]=3m+2$，}
\step{所以 $y=[2x]+[3x]=5m+3$；}
\step{综上所述，当 $n\in[0,1)$ 时，所有函数值的和}
\step{为 $5m+5m+1+5m+2+5m+3=20m+6$，}
\step{又因为 $x\in[-1,2]$，此时 $m=-1,0,1$，且当 $x=2$ 时，$y=[2x]+[3x]=10$，}
\step{所有函数值的和为 $(-14)+6+26+10=28$，即 $A$ 中所有元素之和为 $28$.}
\solend
\fi

\end{enumerate}

\bigsec{二、选择题}

\begin{enumerate}
\setcounter{enumi}{12}

\item 下列写法中，正确的是\mbox{（\quad）}

\keepwithprev
A. $0\in\{\varnothing\}$\qquad B. $\varnothing\in\{\varnothing\}$

C. $\varnothing=\{\varnothing\}$\qquad D. $0\in\varnothing$

\ifanswer
\solbegin
\step{$0\notin\{\varnothing\}$，$\varnothing\in\{\varnothing\}$，故 $A$ 错误，$B$ 正确；}
\step{$\varnothing$ 是不含任何元素的集合，所以 $\varnothing\neq\{\varnothing\}$，$0\notin\varnothing$，故 $CD$ 错误.}
\step{故选 $B$.}
\solend
\fi

% 第 14 题答案“D”直接印在题干括号内“(  D  )”，且该题无【答案】/【解析】段；
% 按 高一37 第 13 题的成例括号留空、答案另提为 \ans{}（见文件头）
\item 若 $a<b<0$，则下列不等式中不能成立的是\mbox{（\quad）}

\keepwithprev
A. $\dfrac1a>\dfrac1b$\qquad B. $a-b>a$

C. $|a|>|b|$\qquad D. $a^3>b^3$

\ifanswer
\ans{$D$}
\fi

\item 已知 $x$、$y$、$z$ 为非零实数，代数式 $\dfrac x{|x|}+\dfrac y{|y|}+\dfrac z{|z|}+\dfrac{xyz}{|xyz|}$
的值所组成的集合是 $M$，则下列判断正确的是\mbox{（\quad）}

\keepwithprev
A. $-4\notin M$\qquad B. $2\in M$

C. $4\notin M$\qquad D. $0\in M$

\ifanswer
\solbegin
\step{当 $x$、$y$、$z$ 均为正数时，$xyz>0$，$\dfrac x{|x|}+\dfrac y{|y|}+\dfrac z{|z|}+\dfrac{xyz}{|xyz|}=1+1+1+1=4$，}
\step{当 $x$、$y$、$z$ 有 2 个正数，1 个负数时，不妨设 $x>0$，$y>0$，$z<0$，}
\step{则 $xyz<0$，$\dfrac x{|x|}+\dfrac y{|y|}+\dfrac z{|z|}+\dfrac{xyz}{|xyz|}=1+1-1-1=0$，}
\step{若 $x$、$y$、$z$ 有 1 个正数，2 个负数时，不妨设 $x>0$，$y<0$，$z<0$，}
\step{则 $xyz>0$，$\dfrac x{|x|}+\dfrac y{|y|}+\dfrac z{|z|}+\dfrac{xyz}{|xyz|}=1-1-1+1=0$，}
\step{当 $x$、$y$、$z$ 均为负数时，$xyz<0$，}
\step{$\dfrac x{|x|}+\dfrac y{|y|}+\dfrac z{|z|}+\dfrac{xyz}{|xyz|}=-1-1-1-1=-4$，}
\step{综上，$M=\{-4,0,4\}$，}
\step{所以 $-4\in M$，$A$ 错误；$2\notin M$，$B$ 错误；$4\in M$，$C$ 错误；}
\step{$0\in M$，$D$ 正确. 故选 $D$.}
\solend
\fi

\item 若 $xy>0$，且 $x+y=9$，则 $\dfrac1x+\dfrac4y$ 的最小值是\mbox{（\quad）}

\keepwithprev
A. $9$\qquad B. $\dfrac59$

C. $\dfrac49$\qquad D. $1$

\ifanswer
\solbegin
\step{$xy>0$，且 $x+y=9$，所以 $x>0$，$y>0$，}
\step{所以 $\dfrac1x+\dfrac4y=\dfrac19\left(\dfrac1x+\dfrac4y\right)(x+y)=\dfrac59+\dfrac19\left(\dfrac yx+\dfrac{4x}y\right)\geqslant\dfrac59+\dfrac29\cdot\sqrt{\dfrac yx\cdot\dfrac{4x}y}=\dfrac59+\dfrac49=1$，}
\step{当且仅当 $\dfrac yx=\dfrac{4x}y$，即 $x=3$，$y=6$ 时取等号，所以 $\dfrac1x+\dfrac4y$ 的最小值为 $1$.}
\step{故选 $D$.}
\solend
\fi

\end{enumerate}

\bigsec{三、解答题}

\begin{enumerate}
\setcounter{enumi}{16}

\item 求解下列不等式：

（1）$\dfrac{x-3}{x+1}\leqslant0$；

（2）$|2x-1|<3$.

\ifanswer
\solbegin
\stepm{(1)}{$\dfrac{x-3}{x+1}\leqslant0$，变形得 $(x-3)\cdot(x+1)\leqslant0$ 且 $x\neq-1$，}
\step{解得 $-1<x\leqslant3$，即不等式的解集为 $(-1,3]$.}
\stepm{(2)}{若 $|2x-1|<3$，则 $-3<2x-1<3$，解得 $-1<x<2$，即 $(-1,2)$.}
\solend
\fi

\ansspace[6]

\item 已知集合 $A=\{x\mid-2<x<15\}$,$B=\{x\mid m-6<x<2m-1,m\in\R\}$.

（1）当 $m=2$ 时，求 $A\cap B$；

（2）若 $A\cup B=A$，求实数 $m$ 的取值范围.

\ifanswer
\solbegin
\stepm{(1)}{当 $m=2$ 时，$A=\{x\mid-2<x<15\}$,$B=\{x\mid-4<x<3\}$，}
\step{此时 $A\cap B=\{x\mid-2<x<3\}$.}
\stepm{(2)}{由 $A\cup B=A$ 得 $B\subseteq A$.}
\step{当 $B=\varnothing$,$m-6\geqslant2m-1$，解得 $m\leqslant-5$；}
\step{当 $B\neq\varnothing$，结合 $B\subseteq A$ 得 $\begin{cases}m-6<2m-1\\ 2m-1\leqslant15\\ m-6\geqslant-2\end{cases}$，解得 $4\leqslant m\leqslant8$.}
\step{综上，$m\leqslant-5$ 或 $4\leqslant m\leqslant8$.}
\solend
\fi

\ansspace[6]

\item 已知 $a\in\R$，命题 $p$：关于 $x$ 的方程 $x^2-2x+a=0$ 无实根；命题 $q$：对任意 $x\in\R$，
不等式 $ax^2-ax+1>0$ 恒成立.

（1）若 $q$ 为真命题，求 $a$ 的取值范围；

（2）若 $p$、$q$ 中有且只有一个为真命题，求 $a$ 的取值范围.

\ifanswer
\solbegin
\stepm{(1)}{当 $q$ 为真命题时，}
\stepm{①}{当 $a=0$ 时，不等式化为 $1>0$，对任意 $x\in\R$ 恒成立，满足条件；}
\stepm{②}{当 $a\neq0$ 时，$ax^2-ax+1>0$ 为一元二次不等式，}
\step{若对任意 $x\in\R$ 恒成立，需满足 $\begin{cases}a>0\\ \Delta=(-a)^2-4a<0\end{cases}$，解得 $0<a<4$；}
\step{综上，$a$ 的取值范围为 $[0,4)$.}
\stepm{(2)}{$p$ 为真时，方程 $x^2-2x+a=0$ 无实根，}
\step{故判别式 $\Delta=4-4a<0$，解得 $a>1$，因此 $p$ 为假时，$a\leqslant1$；}
\step{由（1）得 $q$ 为假时 $a\geqslant4$ 或 $a<0$，}
\step{$p$、$q$ 有且只有一个真命题，分两类：}
\step{当 $p$ 假 $q$ 真时，满足 $\begin{cases}a\leqslant1\\ 0\leqslant a<4\end{cases}$，解得 $0\leqslant a\leqslant1$；}
\step{当 $p$ 真 $q$ 假时，满足 $\begin{cases}a>1\\ a\geqslant4\text{或}a<0\end{cases}$，解得 $a\geqslant4$.}
\step{综上，$a$ 的取值范围是 $[0,1]\cup[4,+\infty)$.}
\solend
\fi

\ansspace[8]

\item 某工厂生产 $x$($x>0$) 件产品的总成本 $C$（单位：元）与生产件数 $x$ 之间满足二次函数
$C(x)=\dfrac14x^2+50x+10000$.

（1）若要求生产总成本不超过 30000 元，求满足生产总成本要求的产品件数的取值范围；

（2）试求该工厂生产多少件产品时，可使每件产品的平均成本最小.

\ifanswer
\solbegin
\stepm{(1)}{若要求生产总成本不超过 30000 元，}
\step{则 $C(x)=\dfrac14x^2+50x+10000\leqslant30000$，}
\step{即 $x^2+200x-80000\leqslant0$，解得 $-400\leqslant x\leqslant200$，又 $x>0$，}
\step{所以满足生产总成本要求的产品件数的取值范围为 $(0,200]$，$x\in\N$；}
\stepm{(2)}{生产件数 $x$ 时每件的平均成本为 $\dfrac{\frac14x^2+50x+10000}x=\dfrac14x+50+\dfrac{10000}x$，}
\step{由基本不等式得 $\dfrac14x+50+\dfrac{10000}x\geqslant2\sqrt{\dfrac14x\cdot\dfrac{10000}x}+50=150$，}
\step{当且仅当 $\dfrac14x=\dfrac{10000}x$，即 $x=200$ 时取等号，}
\step{即该工厂生产 200 件产品时，每件产品的平均成本最小.}
\solend
\fi

\ansspace[8]

% 原卷题干“$(n\geqslant2,n\in N)$”后连印两个全角逗号“，，”，照录未改（见文件头“原卷笔误”①）
\item 已知有限集 $A=\{a_1,a_2,\cdots,a_n\}$($n\geqslant2$,$n\in\N$)，，如果 $A$ 中元素
$a_i$($i=1,2,3,\cdots,n$) 满足 $a_1+a_2+a_3+\cdots+a_n=a_1\times a_2\times\cdots\times a_n$，
就称 $A$ 为“和积平衡集”.

% 原卷 (2) 将“写方程”与“反证法证明”两项任务并作一句（“……表示），可尝试用反证法证明：……”），
% 照录未改（见文件头“原卷笔误”②）
（1）判断集合 $\{-1+\sqrt3,-1-\sqrt3\}$ 是否为“和积平衡集”；

（2）若 $a_1$、$a_2$ 是两个不同的正数，且 $\{a_1,a_2\}$ 是“和积平衡集”，写出以 $a_1$、$a_2$ 为根的
一个一元二次方程（系数可用 $a_1$、$a_2$ 表示），可尝试用反证法证明：$a_1$、$a_2$ 至少有一个大于 2；

（3）若 $a_i\geqslant1$，且 $a_i\in\N$，求所有符合条件的“和积平衡集”.

\ifanswer
\solbegin
\stepm{(1)}{对于集合 $\{-1+\sqrt3,-1-\sqrt3\}$，}
\step{因为 $(-1+\sqrt3)+(-1-\sqrt3)=-2$，$(-1+\sqrt3)\times(-1-\sqrt3)=-2$，}
\step{所以满足 $a_1+a_2=a_1\times a_2$，满足题意，}
\step{所以集合 $\{-1+\sqrt3,-1-\sqrt3\}$ 是和积平衡集.}
\stepm{(2)}{以 $a_1$、$a_2$ 为根的一个一元二次方程可为 $x^2-(a_1+a_2)x+a_1\cdot a_2=0$，}
\step{因为 $a_1$、$a_2$ 是两个不同的正数，且 $\{a_1,a_2\}$ 是和积平衡集，}
\step{所以 $a_1+a_2=a_1\cdot a_2$，}
\step{又 $a_1$ 和 $a_2$ 为 $x^2-(a_1+a_2)x+a_1\cdot a_2=0$ 的两个不同的正根，}
\step{所以 $\Delta=(a_1+a_2)^2-4a_1a_2=(a_1+a_2)^2-4(a_1+a_2)>0$，}
\step{又 $a_1+a_2>0$，则 $a_1+a_2>4$，}
\step{若 $a_1$、$a_2$ 都小于等于 2，则 $a_1+a_2<4$，矛盾，}
\step{所以 $a_1$、$a_2$ 至少有一个大于 2；}
\stepm{(3)}{设 $A$ 中的 $a_1<a_2<a_3<\cdots<a_n$，且 $a_i\in\N^*$，}
\step{由 $a_1+a_2+a_3+\cdots+a_n=a_1\times a_2\times\cdots\times a_n$，得 $n=2$ 时明显不满足，}
\step{所以 $n\geqslant3$，故 $a_1\times a_2\times\cdots\times a_n=a_1+a_2+a_3+\cdots+a_n\leqslant na_n$，}
\step{所以 $a_1\times a_2\times\cdots\times a_{n-1}\leqslant n$，}
\step{当 $n=3$ 时，$a_1\times a_2\leqslant3$，$a_i\in\N^*$，只能有 $a_1=1$，$a_2=2$，}
\step{由 $a_1+a_2+a_3=a_1\times a_2\times a_3$，得 $a_3=3$，此时的和积平衡集为 $\{1,2,3\}$，}
\step{当 $n\geqslant4$ 时，$a_1\times a_2\times\cdots\times a_{n-1}\geqslant(n-1)(n-2)\times\cdots\times2\times1\geqslant(n-1)(n-2)$，}
\step{又 $(n-1)(n-2)-n=n^2-4n+2=n(n-4)+2>0$，即 $(n-1)(n-2)>n$，}
\step{即 $a_1\times a_2\times\cdots\times a_{n-1}>n$，与 $a_1\times a_2\times\cdots\times a_{n-1}\leqslant n$ 矛盾，}
\step{所以 $n\geqslant4$ 不满足条件，}
\step{所以符合条件的和积平衡集有且只有 $n=3$，此时的和积平衡集为 $\{1,2,3\}$.}
\solend
\fi

\ansspace*[10]

\end{enumerate}

\end{document}
"
%
% 原始材料：微信公众号“魔都数学加油站”文章，作者破晓，连载“27学年高一63”；
% 原文标题“27学年高一63——2026-2027年上海市吴淞中学高一上10月月考”；
% 原文链接：https://mp.weixin.qq.com/s/TFQ4TLnt_l66t7t51es4OQ。
% 页面标注发布时间 2026-10-16 21:44（上海）。
% 正文为 10 张 PNG 页（`原图/高一63_吴淞中学_10月月考/01..10.png`）：
% 第 06、10 页 2560×3620，第 01、02、03、04、05、07、08、09 页 4762×6735
% （**同一篇各页分辨率不同**），均按页序读取。10 页全为试卷正文页（卷首标题在第 1 页页顶，
% 卷末解答题第 21 题解析止于第 10 页），无封面页、目录页、二维码或宣传页。
% 粉色斜排“魔都数学加油站”为卷面水印，与公众号页脚一并未录入。
%
% 卷首标题：照卷面印作“2026--2027 年吴淞中学高一上 10 月月考”（公众号标题中的“上海市”
% 卷面标题行没有，本稿照卷面），年份区间按全稿惯例排 en dash。原卷只有一行标题、无副标题。
%
% 篇幅与题号：原卷自第 3 题起编号（第 1、2 题不在本卷内）。填空题 3—12（10 题）、
% 选择题 13—16（4 题）、解答题 17—21（5 题），共 19 题。本稿沿用原节与原题号，
% 只改排为全稿统一的大节样式（原卷小节标题为“一、填空题”“二、选择题”“三、解答题”）。
%
% 讲评版：原卷每题下直接印【答案】或【解析】。**只印【答案】的 2 题**（第 4、7 题），
% 本稿只给 \ans{}；其余 17 题里 16 题都印【解析】（按全稿口径这类题不另提 \ans{}——第 13、15、16 题
% 的结论写在解析末句（“故选 B”等），亦照录在解析内；第 14 题则既无【答案】也无【解析】段。**另有 1 题答案印在题干括号内**——
% 第 14 题（“(  D  )”，且该题无【答案】/【解析】段），本稿按 高一37 第 13 题成例把括号留空、
% 答案另提为 \ans{$D$}。第 16 题题干括号为空但印有【解析】，不算括号内答案。
%
% 副标题：原卷未印副标题。本稿按“知识模块”口径标作“集合与常用逻辑用语、不等式”：
% 19 题中集合与常用逻辑用语 12 题（第 3、4、6、7、9、10、11、12、13、15、18、21 题），
% 不等式 7 题（含等式（韦达定理、恒等式）与不等式性质、基本不等式、二次不等式、恒成立：
% 第 5、8、14、16、17、19、20 题），两类均超过三分之一，故并标。口径同 高一59。
%
% 与原卷的排版差异：
%   ① 原卷每题下直接印【答案】/【解析】，本稿按全稿惯例将答案与解析置于答案版，试卷版保持干净；
%   ② 原卷未印副标题，本稿按“知识模块”口径补出；
%   ③ 原卷填空横线约四个字宽，本稿按全稿惯例统一排 \blankline[4.4em]；
%   ④ 原卷实数集排斜体 $R$（第 9、10、19 题题干与解析）或粗体 $\mathbf{R}$（第 18 题题干），
%      整数集排斜体 $Z$（第 12 题解析），自然数集排斜体 $N$（第 20 题解析、第 21 题题干
%      “$(n\geqslant2,n\in N)$”）或粗体 $\mathbf{N}$（第 21 题 (3) 题干）、斜体 $N^*$
%      （第 21 题 (3) 解析），本稿按全稿惯例统一改用 $\R$、$\Z$、$\N$、$\N^*$
%      （同 高一14 排版差异④：记号归一，粗体／斜体不作区分，$N$ 与 $N^*$ 的区别照原卷保留）；
%      空集记号原卷作 $\varnothing$，照录；
%   ⑤ 原卷在数学式内多处用半角逗号（如 $\{y\mid y=[2x]+[3x],-1\leqslant x\leqslant2\}$、
%      $m=-1,0,1$、$A=\{x\mid-2<x<15\}$,$B=\cdots$），均照录未强改；
%      句末句点按全稿惯例排西文句点；“和积平衡集”自定义名词照录；
%   ⑥ 第 12 题解析里的“法一：”“法二：”标号原卷排**红色**（全卷共 2 处），
%      同 高一53“情况1：”成例，本稿排 \stepm{法一：}{…}、\stepm{法二：}{…}
%      （加粗着色、不改黑色基调），冒号照原卷排全角；
%   ⑦ 第 17—21 题的子问原卷各自另起一行，本稿按全稿惯例用 enumerate 的 (1)(2)(3)；
%      解析中的子问标号一律用 \stepm{(1)}{…}；第 19 题解析中的①②分情形，按全稿惯例作
%      \stepm{①}{…}、\stepm{②}{…}；
%   ⑧ 第 14 题印在题干括号内的答案移入 \ans{}，见上；
%   ⑨ 每道解答题原卷未留作答空白，本稿按全稿惯例加 \ansspace；
%   ⑩ 选择题的作答括号原卷排半角“(　　)”（第 13、15、16 题括号内空白、第 14 题括号内印 D），
%      本稿按全稿惯例排 \mbox{（\quad）}。
%
% 原卷笔误/存疑（均照录未改，正文对应处另有 % 注）：
%   ① 第 21 题题干“$(n\geqslant2$,$n\in N)$”后连印两个全角逗号“，，”（衍字，
%      8 倍放大确认，非水印干扰），本稿照录；
%   ② 第 21 题 (2) 题干将“写出以 $a_1$、$a_2$ 为根的一个一元二次方程”与“用反证法证明：
%      $a_1$、$a_2$ 至少有一个大于 2”两项任务并作一句、之间仅以逗号相接（“……表示），
%      可尝试用反证法证明：……”），读来像抄并两小问，原卷如此，照录未改。
%
% 插图：无（全卷 10 页无任何图形、无“如图”字样）。
%
\documentclass[a4paper,11pt]{article}
\usepackage{common}

\begin{document}

\examtitlest[3]{2026--2027 年吴淞中学高一上 10 月月考}{集合与常用逻辑用语、不等式}

\bigsec{一、填空题}

\begin{enumerate}
\setcounter{enumi}{2}

\item 满足 $A\subseteq\{1,2,3,4\}$ 的集合 $A$ 共有\blankline[4.4em]个.

\ifanswer
\solbegin
\step{由题意，集合 $A$ 是集合 $\{1,2,3,4\}$ 的子集，共有 $2^4=16$ 个.}
\solend
\fi

\item 陈述句“$a>1$ 或 $b>3$”的否定形式为\blankline[4.4em].

\ifanswer
\ans{$a\leqslant1$ 且 $b\leqslant3$}
\fi

\item 已知方程 $3x^2+4x-4=0$ 的两个根为 $x_1$、$x_2$，则 $x_1^2+x_2^2=$\blankline[4.4em].

\ifanswer
\solbegin
\step{对于方程 $3x^2+4x-4=0$，因为 $\Delta=16-4\times3\times(-4)=64>0$，}
\step{所以方程 $3x^2+4x-4=0$ 的两个不相等的实根为 $x_1$、$x_2$，}
\step{由韦达定理可得 $x_1+x_2=-\dfrac43$，$x_1x_2=-\dfrac43$，}
\step{所以 $x_1^2+x_2^2=(x_1+x_2)^2-2x_1x_2=\dfrac{16}9+\dfrac83=\dfrac{40}9$.}
\solend
\fi

\item 设集合 $A=\{y\mid y=x^2+1\}$,$B=\{y\mid y=-x^2+2x+3\}$，则 $A\cap B=$\blankline[4.4em].

\ifanswer
\solbegin
\step{因为 $y=x^2+1\geqslant1$,$y=-x^2+2x+3=-(x-1)^2+4\leqslant4$，}
\step{所以 $A=\{y\mid y\geqslant1\}$,$B=\{y\mid y\leqslant4\}$，得 $A\cap B=\{y\mid1\leqslant y\leqslant4\}$.}
\solend
\fi

\item 若 $x\in\{1,x^2\}$，则 $x=$\blankline[4.4em].

\ifanswer
\ans{$0$}
\fi

\item 已知等式 $2x^2+3x+5=a(2x+1)(x+1)+c$，则 $a+c=$\blankline[4.4em].

\ifanswer
\solbegin
\step{因为 $2x^2+3x+5=a(2x+1)(x+1)+c$ 恒成立，}
\step{所以 $2x^2+3x+5=2ax^2+3ax+a+c$ 恒成立，则 $\begin{cases}2=2a\\ 3=3a\\ 5=a+c\end{cases}$，解得 $a+c=5$.}
\solend
\fi

\item 已知 $a\in\R$，设 $p:x<2$，$q:x<5a-3$，且 $p$ 是 $q$ 的必要非充分条件，则 $a$ 的取值范围是\blankline[4.4em].

\ifanswer
\solbegin
\step{设 $p:x<2$，$q:x<5a-3$，且 $p$ 是 $q$ 的必要非充分条件，}
\step{则 $\{x\mid x<5a-3\}\subset\{x\mid x<2\}$，所以 $5a-3<2$，即 $a<1$.}
\solend
\fi

\item 已知 $a\in\R$，集合 $A=\{x\mid ax^2+3x-1=0\}$，若集合 $A$ 中只有一个元素，则 $a$ 的取值范围是\blankline[4.4em].

\ifanswer
\solbegin
\step{当 $a=0$ 时，则有 $\{x\mid ax^2+3x-1=0\}=\{x\mid3x-1=0\}=\left\{\dfrac13\right\}$，符合题意；}
\step{当 $a\neq0$ 时，由题意得 $\Delta=9+4a=0$，解得 $a=-\dfrac94$. 符合题意；}
\step{综上所述，实数 $a$ 的取值集合为 $\left\{0,-\dfrac94\right\}$.}
\solend
\fi

\item 设集合 $A=\{x\mid x^2+2(m+1)x+m^2-1=0,m\in\R\}$，且 $A\subseteq\{0,-4\}$，则实数 $m$ 的取值范围是\blankline[4.4em].

\ifanswer
\solbegin
\step{由题意得 $A=\varnothing$ 或 $A=\{0\}$ 或 $A=\{-4\}$ 或 $A=\{0,-4\}$，}
\step{当 $A=\varnothing$ 时，方程 $x^2+2(m+1)x+m^2-1=0$ 无实根，}
\step{则 $\Delta=[2(m+1)]^2-4\times1\times(m^2-1)=8m+8<0$，解得 $m<-1$，}
\step{当 $A=\{0\}$ 时，方程 $x^2+2(m+1)x+m^2-1=0$ 有一个实根 $0$，}
\step{则 $\begin{cases}\Delta=8m+8=0\\ m^2-1=0\end{cases}$，解得 $m=-1$，}
\step{当 $A=\{-4\}$ 时，方程 $x^2+2(m+1)x+m^2-1=0$ 有一个实根 $-4$，}
\step{则 $\begin{cases}\Delta=8m+8=0\\ 16-8(m+1)+m^2-1=0\end{cases}$，无解，}
\step{当 $A=\{0,-4\}$ 时，方程 $x^2+2(m+1)x+m^2-1=0$ 有两个实根 $-4$ 和 $0$，}
\step{则 $\begin{cases}\Delta=8m+8>0\\ 16-8(m+1)+m^2-1=0\\ m^2-1=0\end{cases}$，解得 $m=1$，}
\step{综上所述，实数 $m$ 的取值范围是 $\{m\mid m\leqslant-1\text{ 或 }m=1\}$.}
\solend
\fi

\item 对于任意实数 $x$，$[x]$ 表示不超过 $x$ 的最大整数，如 $[1.1]=1$，$[-1.2]=-2$，若
$A=\{y\mid y=[2x]+[3x],-1\leqslant x\leqslant2\}$，则 $A$ 中所有元素的和为\blankline[4.4em].

\ifanswer
\solbegin
% 原卷“法一：”“法二：”标号排红色，本稿按 高一53“情况1：”成例排 \stepm（加粗着色），见文件头
\stepm{法一：}{当 $-1\leqslant x<-\dfrac23$ 时，$[2x]=-2$，$[3x]=-3$，$f(x)=[2x]+[3x]=-5$，}
\step{当 $-\dfrac23\leqslant x<-\dfrac12$ 时，$[2x]=-2$，$[3x]=-2$，$f(x)=[2x]+[3x]=-4$；}
\step{当 $-\dfrac12\leqslant x<-\dfrac13$ 时，$[2x]=-1$，$[3x]=-2$，$f(x)=[2x]+[3x]=-3$；}
\step{当 $-\dfrac13\leqslant x<0$ 时，$[2x]=-1$，$[3x]=-1$，$f(x)=[2x]+[3x]=-2$；}
\step{当 $0\leqslant x<\dfrac13$ 时，$[2x]=0$，$[3x]=0$，$f(x)=[2x]+[3x]=0$；}
\step{当 $\dfrac13\leqslant x<\dfrac12$ 时，$[2x]=0$，$[3x]=1$，$f(x)=[2x]+[3x]=1$；}
\step{当 $\dfrac12\leqslant x<\dfrac23$ 时，$[2x]=1$，$[3x]=1$，$f(x)=[2x]+[3x]=2$；}
\step{当 $\dfrac23\leqslant x<1$ 时，$[2x]=1$，$[3x]=2$，$f(x)=[2x]+[3x]=3$；}
\step{当 $1\leqslant x<\dfrac43$ 时，$[2x]=2$，$[3x]=3$，$f(x)=[2x]+[3x]=5$；}
\step{当 $\dfrac43\leqslant x<\dfrac32$ 时，$[2x]=2$，$[3x]=4$，$f(x)=[2x]+[3x]=6$；}
\step{当 $\dfrac32\leqslant x<\dfrac53$ 时，$[2x]=3$，$[3x]=4$，$f(x)=[2x]+[3x]=7$；}
\step{当 $\dfrac53\leqslant x<2$ 时，$[2x]=3$，$[3x]=5$，$f(x)=[2x]+[3x]=8$；}
\step{当 $x=2$ 时，$[2x]=4$，$[3x]=6$，$f(x)=[2x]+[3x]=10$}
\step{故 $A$ 中所有元素的和为 $-5-4-3-2+0+1+2+3+5+6+7+8+10=28$.}
% 原卷“法二：”标号排红色，同上
\stepm{法二：}{设 $x=m+n$，其中 $m\in\Z$，$n\in[0,1)$，}
\step{若 $0\leqslant n<\dfrac13$，则 $0\leqslant2n<\dfrac23$，$0\leqslant3n<1$，}
\step{得 $[2x]=[2m+2n]=2m$，$[3x]=[3m+3n]=3m$，所以 $y=[2x]+[3x]=5m$；}
\step{若 $\dfrac13\leqslant n<\dfrac12$，则 $\dfrac23\leqslant2n<1$，$1\leqslant3n<\dfrac32$，}
\step{得 $[2x]=[2m+2n]=2m$，$[3x]=[3m+3n]=3m+1$，}
\step{所以 $y=[2x]+[3x]=5m+1$；}
\step{若 $\dfrac12\leqslant n<\dfrac23$，则 $1\leqslant2n<\dfrac43$，$\dfrac32\leqslant3n<2$，}
\step{得 $[2x]=[2m+2n]=2m+1$，$[3x]=[3m+3n]=3m+1$，}
\step{所以 $y=[2x]+[3x]=5m+2$；}
\step{若 $\dfrac23\leqslant n<1$，则 $\dfrac43\leqslant2n<2$，$2\leqslant3n<3$，}
\step{得 $[2x]=[2m+2n]=2m+1$，$[3x]=[3m+3n]=3m+2$，}
\step{所以 $y=[2x]+[3x]=5m+3$；}
\step{综上所述，当 $n\in[0,1)$ 时，所有函数值的和}
\step{为 $5m+5m+1+5m+2+5m+3=20m+6$，}
\step{又因为 $x\in[-1,2]$，此时 $m=-1,0,1$，且当 $x=2$ 时，$y=[2x]+[3x]=10$，}
\step{所有函数值的和为 $(-14)+6+26+10=28$，即 $A$ 中所有元素之和为 $28$.}
\solend
\fi

\end{enumerate}

\bigsec{二、选择题}

\begin{enumerate}
\setcounter{enumi}{12}

\item 下列写法中，正确的是\mbox{（\quad）}

\keepwithprev
A. $0\in\{\varnothing\}$\qquad B. $\varnothing\in\{\varnothing\}$

C. $\varnothing=\{\varnothing\}$\qquad D. $0\in\varnothing$

\ifanswer
\solbegin
\step{$0\notin\{\varnothing\}$，$\varnothing\in\{\varnothing\}$，故 $A$ 错误，$B$ 正确；}
\step{$\varnothing$ 是不含任何元素的集合，所以 $\varnothing\neq\{\varnothing\}$，$0\notin\varnothing$，故 $CD$ 错误.}
\step{故选 $B$.}
\solend
\fi

% 第 14 题答案“D”直接印在题干括号内“(  D  )”，且该题无【答案】/【解析】段；
% 按 高一37 第 13 题的成例括号留空、答案另提为 \ans{}（见文件头）
\item 若 $a<b<0$，则下列不等式中不能成立的是\mbox{（\quad）}

\keepwithprev
A. $\dfrac1a>\dfrac1b$\qquad B. $a-b>a$

C. $|a|>|b|$\qquad D. $a^3>b^3$

\ifanswer
\ans{$D$}
\fi

\item 已知 $x$、$y$、$z$ 为非零实数，代数式 $\dfrac x{|x|}+\dfrac y{|y|}+\dfrac z{|z|}+\dfrac{xyz}{|xyz|}$
的值所组成的集合是 $M$，则下列判断正确的是\mbox{（\quad）}

\keepwithprev
A. $-4\notin M$\qquad B. $2\in M$

C. $4\notin M$\qquad D. $0\in M$

\ifanswer
\solbegin
\step{当 $x$、$y$、$z$ 均为正数时，$xyz>0$，$\dfrac x{|x|}+\dfrac y{|y|}+\dfrac z{|z|}+\dfrac{xyz}{|xyz|}=1+1+1+1=4$，}
\step{当 $x$、$y$、$z$ 有 2 个正数，1 个负数时，不妨设 $x>0$，$y>0$，$z<0$，}
\step{则 $xyz<0$，$\dfrac x{|x|}+\dfrac y{|y|}+\dfrac z{|z|}+\dfrac{xyz}{|xyz|}=1+1-1-1=0$，}
\step{若 $x$、$y$、$z$ 有 1 个正数，2 个负数时，不妨设 $x>0$，$y<0$，$z<0$，}
\step{则 $xyz>0$，$\dfrac x{|x|}+\dfrac y{|y|}+\dfrac z{|z|}+\dfrac{xyz}{|xyz|}=1-1-1+1=0$，}
\step{当 $x$、$y$、$z$ 均为负数时，$xyz<0$，}
\step{$\dfrac x{|x|}+\dfrac y{|y|}+\dfrac z{|z|}+\dfrac{xyz}{|xyz|}=-1-1-1-1=-4$，}
\step{综上，$M=\{-4,0,4\}$，}
\step{所以 $-4\in M$，$A$ 错误；$2\notin M$，$B$ 错误；$4\in M$，$C$ 错误；}
\step{$0\in M$，$D$ 正确. 故选 $D$.}
\solend
\fi

\item 若 $xy>0$，且 $x+y=9$，则 $\dfrac1x+\dfrac4y$ 的最小值是\mbox{（\quad）}

\keepwithprev
A. $9$\qquad B. $\dfrac59$

C. $\dfrac49$\qquad D. $1$

\ifanswer
\solbegin
\step{$xy>0$，且 $x+y=9$，所以 $x>0$，$y>0$，}
\step{所以 $\dfrac1x+\dfrac4y=\dfrac19\left(\dfrac1x+\dfrac4y\right)(x+y)=\dfrac59+\dfrac19\left(\dfrac yx+\dfrac{4x}y\right)\geqslant\dfrac59+\dfrac29\cdot\sqrt{\dfrac yx\cdot\dfrac{4x}y}=\dfrac59+\dfrac49=1$，}
\step{当且仅当 $\dfrac yx=\dfrac{4x}y$，即 $x=3$，$y=6$ 时取等号，所以 $\dfrac1x+\dfrac4y$ 的最小值为 $1$.}
\step{故选 $D$.}
\solend
\fi

\end{enumerate}

\bigsec{三、解答题}

\begin{enumerate}
\setcounter{enumi}{16}

\item 求解下列不等式：

（1）$\dfrac{x-3}{x+1}\leqslant0$；

（2）$|2x-1|<3$.

\ifanswer
\solbegin
\stepm{(1)}{$\dfrac{x-3}{x+1}\leqslant0$，变形得 $(x-3)\cdot(x+1)\leqslant0$ 且 $x\neq-1$，}
\step{解得 $-1<x\leqslant3$，即不等式的解集为 $(-1,3]$.}
\stepm{(2)}{若 $|2x-1|<3$，则 $-3<2x-1<3$，解得 $-1<x<2$，即 $(-1,2)$.}
\solend
\fi

\ansspace[6]

\item 已知集合 $A=\{x\mid-2<x<15\}$,$B=\{x\mid m-6<x<2m-1,m\in\R\}$.

（1）当 $m=2$ 时，求 $A\cap B$；

（2）若 $A\cup B=A$，求实数 $m$ 的取值范围.

\ifanswer
\solbegin
\stepm{(1)}{当 $m=2$ 时，$A=\{x\mid-2<x<15\}$,$B=\{x\mid-4<x<3\}$，}
\step{此时 $A\cap B=\{x\mid-2<x<3\}$.}
\stepm{(2)}{由 $A\cup B=A$ 得 $B\subseteq A$.}
\step{当 $B=\varnothing$,$m-6\geqslant2m-1$，解得 $m\leqslant-5$；}
\step{当 $B\neq\varnothing$，结合 $B\subseteq A$ 得 $\begin{cases}m-6<2m-1\\ 2m-1\leqslant15\\ m-6\geqslant-2\end{cases}$，解得 $4\leqslant m\leqslant8$.}
\step{综上，$m\leqslant-5$ 或 $4\leqslant m\leqslant8$.}
\solend
\fi

\ansspace[6]

\item 已知 $a\in\R$，命题 $p$：关于 $x$ 的方程 $x^2-2x+a=0$ 无实根；命题 $q$：对任意 $x\in\R$，
不等式 $ax^2-ax+1>0$ 恒成立.

（1）若 $q$ 为真命题，求 $a$ 的取值范围；

（2）若 $p$、$q$ 中有且只有一个为真命题，求 $a$ 的取值范围.

\ifanswer
\solbegin
\stepm{(1)}{当 $q$ 为真命题时，}
\stepm{①}{当 $a=0$ 时，不等式化为 $1>0$，对任意 $x\in\R$ 恒成立，满足条件；}
\stepm{②}{当 $a\neq0$ 时，$ax^2-ax+1>0$ 为一元二次不等式，}
\step{若对任意 $x\in\R$ 恒成立，需满足 $\begin{cases}a>0\\ \Delta=(-a)^2-4a<0\end{cases}$，解得 $0<a<4$；}
\step{综上，$a$ 的取值范围为 $[0,4)$.}
\stepm{(2)}{$p$ 为真时，方程 $x^2-2x+a=0$ 无实根，}
\step{故判别式 $\Delta=4-4a<0$，解得 $a>1$，因此 $p$ 为假时，$a\leqslant1$；}
\step{由（1）得 $q$ 为假时 $a\geqslant4$ 或 $a<0$，}
\step{$p$、$q$ 有且只有一个真命题，分两类：}
\step{当 $p$ 假 $q$ 真时，满足 $\begin{cases}a\leqslant1\\ 0\leqslant a<4\end{cases}$，解得 $0\leqslant a\leqslant1$；}
\step{当 $p$ 真 $q$ 假时，满足 $\begin{cases}a>1\\ a\geqslant4\text{或}a<0\end{cases}$，解得 $a\geqslant4$.}
\step{综上，$a$ 的取值范围是 $[0,1]\cup[4,+\infty)$.}
\solend
\fi

\ansspace[8]

\item 某工厂生产 $x$($x>0$) 件产品的总成本 $C$（单位：元）与生产件数 $x$ 之间满足二次函数
$C(x)=\dfrac14x^2+50x+10000$.

（1）若要求生产总成本不超过 30000 元，求满足生产总成本要求的产品件数的取值范围；

（2）试求该工厂生产多少件产品时，可使每件产品的平均成本最小.

\ifanswer
\solbegin
\stepm{(1)}{若要求生产总成本不超过 30000 元，}
\step{则 $C(x)=\dfrac14x^2+50x+10000\leqslant30000$，}
\step{即 $x^2+200x-80000\leqslant0$，解得 $-400\leqslant x\leqslant200$，又 $x>0$，}
\step{所以满足生产总成本要求的产品件数的取值范围为 $(0,200]$，$x\in\N$；}
\stepm{(2)}{生产件数 $x$ 时每件的平均成本为 $\dfrac{\frac14x^2+50x+10000}x=\dfrac14x+50+\dfrac{10000}x$，}
\step{由基本不等式得 $\dfrac14x+50+\dfrac{10000}x\geqslant2\sqrt{\dfrac14x\cdot\dfrac{10000}x}+50=150$，}
\step{当且仅当 $\dfrac14x=\dfrac{10000}x$，即 $x=200$ 时取等号，}
\step{即该工厂生产 200 件产品时，每件产品的平均成本最小.}
\solend
\fi

\ansspace[8]

% 原卷题干“$(n\geqslant2,n\in N)$”后连印两个全角逗号“，，”，照录未改（见文件头“原卷笔误”①）
\item 已知有限集 $A=\{a_1,a_2,\cdots,a_n\}$($n\geqslant2$,$n\in\N$)，，如果 $A$ 中元素
$a_i$($i=1,2,3,\cdots,n$) 满足 $a_1+a_2+a_3+\cdots+a_n=a_1\times a_2\times\cdots\times a_n$，
就称 $A$ 为“和积平衡集”.

% 原卷 (2) 将“写方程”与“反证法证明”两项任务并作一句（“……表示），可尝试用反证法证明：……”），
% 照录未改（见文件头“原卷笔误”②）
（1）判断集合 $\{-1+\sqrt3,-1-\sqrt3\}$ 是否为“和积平衡集”；

（2）若 $a_1$、$a_2$ 是两个不同的正数，且 $\{a_1,a_2\}$ 是“和积平衡集”，写出以 $a_1$、$a_2$ 为根的
一个一元二次方程（系数可用 $a_1$、$a_2$ 表示），可尝试用反证法证明：$a_1$、$a_2$ 至少有一个大于 2；

（3）若 $a_i\geqslant1$，且 $a_i\in\N$，求所有符合条件的“和积平衡集”.

\ifanswer
\solbegin
\stepm{(1)}{对于集合 $\{-1+\sqrt3,-1-\sqrt3\}$，}
\step{因为 $(-1+\sqrt3)+(-1-\sqrt3)=-2$，$(-1+\sqrt3)\times(-1-\sqrt3)=-2$，}
\step{所以满足 $a_1+a_2=a_1\times a_2$，满足题意，}
\step{所以集合 $\{-1+\sqrt3,-1-\sqrt3\}$ 是和积平衡集.}
\stepm{(2)}{以 $a_1$、$a_2$ 为根的一个一元二次方程可为 $x^2-(a_1+a_2)x+a_1\cdot a_2=0$，}
\step{因为 $a_1$、$a_2$ 是两个不同的正数，且 $\{a_1,a_2\}$ 是和积平衡集，}
\step{所以 $a_1+a_2=a_1\cdot a_2$，}
\step{又 $a_1$ 和 $a_2$ 为 $x^2-(a_1+a_2)x+a_1\cdot a_2=0$ 的两个不同的正根，}
\step{所以 $\Delta=(a_1+a_2)^2-4a_1a_2=(a_1+a_2)^2-4(a_1+a_2)>0$，}
\step{又 $a_1+a_2>0$，则 $a_1+a_2>4$，}
\step{若 $a_1$、$a_2$ 都小于等于 2，则 $a_1+a_2<4$，矛盾，}
\step{所以 $a_1$、$a_2$ 至少有一个大于 2；}
\stepm{(3)}{设 $A$ 中的 $a_1<a_2<a_3<\cdots<a_n$，且 $a_i\in\N^*$，}
\step{由 $a_1+a_2+a_3+\cdots+a_n=a_1\times a_2\times\cdots\times a_n$，得 $n=2$ 时明显不满足，}
\step{所以 $n\geqslant3$，故 $a_1\times a_2\times\cdots\times a_n=a_1+a_2+a_3+\cdots+a_n\leqslant na_n$，}
\step{所以 $a_1\times a_2\times\cdots\times a_{n-1}\leqslant n$，}
\step{当 $n=3$ 时，$a_1\times a_2\leqslant3$，$a_i\in\N^*$，只能有 $a_1=1$，$a_2=2$，}
\step{由 $a_1+a_2+a_3=a_1\times a_2\times a_3$，得 $a_3=3$，此时的和积平衡集为 $\{1,2,3\}$，}
\step{当 $n\geqslant4$ 时，$a_1\times a_2\times\cdots\times a_{n-1}\geqslant(n-1)(n-2)\times\cdots\times2\times1\geqslant(n-1)(n-2)$，}
\step{又 $(n-1)(n-2)-n=n^2-4n+2=n(n-4)+2>0$，即 $(n-1)(n-2)>n$，}
\step{即 $a_1\times a_2\times\cdots\times a_{n-1}>n$，与 $a_1\times a_2\times\cdots\times a_{n-1}\leqslant n$ 矛盾，}
\step{所以 $n\geqslant4$ 不满足条件，}
\step{所以符合条件的和积平衡集有且只有 $n=3$，此时的和积平衡集为 $\{1,2,3\}$.}
\solend
\fi

\ansspace*[10]

\end{enumerate}

\end{document}
